% !TeX spellcheck = es_ES
\documentclass[oneside,11pt,parskip=full]{scrbook}

% Incluir sólo algunos capítulos/ficheros. Útil para trabajar sólo en lo que interesa en cada momento.
%\includeonly{teoria-conjuntos-difusos,logica-difusa}

% Utilizar UTF-8 como codificación de entrada
\usepackage[utf8]{inputenc}

% Traducción al español de palabras como "Bibiliografía","Figura" etc.
\usepackage[spanish, es-tabla]{babel}

% Fuentes
\usepackage[T1]{fontenc}
\usepackage{palatino}
\addtokomafont{disposition}{\rmfamily}
\addtokomafont{title}{\rmfamily}

% Márgenes
\usepackage[margin=1in]{geometry}

% Paquetes para generar la página de título
\usepackage[pdftex]{graphicx}
\usepackage{epstopdf,epsfig}

% Texto de pruebas
\usepackage{lipsum}

% Código fuente
\usepackage{listings}
\usepackage{amsthm}
\usepackage{amssymb}

% Índice
\usepackage{makeidx}

% Enumeraciones
\usepackage{enumitem}

% Gráficos
\usepackage{tikz}
\usetikzlibrary{shapes,shadows,arrows}
\usepackage{pgfplots}
\pgfplotsset{compat=newest}
\usepackage{amsmath}

% Tablas
\usepackage{longtable}
\usepackage{multirow}
\usepackage{hhline}

%Algoritmos
\usepackage[linesnumbered,ruled]{algorithm2e}
\renewcommand*{\algorithmcfname}{Algoritmo}
\let\oldnl\nl% Store \nl in \oldnl
\newcommand{\nonl}{\renewcommand{\nl}{\let\nl\oldnl}}% Remove line number for one line

% Float: poner las figuras en un punto exacto
\usepackage{float}
\usepackage{subcaption}

\usepackage[bookmarks]{hyperref}

% Código MATLAB
\usepackage[framed,numbered,autolinebreaks]{mcode}
\usepackage{listingsutf8}
% Transliterar acentos
\lstset{
    literate={á}{{\'a}}1 {é}{{\'e}}1 {í}{{\'i}}1 {ó}{{\'o}}1 {ú}{{\'u}}1 
    			  {Á}{{\'A}}1 {É}{{\'E}}1 {Í}{{\'I}}1 {Ó}{{\'O}}1 {Ú}{{\'U}}1
    			  {ü}{{\"{u}}}1 {Ü}{{\"{U}}}1 
    			  {ñ}{{\~n}}1,
}

% Comandos
\newcommand{\HRule}{\rule{\linewidth}{0.5mm}}

% Nombre de sección en cursiva
\renewcommand*{\headfont}{\normalfont\itshape}
% Nº de página en negrita
\renewcommand*{\pnumfont}{\normalfont\bfseries}

% Nombre para el caption de código fuente
\renewcommand{\lstlistingname}{Código}

% Teoremas, definiciones, ejemplos
\newtheorem{theorem}{Teorema}[chapter]
\newtheorem{example}{Ejemplo}[chapter]
\newtheorem{proposition}{Proposición}[chapter]
\newtheorem{definition}{Definición}[chapter]

% Estilos del header y el footer
\usepackage{scrpage2}
\pagestyle{scrheadings}
\ihead{}
\chead{}
\ohead[]{\headmark}
\ifoot[]{}
\cfoot[]{}
\ofoot[\pagemark]{\pagemark}

% Información general del documento
\title{Documento de ejemplo}
\author{Mikel Pintor}
\date{Febrero 2013}
\titlehead{Un draft de ejemplo escrito en LaTeX}
\publishers{Escuela de ingenieros industriales y de telecomunicaciones\\Universidad de Navarra}

\makeindex

\begin{document}

\begin{titlepage}
\begin{center}

% Logos de la universidad y la escuela
\begin{minipage}[t]{0.15\textwidth}
\includegraphics[width=\linewidth]{./images/logo_upna}
\end{minipage}
\hfill
\begin{minipage}[t]{0.15\textwidth}
\includegraphics[width=\linewidth]{./images/logo_escuela}
\end{minipage}\\[2cm]

\textsc{\LARGE ESCUELA TÉCNICA SUPERIOR DE INGENIEROS INDUSTRIALES Y DE TELECOMUNICACIÓN}\\[1.5cm]

\textsc{\Large Ingeniería Informática}\\[0.5cm]

% Título
\HRule \\[0.4cm]
{ \huge \bfseries Aplicación de un método de interpolación basado en índices de solapamiento a la detección de riesgos ambientales\\[0.4cm] }

\HRule \\[1.5cm]

% Autor y tutor
\vfill
\hfill
\begin{minipage}{0.4\textwidth}
\begin{flushright} \large
\emph{Autor:} Mikel \textsc{Pintor}\\
\emph{Tutores:} Humberto \textsc{Bustince}\\
Fco. Javier \textsc{Fernández}
{\large Pamplona,\today}
\end{flushright}
\end{minipage}
\vfill

\end{center}
\end{titlepage}
\frontmatter
\tableofcontents
\listoftables
\listoffigures

% !TeX spellcheck = es_ES
\chapter{Introducción}
\label{cha:introduccion}
\section{Conjuntos difusos}
En esta sección se presentan algunas definiciones y teoremas que forman la base de toda la teoría de lógica difusa y de los teoremas y métodos presentados en este trabajo.

\begin{definition}
Sea \emph{X} un conjunto clásico de objetos, denominado \emph{universo}, cuyos elementos son denotados como \emph{x}. La pertenencia a un subconjunto clásico \emph{A} de \emph{X} se expresa normalmente por medio de una función de pertenencia $\mu_{A}:X\rightarrow\{0,1\}$ tal que:
\begin{equation}
\mu_{A}(x)=\begin{cases} 1 & \mbox{iff } \mbox{x}\in A, \\ 0 & \mbox{iff } \mbox{x}\notin A \end{cases}
\end{equation}
(se utiliza \emph{iff} como abreviatura de "si y sólo si")
\end{definition}

De esta forma, en un conjunto clásico, se dice que un elemento \emph{x} pertenece al conjunto A ($\mu_{A}(x)=1$) o no pertenece ($\mu_{A}(x)=0$). 
\begin{example}
\normalfont
Un ejemplo de conjunto clásico es el conjunto de todos los números enteros pares, cuya función de pertenencia $\mu_{A}:\mathbb{Z}\rightarrow\{0,1\}$ podría expresarse como:
\begin{equation}
\mu_{A}(x)=\begin{cases} 1 & \mbox{iff } \mbox{x}\bmod 2=0, \\ 0 & \mbox{iff } \mbox{x}\bmod 2\ne 0 \end{cases}
\end{equation}
\end{example}
Si el valor de salida de la función de pertenencia puede ser cualquier número en el intervalo [0,1] (y no sólo los valores discretos \{0,1\}) se dice que el conjunto \emph{A} es \emph{difuso} y que la función $\mu_{A}(x)$ es el \emph{grado de pertenencia} de \emph{x} a \emph{A}. En este tipo de conjuntos, cuanto más próximo sea el valor de $\mu_{A}(x)$ a 1, más pertenece \emph{x} a \emph{A}.

El concepto de conjunto difuso fue introducido por L.A. Zadeh en 1965 \cite{Zadeh65} y se define como:
\begin{definition}
Dado un conjunto de referencia (o universo) \emph{U}, un conjunto difuso \emph{A} sobre \emph{U} es un conjunto tal que:\\
\begin{equation}
\{(u_{i},\mu_{A}(u_{i}))\arrowvert u_{i} \in U\}
\end{equation}
donde \begin{math}\emph{A}:\emph{U}\rightarrow[0,1]\end{math} es la \emph{función de pertenencia}\index{función de pertenencia} (o \emph{grado de pertenencia}) de \emph{A}.
\end{definition}

\section{T-normas y operadores de agregación}\label{sec:t-norms-aggregation-operator}
Dada una función \begin{math}\emph{F}:[0,1]^2\rightarrow[0,1]\end{math} y dos conjuntos difusos \emph{A} y \emph{B} definidos sobre el mismo universo \emph{U}, \begin{math}A,B\in FS(U)\end{math}, denotamos como \begin{math}F(A,B)\end{math} el conjunto difuso sobre \emph{U} cuya función de pertenencia viene dada por:
\begin{equation}
\mu_{F(A,B)}(u_{i}) = F(A(u_{i}),B(u_{i}))
\end{equation}
Una clase importante de este tipo de funciones son las llamadas t-normas y t-conormas (triangular norms) \cite{klement2000triangular}.
\begin{definition}\label{def:tnorma}
Una t-norma es una operación binaria \emph{T} en el intervalo $[0,1]$ que es conmutativa, asociativa, monótona y tiene el valor \emph{1} como elemento neutro. Es decir, una función $\emph{T} : [0,1]^2 \rightarrow [0,1]$ tal que $\forall x,y,z \in [0,1]$:
\begin{enumerate}[label=(T\arabic*),ref=(T\arabic*)]
   \item $T(x,y) = T(y,x)$ (Conmutatividad)\label{T1}
   \item $T(x,T(y,z)) = T(T(x,y),z)$ (Asociatividad)\label{T2}
   \item $T(x,y) \leq T(x,z)$ cuando $y \leq z$ (Monotonía)\label{T3}
   \item $T(x,1) = x$ (Elemento neutro)\label{T4}
  \end{enumerate}
\end{definition}
Estas propiedades son suficientes para garantizar que las t-normas generalizan la conjunción clásica ($x\bigwedge y$) cuando se aplican a valores booleanos ($T(0,0)=0$,$T(0,1)=T(1,0)=0$ y $T(1,1)=1$).
Algunos ejemplos prominentes de t-normas son los siguientes:
\begin{itemize}
	\item \bfseries Mínimo o t-norma de Gödel: $T_{G}(x,y) = min\{x,y\}$
	\item \bfseries Producto: $T_{P}(x,y) = xy$
	\item \bfseries \L{}ukasiewicz: $T_{L}(x,y) = max\{x+y-1,0\}$
\end{itemize}
\begin{figure}[H]
	\centering
	\begin{subfigure}[b]{0.45\textwidth}
		\newlength\figureheight 
		\newlength\figurewidth
		\setlength\figureheight{4.5cm}
		\setlength\figurewidth{6cm}
		% This file was created by matlab2tikz v0.4.7 (commit 4d7ae5c4fd0932fb51051d86111bc23ed23e4580) running on MATLAB 8.0.
% Copyright (c) 2008--2014, Nico Schlömer <nico.schloemer@gmail.com>
% All rights reserved.
% Minimal pgfplots version: 1.3
% 
\begin{tikzpicture}

\begin{axis}[%
width=\figurewidth,
height=\figureheight,
view={-37.5}{30},
scale only axis,
xmin=0,
xmax=1,
xlabel={x},
xmajorgrids,
ymin=0,
ymax=1,
ylabel={y},
ymajorgrids,
zmin=0,
zmax=1,
zlabel={$\text{T}_{\text{G}}\text{(x,y)}$},
zmajorgrids,
axis x line*=bottom,
axis y line*=left,
axis z line*=left
]

\addplot3[%
surf,
shader=faceted,
draw=black,
colormap/jet,
mesh/rows=11]
table[row sep=crcr,header=false] {0	0	0\\
0	0.1	0\\
0	0.2	0\\
0	0.3	0\\
0	0.4	0\\
0	0.5	0\\
0	0.6	0\\
0	0.7	0\\
0	0.8	0\\
0	0.9	0\\
0	1	0\\
0.1	0	0\\
0.1	0.1	0.1\\
0.1	0.2	0.1\\
0.1	0.3	0.1\\
0.1	0.4	0.1\\
0.1	0.5	0.1\\
0.1	0.6	0.1\\
0.1	0.7	0.1\\
0.1	0.8	0.1\\
0.1	0.9	0.1\\
0.1	1	0.1\\
0.2	0	0\\
0.2	0.1	0.1\\
0.2	0.2	0.2\\
0.2	0.3	0.2\\
0.2	0.4	0.2\\
0.2	0.5	0.2\\
0.2	0.6	0.2\\
0.2	0.7	0.2\\
0.2	0.8	0.2\\
0.2	0.9	0.2\\
0.2	1	0.2\\
0.3	0	0\\
0.3	0.1	0.1\\
0.3	0.2	0.2\\
0.3	0.3	0.3\\
0.3	0.4	0.3\\
0.3	0.5	0.3\\
0.3	0.6	0.3\\
0.3	0.7	0.3\\
0.3	0.8	0.3\\
0.3	0.9	0.3\\
0.3	1	0.3\\
0.4	0	0\\
0.4	0.1	0.1\\
0.4	0.2	0.2\\
0.4	0.3	0.3\\
0.4	0.4	0.4\\
0.4	0.5	0.4\\
0.4	0.6	0.4\\
0.4	0.7	0.4\\
0.4	0.8	0.4\\
0.4	0.9	0.4\\
0.4	1	0.4\\
0.5	0	0\\
0.5	0.1	0.1\\
0.5	0.2	0.2\\
0.5	0.3	0.3\\
0.5	0.4	0.4\\
0.5	0.5	0.5\\
0.5	0.6	0.5\\
0.5	0.7	0.5\\
0.5	0.8	0.5\\
0.5	0.9	0.5\\
0.5	1	0.5\\
0.6	0	0\\
0.6	0.1	0.1\\
0.6	0.2	0.2\\
0.6	0.3	0.3\\
0.6	0.4	0.4\\
0.6	0.5	0.5\\
0.6	0.6	0.6\\
0.6	0.7	0.6\\
0.6	0.8	0.6\\
0.6	0.9	0.6\\
0.6	1	0.6\\
0.7	0	0\\
0.7	0.1	0.1\\
0.7	0.2	0.2\\
0.7	0.3	0.3\\
0.7	0.4	0.4\\
0.7	0.5	0.5\\
0.7	0.6	0.6\\
0.7	0.7	0.7\\
0.7	0.8	0.7\\
0.7	0.9	0.7\\
0.7	1	0.7\\
0.8	0	0\\
0.8	0.1	0.1\\
0.8	0.2	0.2\\
0.8	0.3	0.3\\
0.8	0.4	0.4\\
0.8	0.5	0.5\\
0.8	0.6	0.6\\
0.8	0.7	0.7\\
0.8	0.8	0.8\\
0.8	0.9	0.8\\
0.8	1	0.8\\
0.9	0	0\\
0.9	0.1	0.1\\
0.9	0.2	0.2\\
0.9	0.3	0.3\\
0.9	0.4	0.4\\
0.9	0.5	0.5\\
0.9	0.6	0.6\\
0.9	0.7	0.7\\
0.9	0.8	0.8\\
0.9	0.9	0.9\\
0.9	1	0.9\\
1	0	0\\
1	0.1	0.1\\
1	0.2	0.2\\
1	0.3	0.3\\
1	0.4	0.4\\
1	0.5	0.5\\
1	0.6	0.6\\
1	0.7	0.7\\
1	0.8	0.8\\
1	0.9	0.9\\
1	1	1\\
};
\end{axis}
\end{tikzpicture}%
		\caption{Mínimo}
		\label{fig:t-norms-min}
	\end{subfigure}
	\qquad
	\begin{subfigure}[b]{0.45\textwidth}
		\setlength\figureheight{4.5cm}
		\setlength\figurewidth{6cm}
		% This file was created by matlab2tikz v0.4.7 running on MATLAB 8.0.
% Copyright (c) 2008--2014, Nico Schlömer <nico.schloemer@gmail.com>
% All rights reserved.
% Minimal pgfplots version: 1.3
% 
% The latest updates can be retrieved from
%   http://www.mathworks.com/matlabcentral/fileexchange/22022-matlab2tikz
% where you can also make suggestions and rate matlab2tikz.
% 
\begin{tikzpicture}

\begin{axis}[%
width=\figurewidth,
height=\figureheight,
view={-37.5}{30},
scale only axis,
xmin=0,
xmax=1,
xlabel={x},
xmajorgrids,
ymin=0,
ymax=1,
ylabel={y},
ymajorgrids,
zmin=0,
zmax=1,
zlabel={$T_{P}(x,y)$},
zmajorgrids,
axis x line*=bottom,
axis y line*=left,
axis z line*=left
]

\addplot3[%
surf,
shader=faceted,
draw=black,
colormap/jet,
mesh/rows=11]
table[row sep=crcr,header=false] {0	0	0\\
0	0.1	0\\
0	0.2	0\\
0	0.3	0\\
0	0.4	0\\
0	0.5	0\\
0	0.6	0\\
0	0.7	0\\
0	0.8	0\\
0	0.9	0\\
0	1	0\\
0.1	0	0\\
0.1	0.1	0.01\\
0.1	0.2	0.02\\
0.1	0.3	0.03\\
0.1	0.4	0.04\\
0.1	0.5	0.05\\
0.1	0.6	0.06\\
0.1	0.7	0.07\\
0.1	0.8	0.08\\
0.1	0.9	0.09\\
0.1	1	0.1\\
0.2	0	0\\
0.2	0.1	0.02\\
0.2	0.2	0.04\\
0.2	0.3	0.06\\
0.2	0.4	0.08\\
0.2	0.5	0.1\\
0.2	0.6	0.12\\
0.2	0.7	0.14\\
0.2	0.8	0.16\\
0.2	0.9	0.18\\
0.2	1	0.2\\
0.3	0	0\\
0.3	0.1	0.03\\
0.3	0.2	0.06\\
0.3	0.3	0.09\\
0.3	0.4	0.12\\
0.3	0.5	0.15\\
0.3	0.6	0.18\\
0.3	0.7	0.21\\
0.3	0.8	0.24\\
0.3	0.9	0.27\\
0.3	1	0.3\\
0.4	0	0\\
0.4	0.1	0.04\\
0.4	0.2	0.08\\
0.4	0.3	0.12\\
0.4	0.4	0.16\\
0.4	0.5	0.2\\
0.4	0.6	0.24\\
0.4	0.7	0.28\\
0.4	0.8	0.32\\
0.4	0.9	0.36\\
0.4	1	0.4\\
0.5	0	0\\
0.5	0.1	0.05\\
0.5	0.2	0.1\\
0.5	0.3	0.15\\
0.5	0.4	0.2\\
0.5	0.5	0.25\\
0.5	0.6	0.3\\
0.5	0.7	0.35\\
0.5	0.8	0.4\\
0.5	0.9	0.45\\
0.5	1	0.5\\
0.6	0	0\\
0.6	0.1	0.06\\
0.6	0.2	0.12\\
0.6	0.3	0.18\\
0.6	0.4	0.24\\
0.6	0.5	0.3\\
0.6	0.6	0.36\\
0.6	0.7	0.42\\
0.6	0.8	0.48\\
0.6	0.9	0.54\\
0.6	1	0.6\\
0.7	0	0\\
0.7	0.1	0.07\\
0.7	0.2	0.14\\
0.7	0.3	0.21\\
0.7	0.4	0.28\\
0.7	0.5	0.35\\
0.7	0.6	0.42\\
0.7	0.7	0.49\\
0.7	0.8	0.56\\
0.7	0.9	0.63\\
0.7	1	0.7\\
0.8	0	0\\
0.8	0.1	0.08\\
0.8	0.2	0.16\\
0.8	0.3	0.24\\
0.8	0.4	0.32\\
0.8	0.5	0.4\\
0.8	0.6	0.48\\
0.8	0.7	0.56\\
0.8	0.8	0.64\\
0.8	0.9	0.72\\
0.8	1	0.8\\
0.9	0	0\\
0.9	0.1	0.09\\
0.9	0.2	0.18\\
0.9	0.3	0.27\\
0.9	0.4	0.36\\
0.9	0.5	0.45\\
0.9	0.6	0.54\\
0.9	0.7	0.63\\
0.9	0.8	0.72\\
0.9	0.9	0.81\\
0.9	1	0.9\\
1	0	0\\
1	0.1	0.1\\
1	0.2	0.2\\
1	0.3	0.3\\
1	0.4	0.4\\
1	0.5	0.5\\
1	0.6	0.6\\
1	0.7	0.7\\
1	0.8	0.8\\
1	0.9	0.9\\
1	1	1\\
};
\end{axis}
\end{tikzpicture}%
		\caption{Producto}
		\label{fig:t-norms-prod}
	\end{subfigure}
	
	\vspace{1 cm}
	\begin{subfigure}[b]{0.45\textwidth}
		\setlength\figureheight{4.5cm}
		\setlength\figurewidth{6cm}
		% This file was created by matlab2tikz v0.4.7 (commit 4d7ae5c4fd0932fb51051d86111bc23ed23e4580) running on MATLAB 8.0.
% Copyright (c) 2008--2014, Nico Schlömer <nico.schloemer@gmail.com>
% All rights reserved.
% Minimal pgfplots version: 1.3
% 
\begin{tikzpicture}

\begin{axis}[%
width=\figurewidth,
height=\figureheight,
view={-37.5}{30},
scale only axis,
xmin=0,
xmax=1,
xlabel={x},
xmajorgrids,
ymin=0,
ymax=1,
ylabel={y},
ymajorgrids,
zmin=0,
zmax=1,
zlabel={$\text{T}_{\text{L}}\text{(x,y)}$},
zmajorgrids,
axis x line*=bottom,
axis y line*=left,
axis z line*=left
]

\addplot3[%
surf,
shader=faceted,
draw=black,
colormap/jet,
mesh/rows=11]
table[row sep=crcr,header=false] {0	0	0\\
0	0.1	0\\
0	0.2	0\\
0	0.3	0\\
0	0.4	0\\
0	0.5	0\\
0	0.6	0\\
0	0.7	0\\
0	0.8	0\\
0	0.9	0\\
0	1	0\\
0.1	0	0\\
0.1	0.1	0\\
0.1	0.2	0\\
0.1	0.3	0\\
0.1	0.4	0\\
0.1	0.5	0\\
0.1	0.6	0\\
0.1	0.7	0\\
0.1	0.8	0\\
0.1	0.9	0\\
0.1	1	0.1\\
0.2	0	0\\
0.2	0.1	0\\
0.2	0.2	0\\
0.2	0.3	0\\
0.2	0.4	0\\
0.2	0.5	0\\
0.2	0.6	0\\
0.2	0.7	0\\
0.2	0.8	0\\
0.2	0.9	0.1\\
0.2	1	0.2\\
0.3	0	0\\
0.3	0.1	0\\
0.3	0.2	0\\
0.3	0.3	0\\
0.3	0.4	0\\
0.3	0.5	0\\
0.3	0.6	0\\
0.3	0.7	0\\
0.3	0.8	0.1\\
0.3	0.9	0.2\\
0.3	1	0.3\\
0.4	0	0\\
0.4	0.1	0\\
0.4	0.2	0\\
0.4	0.3	0\\
0.4	0.4	0\\
0.4	0.5	0\\
0.4	0.6	0\\
0.4	0.7	0.1\\
0.4	0.8	0.2\\
0.4	0.9	0.3\\
0.4	1	0.4\\
0.5	0	0\\
0.5	0.1	0\\
0.5	0.2	0\\
0.5	0.3	0\\
0.5	0.4	0\\
0.5	0.5	0\\
0.5	0.6	0.1\\
0.5	0.7	0.2\\
0.5	0.8	0.3\\
0.5	0.9	0.4\\
0.5	1	0.5\\
0.6	0	0\\
0.6	0.1	0\\
0.6	0.2	0\\
0.6	0.3	0\\
0.6	0.4	0\\
0.6	0.5	0.1\\
0.6	0.6	0.2\\
0.6	0.7	0.3\\
0.6	0.8	0.4\\
0.6	0.9	0.5\\
0.6	1	0.6\\
0.7	0	0\\
0.7	0.1	0\\
0.7	0.2	0\\
0.7	0.3	0\\
0.7	0.4	0.1\\
0.7	0.5	0.2\\
0.7	0.6	0.3\\
0.7	0.7	0.4\\
0.7	0.8	0.5\\
0.7	0.9	0.6\\
0.7	1	0.7\\
0.8	0	0\\
0.8	0.1	0\\
0.8	0.2	0\\
0.8	0.3	0.1\\
0.8	0.4	0.2\\
0.8	0.5	0.3\\
0.8	0.6	0.4\\
0.8	0.7	0.5\\
0.8	0.8	0.6\\
0.8	0.9	0.7\\
0.8	1	0.8\\
0.9	0	0\\
0.9	0.1	0\\
0.9	0.2	0.1\\
0.9	0.3	0.2\\
0.9	0.4	0.3\\
0.9	0.5	0.4\\
0.9	0.6	0.5\\
0.9	0.7	0.6\\
0.9	0.8	0.7\\
0.9	0.9	0.8\\
0.9	1	0.9\\
1	0	0\\
1	0.1	0.1\\
1	0.2	0.2\\
1	0.3	0.3\\
1	0.4	0.4\\
1	0.5	0.5\\
1	0.6	0.6\\
1	0.7	0.7\\
1	0.8	0.8\\
1	0.9	0.9\\
1	1	1\\
};
\end{axis}
\end{tikzpicture}%
		\caption{\L{}ukasiewicz}
		\label{fig:t-norms-lukasiewicz}
	\end{subfigure}
	\label{fig:t-norms}
	\caption{Representación gráfica de las t-normas \emph{Mínimo} (\ref{fig:t-norms-min}), \emph{Producto} (\ref{fig:t-norms-prod}) y \emph{\L{}ukasiewicz} (\ref{fig:t-norms-lukasiewicz}).}
\end{figure}

\begin{definition}
Una t-conorma es una operación binaria \emph{S} en el intervalo $[0,1]$ que es conmutativa, asociativa, monotona y tiene el valor \emph{0} como elemento neutro. Es decir, una función $\emph{S} : [0,1]^2 \rightarrow [0,1]$ tal que $\forall x,y,z \in [0,1]$:
\begin{enumerate}[label=(S\arabic*),ref=(S\arabic*)]
   \item $S(x,y) = S(y,x)$ (Conmutatividad)\label{S1}
   \item $S(x,S(y,z)) = S(S(x,y),z)$ (Asociatividad)\label{S2}
   \item $S(x,y) \leq S(x,z)$ cuando $y \leq z$ (Monotonía)\label{S3}
   \item $S(x,0) = x$ (Elemento neutro)\label{S4}
  \end{enumerate}
\end{definition}
De la misma forma que las t-normas (definición \ref{def:tnorma}) generalizan la conjunción clásica, las t-conormas generalizan la disyunción ($x\bigvee y$). A continuación se incluyen algunos ejemplos de t-conormas:
\begin{itemize}
	\item \bfseries Máximo: $S_{max}(x,y) = max\{x,y\}$
	\item \bfseries Suma probabilística: $S_{P}(x,y) = x + y - xy$
	\item \bfseries Suma acotada: $S_{B}(x,y) = min\{x+y,1\}$
\end{itemize}
\begin{figure}[H]
	\centering
	\begin{subfigure}[b]{0.45\textwidth}
		\setlength\figureheight{4.5cm}
		\setlength\figurewidth{6cm}
		% This file was created by matlab2tikz v0.4.7 (commit 4d7ae5c4fd0932fb51051d86111bc23ed23e4580) running on MATLAB 8.0.
% Copyright (c) 2008--2014, Nico Schlömer <nico.schloemer@gmail.com>
% All rights reserved.
% Minimal pgfplots version: 1.3
% 
\begin{tikzpicture}

\begin{axis}[%
width=\figurewidth,
height=\figureheight,
view={-37.5}{30},
scale only axis,
xmin=0,
xmax=1,
xlabel={x},
xmajorgrids,
ymin=0,
ymax=1,
ylabel={y},
ymajorgrids,
zmin=0,
zmax=1,
zlabel={$\text{S}_{\text{max}}\text{(x,y)}$},
zmajorgrids,
axis x line*=bottom,
axis y line*=left,
axis z line*=left
]

\addplot3[%
surf,
shader=faceted,
draw=black,
colormap/jet,
mesh/rows=11]
table[row sep=crcr,header=false] {0	0	0\\
0	0.1	0.1\\
0	0.2	0.2\\
0	0.3	0.3\\
0	0.4	0.4\\
0	0.5	0.5\\
0	0.6	0.6\\
0	0.7	0.7\\
0	0.8	0.8\\
0	0.9	0.9\\
0	1	1\\
0.1	0	0.1\\
0.1	0.1	0.1\\
0.1	0.2	0.2\\
0.1	0.3	0.3\\
0.1	0.4	0.4\\
0.1	0.5	0.5\\
0.1	0.6	0.6\\
0.1	0.7	0.7\\
0.1	0.8	0.8\\
0.1	0.9	0.9\\
0.1	1	1\\
0.2	0	0.2\\
0.2	0.1	0.2\\
0.2	0.2	0.2\\
0.2	0.3	0.3\\
0.2	0.4	0.4\\
0.2	0.5	0.5\\
0.2	0.6	0.6\\
0.2	0.7	0.7\\
0.2	0.8	0.8\\
0.2	0.9	0.9\\
0.2	1	1\\
0.3	0	0.3\\
0.3	0.1	0.3\\
0.3	0.2	0.3\\
0.3	0.3	0.3\\
0.3	0.4	0.4\\
0.3	0.5	0.5\\
0.3	0.6	0.6\\
0.3	0.7	0.7\\
0.3	0.8	0.8\\
0.3	0.9	0.9\\
0.3	1	1\\
0.4	0	0.4\\
0.4	0.1	0.4\\
0.4	0.2	0.4\\
0.4	0.3	0.4\\
0.4	0.4	0.4\\
0.4	0.5	0.5\\
0.4	0.6	0.6\\
0.4	0.7	0.7\\
0.4	0.8	0.8\\
0.4	0.9	0.9\\
0.4	1	1\\
0.5	0	0.5\\
0.5	0.1	0.5\\
0.5	0.2	0.5\\
0.5	0.3	0.5\\
0.5	0.4	0.5\\
0.5	0.5	0.5\\
0.5	0.6	0.6\\
0.5	0.7	0.7\\
0.5	0.8	0.8\\
0.5	0.9	0.9\\
0.5	1	1\\
0.6	0	0.6\\
0.6	0.1	0.6\\
0.6	0.2	0.6\\
0.6	0.3	0.6\\
0.6	0.4	0.6\\
0.6	0.5	0.6\\
0.6	0.6	0.6\\
0.6	0.7	0.7\\
0.6	0.8	0.8\\
0.6	0.9	0.9\\
0.6	1	1\\
0.7	0	0.7\\
0.7	0.1	0.7\\
0.7	0.2	0.7\\
0.7	0.3	0.7\\
0.7	0.4	0.7\\
0.7	0.5	0.7\\
0.7	0.6	0.7\\
0.7	0.7	0.7\\
0.7	0.8	0.8\\
0.7	0.9	0.9\\
0.7	1	1\\
0.8	0	0.8\\
0.8	0.1	0.8\\
0.8	0.2	0.8\\
0.8	0.3	0.8\\
0.8	0.4	0.8\\
0.8	0.5	0.8\\
0.8	0.6	0.8\\
0.8	0.7	0.8\\
0.8	0.8	0.8\\
0.8	0.9	0.9\\
0.8	1	1\\
0.9	0	0.9\\
0.9	0.1	0.9\\
0.9	0.2	0.9\\
0.9	0.3	0.9\\
0.9	0.4	0.9\\
0.9	0.5	0.9\\
0.9	0.6	0.9\\
0.9	0.7	0.9\\
0.9	0.8	0.9\\
0.9	0.9	0.9\\
0.9	1	1\\
1	0	1\\
1	0.1	1\\
1	0.2	1\\
1	0.3	1\\
1	0.4	1\\
1	0.5	1\\
1	0.6	1\\
1	0.7	1\\
1	0.8	1\\
1	0.9	1\\
1	1	1\\
};
\end{axis}
\end{tikzpicture}%
		\caption{Máximo}
		\label{fig:t-conorms-max}
	\end{subfigure}
	\qquad
	\begin{subfigure}[b]{0.45\textwidth}
		\setlength\figureheight{4.5cm}
		\setlength\figurewidth{6cm}
		% This file was created by matlab2tikz v0.4.7 (commit 4d7ae5c4fd0932fb51051d86111bc23ed23e4580) running on MATLAB 8.0.
% Copyright (c) 2008--2014, Nico Schlömer <nico.schloemer@gmail.com>
% All rights reserved.
% Minimal pgfplots version: 1.3
% 
\begin{tikzpicture}

\begin{axis}[%
width=\figurewidth,
height=\figureheight,
view={-37.5}{30},
scale only axis,
xmin=0,
xmax=1,
xlabel={x},
xmajorgrids,
ymin=0,
ymax=1,
ylabel={y},
ymajorgrids,
zmin=0,
zmax=1,
zlabel={$\text{S}_{\text{P}}\text{(x,y)}$},
zmajorgrids,
axis x line*=bottom,
axis y line*=left,
axis z line*=left
]

\addplot3[%
surf,
shader=faceted,
draw=black,
colormap/jet,
mesh/rows=11]
table[row sep=crcr,header=false] {0	0	0\\
0	0.1	0.1\\
0	0.2	0.2\\
0	0.3	0.3\\
0	0.4	0.4\\
0	0.5	0.5\\
0	0.6	0.6\\
0	0.7	0.7\\
0	0.8	0.8\\
0	0.9	0.9\\
0	1	1\\
0.1	0	0.1\\
0.1	0.1	0.19\\
0.1	0.2	0.28\\
0.1	0.3	0.37\\
0.1	0.4	0.46\\
0.1	0.5	0.55\\
0.1	0.6	0.64\\
0.1	0.7	0.73\\
0.1	0.8	0.82\\
0.1	0.9	0.91\\
0.1	1	1\\
0.2	0	0.2\\
0.2	0.1	0.28\\
0.2	0.2	0.36\\
0.2	0.3	0.44\\
0.2	0.4	0.52\\
0.2	0.5	0.6\\
0.2	0.6	0.68\\
0.2	0.7	0.76\\
0.2	0.8	0.84\\
0.2	0.9	0.92\\
0.2	1	1\\
0.3	0	0.3\\
0.3	0.1	0.37\\
0.3	0.2	0.44\\
0.3	0.3	0.51\\
0.3	0.4	0.58\\
0.3	0.5	0.65\\
0.3	0.6	0.72\\
0.3	0.7	0.79\\
0.3	0.8	0.86\\
0.3	0.9	0.93\\
0.3	1	1\\
0.4	0	0.4\\
0.4	0.1	0.46\\
0.4	0.2	0.52\\
0.4	0.3	0.58\\
0.4	0.4	0.64\\
0.4	0.5	0.7\\
0.4	0.6	0.76\\
0.4	0.7	0.82\\
0.4	0.8	0.88\\
0.4	0.9	0.94\\
0.4	1	1\\
0.5	0	0.5\\
0.5	0.1	0.55\\
0.5	0.2	0.6\\
0.5	0.3	0.65\\
0.5	0.4	0.7\\
0.5	0.5	0.75\\
0.5	0.6	0.8\\
0.5	0.7	0.85\\
0.5	0.8	0.9\\
0.5	0.9	0.95\\
0.5	1	1\\
0.6	0	0.6\\
0.6	0.1	0.64\\
0.6	0.2	0.68\\
0.6	0.3	0.72\\
0.6	0.4	0.76\\
0.6	0.5	0.8\\
0.6	0.6	0.84\\
0.6	0.7	0.88\\
0.6	0.8	0.92\\
0.6	0.9	0.96\\
0.6	1	1\\
0.7	0	0.7\\
0.7	0.1	0.73\\
0.7	0.2	0.76\\
0.7	0.3	0.79\\
0.7	0.4	0.82\\
0.7	0.5	0.85\\
0.7	0.6	0.88\\
0.7	0.7	0.91\\
0.7	0.8	0.94\\
0.7	0.9	0.97\\
0.7	1	1\\
0.8	0	0.8\\
0.8	0.1	0.82\\
0.8	0.2	0.84\\
0.8	0.3	0.86\\
0.8	0.4	0.88\\
0.8	0.5	0.9\\
0.8	0.6	0.92\\
0.8	0.7	0.94\\
0.8	0.8	0.96\\
0.8	0.9	0.98\\
0.8	1	1\\
0.9	0	0.9\\
0.9	0.1	0.91\\
0.9	0.2	0.92\\
0.9	0.3	0.93\\
0.9	0.4	0.94\\
0.9	0.5	0.95\\
0.9	0.6	0.96\\
0.9	0.7	0.97\\
0.9	0.8	0.98\\
0.9	0.9	0.99\\
0.9	1	1\\
1	0	1\\
1	0.1	1\\
1	0.2	1\\
1	0.3	1\\
1	0.4	1\\
1	0.5	1\\
1	0.6	1\\
1	0.7	1\\
1	0.8	1\\
1	0.9	1\\
1	1	1\\
};
\end{axis}
\end{tikzpicture}%
		\caption{Suma probabilística}
		\label{fig:t-conorms-probabilistic-sum}
	\end{subfigure}
	
	\vspace{1 cm}
	\begin{subfigure}[b]{0.45\textwidth}
		\setlength\figureheight{4.5cm}
		\setlength\figurewidth{6cm}
		% This file was created by matlab2tikz v0.4.7 (commit 4d7ae5c4fd0932fb51051d86111bc23ed23e4580) running on MATLAB 8.0.
% Copyright (c) 2008--2014, Nico Schlömer <nico.schloemer@gmail.com>
% All rights reserved.
% Minimal pgfplots version: 1.3
% 
\begin{tikzpicture}

\begin{axis}[%
width=\figurewidth,
height=\figureheight,
view={-37.5}{30},
scale only axis,
xmin=0,
xmax=1,
xlabel={x},
xmajorgrids,
ymin=0,
ymax=1,
ylabel={y},
ymajorgrids,
zmin=0,
zmax=1,
zlabel={$\text{S}_{\text{B}}\text{(x,y)}$},
zmajorgrids,
axis x line*=bottom,
axis y line*=left,
axis z line*=left
]

\addplot3[%
surf,
shader=faceted,
draw=black,
colormap/jet,
mesh/rows=11]
table[row sep=crcr,header=false] {0	0	0\\
0	0.1	0.1\\
0	0.2	0.2\\
0	0.3	0.3\\
0	0.4	0.4\\
0	0.5	0.5\\
0	0.6	0.6\\
0	0.7	0.7\\
0	0.8	0.8\\
0	0.9	0.9\\
0	1	1\\
0.1	0	0.1\\
0.1	0.1	0.2\\
0.1	0.2	0.3\\
0.1	0.3	0.4\\
0.1	0.4	0.5\\
0.1	0.5	0.6\\
0.1	0.6	0.7\\
0.1	0.7	0.8\\
0.1	0.8	0.9\\
0.1	0.9	1\\
0.1	1	1\\
0.2	0	0.2\\
0.2	0.1	0.3\\
0.2	0.2	0.4\\
0.2	0.3	0.5\\
0.2	0.4	0.6\\
0.2	0.5	0.7\\
0.2	0.6	0.8\\
0.2	0.7	0.9\\
0.2	0.8	1\\
0.2	0.9	1\\
0.2	1	1\\
0.3	0	0.3\\
0.3	0.1	0.4\\
0.3	0.2	0.5\\
0.3	0.3	0.6\\
0.3	0.4	0.7\\
0.3	0.5	0.8\\
0.3	0.6	0.9\\
0.3	0.7	1\\
0.3	0.8	1\\
0.3	0.9	1\\
0.3	1	1\\
0.4	0	0.4\\
0.4	0.1	0.5\\
0.4	0.2	0.6\\
0.4	0.3	0.7\\
0.4	0.4	0.8\\
0.4	0.5	0.9\\
0.4	0.6	1\\
0.4	0.7	1\\
0.4	0.8	1\\
0.4	0.9	1\\
0.4	1	1\\
0.5	0	0.5\\
0.5	0.1	0.6\\
0.5	0.2	0.7\\
0.5	0.3	0.8\\
0.5	0.4	0.9\\
0.5	0.5	1\\
0.5	0.6	1\\
0.5	0.7	1\\
0.5	0.8	1\\
0.5	0.9	1\\
0.5	1	1\\
0.6	0	0.6\\
0.6	0.1	0.7\\
0.6	0.2	0.8\\
0.6	0.3	0.9\\
0.6	0.4	1\\
0.6	0.5	1\\
0.6	0.6	1\\
0.6	0.7	1\\
0.6	0.8	1\\
0.6	0.9	1\\
0.6	1	1\\
0.7	0	0.7\\
0.7	0.1	0.8\\
0.7	0.2	0.9\\
0.7	0.3	1\\
0.7	0.4	1\\
0.7	0.5	1\\
0.7	0.6	1\\
0.7	0.7	1\\
0.7	0.8	1\\
0.7	0.9	1\\
0.7	1	1\\
0.8	0	0.8\\
0.8	0.1	0.9\\
0.8	0.2	1\\
0.8	0.3	1\\
0.8	0.4	1\\
0.8	0.5	1\\
0.8	0.6	1\\
0.8	0.7	1\\
0.8	0.8	1\\
0.8	0.9	1\\
0.8	1	1\\
0.9	0	0.9\\
0.9	0.1	1\\
0.9	0.2	1\\
0.9	0.3	1\\
0.9	0.4	1\\
0.9	0.5	1\\
0.9	0.6	1\\
0.9	0.7	1\\
0.9	0.8	1\\
0.9	0.9	1\\
0.9	1	1\\
1	0	1\\
1	0.1	1\\
1	0.2	1\\
1	0.3	1\\
1	0.4	1\\
1	0.5	1\\
1	0.6	1\\
1	0.7	1\\
1	0.8	1\\
1	0.9	1\\
1	1	1\\
};
\end{axis}
\end{tikzpicture}%
		\caption{Suma acotada}
		\label{fig:t-conorms-bounded-sum}
	\end{subfigure}
	\label{fig:t-conorms}
	\caption{Representación gráfica de las t-conormas \emph{Máximo} (\ref{fig:t-conorms-max}), \emph{Suma probabilística} (\ref{fig:t-conorms-probabilistic-sum}) y \emph{Suma acotada} (\ref{fig:t-conorms-bounded-sum}).}
\end{figure}
Otra clase importante de funciones son los \emph{operadores de agregación} \index{operador de agregación}\cite{calvo2002aggregation, beliakov2007aggregation}.
\begin{definition}
Una función $\emph{M} : [a,b]^{n} \rightarrow [a,b]$ es un operador de agregación si es monótona y no decreciente en cada una de sus componentes y además cumple que $M(a, a, \cdots,a) = a$ y $M(b, b, \cdots,b) = b$.
\end{definition}
\begin{definition}
Una función de agregación \emph{M} se dice que es una media si cumple que:
\begin{equation}\label{def:mean-aggregation-operator}
min(x) = min(x_{1},\cdots,x_{n})\leq \emph{M}(x_{1},\cdots,x_{n}) \leq max(x_{1},\cdots,x_{n}) = max(x)
\end{equation}
\end{definition}
A continuación se incluyen algunos de los operadores de agregación más comunes:
\begin{itemize}
	\item \bfseries Media aritmética: $M(x_{1},x_{2},\cdots,x_{n}) = \frac{1}{n}\sum\limits_{i=1}^{n}x_{i}$
	\item \bfseries Media geométrica: $M(x_{1},x_{2},\cdots,x_{n}) = (\prod\limits_{i=1}^{n}x_{i})^{\frac{1}{n}}$
	\item \bfseries Media ponderada: $M_{w_{1},w_{2},\cdots,w_{n}}(x_{1},x_{2},\cdots,x_{n}) =\sum\limits_{i=1}^{n}(w_{i}\cdot x_{i})$ \normalfont tal que $0 \leq w_{i} \leq 1$ y $\sum\limits_{i=1}^{n}w_{i} = 1$.
	\item \bfseries Mediana: \normalfont Se toma el elemento central del conjunto ordenado de argumentos.
	\item \bfseries Máximo: $M(x_{1},x_{2},\cdots,x_{n}) = max(x_{1},x_{2},\cdots,x_{n})$
	\item \bfseries Mínimo: $M(x_{1},x_{2},\cdots,x_{n}) = min(x_{1},x_{2},\cdots,x_{n})$
\end{itemize}

\section{Funciones de solapamiento}\label{sec:overlap-functions}
En esta sección se presenta el concepto de \emph{función de solapamiento}\index{función de solapamiento} \cite{bustince2010overlapfunctions,bustince2012pairwisecomparisons,jurio2013propertiesoverlap}, que tendrá una gran importancia en el desarrollo de este trabajo, ya que a partir de estas funciones se construyen los índices de solapamiento, que serán utilizados en el método de interpolación presentado.
\begin{definition}
Una función de solapamiento $G_{O} : [0,1]^{2} \rightarrow [0,1]$ cumple:
\begin{enumerate}[label=(G\arabic*),ref=(G\arabic*)]
   \item $G_{O}(x,y) = G_{O}(y,x) \;\; \forall \; x,y \in [0,1]$ \label{G1}
   \item $G_{O}(x,y) = 0$ si y sólo si $x \cdot y = 0$ \label{G2}
   \item $G_{O}(x,y) = 1$ si y sólo si $x \cdot y = 1$ \label{G3}
   \item $G_{O}$ es creciente \label{G4}
   \item $G_{O}$ es continua \label{G5}
\end{enumerate}
\end{definition}
Las funciones de solapamiento generalizan los operadores de intersección tales como el mínimo o, en general, las t-normas (definidas en \ref{def:tnorma}). Además, las funciones de solapamiento son casos particulares de los operadores de agregación sin divisores por cero o divisores por uno. Algunos ejemplos de funciones de solapamiento son:\\
\\
$G_{O}(x,y) = min(x,y)$\\
$G_{O}(x,y) = \sqrt{xy}$\\
$G_{O}(x,y) = min(x^{k}y,xy^{k}), k \in ]0,1[$\\

\section{Índices de solapamiento}
\subsection{Definición y propiedades}
En esta sección se presenta la definición de índice de solapamiento y se estudian algunas de sus propiedades más importantes. Los índices de solapamiento tienen una importancia vital en el desarrollo de este trabajo, ya que son la base del método estudiado. 

El concepto de solapamiento aplicado a conjuntos difusos fue introducido por Zadeh en 1978 \cite{zadeh1978}.

\begin{definition}
Sea $A,B \in FS(U)$. La consistencia entre A y B se define como:
\begin{equation}\label{eq:zadeh-consistency}
O_{Z}(A,B) = sup_{i=1}^{n}(min(A(u_{i}),B(u_{i})))
\end{equation}
\end{definition}

Para conjuntos de referencia finitos, el supremo es en realidad el máximo elemento del conjunto.

En 1982 Dubois y Prade \cite{dubois2000} presentan la siguiente definición para el índice de solapamiento:

\begin{definition}
Un índice de solapamiento es una función $O : FS(U) \times FS(U) \rightarrow [0,1]$ tal que:
\begin{enumerate}[label=(O\arabic*),ref=(O\arabic*)]
\item $O(A,B) = 0$ si y sólo si A y B son completamente disjuntos. \label{DP01}
\item $O(A,B) = 1$, si ($A(u_{i}) = 0$ o $B(u_{i}) = 0$) o ($A(u_{i}) = 1$ o $B(u_{i}) = 1$) \label{DPO2}
\item $O(A,B) = O(B,A)$ \label{DPO3}
\item Si $B \leq C$, entonces $O(A,B) \leq O(A,C)$ \label{DPO4}
\end{enumerate}
\end{definition} \label{def:dubois-overlap-index}
La condición \ref{DPO2} en esta definición presenta la ventaja de que, si \emph{A} no es difuso, entonces $O(A,A) = 1$.

Por esta razón Dubois, Ostasiewicz y Prade imponen en \cite{dubois2000} las siguientes condiciones:
\begin{enumerate}[label=(\arabic*),ref=(\arabic*)]
\item Para todos los conjuntos difusos tales que $A(u_{i}) \le 1$ y $A(u_{i}) \le 1$ para cualquier $u_{i} \in U$, \ref{DPO2} debe ignorarse.
\item El índice ROC (\cite{dubois2000}) no satisface \ref{DPO2}.
\end{enumerate}
Debido a estas consideraciones, normalmente sólo se imponen las condiciones \ref{DP01}, \ref{DPO3}, \ref{DPO4} de la definición \ref{def:dubois-overlap-index} a los índices de solapamiento.
En \cite{bustince2013overlap} se propone la siguiente definición de índice de solapamiento:
\begin{definition}\label{def:bustince-overlap-index}
Un índice de solapamiento es una función $O : FS(U) \times FS(U) \rightarrow [0,1]$ tal que:
\begin{enumerate}[label=(O\arabic*),ref=(O\arabic*)]
\item $O(A,B) = 0$ si y sólo si A y B tienen soportes disjuntos, es decir, $A(u_{i})B(u_{i}) = 0$ para todo $u_{i} \in U$\label{BO1}
\item $O(A,B) = O(B,A)$\label{BO2}
\item Si $B \leq C$, entonces $O(A,B) \leq O(A,C)$\label{BO3}
\end{enumerate}
Un índice de solapamiento normal es un índice \emph{O} tal que:
\begin{enumerate}[label=(O4),ref=(O4)]
\item Si existe un $u_{i} \in U$ tal que $A(u_{i}) = B(u_{i}) = 1$, entonces $O(A,B) = 1$\label{BO4}
\end{enumerate}
\end{definition}
\subsection{Construcción de índices de solapamiento}
En esta sección se presenta el método de construcción de índices de solapamiento propuesto por \cite{bustince2013overlap} que se basa en operadores de agregación (sección \ref{sec:t-norms-aggregation-operator}) y funciones de solapamiento (sección \ref{sec:overlap-functions}).
\begin{theorem}
Sea $M : [0,1]^{2} \rightarrow [0,1]$ una función de agregación tal que $M(x_{1},\cdots,x_{n}) = 0$ si y sólo si $x_{1} = \cdots = x_{n} = 0$. Sea $G_{O} : [0,1]^{2} \rightarrow [0,1]$ una función de solapamiento. Entonces, la función $O: F(U) \times F(U) \rightarrow [0,1]$ definida como:
\begin{equation}\label{eq:construction-overlap-index}
O(A,B) = M(G_{O}(A(u_{1}),B(u_{1})),\cdots,G_{O}(A(u_{n}),B(u_{n})))
\end{equation}
es un índice de solapamiento en el sentido de la definición \ref{def:bustince-overlap-index}. Recíprocamente, si $G_{O}$ es una función de solapamiento y $M:[0,1]^{n} \rightarrow [0,1]$ es un operador de agregación tal que O definido por la ecuación \ref{eq:construction-overlap-index} es un índice de solapamiento, entonces $M(x_{1},\cdots,x_{n}) = 0$ si y sólo si $x_{1} = \cdots = x_{n} = 0$.
\end{theorem}


\mainmatter

% Capítulos
% !TeX spellcheck = es_ES
\chapter{Teoría de conjuntos difusos}
\label{cha:teoria-conjuntos-difusos}

En este primer capítulo se introduce la teoría relacionada con conjuntos difusos necesaria para el desarrollo del proyecto. Los conjuntos difusos forman la base de los métodos estudiados en este proyecto y han tenido una importancia capital en el desarrollo de la inteligencia artificial en las últimas decadas, desde que fueran introducidos por L.A. Zadeh en 1965 \cite{Zadeh65}. Por esta razón, es necesario definirlos de forma apropiada y estudiar algunas de sus propiedades fundamentales.

\section{Conjuntos difusos}

En esta sección se define el concepto de conjunto difuso (sec. \ref{sec:fuzzy-set-definition-and-properties}), así como las operaciones más básicas que se pueden realizar sobre ellos (sec. \ref{sec:fuzzy-set-operations}). En la sección \ref{sec:lang-variables} se presenta el concepto de variable lingüística, una de las aplicaciones más importantes de los conjuntos difusos y clave para el desarrollo de la lógica difusa.

\subsection{Definición y propiedades}
\label{sec:fuzzy-set-definition-and-properties}
\begin{definition}
Sea \emph{X} un conjunto clásico de objetos, denominado \emph{universo}, cuyos elementos son denotados como \emph{x}. La pertenencia a un subconjunto clásico \emph{A} de \emph{X} se expresa normalmente por medio de una función de pertenencia $\mu_{A}:X\rightarrow\{0,1\}$ tal que:
\begin{equation}
\mu_{A}(x)=\begin{cases} 1 & \mbox{iff } \mbox{x}\in A, \\ 0 & \mbox{iff } \mbox{x}\notin A \end{cases}
\end{equation}
(se utiliza \emph{iff} como abreviatura de "si y sólo si")
\end{definition}

De esta forma, en un conjunto clásico, se dice que un elemento \emph{x} pertenece al conjunto A ($\mu_{A}(x)=1$) o no pertenece ($\mu_{A}(x)=0$). 
\begin{example}
\normalfont
Un ejemplo de conjunto clásico es el conjunto de todos los números enteros pares, cuya función de pertenencia $\mu_{A}:\mathbb{Z}\rightarrow\{0,1\}$ podría expresarse como:
\begin{equation}
\mu_{A}(x)=\begin{cases} 1 & \mbox{iff } \mbox{x}\bmod 2=0, \\ 0 & \mbox{iff } \mbox{x}\bmod 2\ne 0 \end{cases}
\end{equation}
\end{example}
Si el valor de salida de la función de pertenencia puede ser cualquier número en el intervalo [0,1] (y no sólo los valores discretos \{0,1\}) se dice que el conjunto \emph{A} es \emph{difuso} y que la función $\mu_{A}(x)$ es el \emph{grado de pertenencia} de \emph{x} a \emph{A}. En este tipo de conjuntos, cuanto más próximo sea el valor de $\mu_{A}(x)$ a 1, más pertenece \emph{x} a \emph{A}.

El concepto de conjunto difuso fue introducido por L.A. Zadeh en 1965 \cite{Zadeh65} y se define como:
\begin{definition}
Dado un conjunto de referencia (o universo) \emph{U}, un conjunto difuso \emph{A} sobre \emph{U} es un conjunto tal que:\\
\begin{equation}
\{(u_{i},\mu_{A}(u_{i}))\arrowvert u_{i} \in U\}
\end{equation}
donde \begin{math}\emph{A}:\emph{U}\rightarrow[0,1]\end{math} es la \emph{función de pertenencia}\index{función de pertenencia} (o \emph{grado de pertenencia}) de \emph{A}.
\end{definition}

\begin{example}
\normalfont
Supongamos que el universo \emph{U} está compuesto por los valores razonables de altura (medida en centímetros) de una persona adulta, tal que:

\begin{equation}
\emph{U} = \{150,151,\ldots,229,230\}
\end{equation}

Se puede definir sobre \emph{U} el conjunto difuso \emph{VT}, para representar el concepto de altura \emph{"Muy alta"}, especificando su función de pertenencia \emph{$\mu_{VT}$}, de forma que los valores de altura \emph{$\mu_{i}$} considerados como \emph{"muy altos"} hagan que el valor de dicha función de pertenencia \emph{$\mu_{VT}(\mu_{i})$} sea más próximo a 1 (pertenecen más al conjunto).\\
\\
Se puede elegir cualquier tipo de función para definir al conjunto difuso, siempre que se ajuste a las necesidades del problema que se pretende resolver. En este ejemplo se va a definir la función de pertenencia \emph{$\mu_{VT}$} como una función sigmoidal tal que:
\begin{equation}
\mu_{VT}(\mu_{i}) = \frac{1}{1 + e^{-0.4(\mu_{i}-190)}}
\end{equation}
La función de pertenencia $\mu_{VT}(\mu_{i})$ se representa en la figura \ref{fig:fuzzyset-verytall-example}. Se puede observar que alturas que no podrían considerarse muy altas como 150, 160, 170,etc. hacen que la función de pertenencia  $\mu_{VT}$ sea cero (no pertenencen al conjunto). Algunos valores del conjunto difuso \emph{VT}:
\begin{equation}
VT = \{(150,0),(160,0),(170,0),\ldots,(190,0.5),(195,0.88),(200,1),(210,1),\ldots\}
\end{equation}
\end{example}

\begin{figure}[t]
	\centering
	\newlength\figureheight 
	\newlength\figurewidth
	\setlength\figureheight{4.5cm}
	\setlength\figurewidth{12cm}
	% This file was created by matlab2tikz v0.4.7 (commit 06bbbef116b887fd29db8b1bac523c404c7f7694) running on MATLAB 8.0.
% Copyright (c) 2008--2014, Nico Schlömer <nico.schloemer@gmail.com>
% All rights reserved.
% Minimal pgfplots version: 1.3
% 
\begin{tikzpicture}

\begin{axis}[%
width=\figurewidth,
height=\figureheight,
scale only axis,
xmin=150,
xmax=230,
xlabel={$\mu{}_{\text{i}}\text{ - Altura en cm}$},
ymin=-0.1,
ymax=1.1,
ylabel={$\mu{}_{\text{VT}}\text{(}\mu{}_{\text{i}}\text{)}$},
legend style={at={(0.97,0.03)},anchor=south east,draw=black,fill=white,legend cell align=left}
]
\addplot [color=blue,solid,line width=1.5pt]
  table[row sep=crcr]{150	1.12535162055095e-07\\
150.1	1.17127808604502e-07\\
150.2	1.21907884570381e-07\\
150.3	1.26883039088593e-07\\
150.4	1.32061233461842e-07\\
150.5	1.37450753899427e-07\\
150.6	1.43060224776898e-07\\
150.7	1.4889862243685e-07\\
150.8	1.54975289552946e-07\\
150.9	1.61299950080123e-07\\
151	1.67882724814952e-07\\
151.1	1.74734147591e-07\\
151.2	1.81865182135148e-07\\
151.3	1.89287239611824e-07\\
151.4	1.97012196883211e-07\\
151.5	2.05052415514691e-07\\
151.6	2.13420761555885e-07\\
151.7	2.22130626128981e-07\\
151.8	2.31195946857281e-07\\
151.9	2.40631230168244e-07\\
152	2.50451574506755e-07\\
152.1	2.6067269449571e-07\\
152.2	2.71310946082621e-07\\
152.3	2.8238335271246e-07\\
152.4	2.93907632568622e-07\\
152.5	3.05902226925625e-07\\
152.6	3.18386329658865e-07\\
152.7	3.313799179587e-07\\
152.8	3.44903784297974e-07\\
152.9	3.58979569704136e-07\\
153	3.73629798389248e-07\\
153.1	3.888779137932e-07\\
153.2	4.04748316097899e-07\\
153.3	4.21266401272407e-07\\
153.4	4.38458601711502e-07\\
153.5	4.56352428532765e-07\\
153.6	4.74976515599755e-07\\
153.7	4.94360665341816e-07\\
153.8	5.14535896443798e-07\\
153.9	5.35534493481965e-07\\
154	5.57390058585609e-07\\
154.1	5.80137565206936e-07\\
154.2	6.03813414085313e-07\\
154.3	6.28455491495431e-07\\
154.4	6.54103229872536e-07\\
154.5	6.80797670911849e-07\\
154.6	7.0858153124299e-07\\
154.7	7.37499270784616e-07\\
154.8	7.67597163888589e-07\\
154.9	7.98923373387501e-07\\
155	8.31528027664132e-07\\
155.1	8.65463300866014e-07\\
155.2	9.00783496393575e-07\\
155.3	9.37545133795406e-07\\
155.4	9.75807039209627e-07\\
155.5	1.01563043949625e-06\\
155.6	1.05707906021093e-06\\
155.7	1.10021922757701e-06\\
155.8	1.14511997461899e-06\\
155.9	1.1918531516272e-06\\
156	1.24049354113058e-06\\
156.1	1.29111897756123e-06\\
156.2	1.3438104718025e-06\\
156.3	1.3986523408198e-06\\
156.4	1.45573234258132e-06\\
156.5	1.51514181648505e-06\\
156.6	1.57697582951616e-06\\
156.7	1.64133332836911e-06\\
156.8	1.70831729777757e-06\\
156.9	1.77803492530552e-06\\
157	1.85059777286345e-06\\
157.1	1.92612195522365e-06\\
157.2	2.00472832582054e-06\\
157.3	2.08654267013323e-06\\
157.4	2.17169590695948e-06\\
157.5	2.26032429790357e-06\\
157.6	2.35256966541272e-06\\
157.7	2.44857961971114e-06\\
157.8	2.54850779499489e-06\\
157.9	2.65251409526506e-06\\
158	2.76076495019305e-06\\
158.1	2.87343358142687e-06\\
158.2	2.9907002797647e-06\\
158.3	3.11275269363916e-06\\
158.4	3.23978612937351e-06\\
158.5	3.37200386369078e-06\\
158.6	3.50961746897491e-06\\
158.7	3.65284715180485e-06\\
158.8	3.80192210530298e-06\\
158.9	3.95708087586131e-06\\
159	4.11857174483263e-06\\
159.1	4.28665312579647e-06\\
159.2	4.46159397803589e-06\\
159.3	4.64367423688635e-06\\
159.4	4.8331852616446e-06\\
159.5	5.03043030175512e-06\\
159.6	5.23572498201847e-06\\
159.7	5.44939780759849e-06\\
159.8	5.67179068963591e-06\\
159.9	5.90325949230864e-06\\
160	6.14417460221472e-06\\
160.1	6.39492152098717e-06\\
160.2	6.65590148208973e-06\\
160.3	6.92753209277934e-06\\
160.4	7.21024800226174e-06\\
160.5	7.50450159711e-06\\
160.6	7.81076372505625e-06\\
160.7	8.12952444831543e-06\\
160.8	8.46129382764526e-06\\
160.9	8.80660273839554e-06\\
161	9.16600371985332e-06\\
161.1	9.54007185923986e-06\\
161.2	9.92940571177426e-06\\
161.3	1.03346282582745e-05\\
161.4	1.07563879018261e-05\\
161.5	1.11953595051133e-05\\
161.6	1.16522454700696e-05\\
161.7	1.21277768615743e-05\\
161.8	1.26227145769914e-05\\
161.9	1.31378505634194e-05\\
162	1.36740090845997e-05\\
162.1	1.42320480395057e-05\\
162.2	1.48128603347223e-05\\
162.3	1.54173753128082e-05\\
162.4	1.60465602389226e-05\\
162.5	1.67014218480952e-05\\
162.6	1.73830079556069e-05\\
162.7	1.8092409133059e-05\\
162.8	1.88307604528061e-05\\
162.9	1.95992433035389e-05\\
163	2.03990872799214e-05\\
163.1	2.12315721492952e-05\\
163.2	2.2098029898597e-05\\
163.3	2.29998468647563e-05\\
163.4	2.39384659519744e-05\\
163.5	2.49153889394299e-05\\
163.6	2.593217888309e-05\\
163.7	2.69904626154658e-05\\
163.8	2.80919333473023e-05\\
163.9	2.92383533753526e-05\\
164	3.04315569005653e-05\\
164.1	3.16734529611747e-05\\
164.2	3.296602848538e-05\\
164.3	3.43113514684834e-05\\
164.4	3.57115742795526e-05\\
164.5	3.71689371028894e-05\\
164.6	3.86857715197865e-05\\
164.7	4.02645042362891e-05\\
164.8	4.19076609629042e-05\\
164.9	4.36178704524425e-05\\
165	4.53978687024344e-05\\
165.1	4.725050332881e-05\\
165.2	4.91787381178212e-05\\
165.3	5.11856577634542e-05\\
165.4	5.32744727978788e-05\\
165.5	5.54485247227949e-05\\
165.6	5.77112913498369e-05\\
165.7	6.00663923585472e-05\\
165.8	6.25175950807636e-05\\
165.9	6.5068820520623e-05\\
166	6.77241496197701e-05\\
166.1	7.04878297777251e-05\\
166.2	7.33642816377876e-05\\
166.3	7.6358106149268e-05\\
166.4	7.94740919172621e-05\\
166.5	8.27172228516664e-05\\
166.6	8.60926861275657e-05\\
166.7	8.96058804696507e-05\\
166.8	9.32624247738134e-05\\
166.9	9.70681670795998e-05\\
167	0.000101029193907773\\
167.1	0.000105151839977773\\
167.2	0.000109442698320505\\
167.3	0.000113908630802462\\
167.4	0.000118556779077877\\
167.5	0.000123394575986232\\
167.6	0.00012842975741318\\
167.7	0.000133670374633634\\
167.8	0.000139124807156559\\
167.9	0.00014480177609175\\
168	0.000150710358059757\\
168.1	0.000156859999666871\\
168.2	0.000163260532568056\\
168.3	0.000169922189141575\\
168.4	0.000176855618800013\\
168.5	0.000184071904963424\\
168.6	0.000191582582721282\\
168.7	0.00019939965721107\\
168.8	0.000207535622742382\\
168.9	0.000216003482696582\\
169	0.000224816770233295\\
169.1	0.000233989569836176\\
169.2	0.000243536539731751\\
169.3	0.000253472935216452\\
169.4	0.000263814632928313\\
169.5	0.000274578156101332\\
169.6	0.000285780700841879\\
169.7	0.000297440163468182\\
169.8	0.000309575168955517\\
169.9	0.000322205100531344\\
170	0.000335350130466478\\
170.1	0.00034903125211006\\
170.2	0.000363270313218076\\
170.3	0.00037809005062704\\
170.4	0.000393514126326484\\
170.5	0.00040956716498605\\
170.6	0.000426274792995\\
170.7	0.000443663679074348\\
170.8	0.000461761576524077\\
170.9	0.000480597367170262\\
171	0.000500201107079564\\
171.1	0.000520604074110941\\
171.2	0.000541838817377267\\
171.3	0.00056393920869224\\
171.4	0.000586940496080804\\
171.5	0.000610879359434401\\
171.6	0.000635793968395226\\
171.7	0.000661724042557041\\
171.8	0.000688710914073259\\
171.9	0.000716797592766405\\
172	0.000746028833836697\\
172.1	0.000776451208270841\\
172.2	0.000808113176056168\\
172.3	0.000841065162308884\\
172.4	0.000875359636429239\\
172.5	0.000911051194400645\\
172.6	0.000948196644353723\\
172.7	0.000986855095520903\\
172.8	0.00102708805071153\\
172.9	0.00106895950244197\\
173	0.00111253603286032\\
173.1	0.00115788691760961\\
173.2	0.00120508423377894\\
173.3	0.00125420297209689\\
173.4	0.00130532115352677\\
173.5	0.00135851995042896\\
173.6	0.00141388381246065\\
173.7	0.00147150059738952\\
173.8	0.00153146170700329\\
173.9	0.00159386222830301\\
174	0.00165880108017442\\
174.1	0.00172638116573715\\
174.2	0.00179670953057844\\
174.3	0.00186989752708406\\
174.4	0.00194606098508556\\
174.5	0.00202532038904988\\
174.6	0.00210780106204325\\
174.7	0.00219363335670874\\
174.8	0.00228295285350307\\
174.9	0.00237590056644499\\
175	0.00247262315663477\\
175.1	0.00257327315381036\\
175.2	0.0026780091862131\\
175.3	0.00278699621904229\\
175.4	0.00290040580178421\\
175.5	0.00301841632470842\\
175.6	0.00314121328482942\\
175.7	0.0032689895616386\\
175.8	0.00340194570291728\\
175.9	0.00354029022094651\\
176	0.00368423989943599\\
176.1	0.00383402011149744\\
176.2	0.00398986514899408\\
176.3	0.00415201856360066\\
176.4	0.00432073351991193\\
176.5	0.00449627316094118\\
176.6	0.00467891098635069\\
176.7	0.00486893124375862\\
176.8	0.00506662933346628\\
176.9	0.00527231222694868\\
177	0.0054862988994504\\
177.1	0.00570892077702333\\
177.2	0.00594052219834034\\
177.3	0.00618146089161117\\
177.4	0.00643210846691865\\
177.5	0.00669285092428486\\
177.6	0.00696408917776207\\
177.7	0.00724623959583139\\
177.8	0.00753973455837343\\
177.9	0.00784502303045565\\
178	0.00816257115315989\\
178.1	0.00849286285164433\\
178.2	0.00883640046060767\\
178.3	0.00919370536728814\\
178.4	0.00956531867209169\\
178.5	0.00995180186690432\\
178.6	0.0103537375310917\\
178.7	0.0107717300451416\\
178.8	0.0112064063218429\\
178.9	0.0116584165548331\\
179	0.0121284349842742\\
179.1	0.0126171606793387\\
179.2	0.0131253183371027\\
179.3	0.0136536590973509\\
179.4	0.0142029613726912\\
179.5	0.0147740316932731\\
179.6	0.0153677055652755\\
179.7	0.0159848483422025\\
179.8	0.0166263561078817\\
179.9	0.0172931565699055\\
180	0.0179862099620916\\
180.1	0.0187065099543546\\
180.2	0.0194550845681929\\
180.3	0.0202329970957855\\
180.4	0.0210413470204683\\
180.5	0.0218812709361305\\
180.6	0.0227539434628039\\
180.7	0.0236605781554611\\
180.8	0.0246024284027395\\
180.9	0.0255807883120077\\
181	0.0265969935768659\\
181.1	0.0276524223228231\\
181.2	0.0287484959265398\\
181.3	0.0298866798036364\\
181.4	0.0310684841596675\\
181.5	0.0322954646984505\\
181.6	0.0335692232814824\\
181.7	0.0348914085317367\\
181.8	0.0362637163746485\\
181.9	0.037687890508606\\
182	0.0391657227967644\\
182.1	0.0406990535714664\\
182.2	0.0422897718420336\\
182.3	0.0439398153961415\\
182.4	0.0456511707844438\\
182.5	0.0474258731775668\\
182.6	0.0492660060840265\\
182.7	0.0511737009170914\\
182.8	0.0531511363980639\\
182.9	0.0552005377829344\\
183	0.0573241758988687\\
183.1	0.0595243659765014\\
183.2	0.0618034662635883\\
183.3	0.0641638764051742\\
183.4	0.0666080355750908\\
183.5	0.0691384203433468\\
183.6	0.0717575422637512\\
183.7	0.0744679451660277\\
183.8	0.0772722021366602\\
183.9	0.0801729121728125\\
184	0.0831726964939224\\
184.1	0.0862741944959165\\
184.2	0.0894800593335611\\
184.3	0.0927929531171574\\
184.4	0.096215541710693\\
184.5	0.0997504891196851\\
184.6	0.103400451458249\\
184.7	0.107168070486528\\
184.8	0.111055966711408\\
184.9	0.11506673204555\\
185	0.119202922022118\\
185.1	0.123467047565224\\
185.2	0.127861566319081\\
185.3	0.132388873542066\\
185.4	0.13705129257546\\
185.5	0.141851064900488\\
185.6	0.146790339801382\\
185.7	0.151871163656659\\
185.8	0.157095468885453\\
185.9	0.162465062580696\\
186	0.167981614866076\\
186.1	0.173646647019005\\
186.2	0.179461519407326\\
186.3	0.185427419292983\\
186.4	0.191545348561468\\
186.5	0.197816111441418\\
186.6	0.204240302284091\\
186.7	0.210818293477746\\
186.8	0.217550223576888\\
186.9	0.224435985730927\\
187	0.231475216500982\\
187.1	0.238667285157089\\
187.2	0.246011283551051\\
187.3	0.253506016662339\\
187.4	0.261149993915751\\
187.5	0.268941421369995\\
187.6	0.27687819487561\\
187.7	0.284957894299009\\
187.8	0.293177778906433\\
187.9	0.301534783997462\\
188	0.310025518872388\\
188.1	0.318646266210974\\
188.2	0.327392982932239\\
188.3	0.336261302595648\\
188.4	0.345246539393681\\
188.5	0.354343693774205\\
188.6	0.363547459718433\\
188.7	0.372852233686803\\
188.8	0.382252125230752\\
188.9	0.391740969253486\\
189	0.401312339887548\\
189.1	0.410959565941334\\
189.2	0.420675747851249\\
189.3	0.430453776060772\\
189.4	0.440286350732808\\
189.5	0.450166002687522\\
189.6	0.460085115444434\\
189.7	0.470035948235427\\
189.8	0.480010659844419\\
189.9	0.490001333120035\\
190	0.5\\
190.1	0.509998666879965\\
190.2	0.519989340155581\\
190.3	0.529964051764573\\
190.4	0.539914884555566\\
190.5	0.549833997312478\\
190.6	0.559713649267192\\
190.7	0.569546223939228\\
190.8	0.579324252148751\\
190.9	0.589040434058666\\
191	0.598687660112452\\
191.1	0.608259030746514\\
191.2	0.617747874769248\\
191.3	0.627147766313197\\
191.4	0.636452540281567\\
191.5	0.645656306225795\\
191.6	0.654753460606319\\
191.7	0.663738697404352\\
191.8	0.672607017067761\\
191.9	0.681353733789026\\
192	0.689974481127613\\
192.1	0.698465216002538\\
192.2	0.706822221093567\\
192.3	0.715042105700991\\
192.4	0.72312180512439\\
192.5	0.731058578630005\\
192.6	0.738850006084249\\
192.7	0.746493983337661\\
192.8	0.753988716448949\\
192.9	0.761332714842911\\
193	0.768524783499018\\
193.1	0.775564014269073\\
193.2	0.782449776423112\\
193.3	0.789181706522254\\
193.4	0.795759697715909\\
193.5	0.802183888558582\\
193.6	0.808454651438532\\
193.7	0.814572580707017\\
193.8	0.820538480592674\\
193.9	0.826353352980995\\
194	0.832018385133924\\
194.1	0.837534937419304\\
194.2	0.842904531114547\\
194.3	0.848128836343341\\
194.4	0.853209660198618\\
194.5	0.858148935099512\\
194.6	0.86294870742454\\
194.7	0.867611126457934\\
194.8	0.872138433680919\\
194.9	0.876532952434776\\
195	0.880797077977882\\
195.1	0.88493326795445\\
195.2	0.888944033288592\\
195.3	0.892831929513472\\
195.4	0.896599548541751\\
195.5	0.900249510880315\\
195.6	0.903784458289307\\
195.7	0.907207046882843\\
195.8	0.910519940666439\\
195.9	0.913725805504084\\
196	0.916827303506078\\
196.1	0.919827087827187\\
196.2	0.92272779786334\\
196.3	0.925532054833972\\
196.4	0.928242457736249\\
196.5	0.930861579656653\\
196.6	0.933391964424909\\
196.7	0.935836123594826\\
196.8	0.938196533736412\\
196.9	0.940475634023499\\
197	0.942675824101131\\
197.1	0.944799462217066\\
197.2	0.946848863601936\\
197.3	0.948826299082909\\
197.4	0.950733993915973\\
197.5	0.952574126822433\\
197.6	0.954348829215556\\
197.7	0.956060184603859\\
197.8	0.957710228157966\\
197.9	0.959300946428534\\
198	0.960834277203236\\
198.1	0.962312109491394\\
198.2	0.963736283625351\\
198.3	0.965108591468263\\
198.4	0.966430776718518\\
198.5	0.96770453530155\\
198.6	0.968931515840333\\
198.7	0.970113320196364\\
198.8	0.97125150407346\\
198.9	0.972347577677177\\
199	0.973403006423134\\
199.1	0.974419211687992\\
199.2	0.97539757159726\\
199.3	0.976339421844539\\
199.4	0.977246056537196\\
199.5	0.978118729063869\\
199.6	0.978958652979532\\
199.7	0.979767002904215\\
199.8	0.980544915431807\\
199.9	0.981293490045645\\
200	0.982013790037908\\
200.1	0.982706843430095\\
200.2	0.983373643892118\\
200.3	0.984015151657797\\
200.4	0.984632294434724\\
200.5	0.985225968306727\\
200.6	0.985797038627309\\
200.7	0.986346340902649\\
200.8	0.986874681662897\\
200.9	0.987382839320661\\
201	0.987871565015726\\
201.1	0.988341583445167\\
201.2	0.988793593678157\\
201.3	0.989228269954858\\
201.4	0.989646262468908\\
201.5	0.990048198133096\\
201.6	0.990434681327908\\
201.7	0.990806294632712\\
201.8	0.991163599539392\\
201.9	0.991507137148356\\
202	0.99183742884684\\
202.1	0.992154976969544\\
202.2	0.992460265441627\\
202.3	0.992753760404169\\
202.4	0.993035910822238\\
202.5	0.993307149075715\\
202.6	0.993567891533081\\
202.7	0.993818539108389\\
202.8	0.99405947780166\\
202.9	0.994291079222977\\
203	0.994513701100549\\
203.1	0.994727687773051\\
203.2	0.994933370666534\\
203.3	0.995131068756241\\
203.4	0.995321089013649\\
203.5	0.995503726839059\\
203.6	0.995679266480088\\
203.7	0.995847981436399\\
203.8	0.996010134851006\\
203.9	0.996165979888503\\
204	0.996315760100564\\
204.1	0.996459709779054\\
204.2	0.996598054297083\\
204.3	0.996731010438361\\
204.4	0.996858786715171\\
204.5	0.996981583675292\\
204.6	0.997099594198216\\
204.7	0.997213003780958\\
204.8	0.997321990813787\\
204.9	0.99742672684619\\
205	0.997527376843365\\
205.1	0.997624099433555\\
205.2	0.997717047146497\\
205.3	0.997806366643291\\
205.4	0.997892198937957\\
205.5	0.99797467961095\\
205.6	0.998053939014914\\
205.7	0.998130102472916\\
205.8	0.998203290469422\\
205.9	0.998273618834263\\
206	0.998341198919826\\
206.1	0.998406137771697\\
206.2	0.998468538292997\\
206.3	0.99852849940261\\
206.4	0.998586116187539\\
206.5	0.998641480049571\\
206.6	0.998694678846473\\
206.7	0.998745797027903\\
206.8	0.998794915766221\\
206.9	0.99884211308239\\
207	0.99888746396714\\
207.1	0.998931040497558\\
207.2	0.998972911949289\\
207.3	0.999013144904479\\
207.4	0.999051803355646\\
207.5	0.999088948805599\\
207.6	0.999124640363571\\
207.7	0.999158934837691\\
207.8	0.999191886823944\\
207.9	0.999223548791729\\
208	0.999253971166163\\
208.1	0.999283202407234\\
208.2	0.999311289085927\\
208.3	0.999338275957443\\
208.4	0.999364206031605\\
208.5	0.999389120640566\\
208.6	0.999413059503919\\
208.7	0.999436060791308\\
208.8	0.999458161182623\\
208.9	0.999479395925889\\
209	0.999499798892921\\
209.1	0.99951940263283\\
209.2	0.999538238423476\\
209.3	0.999556336320926\\
209.4	0.999573725207005\\
209.5	0.999590432835014\\
209.6	0.999606485873673\\
209.7	0.999621909949373\\
209.8	0.999636729686782\\
209.9	0.99965096874789\\
210	0.999664649869534\\
210.1	0.999677794899469\\
210.2	0.999690424831044\\
210.3	0.999702559836532\\
210.4	0.999714219299158\\
210.5	0.999725421843899\\
210.6	0.999736185367072\\
210.7	0.999746527064784\\
210.8	0.999756463460268\\
210.9	0.999766010430164\\
211	0.999775183229767\\
211.1	0.999783996517303\\
211.2	0.999792464377258\\
211.3	0.999800600342789\\
211.4	0.999808417417279\\
211.5	0.999815928095037\\
211.6	0.9998231443812\\
211.7	0.999830077810858\\
211.8	0.999836739467432\\
211.9	0.999843140000333\\
212	0.99984928964194\\
212.1	0.999855198223908\\
212.2	0.999860875192843\\
212.3	0.999866329625366\\
212.4	0.999871570242587\\
212.5	0.999876605424014\\
212.6	0.999881443220922\\
212.7	0.999886091369197\\
212.8	0.99989055730168\\
212.9	0.999894848160022\\
213	0.999898970806092\\
213.1	0.99990293183292\\
213.2	0.999906737575226\\
213.3	0.99991039411953\\
213.4	0.999913907313873\\
213.5	0.999917282777148\\
213.6	0.999920525908083\\
213.7	0.999923641893851\\
213.8	0.999926635718362\\
213.9	0.999929512170222\\
214	0.99993227585038\\
214.1	0.999934931179479\\
214.2	0.999937482404919\\
214.3	0.999939933607641\\
214.4	0.99994228870865\\
214.5	0.999944551475277\\
214.6	0.999946725527202\\
214.7	0.999948814342237\\
214.8	0.999950821261882\\
214.9	0.999952749496671\\
215	0.999954602131298\\
215.1	0.999956382129548\\
215.2	0.999958092339037\\
215.3	0.999959735495764\\
215.4	0.99996131422848\\
215.5	0.999962831062897\\
215.6	0.99996428842572\\
215.7	0.999965688648532\\
215.8	0.999967033971515\\
215.9	0.999968326547039\\
216	0.999969568443099\\
216.1	0.999970761646625\\
216.2	0.999971908066653\\
216.3	0.999973009537385\\
216.4	0.999974067821117\\
216.5	0.999975084611061\\
216.6	0.999976061534048\\
216.7	0.999977000153135\\
216.8	0.999977901970102\\
216.9	0.999978768427851\\
217	0.99997960091272\\
217.1	0.999980400756696\\
217.2	0.999981169239547\\
217.3	0.999981907590867\\
217.4	0.999982616992044\\
217.5	0.999983298578152\\
217.6	0.999983953439761\\
217.7	0.999984582624687\\
217.8	0.999985187139665\\
217.9	0.99998576795196\\
218	0.999986325990915\\
218.1	0.999986862149437\\
218.2	0.999987377285423\\
218.3	0.999987872223139\\
218.4	0.99998834775453\\
218.5	0.999988804640495\\
218.6	0.999989243612098\\
218.7	0.999989665371742\\
218.8	0.999990070594288\\
218.9	0.999990459928141\\
219	0.99999083399628\\
219.1	0.999991193397262\\
219.2	0.999991538706172\\
219.3	0.999991870475552\\
219.4	0.999992189236275\\
219.5	0.999992495498403\\
219.6	0.999992789751998\\
219.7	0.999993072467907\\
219.8	0.999993344098518\\
219.9	0.999993605078479\\
220	0.999993855825398\\
220.1	0.999994096740508\\
220.2	0.99999432820931\\
220.3	0.999994550602192\\
220.4	0.999994764275018\\
220.5	0.999994969569698\\
220.6	0.999995166814738\\
220.7	0.999995356325763\\
220.8	0.999995538406022\\
220.9	0.999995713346874\\
221	0.999995881428255\\
221.1	0.999996042919124\\
221.2	0.999996198077895\\
221.3	0.999996347152848\\
221.4	0.999996490382531\\
221.5	0.999996627996136\\
221.6	0.999996760213871\\
221.7	0.999996887247306\\
221.8	0.99999700929972\\
221.9	0.999997126566419\\
222	0.99999723923505\\
222.1	0.999997347485905\\
222.2	0.999997451492205\\
222.3	0.99999755142038\\
222.4	0.999997647430335\\
222.5	0.999997739675702\\
222.6	0.999997828304093\\
222.7	0.99999791345733\\
222.8	0.999997995271674\\
222.9	0.999998073878045\\
223	0.999998149402227\\
223.1	0.999998221965075\\
223.2	0.999998291682702\\
223.3	0.999998358666672\\
223.4	0.999998423024171\\
223.5	0.999998484858183\\
223.6	0.999998544267658\\
223.7	0.999998601347659\\
223.8	0.999998656189528\\
223.9	0.999998708881022\\
224	0.999998759506459\\
224.1	0.999998808146848\\
224.2	0.999998854880025\\
224.3	0.999998899780772\\
224.4	0.99999894292094\\
224.5	0.999998984369561\\
224.6	0.999999024192961\\
224.7	0.999999062454866\\
224.8	0.999999099216504\\
224.9	0.999999134536699\\
225	0.999999168471972\\
225.1	0.999999201076627\\
225.2	0.999999232402836\\
225.3	0.999999262500729\\
225.4	0.999999291418469\\
225.5	0.999999319202329\\
225.6	0.99999934589677\\
225.7	0.999999371544508\\
225.8	0.999999396186586\\
225.9	0.999999419862435\\
226	0.999999442609941\\
226.1	0.999999464465506\\
226.2	0.999999485464104\\
226.3	0.999999505639335\\
226.4	0.999999525023484\\
226.5	0.999999543647571\\
226.6	0.999999561541398\\
226.7	0.999999578733599\\
226.8	0.999999595251684\\
226.9	0.999999611122086\\
227	0.999999626370202\\
227.1	0.99999964102043\\
227.2	0.999999655096216\\
227.3	0.999999668620082\\
227.4	0.99999968161367\\
227.5	0.999999694097773\\
227.6	0.999999706092367\\
227.7	0.999999717616647\\
227.8	0.999999728689054\\
227.9	0.999999739327305\\
228	0.999999749548426\\
228.1	0.99999975936877\\
228.2	0.999999768804053\\
228.3	0.999999777869374\\
228.4	0.999999786579238\\
228.5	0.999999794947585\\
228.6	0.999999802987803\\
228.7	0.99999981071276\\
228.8	0.999999818134818\\
228.9	0.999999825265853\\
229	0.999999832117275\\
229.1	0.99999983870005\\
229.2	0.999999845024711\\
229.3	0.999999851101378\\
229.4	0.999999856939775\\
229.5	0.999999862549246\\
229.6	0.999999867938767\\
229.7	0.999999873116961\\
229.8	0.999999878092116\\
229.9	0.999999882872191\\
230	0.999999887464838\\
};
\addlegendentry{$\mu{}_{\text{VT}}\text{(}\mu{}_{\text{i}}\text{)}$};

\addplot [color=green,line width=1.5pt,mark size=4.0pt,only marks,mark=asterisk,mark options={solid},forget plot]
  table[row sep=crcr]{195	0.88\\
};
\node[right, inner sep=0mm, text=black]
at (axis cs:195,0.88,0) {$\text{  }\leftarrow\text{ (}\mu{}_{\text{i}}\text{=195; }\mu{}_{\text{VT}}\text{(195)=0.8808)}$};
\end{axis}
\end{tikzpicture}%
	\caption{Función de pertenencia del conjunto difuso \emph{"Muy Alto"} ($\mu_{VT}(\mu_{i})$).}
	\label{fig:fuzzyset-verytall-example}
\end{figure}

\subsection{Variables lingüísticas}
\label{sec:lang-variables}
Una variable lingüística se puede definir de manera informal como una variable que puede tomar palabras del lenguaje natural como valores. Para formular palabras y utilizarlas como valores de variables lingüísticas se utilizan conjuntos difusos (por medio de sus correspondientes funciones de pertenencia). Así pues, una variable lingüística se puede definir como \cite{wang1997}:

\begin{definition}\label{def:informal-lang-variable}
Una \emph{variable lingüística} es un tipo de variable que puede tomar palabras como valores. Las palabras son caracterizadas por conjuntos difusos definidos en el mismo universo en el que la variable es definida.
\end{definition}

\begin{example}\label{ex:lang-variable}
\normalfont
Supongamos que la temperatura medida por un termómetro en grados celsius (ºC) es una variable $x$ que puede tomar valores enteros en el intervalo $[0,100]$. Este intervalo de valores posibles es lo que se define como el \emph{universo de discurso} o \emph{universo de referencia} $U$. Podemos definir sobre $U$ los conjuntos difusos \emph{``baja''},\emph{``media''} y \emph{``alta''} para modelar los conceptos de temperatura alta, media y baja respectivamente. Si consideramos que $x$ es una variable lingüística entonces $x$ puede tomar los valores \emph{``baja''},\emph{``media''} y \emph{``alta''}. Es decir, podemos decir que: \emph{``x es baja''},\emph{``x es media''} o \emph{``x es alta''}, o lo que es lo mismo \emph{``Temperatura es baja''},\emph{``Temperatura es media''} o \emph{``Temperatura es alta''}. En la figura \ref{fig:example-temp-lang-variable} se representan las funciones de pertenencia para los conjuntos difusos \emph{``Temperatura baja''},\emph{``Temperatura media''} y \emph{``Temperatura alta''}.
\end{example}

\begin{figure}[t]
	\centering
	\setlength\figureheight{4.5cm}
	\setlength\figurewidth{12cm}
	% This file was created by matlab2tikz v0.4.7 (commit 05e36af552d2a8253e2edb4c73807088323ff197) running on MATLAB 8.0.
% Copyright (c) 2008--2014, Nico Schlömer <nico.schloemer@gmail.com>
% All rights reserved.
% Minimal pgfplots version: 1.3
% 
\begin{tikzpicture}

\begin{axis}[%
width=\figurewidth,
height=\figureheight,
clip=false,
scale only axis,
xmin=0,
xmax=100,
xlabel={x - Temperatura},
ymin=0,
ymax=1
]
\addplot [color=black,solid,line width=1.5pt,forget plot]
  table[row sep=crcr]{0	1\\
1	0.975\\
2	0.95\\
3	0.925\\
4	0.9\\
5	0.875\\
6	0.85\\
7	0.825\\
8	0.8\\
9	0.775\\
10	0.75\\
11	0.725\\
12	0.7\\
13	0.675\\
14	0.65\\
15	0.625\\
16	0.6\\
17	0.575\\
18	0.55\\
19	0.525\\
20	0.5\\
21	0.475\\
22	0.45\\
23	0.425\\
24	0.4\\
25	0.375\\
26	0.35\\
27	0.325\\
28	0.3\\
29	0.275\\
30	0.25\\
31	0.225\\
32	0.2\\
33	0.175\\
34	0.15\\
35	0.125\\
36	0.1\\
37	0.075\\
38	0.05\\
39	0.025\\
40	0\\
};
\addplot [color=black,solid,line width=1.5pt,forget plot]
  table[row sep=crcr]{10	0\\
11	0.025\\
12	0.05\\
13	0.075\\
14	0.1\\
15	0.125\\
16	0.15\\
17	0.175\\
18	0.2\\
19	0.225\\
20	0.25\\
21	0.275\\
22	0.3\\
23	0.325\\
24	0.35\\
25	0.375\\
26	0.4\\
27	0.425\\
28	0.45\\
29	0.475\\
30	0.5\\
31	0.525\\
32	0.55\\
33	0.575\\
34	0.6\\
35	0.625\\
36	0.65\\
37	0.675\\
38	0.7\\
39	0.725\\
40	0.75\\
41	0.775\\
42	0.8\\
43	0.825\\
44	0.85\\
45	0.875\\
46	0.9\\
47	0.925\\
48	0.95\\
49	0.975\\
50	1\\
51	0.975\\
52	0.95\\
53	0.925\\
54	0.9\\
55	0.875\\
56	0.85\\
57	0.825\\
58	0.8\\
59	0.775\\
60	0.75\\
61	0.725\\
62	0.7\\
63	0.675\\
64	0.65\\
65	0.625\\
66	0.6\\
67	0.575\\
68	0.55\\
69	0.525\\
70	0.5\\
71	0.475\\
72	0.45\\
73	0.425\\
74	0.4\\
75	0.375\\
76	0.35\\
77	0.325\\
78	0.3\\
79	0.275\\
80	0.25\\
81	0.225\\
82	0.2\\
83	0.175\\
84	0.15\\
85	0.125\\
86	0.1\\
87	0.075\\
88	0.05\\
89	0.025\\
90	0\\
};
\addplot [color=black,solid,line width=1.5pt,forget plot]
  table[row sep=crcr]{60	0\\
61	0.025\\
62	0.05\\
63	0.075\\
64	0.1\\
65	0.125\\
66	0.15\\
67	0.175\\
68	0.2\\
69	0.225\\
70	0.25\\
71	0.275\\
72	0.3\\
73	0.325\\
74	0.35\\
75	0.375\\
76	0.4\\
77	0.425\\
78	0.45\\
79	0.475\\
80	0.5\\
81	0.525\\
82	0.55\\
83	0.575\\
84	0.6\\
85	0.625\\
86	0.65\\
87	0.675\\
88	0.7\\
89	0.725\\
90	0.75\\
91	0.775\\
92	0.8\\
93	0.825\\
94	0.85\\
95	0.875\\
96	0.9\\
97	0.925\\
98	0.95\\
99	0.975\\
100	1\\
};
\node[right, inner sep=0mm, text=black]
at (axis cs:1,1.05,0) {$\text{Baja (}\mu{}_{\text{Baja}}\text{(x))}$};
\node[right, inner sep=0mm, text=black]
at (axis cs:40,1.05,0) {$\text{Media (}\mu{}_{\text{Media}}\text{(x))}$};
\node[right, inner sep=0mm, text=black]
at (axis cs:80,1.05,0) {$\text{Alta (}\mu{}_{\text{Alta}}\text{(x))}$};
\end{axis}
\end{tikzpicture}%
	\caption{Ejemplo \ref{ex:lang-variable}: Funciones de pertenencia para los valores \emph{``baja''} ($\mu_{Baja}$),\emph{``media''} ($\mu_{Media}$) y \emph{``alta''} ($\mu_{Alta}$) de la variable lingüística \emph{Temperatura}.}
	\label{fig:example-temp-lang-variable}
\end{figure}

Una definición más formal de \emph{variable lingüística} (equivalente a la definición \ref{def:informal-lang-variable}), dada por Zadeh \cite{Zadeh75}, es la siguiente:

\begin{definition}\label{def:formal-lang-variable}
\normalfont
Una \emph{variable lingüística} se caracteriza por la tupla $(X,T,U,M)$, donde:
\begin{itemize}
  \item $X$ es el nombre de la variable lingüística. En el ejemplo \ref{ex:lang-variable},``Temperatura en ºC''.
  \item $T$ es el conjunto de valores lingüísticos que $X$ puede tomar. En el ejemplo \ref{ex:lang-variable}, $T$ = \{baja, media, alta\}
  \item $U$ es el universo de referencia del que cada valor lingüístico toma sus valores cuantitativos (escalares). En el ejemplo \ref{ex:lang-variable}, $U = [0,100]$
  \item $M$ es la regla semántica que relaciona cada valor lingüístico en $T$ con un conjunto difuso en $U$. En el ejemplo \ref{ex:lang-variable}, $M$ relaciona los valores ``baja'', ``media'' y ``alta'' con las funciones de pertenencia representadas en la figura \ref{fig:example-temp-lang-variable}.
\end{itemize}
\end{definition}
\subsection{Operaciones sobre conjuntos difusos}
\label{sec:fuzzy-set-operations}
En esta sección se definen algunas operaciones básicas sobre conjuntos difusos. Un conjunto difuso compuesto a partir de dos conjuntos difusos sobre el mismo universo \emph{U} puede definirse por su función de pertenencia tal que:
\begin{definition}
Dada una función \begin{math}\emph{F}:[0,1]^2\rightarrow[0,1]\end{math} y dos conjuntos difusos \emph{A} y \emph{B} definidos sobre el mismo universo \emph{U}, \begin{math}A,B\in FS(U)\end{math}, denotamos como \begin{math}F(A,B)\end{math} el conjunto difuso sobre \emph{U} cuya función de pertenencia viene dada por:
\begin{equation}
\mu_{F(A,B)}(u_{i}) = F(A(u_{i}),B(u_{i}))
\end{equation}
\end{definition}
De esta forma, se pueden definir las operaciones de unión ($\cup$) y de intersección ($\cap$) clásicas sobre los subconjuntos difusos \emph{A} y \emph{B} sobre \emph{U} \cite{dubois1980} tal que:
\begin{equation}
\forall\mu_{i}\in U,\quad\mu_{A\cup B}(u_{i}) = max(\mu_{A}(u_{i}),\mu_{B}(u_{i}))
\end{equation}
\begin{equation}
\forall\mu_{i}\in U,\quad\mu_{A\cap B}(u_{i}) = min(\mu_{A}(u_{i}),\mu_{B}(u_{i}))
\end{equation}
donde $\mu_{A\cup B}$ y $\mu_{A\cap B}$ son las funciones de pertenencia de $A\cup B$ (unión) y $A\cap B$ (intersección) respectivamente.
\begin{definition}
El \emph{complementario} \emph{$\overline{A}$} de un conjunto difuso \emph{A} sobre \emph{U} puede definirse por su función de pertenencia \cite{Zadeh65}, tal que:
\begin{equation}
\forall \mu_{i} \in U, \quad \mu_{\overline{A}}(\mu_{i}) = 1 - \mu_{A}(\mu_{i})
\end{equation}
\end{definition}


\section{T-normas y operadores de agregación}\label{sec:t-norms-aggregation-operator}
Dada una función \begin{math}\emph{F}:[0,1]^2\rightarrow[0,1]\end{math} y dos conjuntos difusos \emph{A} y \emph{B} definidos sobre el mismo universo \emph{U}, \begin{math}A,B\in FS(U)\end{math}, denotamos como \begin{math}F(A,B)\end{math} el conjunto difuso sobre \emph{U} cuya función de pertenencia viene dada por:
\begin{equation}
\mu_{F(A,B)}(u_{i}) = F(A(u_{i}),B(u_{i}))
\end{equation}
Una clase importante de este tipo de funciones son las llamadas t-normas y t-conormas (normas triangulares) \cite{klement2000triangular}.
\begin{definition}\label{def:tnorma}
Una t-norma es una operación binaria \emph{T} en el intervalo $[0,1]$ que es conmutativa, asociativa, monótona y tiene el valor \emph{1} como elemento neutro. Es decir, una función $\emph{T} : [0,1]^2 \rightarrow [0,1]$ tal que $\forall x,y,z \in [0,1]$:
\begin{enumerate}[label=(T\arabic*),ref=(T\arabic*)]
   \item $T(x,y) = T(y,x)$ (Conmutatividad)\label{T1}
   \item $T(x,T(y,z)) = T(T(x,y),z)$ (Asociatividad)\label{T2}
   \item $T(x,y) \leq T(x,z)$ cuando $y \leq z$ (Monotonía)\label{T3}
   \item $T(x,1) = x$ (Elemento neutro)\label{T4}
  \end{enumerate}
\end{definition}
Estas propiedades son suficientes para garantizar que las t-normas generalizan la conjunción clásica ($x\bigwedge y$) cuando se aplican a valores booleanos ($T(0,0)=0$,$T(0,1)=T(1,0)=0$ y $T(1,1)=1$).
Algunos ejemplos prominentes de t-normas son los siguientes:
\begin{itemize}
	\item \bfseries Mínimo o t-norma de Gödel: $T_{G}(x,y) = min\{x,y\}$
	\item \bfseries Producto: $T_{P}(x,y) = xy$
	\item \bfseries \L{}ukasiewicz: $T_{L}(x,y) = max\{x+y-1,0\}$
	\item \bfseries T-norma drástica: \mdseries $T_{D}(x,y) = \begin{cases} y & if\quad x=1, \\ x & if\quad y=1, \\0 & \text{en cualquier otro caso}. \end{cases}$
\end{itemize}
\begin{figure}[H]
	\centering
	\begin{subfigure}[b]{0.45\textwidth}
		\setlength\figureheight{4.5cm}
		\setlength\figurewidth{6cm}
		% This file was created by matlab2tikz v0.4.7 (commit 4d7ae5c4fd0932fb51051d86111bc23ed23e4580) running on MATLAB 8.0.
% Copyright (c) 2008--2014, Nico Schlömer <nico.schloemer@gmail.com>
% All rights reserved.
% Minimal pgfplots version: 1.3
% 
\begin{tikzpicture}

\begin{axis}[%
width=\figurewidth,
height=\figureheight,
view={-37.5}{30},
scale only axis,
xmin=0,
xmax=1,
xlabel={x},
xmajorgrids,
ymin=0,
ymax=1,
ylabel={y},
ymajorgrids,
zmin=0,
zmax=1,
zlabel={$\text{T}_{\text{G}}\text{(x,y)}$},
zmajorgrids,
axis x line*=bottom,
axis y line*=left,
axis z line*=left
]

\addplot3[%
surf,
shader=faceted,
draw=black,
colormap/jet,
mesh/rows=11]
table[row sep=crcr,header=false] {0	0	0\\
0	0.1	0\\
0	0.2	0\\
0	0.3	0\\
0	0.4	0\\
0	0.5	0\\
0	0.6	0\\
0	0.7	0\\
0	0.8	0\\
0	0.9	0\\
0	1	0\\
0.1	0	0\\
0.1	0.1	0.1\\
0.1	0.2	0.1\\
0.1	0.3	0.1\\
0.1	0.4	0.1\\
0.1	0.5	0.1\\
0.1	0.6	0.1\\
0.1	0.7	0.1\\
0.1	0.8	0.1\\
0.1	0.9	0.1\\
0.1	1	0.1\\
0.2	0	0\\
0.2	0.1	0.1\\
0.2	0.2	0.2\\
0.2	0.3	0.2\\
0.2	0.4	0.2\\
0.2	0.5	0.2\\
0.2	0.6	0.2\\
0.2	0.7	0.2\\
0.2	0.8	0.2\\
0.2	0.9	0.2\\
0.2	1	0.2\\
0.3	0	0\\
0.3	0.1	0.1\\
0.3	0.2	0.2\\
0.3	0.3	0.3\\
0.3	0.4	0.3\\
0.3	0.5	0.3\\
0.3	0.6	0.3\\
0.3	0.7	0.3\\
0.3	0.8	0.3\\
0.3	0.9	0.3\\
0.3	1	0.3\\
0.4	0	0\\
0.4	0.1	0.1\\
0.4	0.2	0.2\\
0.4	0.3	0.3\\
0.4	0.4	0.4\\
0.4	0.5	0.4\\
0.4	0.6	0.4\\
0.4	0.7	0.4\\
0.4	0.8	0.4\\
0.4	0.9	0.4\\
0.4	1	0.4\\
0.5	0	0\\
0.5	0.1	0.1\\
0.5	0.2	0.2\\
0.5	0.3	0.3\\
0.5	0.4	0.4\\
0.5	0.5	0.5\\
0.5	0.6	0.5\\
0.5	0.7	0.5\\
0.5	0.8	0.5\\
0.5	0.9	0.5\\
0.5	1	0.5\\
0.6	0	0\\
0.6	0.1	0.1\\
0.6	0.2	0.2\\
0.6	0.3	0.3\\
0.6	0.4	0.4\\
0.6	0.5	0.5\\
0.6	0.6	0.6\\
0.6	0.7	0.6\\
0.6	0.8	0.6\\
0.6	0.9	0.6\\
0.6	1	0.6\\
0.7	0	0\\
0.7	0.1	0.1\\
0.7	0.2	0.2\\
0.7	0.3	0.3\\
0.7	0.4	0.4\\
0.7	0.5	0.5\\
0.7	0.6	0.6\\
0.7	0.7	0.7\\
0.7	0.8	0.7\\
0.7	0.9	0.7\\
0.7	1	0.7\\
0.8	0	0\\
0.8	0.1	0.1\\
0.8	0.2	0.2\\
0.8	0.3	0.3\\
0.8	0.4	0.4\\
0.8	0.5	0.5\\
0.8	0.6	0.6\\
0.8	0.7	0.7\\
0.8	0.8	0.8\\
0.8	0.9	0.8\\
0.8	1	0.8\\
0.9	0	0\\
0.9	0.1	0.1\\
0.9	0.2	0.2\\
0.9	0.3	0.3\\
0.9	0.4	0.4\\
0.9	0.5	0.5\\
0.9	0.6	0.6\\
0.9	0.7	0.7\\
0.9	0.8	0.8\\
0.9	0.9	0.9\\
0.9	1	0.9\\
1	0	0\\
1	0.1	0.1\\
1	0.2	0.2\\
1	0.3	0.3\\
1	0.4	0.4\\
1	0.5	0.5\\
1	0.6	0.6\\
1	0.7	0.7\\
1	0.8	0.8\\
1	0.9	0.9\\
1	1	1\\
};
\end{axis}
\end{tikzpicture}%
		\caption{Mínimo}
		\label{fig:t-norms-min}
	\end{subfigure}
	\qquad
	\begin{subfigure}[b]{0.45\textwidth}
		\setlength\figureheight{4.5cm}
		\setlength\figurewidth{6cm}
		% This file was created by matlab2tikz v0.4.7 (commit 4d7ae5c4fd0932fb51051d86111bc23ed23e4580) running on MATLAB 8.0.
% Copyright (c) 2008--2014, Nico Schlömer <nico.schloemer@gmail.com>
% All rights reserved.
% Minimal pgfplots version: 1.3
% 
\begin{tikzpicture}

\begin{axis}[%
width=\figurewidth,
height=\figureheight,
view={-37.5}{30},
scale only axis,
xmin=0,
xmax=1,
xlabel={x},
xmajorgrids,
ymin=0,
ymax=1,
ylabel={y},
ymajorgrids,
zmin=0,
zmax=1,
zlabel={$\text{T}_{\text{P}}\text{(x,y)}$},
zmajorgrids,
axis x line*=bottom,
axis y line*=left,
axis z line*=left
]

\addplot3[%
surf,
shader=faceted,
draw=black,
colormap/jet,
mesh/rows=11]
table[row sep=crcr,header=false] {0	0	0\\
0	0.1	0\\
0	0.2	0\\
0	0.3	0\\
0	0.4	0\\
0	0.5	0\\
0	0.6	0\\
0	0.7	0\\
0	0.8	0\\
0	0.9	0\\
0	1	0\\
0.1	0	0\\
0.1	0.1	0.01\\
0.1	0.2	0.02\\
0.1	0.3	0.03\\
0.1	0.4	0.04\\
0.1	0.5	0.05\\
0.1	0.6	0.06\\
0.1	0.7	0.07\\
0.1	0.8	0.08\\
0.1	0.9	0.09\\
0.1	1	0.1\\
0.2	0	0\\
0.2	0.1	0.02\\
0.2	0.2	0.04\\
0.2	0.3	0.06\\
0.2	0.4	0.08\\
0.2	0.5	0.1\\
0.2	0.6	0.12\\
0.2	0.7	0.14\\
0.2	0.8	0.16\\
0.2	0.9	0.18\\
0.2	1	0.2\\
0.3	0	0\\
0.3	0.1	0.03\\
0.3	0.2	0.06\\
0.3	0.3	0.09\\
0.3	0.4	0.12\\
0.3	0.5	0.15\\
0.3	0.6	0.18\\
0.3	0.7	0.21\\
0.3	0.8	0.24\\
0.3	0.9	0.27\\
0.3	1	0.3\\
0.4	0	0\\
0.4	0.1	0.04\\
0.4	0.2	0.08\\
0.4	0.3	0.12\\
0.4	0.4	0.16\\
0.4	0.5	0.2\\
0.4	0.6	0.24\\
0.4	0.7	0.28\\
0.4	0.8	0.32\\
0.4	0.9	0.36\\
0.4	1	0.4\\
0.5	0	0\\
0.5	0.1	0.05\\
0.5	0.2	0.1\\
0.5	0.3	0.15\\
0.5	0.4	0.2\\
0.5	0.5	0.25\\
0.5	0.6	0.3\\
0.5	0.7	0.35\\
0.5	0.8	0.4\\
0.5	0.9	0.45\\
0.5	1	0.5\\
0.6	0	0\\
0.6	0.1	0.06\\
0.6	0.2	0.12\\
0.6	0.3	0.18\\
0.6	0.4	0.24\\
0.6	0.5	0.3\\
0.6	0.6	0.36\\
0.6	0.7	0.42\\
0.6	0.8	0.48\\
0.6	0.9	0.54\\
0.6	1	0.6\\
0.7	0	0\\
0.7	0.1	0.07\\
0.7	0.2	0.14\\
0.7	0.3	0.21\\
0.7	0.4	0.28\\
0.7	0.5	0.35\\
0.7	0.6	0.42\\
0.7	0.7	0.49\\
0.7	0.8	0.56\\
0.7	0.9	0.63\\
0.7	1	0.7\\
0.8	0	0\\
0.8	0.1	0.08\\
0.8	0.2	0.16\\
0.8	0.3	0.24\\
0.8	0.4	0.32\\
0.8	0.5	0.4\\
0.8	0.6	0.48\\
0.8	0.7	0.56\\
0.8	0.8	0.64\\
0.8	0.9	0.72\\
0.8	1	0.8\\
0.9	0	0\\
0.9	0.1	0.09\\
0.9	0.2	0.18\\
0.9	0.3	0.27\\
0.9	0.4	0.36\\
0.9	0.5	0.45\\
0.9	0.6	0.54\\
0.9	0.7	0.63\\
0.9	0.8	0.72\\
0.9	0.9	0.81\\
0.9	1	0.9\\
1	0	0\\
1	0.1	0.1\\
1	0.2	0.2\\
1	0.3	0.3\\
1	0.4	0.4\\
1	0.5	0.5\\
1	0.6	0.6\\
1	0.7	0.7\\
1	0.8	0.8\\
1	0.9	0.9\\
1	1	1\\
};
\end{axis}
\end{tikzpicture}%
		\caption{Producto}
		\label{fig:t-norms-prod}
	\end{subfigure}
	
	\vspace{1 cm}
	\begin{subfigure}[b]{0.45\textwidth}
		\setlength\figureheight{4.5cm}
		\setlength\figurewidth{6cm}
		% This file was created by matlab2tikz v0.4.7 (commit 4d7ae5c4fd0932fb51051d86111bc23ed23e4580) running on MATLAB 8.0.
% Copyright (c) 2008--2014, Nico Schlömer <nico.schloemer@gmail.com>
% All rights reserved.
% Minimal pgfplots version: 1.3
% 
\begin{tikzpicture}

\begin{axis}[%
width=\figurewidth,
height=\figureheight,
view={-37.5}{30},
scale only axis,
xmin=0,
xmax=1,
xlabel={x},
xmajorgrids,
ymin=0,
ymax=1,
ylabel={y},
ymajorgrids,
zmin=0,
zmax=1,
zlabel={$\text{T}_{\text{L}}\text{(x,y)}$},
zmajorgrids,
axis x line*=bottom,
axis y line*=left,
axis z line*=left
]

\addplot3[%
surf,
shader=faceted,
draw=black,
colormap/jet,
mesh/rows=11]
table[row sep=crcr,header=false] {0	0	0\\
0	0.1	0\\
0	0.2	0\\
0	0.3	0\\
0	0.4	0\\
0	0.5	0\\
0	0.6	0\\
0	0.7	0\\
0	0.8	0\\
0	0.9	0\\
0	1	0\\
0.1	0	0\\
0.1	0.1	0\\
0.1	0.2	0\\
0.1	0.3	0\\
0.1	0.4	0\\
0.1	0.5	0\\
0.1	0.6	0\\
0.1	0.7	0\\
0.1	0.8	0\\
0.1	0.9	0\\
0.1	1	0.1\\
0.2	0	0\\
0.2	0.1	0\\
0.2	0.2	0\\
0.2	0.3	0\\
0.2	0.4	0\\
0.2	0.5	0\\
0.2	0.6	0\\
0.2	0.7	0\\
0.2	0.8	0\\
0.2	0.9	0.1\\
0.2	1	0.2\\
0.3	0	0\\
0.3	0.1	0\\
0.3	0.2	0\\
0.3	0.3	0\\
0.3	0.4	0\\
0.3	0.5	0\\
0.3	0.6	0\\
0.3	0.7	0\\
0.3	0.8	0.1\\
0.3	0.9	0.2\\
0.3	1	0.3\\
0.4	0	0\\
0.4	0.1	0\\
0.4	0.2	0\\
0.4	0.3	0\\
0.4	0.4	0\\
0.4	0.5	0\\
0.4	0.6	0\\
0.4	0.7	0.1\\
0.4	0.8	0.2\\
0.4	0.9	0.3\\
0.4	1	0.4\\
0.5	0	0\\
0.5	0.1	0\\
0.5	0.2	0\\
0.5	0.3	0\\
0.5	0.4	0\\
0.5	0.5	0\\
0.5	0.6	0.1\\
0.5	0.7	0.2\\
0.5	0.8	0.3\\
0.5	0.9	0.4\\
0.5	1	0.5\\
0.6	0	0\\
0.6	0.1	0\\
0.6	0.2	0\\
0.6	0.3	0\\
0.6	0.4	0\\
0.6	0.5	0.1\\
0.6	0.6	0.2\\
0.6	0.7	0.3\\
0.6	0.8	0.4\\
0.6	0.9	0.5\\
0.6	1	0.6\\
0.7	0	0\\
0.7	0.1	0\\
0.7	0.2	0\\
0.7	0.3	0\\
0.7	0.4	0.1\\
0.7	0.5	0.2\\
0.7	0.6	0.3\\
0.7	0.7	0.4\\
0.7	0.8	0.5\\
0.7	0.9	0.6\\
0.7	1	0.7\\
0.8	0	0\\
0.8	0.1	0\\
0.8	0.2	0\\
0.8	0.3	0.1\\
0.8	0.4	0.2\\
0.8	0.5	0.3\\
0.8	0.6	0.4\\
0.8	0.7	0.5\\
0.8	0.8	0.6\\
0.8	0.9	0.7\\
0.8	1	0.8\\
0.9	0	0\\
0.9	0.1	0\\
0.9	0.2	0.1\\
0.9	0.3	0.2\\
0.9	0.4	0.3\\
0.9	0.5	0.4\\
0.9	0.6	0.5\\
0.9	0.7	0.6\\
0.9	0.8	0.7\\
0.9	0.9	0.8\\
0.9	1	0.9\\
1	0	0\\
1	0.1	0.1\\
1	0.2	0.2\\
1	0.3	0.3\\
1	0.4	0.4\\
1	0.5	0.5\\
1	0.6	0.6\\
1	0.7	0.7\\
1	0.8	0.8\\
1	0.9	0.9\\
1	1	1\\
};
\end{axis}
\end{tikzpicture}%
		\caption{\L{}ukasiewicz}
		\label{fig:t-norms-lukasiewicz}
	\end{subfigure}
	\qquad
	\begin{subfigure}[b]{0.45\textwidth}
			\setlength\figureheight{4.5cm}
			\setlength\figurewidth{6cm}
			% This file was created by matlab2tikz v0.4.7 (commit 730e9c094bba45cf576cc1ba69976f02cfab8e0c) running on MATLAB 8.0.
% Copyright (c) 2008--2014, Nico Schlömer <nico.schloemer@gmail.com>
% All rights reserved.
% Minimal pgfplots version: 1.3
% 
\begin{tikzpicture}

\begin{axis}[%
width=\figurewidth,
height=\figureheight,
view={-37.5}{30},
scale only axis,
xmin=0,
xmax=1,
xlabel={x},
xmajorgrids,
ymin=0,
ymax=1,
ylabel={y},
ymajorgrids,
zmin=0,
zmax=1,
zlabel={$\text{T}_{\text{D}}\text{(x,y)}$},
zmajorgrids,
axis x line*=bottom,
axis y line*=left,
axis z line*=left
]

\addplot3[%
surf,
shader=faceted,
draw=black,
colormap/jet,
mesh/rows=11]
table[row sep=crcr,header=false] {0	0	0\\
0	0.1	0\\
0	0.2	0\\
0	0.3	0\\
0	0.4	0\\
0	0.5	0\\
0	0.6	0\\
0	0.7	0\\
0	0.8	0\\
0	0.9	0\\
0	1	0\\
0.1	0	0\\
0.1	0.1	0\\
0.1	0.2	0\\
0.1	0.3	0\\
0.1	0.4	0\\
0.1	0.5	0\\
0.1	0.6	0\\
0.1	0.7	0\\
0.1	0.8	0\\
0.1	0.9	0\\
0.1	1	0.1\\
0.2	0	0\\
0.2	0.1	0\\
0.2	0.2	0\\
0.2	0.3	0\\
0.2	0.4	0\\
0.2	0.5	0\\
0.2	0.6	0\\
0.2	0.7	0\\
0.2	0.8	0\\
0.2	0.9	0\\
0.2	1	0.2\\
0.3	0	0\\
0.3	0.1	0\\
0.3	0.2	0\\
0.3	0.3	0\\
0.3	0.4	0\\
0.3	0.5	0\\
0.3	0.6	0\\
0.3	0.7	0\\
0.3	0.8	0\\
0.3	0.9	0\\
0.3	1	0.3\\
0.4	0	0\\
0.4	0.1	0\\
0.4	0.2	0\\
0.4	0.3	0\\
0.4	0.4	0\\
0.4	0.5	0\\
0.4	0.6	0\\
0.4	0.7	0\\
0.4	0.8	0\\
0.4	0.9	0\\
0.4	1	0.4\\
0.5	0	0\\
0.5	0.1	0\\
0.5	0.2	0\\
0.5	0.3	0\\
0.5	0.4	0\\
0.5	0.5	0\\
0.5	0.6	0\\
0.5	0.7	0\\
0.5	0.8	0\\
0.5	0.9	0\\
0.5	1	0.5\\
0.6	0	0\\
0.6	0.1	0\\
0.6	0.2	0\\
0.6	0.3	0\\
0.6	0.4	0\\
0.6	0.5	0\\
0.6	0.6	0\\
0.6	0.7	0\\
0.6	0.8	0\\
0.6	0.9	0\\
0.6	1	0.6\\
0.7	0	0\\
0.7	0.1	0\\
0.7	0.2	0\\
0.7	0.3	0\\
0.7	0.4	0\\
0.7	0.5	0\\
0.7	0.6	0\\
0.7	0.7	0\\
0.7	0.8	0\\
0.7	0.9	0\\
0.7	1	0.7\\
0.8	0	0\\
0.8	0.1	0\\
0.8	0.2	0\\
0.8	0.3	0\\
0.8	0.4	0\\
0.8	0.5	0\\
0.8	0.6	0\\
0.8	0.7	0\\
0.8	0.8	0\\
0.8	0.9	0\\
0.8	1	0.8\\
0.9	0	0\\
0.9	0.1	0\\
0.9	0.2	0\\
0.9	0.3	0\\
0.9	0.4	0\\
0.9	0.5	0\\
0.9	0.6	0\\
0.9	0.7	0\\
0.9	0.8	0\\
0.9	0.9	0\\
0.9	1	0.9\\
1	0	0\\
1	0.1	0.1\\
1	0.2	0.2\\
1	0.3	0.3\\
1	0.4	0.4\\
1	0.5	0.5\\
1	0.6	0.6\\
1	0.7	0.7\\
1	0.8	0.8\\
1	0.9	0.9\\
1	1	1\\
};
\end{axis}
\end{tikzpicture}%
			\caption{T-norma drástica}
			\label{fig:t-norms-drastic}
	\end{subfigure}
	\label{fig:t-norms}
	\caption{Representación gráfica de las t-normas \emph{Mínimo} (\ref{fig:t-norms-min}), \emph{Producto} (\ref{fig:t-norms-prod}), \emph{\L{}ukasiewicz} (\ref{fig:t-norms-lukasiewicz}) y la \emph{t-norma drástica} (\ref{fig:t-norms-drastic}).}
\end{figure}

\begin{definition}
Una t-conorma es una operación binaria \emph{S} en el intervalo $[0,1]$ que es conmutativa, asociativa, monotona y tiene el valor \emph{0} como elemento neutro. Es decir, una función $\emph{S} : [0,1]^2 \rightarrow [0,1]$ tal que $\forall x,y,z \in [0,1]$:
\begin{enumerate}[label=(S\arabic*),ref=(S\arabic*)]
   \item $S(x,y) = S(y,x)$ (Conmutatividad)\label{S1}
   \item $S(x,S(y,z)) = S(S(x,y),z)$ (Asociatividad)\label{S2}
   \item $S(x,y) \leq S(x,z)$ cuando $y \leq z$ (Monotonía)\label{S3}
   \item $S(x,0) = x$ (Elemento neutro)\label{S4}
  \end{enumerate}
\end{definition}
De la misma forma que las t-normas (definición \ref{def:tnorma}) generalizan la conjunción clásica, las t-conormas generalizan la disyunción ($x\bigvee y$). A continuación se incluyen algunos ejemplos de t-conormas:
\begin{itemize}
	\item \bfseries Máximo: $S_{max}(x,y) = max\{x,y\}$
	\item \bfseries Suma probabilística: $S_{P}(x,y) = x + y - xy$
	\item \bfseries Suma acotada: $S_{B}(x,y) = min\{x+y,1\}$
	\item \bfseries T-conorma drástica: \mdseries $S_{D}(x,y) = \begin{cases} y & if\quad x=0, \\ x & if\quad y=0, \\1 & \text{en cualquier otro caso}. \end{cases}$
\end{itemize}
\begin{figure}[H]
	\centering
	\begin{subfigure}[b]{0.45\textwidth}
		\setlength\figureheight{4.5cm}
		\setlength\figurewidth{6cm}
		% This file was created by matlab2tikz v0.4.7 (commit 4d7ae5c4fd0932fb51051d86111bc23ed23e4580) running on MATLAB 8.0.
% Copyright (c) 2008--2014, Nico Schlömer <nico.schloemer@gmail.com>
% All rights reserved.
% Minimal pgfplots version: 1.3
% 
\begin{tikzpicture}

\begin{axis}[%
width=\figurewidth,
height=\figureheight,
view={-37.5}{30},
scale only axis,
xmin=0,
xmax=1,
xlabel={x},
xmajorgrids,
ymin=0,
ymax=1,
ylabel={y},
ymajorgrids,
zmin=0,
zmax=1,
zlabel={$\text{S}_{\text{max}}\text{(x,y)}$},
zmajorgrids,
axis x line*=bottom,
axis y line*=left,
axis z line*=left
]

\addplot3[%
surf,
shader=faceted,
draw=black,
colormap/jet,
mesh/rows=11]
table[row sep=crcr,header=false] {0	0	0\\
0	0.1	0.1\\
0	0.2	0.2\\
0	0.3	0.3\\
0	0.4	0.4\\
0	0.5	0.5\\
0	0.6	0.6\\
0	0.7	0.7\\
0	0.8	0.8\\
0	0.9	0.9\\
0	1	1\\
0.1	0	0.1\\
0.1	0.1	0.1\\
0.1	0.2	0.2\\
0.1	0.3	0.3\\
0.1	0.4	0.4\\
0.1	0.5	0.5\\
0.1	0.6	0.6\\
0.1	0.7	0.7\\
0.1	0.8	0.8\\
0.1	0.9	0.9\\
0.1	1	1\\
0.2	0	0.2\\
0.2	0.1	0.2\\
0.2	0.2	0.2\\
0.2	0.3	0.3\\
0.2	0.4	0.4\\
0.2	0.5	0.5\\
0.2	0.6	0.6\\
0.2	0.7	0.7\\
0.2	0.8	0.8\\
0.2	0.9	0.9\\
0.2	1	1\\
0.3	0	0.3\\
0.3	0.1	0.3\\
0.3	0.2	0.3\\
0.3	0.3	0.3\\
0.3	0.4	0.4\\
0.3	0.5	0.5\\
0.3	0.6	0.6\\
0.3	0.7	0.7\\
0.3	0.8	0.8\\
0.3	0.9	0.9\\
0.3	1	1\\
0.4	0	0.4\\
0.4	0.1	0.4\\
0.4	0.2	0.4\\
0.4	0.3	0.4\\
0.4	0.4	0.4\\
0.4	0.5	0.5\\
0.4	0.6	0.6\\
0.4	0.7	0.7\\
0.4	0.8	0.8\\
0.4	0.9	0.9\\
0.4	1	1\\
0.5	0	0.5\\
0.5	0.1	0.5\\
0.5	0.2	0.5\\
0.5	0.3	0.5\\
0.5	0.4	0.5\\
0.5	0.5	0.5\\
0.5	0.6	0.6\\
0.5	0.7	0.7\\
0.5	0.8	0.8\\
0.5	0.9	0.9\\
0.5	1	1\\
0.6	0	0.6\\
0.6	0.1	0.6\\
0.6	0.2	0.6\\
0.6	0.3	0.6\\
0.6	0.4	0.6\\
0.6	0.5	0.6\\
0.6	0.6	0.6\\
0.6	0.7	0.7\\
0.6	0.8	0.8\\
0.6	0.9	0.9\\
0.6	1	1\\
0.7	0	0.7\\
0.7	0.1	0.7\\
0.7	0.2	0.7\\
0.7	0.3	0.7\\
0.7	0.4	0.7\\
0.7	0.5	0.7\\
0.7	0.6	0.7\\
0.7	0.7	0.7\\
0.7	0.8	0.8\\
0.7	0.9	0.9\\
0.7	1	1\\
0.8	0	0.8\\
0.8	0.1	0.8\\
0.8	0.2	0.8\\
0.8	0.3	0.8\\
0.8	0.4	0.8\\
0.8	0.5	0.8\\
0.8	0.6	0.8\\
0.8	0.7	0.8\\
0.8	0.8	0.8\\
0.8	0.9	0.9\\
0.8	1	1\\
0.9	0	0.9\\
0.9	0.1	0.9\\
0.9	0.2	0.9\\
0.9	0.3	0.9\\
0.9	0.4	0.9\\
0.9	0.5	0.9\\
0.9	0.6	0.9\\
0.9	0.7	0.9\\
0.9	0.8	0.9\\
0.9	0.9	0.9\\
0.9	1	1\\
1	0	1\\
1	0.1	1\\
1	0.2	1\\
1	0.3	1\\
1	0.4	1\\
1	0.5	1\\
1	0.6	1\\
1	0.7	1\\
1	0.8	1\\
1	0.9	1\\
1	1	1\\
};
\end{axis}
\end{tikzpicture}%
		\caption{Máximo}
		\label{fig:t-conorms-max}
	\end{subfigure}
	\qquad
	\begin{subfigure}[b]{0.45\textwidth}
		\setlength\figureheight{4.5cm}
		\setlength\figurewidth{6cm}
		% This file was created by matlab2tikz v0.4.7 (commit 4d7ae5c4fd0932fb51051d86111bc23ed23e4580) running on MATLAB 8.0.
% Copyright (c) 2008--2014, Nico Schlömer <nico.schloemer@gmail.com>
% All rights reserved.
% Minimal pgfplots version: 1.3
% 
\begin{tikzpicture}

\begin{axis}[%
width=\figurewidth,
height=\figureheight,
view={-37.5}{30},
scale only axis,
xmin=0,
xmax=1,
xlabel={x},
xmajorgrids,
ymin=0,
ymax=1,
ylabel={y},
ymajorgrids,
zmin=0,
zmax=1,
zlabel={$\text{S}_{\text{P}}\text{(x,y)}$},
zmajorgrids,
axis x line*=bottom,
axis y line*=left,
axis z line*=left
]

\addplot3[%
surf,
shader=faceted,
draw=black,
colormap/jet,
mesh/rows=11]
table[row sep=crcr,header=false] {0	0	0\\
0	0.1	0.1\\
0	0.2	0.2\\
0	0.3	0.3\\
0	0.4	0.4\\
0	0.5	0.5\\
0	0.6	0.6\\
0	0.7	0.7\\
0	0.8	0.8\\
0	0.9	0.9\\
0	1	1\\
0.1	0	0.1\\
0.1	0.1	0.19\\
0.1	0.2	0.28\\
0.1	0.3	0.37\\
0.1	0.4	0.46\\
0.1	0.5	0.55\\
0.1	0.6	0.64\\
0.1	0.7	0.73\\
0.1	0.8	0.82\\
0.1	0.9	0.91\\
0.1	1	1\\
0.2	0	0.2\\
0.2	0.1	0.28\\
0.2	0.2	0.36\\
0.2	0.3	0.44\\
0.2	0.4	0.52\\
0.2	0.5	0.6\\
0.2	0.6	0.68\\
0.2	0.7	0.76\\
0.2	0.8	0.84\\
0.2	0.9	0.92\\
0.2	1	1\\
0.3	0	0.3\\
0.3	0.1	0.37\\
0.3	0.2	0.44\\
0.3	0.3	0.51\\
0.3	0.4	0.58\\
0.3	0.5	0.65\\
0.3	0.6	0.72\\
0.3	0.7	0.79\\
0.3	0.8	0.86\\
0.3	0.9	0.93\\
0.3	1	1\\
0.4	0	0.4\\
0.4	0.1	0.46\\
0.4	0.2	0.52\\
0.4	0.3	0.58\\
0.4	0.4	0.64\\
0.4	0.5	0.7\\
0.4	0.6	0.76\\
0.4	0.7	0.82\\
0.4	0.8	0.88\\
0.4	0.9	0.94\\
0.4	1	1\\
0.5	0	0.5\\
0.5	0.1	0.55\\
0.5	0.2	0.6\\
0.5	0.3	0.65\\
0.5	0.4	0.7\\
0.5	0.5	0.75\\
0.5	0.6	0.8\\
0.5	0.7	0.85\\
0.5	0.8	0.9\\
0.5	0.9	0.95\\
0.5	1	1\\
0.6	0	0.6\\
0.6	0.1	0.64\\
0.6	0.2	0.68\\
0.6	0.3	0.72\\
0.6	0.4	0.76\\
0.6	0.5	0.8\\
0.6	0.6	0.84\\
0.6	0.7	0.88\\
0.6	0.8	0.92\\
0.6	0.9	0.96\\
0.6	1	1\\
0.7	0	0.7\\
0.7	0.1	0.73\\
0.7	0.2	0.76\\
0.7	0.3	0.79\\
0.7	0.4	0.82\\
0.7	0.5	0.85\\
0.7	0.6	0.88\\
0.7	0.7	0.91\\
0.7	0.8	0.94\\
0.7	0.9	0.97\\
0.7	1	1\\
0.8	0	0.8\\
0.8	0.1	0.82\\
0.8	0.2	0.84\\
0.8	0.3	0.86\\
0.8	0.4	0.88\\
0.8	0.5	0.9\\
0.8	0.6	0.92\\
0.8	0.7	0.94\\
0.8	0.8	0.96\\
0.8	0.9	0.98\\
0.8	1	1\\
0.9	0	0.9\\
0.9	0.1	0.91\\
0.9	0.2	0.92\\
0.9	0.3	0.93\\
0.9	0.4	0.94\\
0.9	0.5	0.95\\
0.9	0.6	0.96\\
0.9	0.7	0.97\\
0.9	0.8	0.98\\
0.9	0.9	0.99\\
0.9	1	1\\
1	0	1\\
1	0.1	1\\
1	0.2	1\\
1	0.3	1\\
1	0.4	1\\
1	0.5	1\\
1	0.6	1\\
1	0.7	1\\
1	0.8	1\\
1	0.9	1\\
1	1	1\\
};
\end{axis}
\end{tikzpicture}%
		\caption{Suma probabilística}
		\label{fig:t-conorms-probabilistic-sum}
	\end{subfigure}
	
	\vspace{1 cm}
	\begin{subfigure}[b]{0.45\textwidth}
		\setlength\figureheight{4.5cm}
		\setlength\figurewidth{6cm}
		% This file was created by matlab2tikz v0.4.7 (commit 4d7ae5c4fd0932fb51051d86111bc23ed23e4580) running on MATLAB 8.0.
% Copyright (c) 2008--2014, Nico Schlömer <nico.schloemer@gmail.com>
% All rights reserved.
% Minimal pgfplots version: 1.3
% 
\begin{tikzpicture}

\begin{axis}[%
width=\figurewidth,
height=\figureheight,
view={-37.5}{30},
scale only axis,
xmin=0,
xmax=1,
xlabel={x},
xmajorgrids,
ymin=0,
ymax=1,
ylabel={y},
ymajorgrids,
zmin=0,
zmax=1,
zlabel={$\text{S}_{\text{B}}\text{(x,y)}$},
zmajorgrids,
axis x line*=bottom,
axis y line*=left,
axis z line*=left
]

\addplot3[%
surf,
shader=faceted,
draw=black,
colormap/jet,
mesh/rows=11]
table[row sep=crcr,header=false] {0	0	0\\
0	0.1	0.1\\
0	0.2	0.2\\
0	0.3	0.3\\
0	0.4	0.4\\
0	0.5	0.5\\
0	0.6	0.6\\
0	0.7	0.7\\
0	0.8	0.8\\
0	0.9	0.9\\
0	1	1\\
0.1	0	0.1\\
0.1	0.1	0.2\\
0.1	0.2	0.3\\
0.1	0.3	0.4\\
0.1	0.4	0.5\\
0.1	0.5	0.6\\
0.1	0.6	0.7\\
0.1	0.7	0.8\\
0.1	0.8	0.9\\
0.1	0.9	1\\
0.1	1	1\\
0.2	0	0.2\\
0.2	0.1	0.3\\
0.2	0.2	0.4\\
0.2	0.3	0.5\\
0.2	0.4	0.6\\
0.2	0.5	0.7\\
0.2	0.6	0.8\\
0.2	0.7	0.9\\
0.2	0.8	1\\
0.2	0.9	1\\
0.2	1	1\\
0.3	0	0.3\\
0.3	0.1	0.4\\
0.3	0.2	0.5\\
0.3	0.3	0.6\\
0.3	0.4	0.7\\
0.3	0.5	0.8\\
0.3	0.6	0.9\\
0.3	0.7	1\\
0.3	0.8	1\\
0.3	0.9	1\\
0.3	1	1\\
0.4	0	0.4\\
0.4	0.1	0.5\\
0.4	0.2	0.6\\
0.4	0.3	0.7\\
0.4	0.4	0.8\\
0.4	0.5	0.9\\
0.4	0.6	1\\
0.4	0.7	1\\
0.4	0.8	1\\
0.4	0.9	1\\
0.4	1	1\\
0.5	0	0.5\\
0.5	0.1	0.6\\
0.5	0.2	0.7\\
0.5	0.3	0.8\\
0.5	0.4	0.9\\
0.5	0.5	1\\
0.5	0.6	1\\
0.5	0.7	1\\
0.5	0.8	1\\
0.5	0.9	1\\
0.5	1	1\\
0.6	0	0.6\\
0.6	0.1	0.7\\
0.6	0.2	0.8\\
0.6	0.3	0.9\\
0.6	0.4	1\\
0.6	0.5	1\\
0.6	0.6	1\\
0.6	0.7	1\\
0.6	0.8	1\\
0.6	0.9	1\\
0.6	1	1\\
0.7	0	0.7\\
0.7	0.1	0.8\\
0.7	0.2	0.9\\
0.7	0.3	1\\
0.7	0.4	1\\
0.7	0.5	1\\
0.7	0.6	1\\
0.7	0.7	1\\
0.7	0.8	1\\
0.7	0.9	1\\
0.7	1	1\\
0.8	0	0.8\\
0.8	0.1	0.9\\
0.8	0.2	1\\
0.8	0.3	1\\
0.8	0.4	1\\
0.8	0.5	1\\
0.8	0.6	1\\
0.8	0.7	1\\
0.8	0.8	1\\
0.8	0.9	1\\
0.8	1	1\\
0.9	0	0.9\\
0.9	0.1	1\\
0.9	0.2	1\\
0.9	0.3	1\\
0.9	0.4	1\\
0.9	0.5	1\\
0.9	0.6	1\\
0.9	0.7	1\\
0.9	0.8	1\\
0.9	0.9	1\\
0.9	1	1\\
1	0	1\\
1	0.1	1\\
1	0.2	1\\
1	0.3	1\\
1	0.4	1\\
1	0.5	1\\
1	0.6	1\\
1	0.7	1\\
1	0.8	1\\
1	0.9	1\\
1	1	1\\
};
\end{axis}
\end{tikzpicture}%
		\caption{Suma acotada}
		\label{fig:t-conorms-bounded-sum}
	\end{subfigure}
	\qquad
	\begin{subfigure}[b]{0.45\textwidth}
			\setlength\figureheight{4.5cm}
			\setlength\figurewidth{6cm}
			% This file was created by matlab2tikz v0.4.7 (commit 730e9c094bba45cf576cc1ba69976f02cfab8e0c) running on MATLAB 8.0.
% Copyright (c) 2008--2014, Nico Schlömer <nico.schloemer@gmail.com>
% All rights reserved.
% Minimal pgfplots version: 1.3
% 
\begin{tikzpicture}

\begin{axis}[%
width=\figurewidth,
height=\figureheight,
view={-37.5}{30},
scale only axis,
xmin=0,
xmax=1,
xlabel={x},
xmajorgrids,
ymin=0,
ymax=1,
ylabel={y},
ymajorgrids,
zmin=0,
zmax=1,
zlabel={$\text{S}_{\text{D}}\text{(x,y)}$},
zmajorgrids,
axis x line*=bottom,
axis y line*=left,
axis z line*=left
]

\addplot3[%
surf,
shader=faceted,
draw=black,
colormap/jet,
mesh/rows=11]
table[row sep=crcr,header=false] {0	0	0\\
0	0.1	0.1\\
0	0.2	0.2\\
0	0.3	0.3\\
0	0.4	0.4\\
0	0.5	0.5\\
0	0.6	0.6\\
0	0.7	0.7\\
0	0.8	0.8\\
0	0.9	0.9\\
0	1	1\\
0.1	0	0.1\\
0.1	0.1	1\\
0.1	0.2	1\\
0.1	0.3	1\\
0.1	0.4	1\\
0.1	0.5	1\\
0.1	0.6	1\\
0.1	0.7	1\\
0.1	0.8	1\\
0.1	0.9	1\\
0.1	1	1\\
0.2	0	0.2\\
0.2	0.1	1\\
0.2	0.2	1\\
0.2	0.3	1\\
0.2	0.4	1\\
0.2	0.5	1\\
0.2	0.6	1\\
0.2	0.7	1\\
0.2	0.8	1\\
0.2	0.9	1\\
0.2	1	1\\
0.3	0	0.3\\
0.3	0.1	1\\
0.3	0.2	1\\
0.3	0.3	1\\
0.3	0.4	1\\
0.3	0.5	1\\
0.3	0.6	1\\
0.3	0.7	1\\
0.3	0.8	1\\
0.3	0.9	1\\
0.3	1	1\\
0.4	0	0.4\\
0.4	0.1	1\\
0.4	0.2	1\\
0.4	0.3	1\\
0.4	0.4	1\\
0.4	0.5	1\\
0.4	0.6	1\\
0.4	0.7	1\\
0.4	0.8	1\\
0.4	0.9	1\\
0.4	1	1\\
0.5	0	0.5\\
0.5	0.1	1\\
0.5	0.2	1\\
0.5	0.3	1\\
0.5	0.4	1\\
0.5	0.5	1\\
0.5	0.6	1\\
0.5	0.7	1\\
0.5	0.8	1\\
0.5	0.9	1\\
0.5	1	1\\
0.6	0	0.6\\
0.6	0.1	1\\
0.6	0.2	1\\
0.6	0.3	1\\
0.6	0.4	1\\
0.6	0.5	1\\
0.6	0.6	1\\
0.6	0.7	1\\
0.6	0.8	1\\
0.6	0.9	1\\
0.6	1	1\\
0.7	0	0.7\\
0.7	0.1	1\\
0.7	0.2	1\\
0.7	0.3	1\\
0.7	0.4	1\\
0.7	0.5	1\\
0.7	0.6	1\\
0.7	0.7	1\\
0.7	0.8	1\\
0.7	0.9	1\\
0.7	1	1\\
0.8	0	0.8\\
0.8	0.1	1\\
0.8	0.2	1\\
0.8	0.3	1\\
0.8	0.4	1\\
0.8	0.5	1\\
0.8	0.6	1\\
0.8	0.7	1\\
0.8	0.8	1\\
0.8	0.9	1\\
0.8	1	1\\
0.9	0	0.9\\
0.9	0.1	1\\
0.9	0.2	1\\
0.9	0.3	1\\
0.9	0.4	1\\
0.9	0.5	1\\
0.9	0.6	1\\
0.9	0.7	1\\
0.9	0.8	1\\
0.9	0.9	1\\
0.9	1	1\\
1	0	1\\
1	0.1	1\\
1	0.2	1\\
1	0.3	1\\
1	0.4	1\\
1	0.5	1\\
1	0.6	1\\
1	0.7	1\\
1	0.8	1\\
1	0.9	1\\
1	1	1\\
};
\end{axis}
\end{tikzpicture}%
			\caption{T-conorma drástica}
			\label{fig:t-conorms-drastic}
	\end{subfigure}
	\label{fig:t-conorms}
	\caption{Representación gráfica de las t-conormas \emph{Máximo} (\ref{fig:t-conorms-max}), \emph{Suma probabilística} (\ref{fig:t-conorms-probabilistic-sum}), \emph{Suma acotada} (\ref{fig:t-conorms-bounded-sum})  y la \emph{t-conorma drástica} (\ref{fig:t-conorms-drastic}).}
\end{figure}
Otra clase importante de funciones son los \emph{operadores de agregación} \index{operador de agregación}\cite{calvo2002aggregation, beliakov2007aggregation}.
\begin{definition}
Una función $\emph{M} : [a,b]^{n} \rightarrow [a,b]$ es un operador de agregación si es monótona y no decreciente en cada una de sus componentes y además cumple que $M(a, a, \cdots,a) = a$ y $M(b, b, \cdots,b) = b$.
\end{definition}
\begin{definition}
Una función de agregación \emph{M} se dice que es una media si cumple que:
\begin{equation}\label{def:mean-aggregation-operator}
min(x) = min(x_{1},\cdots,x_{n})\leq \emph{M}(x_{1},\cdots,x_{n}) \leq max(x_{1},\cdots,x_{n}) = max(x)
\end{equation}
\end{definition}
A continuación se incluyen algunos de los operadores de agregación más comunes:
\begin{itemize}
	\item \bfseries Media aritmética: $M(x_{1},x_{2},\cdots,x_{n}) = \frac{1}{n}\sum\limits_{i=1}^{n}x_{i}$
	\item \bfseries Media geométrica: $M(x_{1},x_{2},\cdots,x_{n}) = (\prod\limits_{i=1}^{n}x_{i})^{\frac{1}{n}}$
	\item \bfseries Media ponderada: $M_{w_{1},w_{2},\cdots,w_{n}}(x_{1},x_{2},\cdots,x_{n}) =\sum\limits_{i=1}^{n}(w_{i}\cdot x_{i})$ \normalfont tal que $0 \leq w_{i} \leq 1$ y $\sum\limits_{i=1}^{n}w_{i} = 1$.
	\item \bfseries Mediana: \normalfont Se toma el elemento central del conjunto ordenado de argumentos.
	\item \bfseries Máximo: $M(x_{1},x_{2},\cdots,x_{n}) = max(x_{1},x_{2},\cdots,x_{n})$
	\item \bfseries Mínimo: $M(x_{1},x_{2},\cdots,x_{n}) = min(x_{1},x_{2},\cdots,x_{n})$
\end{itemize}

\section{Funciones de solapamiento}\label{sec:overlap-functions}
En esta sección se presenta el concepto de \emph{función de solapamiento}\index{función de solapamiento} \cite{bustince2009overlap,bustince2010overlapfunctions,bustince2012pairwisecomparisons,jurio2013propertiesoverlap}, que tendrá una gran importancia en el desarrollo de este trabajo, ya que a partir de estas funciones se construyen los índices de solapamiento, que serán utilizados en el método de interpolación presentado.
\subsection{Definición de función de solapamiento y teorema de construcción.}
\begin{definition}
Una función de solapamiento $G_{O} : [0,1]^{2} \rightarrow [0,1]$ cumple:
\begin{enumerate}[label=(G\arabic*),ref=(G\arabic*)]
   \item $G_{O}(x,y) = G_{O}(y,x) \;\; \forall \; x,y \in [0,1]$ \label{G1}
   \item $G_{O}(x,y) = 0$ si y sólo si $x \cdot y = 0$ \label{G2}
   \item $G_{O}(x,y) = 1$ si y sólo si $x \cdot y = 1$ \label{G3}
   \item $G_{O}$ es creciente \label{G4}
   \item $G_{O}$ es continua \label{G5}
\end{enumerate}
\end{definition}
Las funciones de solapamiento generalizan los operadores de intersección tales como el mínimo o, en general, las t-normas (definidas en \ref{def:tnorma}). Además, las funciones de solapamiento son casos particulares de los operadores de agregación sin divisores por cero o divisores por uno. Algunos ejemplos de funciones de solapamiento son:\\
\\
$G_{O}(x,y) = min(x,y)$\\
$G_{O}(x,y) = \sqrt{xy}$\\
$G_{O}(x,y) = min(x^{k}y,xy^{k}), k \in ]0,1[$\\
\\
En \cite{bustince2009overlap} se presenta un teorema que proporciona tanto una caracterización como un método de construcción para funciones de solapamiento:
\begin{theorem}\label{th:overlap-function-construction}
Una función $G_{O} : [0,1]^{2} \rightarrow [0,1]$ es una función de solapamiento si y sólo si puede ser expresada como:
\begin{equation}
G_{O}(x,y) = \frac{f(x,y)}{f(x,y) + h(x,y)}
\end{equation}
para cualesquiera $f,h: [0,1]^{2} \rightarrow [0,1]$ tales que:
\begin{enumerate}[label=(\arabic*),ref=(\arabic*)]
   \item f y h son simétricas;
   \item f es no decreciente y h es no creciente;
   \item $f(x,y) = 0$ si y sólo si $min(x.y) = 0$;
   \item $h(x,y) = 0$ si y sólo si $min(x.y) = 1$;
   \item f y h son continuas
\end{enumerate}
\end{theorem}
\begin{example}
A continuación se presentan algunos ejemplos de construcción de funciones de solapamiento utilizando el teorema  \ref{th:overlap-function-construction}:
\begin{enumerate}[label=(\arabic*),ref=(\arabic*)]
	\item Si $f(x,y)=min(x,y)$ y $h(x,y) = max(1-x,1-y)$ entonces: 
	\begin{equation}
	G_{O}(x,y) = min(x,y)
	\end{equation} es una función de solapamiento.
	\item Si $f(x,y)=\sqrt{x \times y}$ y $h(x,y) = max(1-x,1-y)$ entonces: 
	\begin{equation}G_{O}(x,y) = \frac{\sqrt{x \times y}}{\sqrt{x \times y} + max(1-x,1-y)}
	\end{equation} es una función de solapamiento.
	\item Si $f(x,y)=\sqrt{x \times y}$ y $h(x,y) = 1 - x \times y $ entonces: 
	\begin{equation}G_{O}(x,y) = \frac{\sqrt{x \times y}}{\sqrt{x \times y} + 1 - x \times y}
	\end{equation} es una función de solapamiento.
\end{enumerate}
\end{example}
\begin{figure}[t]
	\centering
	\begin{subfigure}[b]{0.45\textwidth}
		\setlength\figureheight{4.5cm}
		\setlength\figurewidth{6cm}
		% This file was created by matlab2tikz v0.4.7 (commit 730e9c094bba45cf576cc1ba69976f02cfab8e0c) running on MATLAB 8.0.
% Copyright (c) 2008--2014, Nico Schlömer <nico.schloemer@gmail.com>
% All rights reserved.
% Minimal pgfplots version: 1.3
% 
\begin{tikzpicture}

\begin{axis}[%
width=\figurewidth,
height=\figureheight,
view={-37.5}{30},
scale only axis,
xmin=0,
xmax=1,
xlabel={x},
xmajorgrids,
ymin=0,
ymax=1,
ylabel={y},
ymajorgrids,
zmin=0,
zmax=1,
zlabel={$\text{G}_{\text{O}}\text{(x,y)}$},
zmajorgrids,
axis x line*=bottom,
axis y line*=left,
axis z line*=left
]

\addplot3[%
surf,
shader=faceted,
draw=black,
colormap/jet,
mesh/rows=11]
table[row sep=crcr,header=false] {0	0	0\\
0	0.1	0\\
0	0.2	0\\
0	0.3	0\\
0	0.4	0\\
0	0.5	0\\
0	0.6	0\\
0	0.7	0\\
0	0.8	0\\
0	0.9	0\\
0	1	0\\
0.1	0	0\\
0.1	0.1	0.1\\
0.1	0.2	0.1\\
0.1	0.3	0.1\\
0.1	0.4	0.1\\
0.1	0.5	0.1\\
0.1	0.6	0.1\\
0.1	0.7	0.1\\
0.1	0.8	0.1\\
0.1	0.9	0.1\\
0.1	1	0.1\\
0.2	0	0\\
0.2	0.1	0.1\\
0.2	0.2	0.2\\
0.2	0.3	0.2\\
0.2	0.4	0.2\\
0.2	0.5	0.2\\
0.2	0.6	0.2\\
0.2	0.7	0.2\\
0.2	0.8	0.2\\
0.2	0.9	0.2\\
0.2	1	0.2\\
0.3	0	0\\
0.3	0.1	0.1\\
0.3	0.2	0.2\\
0.3	0.3	0.3\\
0.3	0.4	0.3\\
0.3	0.5	0.3\\
0.3	0.6	0.3\\
0.3	0.7	0.3\\
0.3	0.8	0.3\\
0.3	0.9	0.3\\
0.3	1	0.3\\
0.4	0	0\\
0.4	0.1	0.1\\
0.4	0.2	0.2\\
0.4	0.3	0.3\\
0.4	0.4	0.4\\
0.4	0.5	0.4\\
0.4	0.6	0.4\\
0.4	0.7	0.4\\
0.4	0.8	0.4\\
0.4	0.9	0.4\\
0.4	1	0.4\\
0.5	0	0\\
0.5	0.1	0.1\\
0.5	0.2	0.2\\
0.5	0.3	0.3\\
0.5	0.4	0.4\\
0.5	0.5	0.5\\
0.5	0.6	0.5\\
0.5	0.7	0.5\\
0.5	0.8	0.5\\
0.5	0.9	0.5\\
0.5	1	0.5\\
0.6	0	0\\
0.6	0.1	0.1\\
0.6	0.2	0.2\\
0.6	0.3	0.3\\
0.6	0.4	0.4\\
0.6	0.5	0.5\\
0.6	0.6	0.6\\
0.6	0.7	0.6\\
0.6	0.8	0.6\\
0.6	0.9	0.6\\
0.6	1	0.6\\
0.7	0	0\\
0.7	0.1	0.1\\
0.7	0.2	0.2\\
0.7	0.3	0.3\\
0.7	0.4	0.4\\
0.7	0.5	0.5\\
0.7	0.6	0.6\\
0.7	0.7	0.7\\
0.7	0.8	0.7\\
0.7	0.9	0.7\\
0.7	1	0.7\\
0.8	0	0\\
0.8	0.1	0.1\\
0.8	0.2	0.2\\
0.8	0.3	0.3\\
0.8	0.4	0.4\\
0.8	0.5	0.5\\
0.8	0.6	0.6\\
0.8	0.7	0.7\\
0.8	0.8	0.8\\
0.8	0.9	0.8\\
0.8	1	0.8\\
0.9	0	0\\
0.9	0.1	0.1\\
0.9	0.2	0.2\\
0.9	0.3	0.3\\
0.9	0.4	0.4\\
0.9	0.5	0.5\\
0.9	0.6	0.6\\
0.9	0.7	0.7\\
0.9	0.8	0.8\\
0.9	0.9	0.9\\
0.9	1	0.9\\
1	0	0\\
1	0.1	0.1\\
1	0.2	0.2\\
1	0.3	0.3\\
1	0.4	0.4\\
1	0.5	0.5\\
1	0.6	0.6\\
1	0.7	0.7\\
1	0.8	0.8\\
1	0.9	0.9\\
1	1	1\\
};
\end{axis}
\end{tikzpicture}%
		\caption{$G_{O}(x,y) = min(x,y)$}
		\label{fig:overlap-functions-min}
	\end{subfigure}
	\qquad
	\begin{subfigure}[b]{0.45\textwidth}
		\setlength\figureheight{4.5cm}
		\setlength\figurewidth{6cm}
		% This file was created by matlab2tikz v0.4.7 (commit 921463bb921f990ffc6a7f597361910de95829d7) running on MATLAB 8.0.
% Copyright (c) 2008--2014, Nico Schlömer <nico.schloemer@gmail.com>
% All rights reserved.
% Minimal pgfplots version: 1.3
% 
\begin{tikzpicture}

\begin{axis}[%
width=\figurewidth,
height=\figureheight,
view={-37.5}{30},
scale only axis,
xmin=0,
xmax=1,
xlabel={x},
xmajorgrids,
ymin=0,
ymax=1,
ylabel={y},
ymajorgrids,
zmin=0,
zmax=1,
zlabel={$\text{G}_{\text{O}}\text{(x,y)}$},
zmajorgrids,
axis x line*=bottom,
axis y line*=left,
axis z line*=left
]

\addplot3[%
surf,
shader=faceted,
draw=black,
colormap/jet,
mesh/rows=11]
table[row sep=crcr,header=false] {0	0	0\\
0	0.1	0\\
0	0.2	0\\
0	0.3	0\\
0	0.4	0\\
0	0.5	0\\
0	0.6	0\\
0	0.7	0\\
0	0.8	0\\
0	0.9	0\\
0	1	0\\
0.1	0	0\\
0.1	0.1	0.1\\
0.1	0.2	0.14142135623731\\
0.1	0.3	0.173205080756888\\
0.1	0.4	0.2\\
0.1	0.5	0.223606797749979\\
0.1	0.6	0.244948974278318\\
0.1	0.7	0.264575131106459\\
0.1	0.8	0.282842712474619\\
0.1	0.9	0.3\\
0.1	1	0.316227766016838\\
0.2	0	0\\
0.2	0.1	0.14142135623731\\
0.2	0.2	0.2\\
0.2	0.3	0.244948974278318\\
0.2	0.4	0.282842712474619\\
0.2	0.5	0.316227766016838\\
0.2	0.6	0.346410161513775\\
0.2	0.7	0.374165738677394\\
0.2	0.8	0.4\\
0.2	0.9	0.424264068711929\\
0.2	1	0.447213595499958\\
0.3	0	0\\
0.3	0.1	0.173205080756888\\
0.3	0.2	0.244948974278318\\
0.3	0.3	0.3\\
0.3	0.4	0.346410161513776\\
0.3	0.5	0.387298334620742\\
0.3	0.6	0.424264068711929\\
0.3	0.7	0.458257569495584\\
0.3	0.8	0.489897948556636\\
0.3	0.9	0.519615242270663\\
0.3	1	0.547722557505166\\
0.4	0	0\\
0.4	0.1	0.2\\
0.4	0.2	0.282842712474619\\
0.4	0.3	0.346410161513776\\
0.4	0.4	0.4\\
0.4	0.5	0.447213595499958\\
0.4	0.6	0.489897948556636\\
0.4	0.7	0.529150262212918\\
0.4	0.8	0.565685424949238\\
0.4	0.9	0.6\\
0.4	1	0.632455532033676\\
0.5	0	0\\
0.5	0.1	0.223606797749979\\
0.5	0.2	0.316227766016838\\
0.5	0.3	0.387298334620742\\
0.5	0.4	0.447213595499958\\
0.5	0.5	0.5\\
0.5	0.6	0.547722557505166\\
0.5	0.7	0.591607978309962\\
0.5	0.8	0.632455532033676\\
0.5	0.9	0.670820393249937\\
0.5	1	0.707106781186548\\
0.6	0	0\\
0.6	0.1	0.244948974278318\\
0.6	0.2	0.346410161513775\\
0.6	0.3	0.424264068711929\\
0.6	0.4	0.489897948556636\\
0.6	0.5	0.547722557505166\\
0.6	0.6	0.6\\
0.6	0.7	0.648074069840786\\
0.6	0.8	0.692820323027551\\
0.6	0.9	0.734846922834953\\
0.6	1	0.774596669241483\\
0.7	0	0\\
0.7	0.1	0.264575131106459\\
0.7	0.2	0.374165738677394\\
0.7	0.3	0.458257569495584\\
0.7	0.4	0.529150262212918\\
0.7	0.5	0.591607978309962\\
0.7	0.6	0.648074069840786\\
0.7	0.7	0.7\\
0.7	0.8	0.748331477354788\\
0.7	0.9	0.793725393319377\\
0.7	1	0.836660026534076\\
0.8	0	0\\
0.8	0.1	0.282842712474619\\
0.8	0.2	0.4\\
0.8	0.3	0.489897948556636\\
0.8	0.4	0.565685424949238\\
0.8	0.5	0.632455532033676\\
0.8	0.6	0.692820323027551\\
0.8	0.7	0.748331477354788\\
0.8	0.8	0.8\\
0.8	0.9	0.848528137423857\\
0.8	1	0.894427190999916\\
0.9	0	0\\
0.9	0.1	0.3\\
0.9	0.2	0.424264068711929\\
0.9	0.3	0.519615242270663\\
0.9	0.4	0.6\\
0.9	0.5	0.670820393249937\\
0.9	0.6	0.734846922834953\\
0.9	0.7	0.793725393319377\\
0.9	0.8	0.848528137423857\\
0.9	0.9	0.9\\
0.9	1	0.948683298050514\\
1	0	0\\
1	0.1	0.316227766016838\\
1	0.2	0.447213595499958\\
1	0.3	0.547722557505166\\
1	0.4	0.632455532033676\\
1	0.5	0.707106781186548\\
1	0.6	0.774596669241483\\
1	0.7	0.836660026534076\\
1	0.8	0.894427190999916\\
1	0.9	0.948683298050514\\
1	1	1\\
};
\end{axis}
\end{tikzpicture}%
		\caption{$G_{O}(x,y) = \sqrt{x \cdot y}$}
		\label{fig:overlap-functions-sqrt}
	\end{subfigure}	
	
	\vspace{1 cm}
	\begin{subfigure}[b]{0.45\textwidth}
		\setlength\figureheight{4.5cm}
		\setlength\figurewidth{6cm}
		% This file was created by matlab2tikz v0.4.7 (commit 730e9c094bba45cf576cc1ba69976f02cfab8e0c) running on MATLAB 8.0.
% Copyright (c) 2008--2014, Nico Schlömer <nico.schloemer@gmail.com>
% All rights reserved.
% Minimal pgfplots version: 1.3
% 
\begin{tikzpicture}

\begin{axis}[%
width=\figurewidth,
height=\figureheight,
view={-37.5}{30},
scale only axis,
xmin=0,
xmax=1,
xlabel={x},
xmajorgrids,
ymin=0,
ymax=1,
ylabel={y},
ymajorgrids,
zmin=0,
zmax=1,
zlabel={$\text{G}_{\text{O}}\text{(x,y)}$},
zmajorgrids,
axis x line*=bottom,
axis y line*=left,
axis z line*=left
]

\addplot3[%
surf,
shader=faceted,
draw=black,
colormap/jet,
mesh/rows=11]
table[row sep=crcr,header=false] {0	0	0\\
0	0.1	0\\
0	0.2	0\\
0	0.3	0\\
0	0.4	0\\
0	0.5	0\\
0	0.6	0\\
0	0.7	0\\
0	0.8	0\\
0	0.9	0\\
0	1	0\\
0.1	0	0\\
0.1	0.1	0.1\\
0.1	0.2	0.13579648178934\\
0.1	0.3	0.161390477796409\\
0.1	0.4	0.181818181818182\\
0.1	0.5	0.19900804996708\\
0.1	0.6	0.213938769133981\\
0.1	0.7	0.227185970264612\\
0.1	0.8	0.239121152365969\\
0.1	0.9	0.25\\
0.1	1	0.260007027345288\\
0.2	0	0\\
0.2	0.1	0.13579648178934\\
0.2	0.2	0.2\\
0.2	0.3	0.234412378314921\\
0.2	0.4	0.261203874963741\\
0.2	0.5	0.283300394099019\\
0.2	0.6	0.302169479251962\\
0.2	0.7	0.318665181883831\\
0.2	0.8	0.333333333333333\\
0.2	0.9	0.346546206455528\\
0.2	1	0.358570173636287\\
0.3	0	0\\
0.3	0.1	0.161390477796409\\
0.3	0.2	0.234412378314921\\
0.3	0.3	0.3\\
0.3	0.4	0.331046251512548\\
0.3	0.5	0.356202453630939\\
0.3	0.6	0.377370477736613\\
0.3	0.7	0.39564392373896\\
0.3	0.8	0.41171425595858\\
0.3	0.9	0.426048498133928\\
0.3	1	0.438977843440086\\
0.4	0	0\\
0.4	0.1	0.181818181818182\\
0.4	0.2	0.261203874963741\\
0.4	0.3	0.331046251512548\\
0.4	0.4	0.4\\
0.4	0.5	0.427050983124842\\
0.4	0.6	0.449489742783178\\
0.4	0.7	0.468626966596886\\
0.4	0.8	0.48528137423857\\
0.4	0.9	0.5\\
0.4	1	0.513167019494862\\
0.5	0	0\\
0.5	0.1	0.19900804996708\\
0.5	0.2	0.283300394099019\\
0.5	0.3	0.356202453630939\\
0.5	0.4	0.427050983124842\\
0.5	0.5	0.5\\
0.5	0.6	0.522774424948339\\
0.5	0.7	0.541960108450192\\
0.5	0.8	0.558481559887747\\
0.5	0.9	0.572949016875158\\
0.5	1	0.585786437626905\\
0.6	0	0\\
0.6	0.1	0.213938769133981\\
0.6	0.2	0.302169479251962\\
0.6	0.3	0.377370477736613\\
0.6	0.4	0.449489742783178\\
0.6	0.5	0.522774424948339\\
0.6	0.6	0.6\\
0.6	0.7	0.618347584860329\\
0.6	0.8	0.633974596215561\\
0.6	0.9	0.647529554910575\\
0.6	1	0.659457573416833\\
0.7	0	0\\
0.7	0.1	0.227185970264612\\
0.7	0.2	0.318665181883831\\
0.7	0.3	0.39564392373896\\
0.7	0.4	0.468626966596886\\
0.7	0.5	0.541960108450192\\
0.7	0.6	0.618347584860329\\
0.7	0.7	0.7\\
0.7	0.8	0.713830971901199\\
0.7	0.9	0.725708114822568\\
0.7	1	0.736068839409471\\
0.8	0	0\\
0.8	0.1	0.239121152365969\\
0.8	0.2	0.333333333333333\\
0.8	0.3	0.41171425595858\\
0.8	0.4	0.48528137423857\\
0.8	0.5	0.558481559887747\\
0.8	0.6	0.633974596215561\\
0.8	0.7	0.713830971901199\\
0.8	0.8	0.8\\
0.8	0.9	0.809256430169454\\
0.8	1	0.817256002368443\\
0.9	0	0\\
0.9	0.1	0.25\\
0.9	0.2	0.346546206455528\\
0.9	0.3	0.426048498133928\\
0.9	0.4	0.5\\
0.9	0.5	0.572949016875158\\
0.9	0.6	0.647529554910575\\
0.9	0.7	0.725708114822568\\
0.9	0.8	0.809256430169454\\
0.9	0.9	0.9\\
0.9	1	0.90464232606174\\
1	0	0\\
1	0.1	0.260007027345288\\
1	0.2	0.358570173636287\\
1	0.3	0.438977843440086\\
1	0.4	0.513167019494862\\
1	0.5	0.585786437626905\\
1	0.6	0.659457573416833\\
1	0.7	0.736068839409471\\
1	0.8	0.817256002368443\\
1	0.9	0.90464232606174\\
1	1	1\\
};
\end{axis}
\end{tikzpicture}%
		\caption{$G_{O}(x,y) = \frac{\sqrt{x \times y}}{\sqrt{x \times y} + max(1-x,1-y)}$}
		\label{fig:overlap-functions-sqrt-max}
	\end{subfigure}
	\qquad
	\begin{subfigure}[b]{0.45\textwidth}
		\setlength\figureheight{4.5cm}
		\setlength\figurewidth{6cm}
		% This file was created by matlab2tikz v0.4.7 (commit 730e9c094bba45cf576cc1ba69976f02cfab8e0c) running on MATLAB 8.0.
% Copyright (c) 2008--2014, Nico Schlömer <nico.schloemer@gmail.com>
% All rights reserved.
% Minimal pgfplots version: 1.3
% 
\begin{tikzpicture}

\begin{axis}[%
width=\figurewidth,
height=\figureheight,
view={-37.5}{30},
scale only axis,
xmin=0,
xmax=1,
xlabel={x},
xmajorgrids,
ymin=0,
ymax=1,
ylabel={y},
ymajorgrids,
zmin=0,
zmax=1,
zlabel={$\text{G}_{\text{O}}\text{(x,y)}$},
zmajorgrids,
axis x line*=bottom,
axis y line*=left,
axis z line*=left
]

\addplot3[%
surf,
shader=faceted,
draw=black,
colormap/jet,
mesh/rows=11]
table[row sep=crcr,header=false] {0	0	0\\
0	0.1	0\\
0	0.2	0\\
0	0.3	0\\
0	0.4	0\\
0	0.5	0\\
0	0.6	0\\
0	0.7	0\\
0	0.8	0\\
0	0.9	0\\
0	1	0\\
0.1	0	0\\
0.1	0.1	0.0917431192660551\\
0.1	0.2	0.12610902712948\\
0.1	0.3	0.151508319611572\\
0.1	0.4	0.172413793103448\\
0.1	0.5	0.190529569340153\\
0.1	0.6	0.206716896335137\\
0.1	0.7	0.221480528279038\\
0.1	0.8	0.235145218523812\\
0.1	0.9	0.247933884297521\\
0.1	1	0.260007027345288\\
0.2	0	0\\
0.2	0.1	0.12610902712948\\
0.2	0.2	0.172413793103448\\
0.2	0.3	0.206716896335137\\
0.2	0.4	0.235145218523812\\
0.2	0.5	0.260007027345288\\
0.2	0.6	0.282458652402388\\
0.2	0.7	0.30317300744256\\
0.2	0.8	0.32258064516129\\
0.2	0.9	0.340975906465844\\
0.2	1	0.358570173636287\\
0.3	0	0\\
0.3	0.1	0.151508319611572\\
0.3	0.2	0.206716896335137\\
0.3	0.3	0.247933884297521\\
0.3	0.4	0.282458652402388\\
0.3	0.5	0.313019361445643\\
0.3	0.6	0.340975906465844\\
0.3	0.7	0.367117797395584\\
0.3	0.8	0.391950358125128\\
0.3	0.9	0.415820185840944\\
0.3	1	0.438977843440086\\
0.4	0	0\\
0.4	0.1	0.172413793103448\\
0.4	0.2	0.235145218523812\\
0.4	0.3	0.282458652402388\\
0.4	0.4	0.32258064516129\\
0.4	0.5	0.358570173636287\\
0.4	0.6	0.391950358125128\\
0.4	0.7	0.423608174468545\\
0.4	0.8	0.454115793296923\\
0.4	0.9	0.483870967741936\\
0.4	1	0.513167019494862\\
0.5	0	0\\
0.5	0.1	0.190529569340153\\
0.5	0.2	0.260007027345288\\
0.5	0.3	0.313019361445643\\
0.5	0.4	0.358570173636287\\
0.5	0.5	0.4\\
0.5	0.6	0.438977843440086\\
0.5	0.7	0.476485322778966\\
0.5	0.8	0.513167019494862\\
0.5	0.9	0.549483279407015\\
0.5	1	0.585786437626905\\
0.6	0	0\\
0.6	0.1	0.206716896335137\\
0.6	0.2	0.282458652402388\\
0.6	0.3	0.340975906465844\\
0.6	0.4	0.391950358125128\\
0.6	0.5	0.438977843440086\\
0.6	0.6	0.483870967741935\\
0.6	0.7	0.527715783401245\\
0.6	0.8	0.571247290198824\\
0.6	0.9	0.615013445480881\\
0.6	1	0.659457573416833\\
0.7	0	0\\
0.7	0.1	0.221480528279038\\
0.7	0.2	0.30317300744256\\
0.7	0.3	0.367117797395584\\
0.7	0.4	0.423608174468545\\
0.7	0.5	0.476485322778966\\
0.7	0.6	0.527715783401245\\
0.7	0.7	0.578512396694215\\
0.7	0.8	0.629732942041193\\
0.7	0.9	0.682055575890956\\
0.7	1	0.736068839409471\\
0.8	0	0\\
0.8	0.1	0.235145218523812\\
0.8	0.2	0.32258064516129\\
0.8	0.3	0.391950358125128\\
0.8	0.4	0.454115793296923\\
0.8	0.5	0.513167019494862\\
0.8	0.6	0.571247290198824\\
0.8	0.7	0.629732942041193\\
0.8	0.8	0.689655172413793\\
0.8	0.9	0.751889216834975\\
0.8	1	0.817256002368443\\
0.9	0	0\\
0.9	0.1	0.247933884297521\\
0.9	0.2	0.340975906465844\\
0.9	0.3	0.415820185840944\\
0.9	0.4	0.483870967741936\\
0.9	0.5	0.549483279407015\\
0.9	0.6	0.615013445480881\\
0.9	0.7	0.682055575890956\\
0.9	0.8	0.751889216834975\\
0.9	0.9	0.825688073394496\\
0.9	1	0.90464232606174\\
1	0	0\\
1	0.1	0.260007027345288\\
1	0.2	0.358570173636287\\
1	0.3	0.438977843440086\\
1	0.4	0.513167019494862\\
1	0.5	0.585786437626905\\
1	0.6	0.659457573416833\\
1	0.7	0.736068839409471\\
1	0.8	0.817256002368443\\
1	0.9	0.90464232606174\\
1	1	1\\
};
\end{axis}
\end{tikzpicture}%
		\caption{$G_{O}(x,y) = \frac{\sqrt{x \times y}}{\sqrt{x \times y} + 1 - x \times y}$}
		\label{fig:overlap-functions-sqrt-2}
	\end{subfigure}
	\caption{Representación gráfica de algunas funciones de solapamiento.}
	\label{fig:overlap-functions-examples}
\end{figure}
En la figura \ref{fig:overlap-functions-examples} se pueden ver las representaciones gráficas de las funciones de solapamiento del ejemplo anterior.

\subsection{Caso particular: t-normas.}
Por la definición \ref{def:tnorma} sabemos que una t-norma es una función conmutativa, asociativa, y creciente $T:[0,1]^2 \rightarrow [0,1]$ tal que $T(x,1) = x$ para todo $x \in [0,1]$. Bajo ciertas condiciones, las t-normas también cumplen la definición de función de solapamiento \cite{bustince2009overlap}. 
\begin{theorem}
Si una t-norma \emph{T} es una función de solapamiento, entonces \emph{T} es uno de los siguientes 3 tipos:
\begin{enumerate}[label=(\arabic*),ref=(\arabic*)]
	\item $T = \min$
	\item T es estricta (continua y estrictamente monótona)
	\item T es la suma ordinal de la familia $\{([a_{m},b_{m}],T_{m})\}$, siendo $T_{m}$ todas las t-normas continuas y arquimedeanas tales que si para algún $m_{0}$ tenemos que $a_{m_{0}} = 0$ entonces $T_{m_{0}}$ es necesariamente una t-norma estricta.
\end{enumerate}
\end{theorem}
\begin{proposition}
Sea $G_{O}$ una función de solapamiento tal que $G_{O}(x,G_{O}(y,z)) = G_{O}(y,G_{O}(x,z))$, entonces \emph{T} es una t-norma \cite{bustince2013overlap}.
\end{proposition}
\section{Índices de solapamiento}\label{sec:overlap-indexes}
\subsection{Definición y propiedades}
En esta sección se presenta la definición de índice de solapamiento y se estudian algunas de sus propiedades más importantes. Los índices de solapamiento tienen una importancia vital en el desarrollo de este trabajo, ya que son la base del método estudiado. 

El concepto de solapamiento aplicado a conjuntos difusos fue introducido por Zadeh en 1978 \cite{zadeh1978}.

\begin{definition}
Sea $A,B \in FS(U)$. La consistencia entre A y B se define como:
\begin{equation}\label{eq:zadeh-consistency}
O_{Z}(A,B) = sup_{i=1}^{n}(min(A(u_{i}),B(u_{i})))
\end{equation}
\end{definition}

Para conjuntos de referencia finitos, el supremo es en realidad el máximo elemento del conjunto.

En 1982 Dubois y Prade \cite{dubois2000} presentan la siguiente definición para el índice de solapamiento:

\begin{definition}
Un índice de solapamiento es una función $O : FS(U) \times FS(U) \rightarrow [0,1]$ tal que:
\begin{enumerate}[label=(O\arabic*),ref=(O\arabic*)]
\item $O(A,B) = 0$ si y sólo si A y B son completamente disjuntos. \label{DP01}
\item $O(A,B) = 1$, si ($A(u_{i}) = 0$ o $B(u_{i}) = 0$) o ($A(u_{i}) = 1$ o $B(u_{i}) = 1$) \label{DPO2}
\item $O(A,B) = O(B,A)$ \label{DPO3}
\item Si $B \leq C$, entonces $O(A,B) \leq O(A,C)$ \label{DPO4}
\end{enumerate}
\label{def:dubois-overlap-index}
\end{definition}
La condición \ref{DPO2} en esta definición presenta la ventaja de que, si \emph{A} no es difuso, entonces $O(A,A) = 1$.

Por esta razón Dubois, Ostasiewicz y Prade imponen en \cite{dubois2000} las siguientes condiciones:
\begin{enumerate}[label=(\arabic*),ref=(\arabic*)]
\item Para todos los conjuntos difusos tales que $A(u_{i}) \le 1$ y $A(u_{i}) \le 1$ para cualquier $u_{i} \in U$, \ref{DPO2} debe ignorarse.
\item El índice ROC (\cite{dubois2000}) no satisface \ref{DPO2}.
\end{enumerate}
Debido a estas consideraciones, normalmente sólo se imponen las condiciones \ref{DP01}, \ref{DPO3}, \ref{DPO4} de la definición \ref{def:dubois-overlap-index} a los índices de solapamiento.
En \cite{bustince2013overlap} se propone la siguiente definición de índice de solapamiento:
\begin{definition}\label{def:bustince-overlap-index}
Un índice de solapamiento es una función $O : FS(U) \times FS(U) \rightarrow [0,1]$ tal que:
\begin{enumerate}[label=(O\arabic*),ref=(O\arabic*)]
\item $O(A,B) = 0$ si y sólo si A y B tienen soportes disjuntos, es decir, $A(u_{i})B(u_{i}) = 0$ para todo $u_{i} \in U$\label{BO1}
\item $O(A,B) = O(B,A)$\label{BO2}
\item Si $B \leq C$, entonces $O(A,B) \leq O(A,C)$\label{BO3}
\end{enumerate}
Un índice de solapamiento normal es un índice \emph{O} tal que:
\begin{enumerate}[label=(O4),ref=(O4)]
\item Si existe un $u_{i} \in U$ tal que $A(u_{i}) = B(u_{i}) = 1$, entonces $O(A,B) = 1$\label{BO4}
\end{enumerate}
\end{definition}
\subsection{Construcción de índices de solapamiento}\label{sec:overlap-index-construction}
En esta sección se presenta el método de construcción de índices de solapamiento propuesto por \cite{bustince2013overlap} que se basa en operadores de agregación (sección \ref{sec:t-norms-aggregation-operator}) y funciones de solapamiento (sección \ref{sec:overlap-functions}).
\begin{theorem}
Sea $M : [0,1]^{2} \rightarrow [0,1]$ una función de agregación tal que $M(x_{1},\cdots,x_{n}) = 0$ si y sólo si $x_{1} = \cdots = x_{n} = 0$. Sea $G_{O} : [0,1]^{2} \rightarrow [0,1]$ una función de solapamiento. Entonces, la función $O: F(U) \times F(U) \rightarrow [0,1]$ definida como:
\begin{equation}\label{eq:construction-overlap-index}
O(A,B) = M(G_{O}(A(u_{1}),B(u_{1})),\cdots,G_{O}(A(u_{n}),B(u_{n})))
\end{equation}
es un índice de solapamiento en el sentido de la definición \ref{def:bustince-overlap-index}. Recíprocamente, si $G_{O}$ es una función de solapamiento y $M:[0,1]^{n} \rightarrow [0,1]$ es un operador de agregación tal que O definido por la ecuación \ref{eq:construction-overlap-index} es un índice de solapamiento, entonces $M(x_{1},\cdots,x_{n}) = 0$ si y sólo si $x_{1} = \cdots = x_{n} = 0$.
\end{theorem}


% !TeX spellcheck = es_ES
\chapter{Lógica difusa}
\label{cha:logica-difusa}
En esta sección se introducen las bases de la lógica difusa y los mecanismos de inferencia difusa. Además se describen algunos de los métodos clásicos de inferencia difusa basada en reglas.

\section{¿Qué es la lógica difusa?}
La lógica difusa fue introducida por Lofti A. Zadeh en 1965. Desde entonces se ha aplicado en multitud de escenarios tales como control de procesos industriales, medicina, electrónica y otros tipos de sistemas expertos. En general, la mayoría de aplicaciones de la lógica difusa son el área del control.

La motivación para el desarrollo de una lógica difusa fue expresada por Zadeh (1984, \cite{Zadeh1984}) de la siguiente manera:
\begin{quote}
\emph{``La habilidad de la mente humana para razonar en términos difusos es realmente una gran ventaja. Aunque una gran cantidad de información es captada por los sentidos en una situación determinada, de alguna manera la mente humana es capaz de descartar la mayoría de esa información y concentrarse sólo en aquello que es relevante.''}
\end{quote}
El objetivo principal de la lógica difusa es intentar dotar a las máquinas de un sistema de razonamiento aproximado similar al de los humanos, que pueda lidiar con imprecisiones y términos inexactos. Para ello, la lógica difusa se basa en la utilización de variables lingüísticas expresadas en lenguaje natural, que forman la base del razonamiento aproximado. 

Por medio del uso de variables lingüísticas se pueden modelar conceptos como \emph{``muy alto''} o \emph{``bastante caliente''} y establecer relaciones entre ellos, expresadas de forma matemática y algorítmica. Los sistemas de lógica difusa tratan la imprecisión de las entradas mediante el uso de variables lingüísticas (expresadas como conjuntos difusos), que transforman dichos valores de entrada en números difusos.

\section{Razonamiento aproximado}
El razonamiento aproximado utiliza conjuntos difusos y lógica difusa para modelar la forma de discurrir y razonar del ser humano \cite{bojadziev1995}. El razonamiento aproximado no posee la precisión y exactitud de la lógica clásica, pero es más efectivo a la hora de lidiar con imprecisiones, conceptos vagos o sistemas complejos y/o pobremente definidos.

En este proyecto se va a implementar una forma de razonamiento aproximado que utiliza un mecanismo de inferencia difusa basado en el modus ponens generalizado, que es una extensión del modus ponens clásico.

\subsection{Modus ponens clásico}
El modus ponens clásico es una proposición compuesta muy conocida en la lógica clásica. Tiene la forma:

\begin{equation}
((p \wedge (p \rightarrow q)) \rightarrow q)
\end{equation}

Para inferir el valor de verdad de \emph{q} a partir del de \emph{p} utilizando el modus ponens clásico, se utiliza una \emph{regla} de inferencia, expresada de forma simbólica utilizando el \emph{silogismo}:

\begin{equation}
\begin{array}{rl}
\text{Premisa 1:} & p \\
\text{Premisa 2:} & p \to q \\
\cline{2-2}
\text{Conclusión:} & q
\end{array}
\label{eq:inference-rule}
\end{equation}

que puede expresarse como: si la proposición \emph{p} es verdadera (premisa 1) y la proposición $p \to q$ es verdadera (premisa 2), entonces \emph{q} es verdadera (conclusión). Esto permite obtener el valor de verdad de \emph{q} a partir del de \emph{p} utilizando una regla de inferencia. Es decir, por medio de esta regla de inferencia sabemos que, si \emph{p} es verdadero, entonces \emph{q} es también verdadero.

La regla de inferencia \ref{eq:inference-rule} puede expresarse de manera más detallada como:

\begin{equation}
\begin{array}{rl}
\text{Premisa 1:} & x \text{ es } A \\
\text{Premisa 2:} & \text{IF } x \text{ es } A \text{ THEN } y \text{ es } B \\
\cline{2-2}
\text{Conclusión:} & y \text{ es } B
\end{array}
\label{eq:inference-rule-detail}
\end{equation}

donde $p = x \text{ es } A$ y $q = y \text{ es } B$.

\begin{example}
Un ejemplo de la regla de inferencia \ref{eq:inference-rule-detail} es la regla de divisibilidad por 3: \emph{``Ùn número entero es divisible por 3 si la suma de sus dígitos es múltiplo de 3''}. Puede expresarse mediante el siguiente silogismo:

\begin{equation}
\normalfont
\begin{array}{rl}
\text{Premisa 1:} & \text{\emph{La suma de los dígitos de un número es múltiplo de 3}} \\
\text{Premisa 2:} & \text{IF \emph{La suma de los dígitos de un número es múltiplo de 3} THEN \emph{El número es divisible por 3}} \\
\cline{2-2}
\text{Conclusión:} & \text{\emph{El número es divisible por 3}}
\end{array}
\label{eq:inference-rule-example}
\end{equation}

En este ejemplo se utilizan conceptos exactos (lógica clásica). Un número pertenece o no pertence al conjunto de números cuya suma de cifras es múltiplo de 3. Por lo tanto la proposición \emph{``La suma de los dígitos de un número es múltiplo de 3''} sólo puede tomar los valores \emph{Verdadero} (1) o \emph{Falso} (0). 
\end{example}

\subsection{Modus ponens generalizado}

El \emph{modus ponens generalizado} es, como su nombre índica, una generalización del modus ponens clásico (definido en la sección anterior) utilizado en la lógica difusa. Tiene la forma:

\begin{equation}
\begin{array}{rl}
\text{Premisa 1:} & p' \\
\text{Premisa 2:} & p \to q \\
\cline{2-2}
\text{Conclusión:} & q'
\end{array}
\label{eq:fuzzy-inference-rule}
\end{equation}

o lo que es lo mismo:

\begin{equation}
\begin{array}{rl}
\text{Premisa 1:} & x \text{ es } A' \\
\text{Premisa 2:} & \text{IF } x \text{ es } A \text{ THEN } y \text{ es } B \\
\cline{2-2}
\text{Conclusión:} & y \text{ es } B'
\end{array}
\label{eq:fuzzy-inference-rule-detail}
\end{equation}

En este caso las proposiciones $p'=x \text{ es } A'$, $p=x \text{ es } A$, $q=y \text{ es } B$, $q'=y \text{ es } B'$ vienen caracterizadas por los conjuntos difusos $A'$,$A$,$B$ y $B'$,que representan conceptos difusos. Esta regla de inferencia puede interpretarse de la siguiente manera: si $p \to q$ y tenemos $p'$ (aproximadamente $p$) entonces tenemos $q'$ (aproximadamente q). El modus ponens generalizado permite, por tanto, utilizar conceptos difusos y obtener el ``grado de verdad'' de la conclusión a partir de las premisas.

En la siguiente sección se define el concepto de \emph{regla difusa} (p.e. $\text{IF } x \text{ es } A \text{ THEN } y \text{ es } B$), que constituye una de las premisas en el modus ponens generalizado. En las secciones \ref{sec:mamdani-controller} y \ref{sec:interpolation-method} se definen los métodos de inferencia difusa que se van a estudiar en este proyecto. Estos métodos utilizan reglas difusas y el modus ponens generalizado para obtener las salidas ($B'$) a partir de las premisas ($A'$).

\section{Reglas difusas}
Una \emph{regla difusa IF-THEN}\footnote{Podría traducirse como SI-ENTONCES.} puede definirse de manera informal como:

\begin{definition}
Una regla difusa IF-THEN es una estructura condicional de la forma:
\begin{equation}
\text{IF [antecedentes] THEN [consecuentes]}
\end{equation},
donde [antecedentes] y [consecuentes] son proposiciones difusas.
\end{definition}

Las proposiciones difusas pueden ser \emph{simples} o \emph{compuestas}. Una proposición difusa simple es un predicado único formado por una variable lingüística \emph{x} y un valor \emph{A} de dicha variable (un conjunto difuso). Es decir, \emph{A} es un conjunto difuso (valor lingüístico) definido sobre el mismo universo de referencia \emph{U} de \emph{x}. Por ejemplo si la variable \emph{x} es la \emph{``Temperatura medida en ºC''} del ejemplo \ref{ex:lang-variable}, entonces las siguientes son proposiciones difusas sobre \emph{x}:

\begin{equation}
\text{$x$ es B}
\end{equation}
\begin{equation}
\text{$x$ es M}
\end{equation}
\begin{equation}
\text{$x$ es A}
\end{equation}

donde \emph{B}, \emph{M} y \emph{A} son los valores lingüísticos ``baja'', ``media'' y ``alta'' respectivamente. Estos valores lingüísticos vienen definidos por conjuntos difusos con funciones de pertenencia $\mu_{Baja}$, $\mu_{Media}$ y $\mu_{Alta}$.

Las proposiciones difusas simples se pueden combinar para formar predicados más complejos. Para ello se utilizan normalmente los conectivos \emph{AND} (intersección difusa) y \emph{OR} (unión difusa). Por ejemplo, se puede definir la condición \emph{``Temperatura es Alta Ó Temperatura es Media''} con la siguiente proposición difusa:

\begin{equation}
\text{$x$ es $A$ OR $x$ es $M$}
\end{equation}

donde $A$ y $M$ son los valores lingüísticos \emph{``alta''} ($\mu_{Alta}$) y \emph{``media''} ($\mu_{Media}$) respectivamente. En la misma proposición difusa se pueden utilizar variables diferentes (generalmente es así) definidas normalmente sobre universos de referencia diferentes. El consecuente generalmente suele ser una proposición difusa simple formado por una variable \emph{y} sobre \emph{V}. Por ello, se puede definir el concepto de \emph{regla difusa IF-THEN} de manera más formal como:

\begin{definition}
Una \emph{regla difusa IF-THEN} es una estructura condicional de la forma:

\begin{equation}
\text{IF $x_{1}$ es $A_{1}$ $\odot$ $x_{2}$ es $A_{2}$ $\odot$ \ldots $\odot$ $x_{n}$ es $A_{n}$ THEN $y$ es $B$}
\end{equation}
donde $x_{1} \in U_{i}$,\ldots, $x_{n} \in U_{n}$,$y \in V$ son variables lingüísticas y $\odot$ algún conectivo difuso (AND,OR,AND NOT,etc.).
\end{definition}

\section{Sistemas difusos basados en reglas}\label{sec:sistemas-difusos-basados-en-reglas}
Un enfoque muy utilizado a la hora de modelar un sistema de lógica difusa es la utilización de \emph{reglas} expresadas en un lenguaje muy próximo al natural, mediante variables lingüísticas. Estas reglas habitualmente son formuladas por un experto humano aunque pueden ser derivadas también de datos numéricos. A este tipo de sistemas difusos se les denomina \emph{sistemas difusos basados en reglas}. Los métodos estudiados en este trabajo pertenecen a este tipo de sistemas.

Existen principalmente 3 tipos de sistemas difusos \cite{wang1997}: sistemas difusos puros, sistemas difusos tipo Takagi-Sugeno-Kang (TKG) y sistemas difusos con fusificador y defusificador. A continuación se realiza una breve descripción de las características de estos sistemas.

En la figura \ref{fig:pure-fuzzy-system} se representa el esquema básico de un sistema difuso puro. En los sistemas difusos puros tanto las entradas como la salida del sistema son conjuntos difusos. El sistema difuso está dotado de un \emph{conjunto de reglas difusas} \emph{IF-THEN}, que forman la base de conocimiento sobre el dominio del problema. El \emph{sistema de inferencia difusa} combina estas reglas \emph{IF-THEN} y transforma los conjuntos difusos de entrada sobre el universo $U \subset R^{n}$ ($x_{1}.\ldots,x_{n}$) en un conjunto difuso sobre $V \subset R$ ($y$), utilizando principios de lógica difusa (modus ponens).

\begin{figure}[tb]
	\centering
	\tikzstyle{block} = [draw, fill=blue!20, rectangle, minimum height=3em, minimum width=6em]
\tikzstyle{input} = [coordinate]
\tikzstyle{output} = [coordinate]

\tikzset{
    state/.style={
    	   fill=white,
           rectangle,
           rounded corners,
           draw=black, very thick,
           minimum height=4em,
           minimum width=15em,
           inner sep=2pt,
           text centered,
           drop shadow
           },
}

\tikzstyle{materia}=[draw, fill=blue!20, text width=6.0em, text centered,
  minimum height=1.5em,drop shadow]


\shorthandoff{>}
\begin{tikzpicture}[auto, node distance=2cm,>=latex']
\node [input, name=input] {};
\node [state, right of=input,node distance=7cm,align=left] (inference-engine) {Sistema de inferencia difusa};
\node [state, above of=inference-engine] (rule-base) {Conjunto de reglas difusas};
\node [output, right of=inference-engine,node distance=7cm] (output) {};

\draw [->] (input) -- node {$x \in FS(U)$} (inference-engine);
\draw [->] (rule-base) -- node {} (inference-engine);
\draw [->] (inference-engine) -- node {$y \in FS(V)$} (output);
\end{tikzpicture}
	\caption{Diagrama de un sistema difuso puro.}
	\label{fig:pure-fuzzy-system}
\end{figure}

El principal problema de los sistemas difusos puros es que tanto las entradas como las salidas son conjuntos difusos, mientras que en aplicaciones de control industrial se utilizan generalmente valores escalares reales. Para solucionar este problema, Takagi, Sugeno y Kang propusieron un nuevo método simplificado que utiliza entradas y salidas escalares \cite{takagisugeno1985}\cite{sugenokang1988}.

El método de Takagi-Sugeno-Kang utiliza reglas de la forma:

\begin{equation}
R_{i}: \text{IF }x_{1}\text{ es }A_{1i}\text{, }x_{2}\text{ es }A_{2i}\text{ , \ldots , }x_{n}\text{ es }A_{ni}\text{ THEN } y = f_{i}(x_{1},x_{2},\ldots,x_{n})
\end{equation}

donde $A_{ij}$ representa las funciones de pertenencia asociadas a las variables lingüísticas utilizadas en los antecedentes y $f_{i}(x_{1},x_{2},\ldots,x_{n})$ representa los consecuentes. Normalmente $f_{i}$ es un polinomio en las variables de entrada, pero puede ser cualquier tipo de función mientras pueda describir la salida del modelo de forma apropiada, teniendo en cuenta las entradas. 

Como se puede ver, la principal diferencia entre un sistema difuso puro y el sistema de Takagi-Sugeno-Kang es que el consecuente de la regla (\emph{THEN}) es sustituido por una simple fórmula matemática. Este cambio supone que el método para combinar las diferentes reglas para proporcionar una salida es más sencillo. De hecho, en el método de Takagi-Sugeno-Kang se utiliza la media ponderada de los consecuentes obtenidos de las diferentes reglas.

El principal problema del método de Takagi-Sugeno-Kang es que el consecuente de las reglas es una función matemática y por lo tanto se pierde cierta capacidad de modelar el sistema utilizando el lenguaje natural. Para solucionar estos problemas, se puede modificar el sistema difuso puro, añadiendo un fusificador a la entrada de las variables y un defusificador a la salida.

\begin{figure}[tb]
	\centering
	\tikzstyle{block} = [draw, fill=blue!20, rectangle, minimum height=3em, minimum width=6em]
\tikzstyle{input} = [coordinate]
\tikzstyle{output} = [coordinate]

\tikzset{
    state/.style={
    	   fill=white,
           rectangle,
           rounded corners,
           draw=black, very thick,
           minimum height=4em,
           minimum width=15em,
           inner sep=2pt,
           text centered,
           drop shadow
           },
}

\tikzset{
    state-small/.style={
    	   fill=white,
           rectangle,
           rounded corners,
           draw=black, very thick,
           minimum height=4em,
           minimum width=5em,
           inner sep=2pt,
           text centered,
           drop shadow
           },
}

\tikzstyle{materia}=[draw, fill=blue!20, text width=6.0em, text centered,
  minimum height=1.5em,drop shadow]


\shorthandoff{>}
\begin{tikzpicture}[auto, node distance=2cm,>=latex']
\node [input, name=input] {};
\node [state-small, right of=input,node distance=2.5cm,align=left] (fuzzifier) {Fusificar};
\node [state, right of=fuzzifier,node distance=6cm,align=left] (inference-engine) {Sistema de inferencia difusa};
\node [state, above of=inference-engine] (rule-base) {Conjunto de reglas difusas};
\node [state-small, right of=inference-engine,node distance=6cm,align=left] (defuzzifier) {Defusificar};
\node [output, right of=defuzzifier,node distance=2.5cm] (output) {};

\draw [->] (input) -- node {$x \in U$} (fuzzifier);
\draw [->] (fuzzifier) -- node {$x \in FS(U)$} (inference-engine);
\draw [->] (rule-base) -- node {} (inference-engine);
\draw [->] (inference-engine) -- node {$y \in FS(V)$} (defuzzifier);
\draw [->] (defuzzifier) -- node {$y \in V$} (output);
\end{tikzpicture}
	\caption{Diagrama de un sistema difuso con fusificador y defusificador.}
	\label{fig:fuzzy-system}
\end{figure}

En la figura \ref{fig:fuzzy-system} se representa el esquema básico de un sistema difuso con fusificador y defusificador. El fusificador transforma las entradas escalares en $U$ del sistema en conjuntos difusos. Una vez realizada la inferencia difusa se obtiene un conjunto difuso como salida. El defusificador transforma el conjunto difuso de salida en una variable escalar en $V$. Existen multitud de métodos para realizar estas transformaciones y en siguientes secciones se detallarán algunos de ellos. El sistema difuso presentado en este trabajo se basa en este último tipo de sistema.

\subsection{Fusificadores}\label{sec:fusificadores}

En la sección \ref{sec:sistemas-difusos-basados-en-reglas} se ha definido un tipo de sistema difuso que utiliza un concepto llamado \emph{fusificador} para transformar las entradas escalares del sistema en conjuntos difusos. Un \emph{fusificador} puede definirse formalmente como:

\begin{definition}\label{def:fusificador}
Un \emph{fusificador} es una función $F:\mathbb{R} \rightarrow FS(U)$ que toma un valor escalar real $x^{*}$ en el universo $U$, $x^{*}\in U \subset \mathbb{R}$, y lo transforma en un conjunto difuso $A'$ sobre $U$.
\end{definition}

Para que el conjunto de salida represente de manera lo más fiel posible al valor de entrada, es necesario imponer algunas restricciones a la hora de construir el fusificador. La principal es que la función de pertenencia de $A'$ debe tener un valor alto en el punto $x^*$ (generalmente máximo, $\mu_{A'}(x^*) = \max(\mu_{A'}(\mu_{i})),\forall \mu_{i} \in U$). Además, si se considera que existe ruido o imprecisión en los valores entrada (por ejemplo imprecisiones provocadas por instrumentos de medida) se puede utilizar una función de pertenencia cuyo valor sea alto en los puntos próximos a $x^*$ y bajo (o incluso cero) en los más alejados.

Algunos de los fusificadores más utilizados son los siguientes \cite{wang1997}:

\begin{itemize}
\item\bfseries Fusificador singleton: \normalfont El \emph{fusificador singleton} (fig. \ref{fig:fuzzifier-singleton}) transforma un valor escalar real $x^* \in U$ en un conjunto difuso $A'$ en $U$, de forma que la función de pertenencia de $A'$ tiene valor 1 en el punto $x^*$ y 0 en el resto de puntos de $U$:
\begin{equation}
\mu_{A'}(x) = \begin{cases} 1 & \mbox{si } \mbox{x}=x^*, \\ 0 & \mbox{en cualquier otro caso} \end{cases}
\end{equation}
Este tipo de fusificador tiene la ventaja de ser simple y facilitar los cálculos realizados en el motor de inferencia. Sin embargo, tiene la desventaja de que no puede lidiar con ruido o imprecisiones en los valores de entrada. 
\item\bfseries Fusificador gaussiano: \normalfont El \emph{fusificador gaussiano} (fig. \ref{fig:fuzzifier-gaussian}) transforma un valor escalar real $x^* \in U$ en un conjunto difuso $A'$ en $U$, utilizando una función de pertenencia Gaussiana:
\begin{equation}
\mu_{A'}(x) = e^{-(\frac{x_1 - x_1^*}{a_1})^2}\star\ldots\star e^{-(\frac{x_n - x_n^*}{a_n})^2}
\end{equation}
donde $a_i$ son parámetros positivos y $\star$ es una t-norma, generalmente el producto algebraico o el mínimo. Este tipo de fusificador es adecuado cuando existe ruido o imprecisiones en los datos de entrada.
\item\bfseries Fusificador triangular: \normalfont El \emph{fusificador triangular} (fig. \ref{fig:fuzzifier-triangular}) transforma un valor escalar real $x^* \in U$ en un conjunto difuso $A'$ en $U$, utilizando una función de pertenencia triangular:
\begin{equation}
\mu_{A'}(x) = \begin{cases} (1 - \frac{|x_1 - x_1^*|}{b_1})\star\ldots\star  (1 - \frac{|x_n - x_n^*|}{b_n}) & \mbox{si } |x_i - x_i^*| \leq b_i, i = 1,2,\ldots,n, \\ 0 & \mbox{en cualquier otro caso} \end{cases}
\end{equation}
donde $b_i$ son parámetros positivos y $\star$ es una t-norma, generalmente el producto algebraico o el mínimo. Al igual que el fusificador gaussiano este tipo de fusificador es adecuado cuando existe ruido o imprecisiones en los datos de entrada.
\end{itemize}

\begin{figure}[H]
	\centering
	\begin{subfigure}[t]{\textwidth}
		\setlength\figureheight{4cm}
		\setlength\figurewidth{12cm}
		% This file was created by matlab2tikz v0.4.7 (commit b040dd9e3d20b4d02324400633a5d15752e58612) running on MATLAB 8.0.
% Copyright (c) 2008--2014, Nico Schlömer <nico.schloemer@gmail.com>
% All rights reserved.
% Minimal pgfplots version: 1.3
% 
\begin{tikzpicture}

\begin{axis}[%
width=\figurewidth,
height=\figureheight,
scale only axis,
xmin=0,
xmax=100,
xtick={0,20,40,60,80,100},
xticklabels={\empty},
xlabel={x},
ymin=0,
ymax=1,
ylabel={$\mu{}_{\text{A'}}\text{(x)}$},
axis x line*=bottom,
axis y line*=left
]
\addplot [color=black,solid,line width=1.5pt,forget plot]
  table[row sep=crcr]{0	0\\
1	0\\
2	0\\
3	0\\
4	0\\
5	0\\
6	0\\
7	0\\
8	0\\
9	0\\
10	0\\
11	0\\
12	0\\
13	0\\
14	0\\
15	0\\
16	0\\
17	0\\
18	0\\
19	0\\
20	0\\
21	0\\
22	0\\
23	0\\
24	0\\
25	0\\
26	0\\
27	0\\
28	0\\
29	0\\
30	0\\
31	0\\
32	0\\
33	0\\
34	0\\
35	0\\
36	0\\
37	0\\
38	0\\
39	0\\
40	0\\
41	0\\
42	0\\
43	0\\
44	0\\
45	0\\
46	0\\
47	0\\
48	0\\
49	0\\
50	0\\
51	0\\
52	0\\
53	0\\
54	0\\
55	0\\
56	0\\
57	0\\
58	0\\
59	0\\
60	0\\
61	0\\
62	0\\
63	0\\
64	0\\
65	0\\
66	0\\
67	0\\
68	0\\
69	0\\
70	0\\
71	0\\
72	0\\
73	0\\
74	0\\
75	0\\
76	0\\
77	0\\
78	0\\
79	0\\
80	0\\
81	0\\
82	0\\
83	0\\
84	0\\
85	0\\
86	0\\
87	0\\
88	0\\
89	0\\
90	0\\
91	0\\
92	0\\
93	0\\
94	0\\
95	0\\
96	0\\
97	0\\
98	0\\
99	0\\
100	0\\
};
\addplot [color=blue,line width=1.5pt,mark size=4.0pt,only marks,mark=asterisk,mark options={solid},forget plot]
  table[row sep=crcr]{40	1\\
};
\addplot [color=black,dotted,line width=1.5pt,forget plot]
  table[row sep=crcr]{40	0\\
40	1\\
};
\node[right, inner sep=0mm, text=black]
at (axis cs:42,0.95,0) {$\text{(x*,}\mu{}_{\text{A'}}\text{(x*))}$};
\end{axis}
\end{tikzpicture}%
		\caption{Fusificador singleton.}
		\label{fig:fuzzifier-singleton}
	\end{subfigure}
	
	\vspace{1 cm}
	\begin{subfigure}[t]{\textwidth}
		\setlength\figureheight{4cm}
		\setlength\figurewidth{12cm}
		% This file was created by matlab2tikz v0.4.7 (commit b040dd9e3d20b4d02324400633a5d15752e58612) running on MATLAB 8.0.
% Copyright (c) 2008--2014, Nico Schlömer <nico.schloemer@gmail.com>
% All rights reserved.
% Minimal pgfplots version: 1.3
% 
\begin{tikzpicture}

\begin{axis}[%
width=\figurewidth,
height=\figureheight,
scale only axis,
xmin=0,
xmax=100,
xtick={0,20,40,60,80,100},
xticklabels={\empty},
xlabel={x},
ymin=0,
ymax=1,
ylabel={$\mu{}_{\text{A'}}\text{(x)}$},
axis x line*=bottom,
axis y line*=left
]
\addplot [color=black,solid,line width=1.5pt,forget plot]
  table[row sep=crcr]{0	0.000335462627902512\\
1	0.000497955421503273\\
2	0.000731802418880473\\
3	0.00106476623666792\\
4	0.00153381067932446\\
5	0.00218749111818289\\
6	0.00308871540823677\\
7	0.00431784000763308\\
8	0.00597602289500594\\
9	0.00818870101437408\\
10	0.0111089965382423\\
11	0.0149207860690678\\
12	0.0198410947443703\\
13	0.0261214098539182\\
14	0.0340474547345993\\
15	0.0439369336234074\\
16	0.0561347628341337\\
17	0.071005353739637\\
18	0.0889216174593863\\
19	0.110250525304485\\
20	0.135335283236613\\
21	0.164474456577155\\
22	0.197898699083615\\
23	0.235746076555864\\
24	0.278037300453194\\
25	0.32465246735835\\
26	0.3753110988514\\
27	0.429557358210739\\
28	0.486752255959972\\
29	0.546074426639709\\
30	0.606530659712633\\
31	0.666976810858474\\
32	0.726149037073691\\
33	0.782704538241868\\
34	0.835270211411272\\
35	0.882496902584595\\
36	0.923116346386636\\
37	0.9559974818331\\
38	0.980198673306755\\
39	0.995012479192682\\
40	1\\
41	0.995012479192682\\
42	0.980198673306755\\
43	0.9559974818331\\
44	0.923116346386636\\
45	0.882496902584595\\
46	0.835270211411272\\
47	0.782704538241868\\
48	0.726149037073691\\
49	0.666976810858474\\
50	0.606530659712633\\
51	0.546074426639709\\
52	0.486752255959972\\
53	0.429557358210739\\
54	0.3753110988514\\
55	0.32465246735835\\
56	0.278037300453194\\
57	0.235746076555864\\
58	0.197898699083615\\
59	0.164474456577155\\
60	0.135335283236613\\
61	0.110250525304485\\
62	0.0889216174593863\\
63	0.071005353739637\\
64	0.0561347628341337\\
65	0.0439369336234074\\
66	0.0340474547345993\\
67	0.0261214098539182\\
68	0.0198410947443703\\
69	0.0149207860690678\\
70	0.0111089965382423\\
71	0.00818870101437408\\
72	0.00597602289500594\\
73	0.00431784000763308\\
74	0.00308871540823677\\
75	0.00218749111818289\\
76	0.00153381067932446\\
77	0.00106476623666792\\
78	0.000731802418880473\\
79	0.000497955421503273\\
80	0.000335462627902512\\
81	0.000223745793720621\\
82	0.000147748360232034\\
83	9.6593413722184e-05\\
84	6.25215037748203e-05\\
85	4.00652973929511e-05\\
86	2.54193465161992e-05\\
87	1.59667838978047e-05\\
88	9.92950430585108e-06\\
89	6.1135679663714e-06\\
90	3.72665317207867e-06\\
91	2.24905596703234e-06\\
92	1.34381227763152e-06\\
93	7.94939361534914e-07\\
94	4.65571571578309e-07\\
95	2.69957850336301e-07\\
96	1.5497531357029e-07\\
97	8.80817919646056e-08\\
98	4.9564053191725e-08\\
99	2.76124245682803e-08\\
100	1.52299797447126e-08\\
};
\addplot [color=blue,line width=1.5pt,mark size=4.0pt,only marks,mark=asterisk,mark options={solid},forget plot]
  table[row sep=crcr]{40	1\\
};
\addplot [color=black,dotted,line width=1.5pt,forget plot]
  table[row sep=crcr]{40	0\\
40	1\\
};
\node[right, inner sep=0mm, text=black]
at (axis cs:44,0.96,0) {$\text{(x*,}\mu{}_{\text{A'}}\text{(x*))}$};
\end{axis}
\end{tikzpicture}%
		\caption{Fusificador gaussiano.}
		\label{fig:fuzzifier-gaussian}
	\end{subfigure}
	
	\vspace{1 cm}
	\begin{subfigure}[t]{\textwidth}
		\setlength\figureheight{4cm}
		\setlength\figurewidth{12cm}
		% This file was created by matlab2tikz v0.4.7 (commit b040dd9e3d20b4d02324400633a5d15752e58612) running on MATLAB 8.0.
% Copyright (c) 2008--2014, Nico Schlömer <nico.schloemer@gmail.com>
% All rights reserved.
% Minimal pgfplots version: 1.3
% 
\begin{tikzpicture}

\begin{axis}[%
width=\figurewidth,
height=\figureheight,
scale only axis,
xmin=0,
xmax=100,
xtick={0,20,40,60,80,100},
xticklabels={\empty},
xlabel={x},
ymin=0,
ymax=1,
ylabel={$\mu{}_{\text{A'}}\text{(x)}$},
axis x line*=bottom,
axis y line*=left
]
\addplot [color=black,solid,line width=1.5pt,forget plot]
  table[row sep=crcr]{0	0\\
1	0\\
2	0\\
3	0\\
4	0\\
5	0\\
6	0\\
7	0\\
8	0\\
9	0\\
10	0\\
11	0\\
12	0\\
13	0\\
14	0\\
15	0\\
16	0\\
17	0\\
18	0\\
19	0\\
20	0\\
21	0.05\\
22	0.1\\
23	0.15\\
24	0.2\\
25	0.25\\
26	0.3\\
27	0.35\\
28	0.4\\
29	0.45\\
30	0.5\\
31	0.55\\
32	0.6\\
33	0.65\\
34	0.7\\
35	0.75\\
36	0.8\\
37	0.85\\
38	0.9\\
39	0.95\\
40	1\\
41	0.95\\
42	0.9\\
43	0.85\\
44	0.8\\
45	0.75\\
46	0.7\\
47	0.65\\
48	0.6\\
49	0.55\\
50	0.5\\
51	0.45\\
52	0.4\\
53	0.35\\
54	0.3\\
55	0.25\\
56	0.2\\
57	0.15\\
58	0.1\\
59	0.05\\
60	0\\
61	0\\
62	0\\
63	0\\
64	0\\
65	0\\
66	0\\
67	0\\
68	0\\
69	0\\
70	0\\
71	0\\
72	0\\
73	0\\
74	0\\
75	0\\
76	0\\
77	0\\
78	0\\
79	0\\
80	0\\
81	0\\
82	0\\
83	0\\
84	0\\
85	0\\
86	0\\
87	0\\
88	0\\
89	0\\
90	0\\
91	0\\
92	0\\
93	0\\
94	0\\
95	0\\
96	0\\
97	0\\
98	0\\
99	0\\
100	0\\
};
\addplot [color=blue,line width=1.5pt,mark size=4.0pt,only marks,mark=asterisk,mark options={solid},forget plot]
  table[row sep=crcr]{40	1\\
};
\addplot [color=black,dotted,line width=1.5pt,forget plot]
  table[row sep=crcr]{40	0\\
40	1\\
};
\node[right, inner sep=0mm, text=black]
at (axis cs:44,0.96,0) {$\text{(x*,}\mu{}_{\text{A'}}\text{(x*))}$};
\end{axis}
\end{tikzpicture}%
		\caption{Fusificador triangular.}
		\label{fig:fuzzifier-triangular}
	\end{subfigure}
		\caption{Representación de las funciones de pertenencia obtenidas con los fusificadores \emph{singleton} (\ref{fig:fuzzifier-singleton}), \emph{gaussiano} (\ref{fig:fuzzifier-gaussian}) y \emph{triangular} (\ref{fig:fuzzifier-triangular}).}
		\label{fig:fuzzifiers}
\end{figure}

\subsection{Defusificadores}\label{sec:defusificadores}
El último paso en el sistema difuso definido en \ref{sec:sistemas-difusos-basados-en-reglas} es la \emph{defusificación}, es decir, transformar el conjunto difuso, obtenido al aplicar las reglas de inferencia, en un valor escalar. Puede decirse que un defusificador realiza la operación complementaria a un fusificador (descritos en la sección \ref{sec:fusificadores}). Así pues, un defusificador puede definirse como \cite{wang1997}:

\begin{definition}
Un \emph{defusificador} es una función que transforma un conjunto difuso $B'$ en $V \subset \mathbb{R}$ (la salida del sistema de inferencia difusa) en un valor escalar $y^* \in V$.
\end{definition}

La misión de un defusificador es obtener el punto $y^*$ en $V$ que mejor representa al conjunto difuso $B'$. En este sentido, los defusificadores cumplen un rol similar a los operadores de agregación (por ejemplo la media aritmética). Se pueden utilizar diversas funciones como defusificadores, aunque generalmente se eligen aquellas que satisfacen las siguientes condiciones \cite{hellendoorn93}:

\begin{itemize}
\item\bfseries Continuidad: \normalfont Un cambio pequeño en la entrada del sistema difuso no debería causar una gran variación en la salida.
\item\bfseries Desambiguación: \normalfont El método de defusificación debe devolver un único valor para $y^*$. Es decir, no debe existir ambigüedad a la hora de seleccionar el valor para $y^*$.
\item\bfseries Plausibilidad: \normalfont El punto $y^*$ debe representar al conjunto difuso $B'$ de forma intuitiva. Es decir, el punto $y^*$ debe estar aproximadamente en el centro del soporte de $B'$ y tener un grado de pertenencia alto en $B'$.
\item\bfseries Simplicidad computacional: \normalfont La función debe ser lo más sencilla posible de calcular. Esto es especialmente importante en controladores difusos que operan en tiempo real.
\item\bfseries Método de ponderado: \normalfont El método seleccionado debe ponderar los diferentes conjuntos de salida. Este criterio depende del ámbito del problema.
\end{itemize}
A continuación se presentan algunos de los métodos de defusificación más utilizados:

\begin{itemize}
  \item\bfseries Centroide: \normalfont También conocido como \emph{centro de gravedad}, el método del \emph{centroide} fue propuesto por Sugeno en 1985 y es uno de los métodos más utilizados \cite{lee1990}. Este método obtiene el valor $y^*$ como el centro del área cubierta por la función de pertenencia de $B'$. El método del centro de gravedad puede definirse como:
  \begin{equation}
	y^* = \frac{\sum y \mu_{B'}(y)}{\sum \mu_{B'}(y)}
  \end{equation}
  Este método tiene la ventaja de que proporciona un valor plausible de forma intuitiva.
  \item\bfseries Bisector: \normalfont El método del bisector obtiene el punto por el que pasa la línea que divide la región delimitada por la función de pertenencia de $B'$ y el eje de abscisas, en dos subsecciones de igual área:
  \begin{equation}
	y^*\text{ tal que: } \sum\limits_{y=y_0}^{y=y^*}\mu_{B'}(y) = \sum\limits_{y=y^*}^{y=y_n}\mu_{B'}(y)
  \end{equation}
  El valor obtenido con este método coincide en ocasiones con el método del \emph{centroide} y también es computacionalmente costoso.
  \item\bfseries Máximo: \normalfont El defusificador del máximo obtiene el punto $y^*$ en el que la función de pertenencia $\mu_{B'}(y)$ alcanza su máximo valor. Es decir:
  \begin{equation}
	y^*\text{ tal que: } \mu_{B'}(y^*) = \max(\mu_{B'}(y_{i}))\qquad ,\forall y_{i} \in V
  \end{equation}
  Este método tiene la ventaja de ser fácil de implementar y poco costoso de calcular. Sin embargo, puede producir resultados ambiguos, puesto que pueden existir varios puntos en los que $\mu_{B'}(y)$ alcanza un valor máximo. En el caso de que existan varios puntos con valores máximos es necesario imponer algún criterio para seleccionar uno de ellos. Si se construye el conjunto:
  \begin{equation}
	hgt(B') = \{y \in V | \mu_{B'}(y) = \sup\limits_{y \in V}\mu_{B'}(y)\}
  \end{equation}
  es decir, el conjunto de todos los puntos en $V$ donde $\mu_{B'}(y)$ alcanza su valor máximo. El defusificador del máximo obtiene un elemento arbitrario $y^*$ del conjunto $hgt(B')$, es decir:
  \begin{equation}
	y^* = \text{ cualquier punto en } hgt(B')
  \end{equation}
  \item\bfseries Menor de máximos (SOM, \emph{smallest of maximum}): \normalfont Dado el conjunto $hgt(B')$ con todos los puntos donde $\mu_{B'}(y)$ alcanza su valor máximo, el defusificador \emph{menor de máximos} obtiene el menor punto $y^*$ del conjunto tal que:
  \begin{equation}
	y^* = \inf\{y \in hgt(B')\}
  \end{equation}
  \item\bfseries Mayor de máximos (LOM, \emph{largest of maximum}): \normalfont Dado el conjunto $hgt(B')$ con todos los puntos donde $\mu_{B'}(y)$ alcanza su valor máximo, el defusificador \emph{mayor de máximos} obtiene el mayor punto $y^*$ del conjunto tal que:
   \begin{equation}
  	y^* = \sup\{y \in hgt(B')\}
    \end{equation}
   \item\bfseries Media de máximos (MOM, \emph{mean of maximum}): \normalfont Dado el conjunto $hgt(B')$ con todos los puntos donde $\mu_{B'}(y)$ alcanza su valor máximo, el defusificador \emph{media de máximos} obtiene el punto medio $y^*$ del conjunto tal que:
   \begin{equation}
     	y^* = \frac{ \inf\{y \in hgt(B')\} + \sup\{y \in hgt(B')\}}{2}
   \end{equation}
\end{itemize}
En la figura \ref{fig:defuzzifiers} se presenta una función de pertenencia y los resultados de aplicar los defusificadores definidos anteriormente. En este caso, los resultados obtenidos con el método del \emph{centroide} y el \emph{bisector} coinciden, aunque no siempre tiene por qué ser así.
\begin{figure}[t]
	\centering
	\setlength\figureheight{5.5cm}
	\setlength\figurewidth{12cm}
	% This file was created by matlab2tikz v0.4.7 (commit e06548f86e222baeac9b1532d97c4b747a61abe5) running on MATLAB 8.0.
% Copyright (c) 2008--2014, Nico Schlömer <nico.schloemer@gmail.com>
% All rights reserved.
% Minimal pgfplots version: 1.3
% 
%
% defining custom colors
\definecolor{mycolor1}{rgb}{0.00000,0.75000,0.75000}%
\definecolor{mycolor2}{rgb}{0.75000,0.00000,0.75000}%
%
\begin{tikzpicture}

\begin{axis}[%
width=\figurewidth,
height=\figureheight,
scale only axis,
xmin=0,
xmax=100,
xtick={0,10,20,30,40,50,60,70,80,90,100},
xticklabels={\empty},
ymin=0,
ymax=1,
axis x line*=bottom,
axis y line*=left,
legend style={draw=black,fill=white,legend cell align=left}
]
\addplot [color=black,solid,line width=1.2pt]
  table[row sep=crcr]{0	0\\
1	0.06\\
2	0.12\\
3	0.18\\
4	0.24\\
5	0.3\\
6	0.36\\
7	0.42\\
8	0.48\\
9	0.54\\
10	0.6\\
11	0.6\\
12	0.6\\
13	0.6\\
14	0.6\\
15	0.6\\
16	0.6\\
17	0.6\\
18	0.6\\
19	0.6\\
20	0.6\\
21	0.6\\
22	0.6\\
23	0.6\\
24	0.6\\
25	0.6\\
26	0.6\\
27	0.6\\
28	0.6\\
29	0.633333333333333\\
30	0.666666666666667\\
31	0.7\\
32	0.733333333333333\\
33	0.766666666666667\\
34	0.8\\
35	0.833333333333333\\
36	0.866666666666667\\
37	0.9\\
38	0.933333333333333\\
39	0.966666666666667\\
40	1\\
41	1\\
42	1\\
43	1\\
44	1\\
45	1\\
46	1\\
47	1\\
48	1\\
49	1\\
50	1\\
51	1\\
52	1\\
53	1\\
54	1\\
55	1\\
56	1\\
57	1\\
58	1\\
59	1\\
60	1\\
61	0.95\\
62	0.9\\
63	0.85\\
64	0.8\\
65	0.75\\
66	0.7\\
67	0.65\\
68	0.6\\
69	0.55\\
70	0.5\\
71	0.45\\
72	0.4\\
73	0.4\\
74	0.4\\
75	0.4\\
76	0.4\\
77	0.4\\
78	0.4\\
79	0.4\\
80	0.4\\
81	0.4\\
82	0.4\\
83	0.4\\
84	0.4\\
85	0.4\\
86	0.4\\
87	0.4\\
88	0.4\\
89	0.4\\
90	0.4\\
91	0.36\\
92	0.32\\
93	0.28\\
94	0.24\\
95	0.2\\
96	0.16\\
97	0.12\\
98	0.08\\
99	0.04\\
100	0\\
};
\addlegendentry{$\mu\text{(x)}$};

\addplot [color=blue,line width=1.2pt,mark size=4.5pt,only marks,mark=asterisk,mark options={solid}]
  table[row sep=crcr]{47	1\\
};
\addlegendentry{centroid};

\addplot [color=black!50!green,line width=1.2pt,mark size=4.5pt,only marks,mark=+,mark options={solid}]
  table[row sep=crcr]{47	1\\
};
\addlegendentry{bisector};

\addplot [color=red,line width=1.2pt,mark size=3.2pt,only marks,mark=square,mark options={solid}]
  table[row sep=crcr]{50	1\\
};
\addlegendentry{mom};

\addplot [color=mycolor1,line width=1.2pt,mark size=3.0pt,only marks,mark=triangle,mark options={solid,rotate=180}]
  table[row sep=crcr]{40	1\\
};
\addlegendentry{som};

\addplot [color=mycolor2,line width=1.2pt,mark size=3.0pt,only marks,mark=triangle,mark options={solid}]
  table[row sep=crcr]{60	1\\
};
\addlegendentry{lom};

\end{axis}
\end{tikzpicture}%
	\caption{Resultado de aplicar los defusificadores \emph{centroide} (\textasteriskcentered), \emph{bisector}(+), \emph{media de máximos} ($\square$), \emph{menor de máximos}($\triangledown$) y \emph{mayor de máximos}($\vartriangle$).}
	\label{fig:defuzzifiers}
\end{figure}

\section{El controlador de Mamdani}\label{sec:mamdani-controller}
Uno de los métodos de inferencia difusa más utilizado es el llamado \emph{controlador de Mamdani} \cite{Mamdani1975}, propuesto por Mamdani y Assilian en 1975 para realizar el control de un motor de vapor a partir de un conjunto de reglas obtenidas de operadores humanos experimentados.

El método de Mamdani utiliza reglas difusas de la forma:

\begin{equation}
R_{i}: \text{IF }\chi_{1}\text{ es }A_{i1}\text{, }\chi_{2}\text{ es }A_{i2}\text{ , \ldots , }\chi_{n}\text{ es }A_{im}\text{ THEN } y \text{ es } B_i
\end{equation}

donde $\chi_1,\ldots,\chi_m$ son variables lingüísticas cuyos valores son conjuntos difusos $A_{ij} \subset FS(U_j)$ con $j\in\{1,\ldots,m\}$. El consecuente $y$ es también una variable lingüística sobre $Y$ ($B_i \subset FS(Y)$). Las entradas del sistema son valores escalares $x_1 \in U_1,\ldots,x_m \in U_m$. 

Cada regla se evalúa por separado y se obtiene el grado de pertenencia de cada entrada al conjunto difuso de su correspondiente antecedente, utilizando el operador mínimo para realizar la conjunción AND de los antecedentes ($k_i = \min(\mu_{Ai1}(x_1),\ldots,\mu_{Aim}(x_m))$). A este valor $k_i$ se le denomina normalmente \emph{grado de activación} de la regla. El resultado de evaluar cada regla es un conjunto difuso sobre $Y$ ($B_i'$), que se calcula ``truncando'' el consecuente con el grado de activación ($B_i' = \{(y,\min(B_i(y),k_i))|y \in Y\}$). En el caso de múltiples reglas, el conjunto difuso de salida se calcula combinando los conjuntos difusos obtenidos en cada regla utilizando el máximo ($B'(y) = \max\limits_{i=1}^{n}(B_i')$).

\begin{algorithm}
\caption{Método de Mamdani}
\label{algo:mamdani}
\DontPrintSemicolon
\KwIn{Un conjunto de reglas $R_{i}$ ($i \in \{1,\ldots,n\}$) con varios antecedentes $A_{ij}$ ($j \in \{1,\ldots,m\}$)  y entradas escalares $x_1,\ldots,x_m$.}
\KwOut{\emph{B'}.}
\For{$i \in \{1,\ldots,n\}$} {
Calcular $k_i = \min(\mu_{Ai1}(x_1),\ldots,\mu_{Aim}(x_m))$ \\
Calcular $B_i' = \{(y,\min(B_i(y),k_i))|y \in Y\}$
}
Construir $B' = \{(y, B'(y))|y \in Y\}$ dado por: \\
\centering
\nonl $B'(y) = \max\limits_{i=1}^{n}(B_i')$.\\
\Return{$B'$}\;
\end{algorithm}

La principal ventaja del método de Mamdani es que es relativamente sencillo de implementar y computacionalmente eficiente. El método de Larsen es similar, pero utilizando el producto en vez del mínimo como operador de implicación \cite{larsen1980}.

\section{Método de interpolación basado en índices de solapamiento}\label{sec:interpolation-method}
En esta sección se introduce un nuevo método para resolver sistemas basados en reglas que generaliza el método clásico de interpolación , utilizando para ello índices de solapamiento \cite{bustince2013overlap}.\\
\\
Históricamente el método más utilizado para resolver sistemas basados en reglas, es decir, para calcular el consecuente \emph{B'}, era el método de interpolación \cite{klir1987}. En este método se utiliza la consistencia de Zadeh \emph{$O_{Z}$} \cite{zadeh1978}, dada en la ecuación \ref{eq:zadeh-consistency}. Los pasos a seguir son:

\begin{algorithm}
\DontPrintSemicolon
\KwIn{Un conjunto de reglas $R_{j}$, con $j \in \{1,\ldots,n\}$, un hecho \emph{A'} y el índice de consistencia $O_{Z}$ (ec. \ref{eq:zadeh-consistency}).}
\KwOut{\emph{B'}.}
\vspace{0.4 cm}
\For{$j \in \{1,\ldots,n\}$} {
Calcular $O_{Z}(A',A_{j}) = \max\limits_{x \in X}(\min(A'(x),A_{j}(x))) $
}
Construir $B' = \{(y, B'(y))|y \in Y\}$ dado por: \\
\centering
\nonl $B'(y) = \max\limits_{j=1}^{n}(\min(B_{j}(y),O_{Z}(A', A_{j})))$.\\
\Return{$B'$}\;
\caption{Método de interpolación}
\label{algo:interpolation-method}
\end{algorithm}
Dado que el índice de consistencia $O_{Z}$ es un índice de solapamiento que cumple la definición \ref{def:dubois-overlap-index}, el algoritmo \ref{algo:interpolation-method} se puede generalizar para utilizar cualquier índice de solapamiento:

\begin{algorithm}
\DontPrintSemicolon
\KwIn{Un conjunto de reglas $R_{j}$, con $j \in \{1,\ldots,n\}$ y un hecho \emph{A'}.}
\KwOut{\emph{B'}.}
\vspace{0.4 cm}
Seleccionar un operador de agregación $M_{1}$, una función de solapamiento $G_{O}$ y un índice de solapamiento O.\\
\For{$j \in \{1,\ldots,n\}$} {
Calcular $O(A', A_{j})$
}
Construir $B' = \{(y, B'(y))|y \in Y\}$ dado por: \\
\centering
\nonl $B'(y) = \overset{n}{\underset{j=1}{M_{1}}}(G_{O}(B_{j}(y),O(A', A_{j})))$.\\
\Return{$B'$}\;
\caption{Método de interpolación generalizado}
\label{algo:overlap-interpolation-method}
\end{algorithm}

Si en el algoritmo \ref{algo:overlap-interpolation-method} utilizamos $M_{1} = \max$, $G_{O} = \min$ y $O = O_{Z}$ entonces se recupera el algoritmo \ref{algo:interpolation-method}.

En el caso de que cada regla tenga varios antecedentes, se puede generalizar el algoritmo \ref{algo:overlap-interpolation-method} de la siguiente manera:

\begin{algorithm}
\DontPrintSemicolon
\KwIn{Un conjunto de reglas $R_{j}$ con varios antecedentes, con $j \in \{1,\ldots,n\}$ y un hecho \emph{A'}.}
\KwOut{\emph{B'}.}
\vspace{0.4 cm}
Seleccionar un operador de agregación $M$, una t-norma \emph{T} y un índice de solapamiento \emph{O}.\\
\For{$i =1 \to n$} {
Calcular $O(A'_{i}, A_{i1}),\ldots,O(A'_{m}, A_{im})$\\
Calcular $k_{i} = T(O(A'_{i}, A_{i1}),\ldots,O(A'_{m}, A_{im}))$\\
Construir sobre el universo de referencia \emph{Y} el conjunto $K_{i} = \{(y,k_{i})|y \in Y\}$
}
Construir $B' = \{(y, B'(y))|y \in Y\}$ dado por: \\
\centering
\nonl $B'(y) = \overset{n}{\underset{i=1}{M}}(\min(K_{i},B_{i}))$.\\
\Return{$B'$}\;
\caption{Método de interpolación generalizado para reglas con varios antecedentes}
\label{algo:multi-overlap-interpolation-method}
\end{algorithm}

Como en el algoritmo anterior, se toma un conjunto de reglas, en este caso con varios antecedentes cada una, y un hecho $A'$. Se debe seleccionar también un operador de agregación $M$, una t-norma $T$ y un índice de solapamiento $O$. En caso de que las reglas tengan más de un antecedente, el hecho $A'$ debe tener un valor para cada uno de ellos ($A'_1,\cdots,A'_n$).

Cada regla se evalúa por separado y para cada una se calculan los índices de solapamiento de cada valor del hecho $A'$ con su correspondiente antecedente en la regla (línea 3). Estos índices de solapamiento se agregan utilizando la t-norma $T$, que modela la conjunción AND (línea 4). El valor obtenido es el ``grado de activación'' de la regla ($k_i$). Este valor se utiliza para construir sobre el universo de salida $Y$ el conjunto difuso $K_i$, tal que todos los elementos $y \in Y$ tienen como grado de pertenencia el valor $k_i$ (línea 5).

Por último, para cada regla se calcula el conjunto de salida como el mínimo entre su correspondiente conjunto $K_i$ (calculado en pasos previos) y el conjunto difuso de salida de la regla $B_i$. Los conjuntos de salida obtenidos para cada regla se agregan utilizando el operador de agregación $M$ para obtener el conjunto de salida $B'$ (línea 7). El conjunto $B'$ es el resultado final de aplicar el método de inferencia a la entrada $A'$.



% !TeX spellcheck = es_ES
\chapter{Detección de incendios forestales}
\label{cha:deteccion-incendios-forestales}
En este capítulo se presenta una aplicación práctica de los métodos de inferencia difusa basados en reglas, descritos en los capítulos anteriores (REFERENCIA). El objetivo es definir un sistema basado en reglas capaz de determinar el riesgo de incendios forestales a partir de mediciones de algunas magnitudes tales como la temperatura, el humo, la humedad etc. Estas magnitudes constituirán las entradas del sistema de inferencia que obtendrá una estimación cuantitativa del riesgo del incendio forestal.

\section{Red de sensores inalámbricos}
El primer paso para la determinación del riesgo de un incendio forestal es la medición y toma de datos ambientales relacionados con dicho incendio. Para ello se puede desplegar una red de sensores inalámbricos (\emph{Wireless Sensor Networks}, WSN), que han sido desarrolladas y utilizadas en una gran variedad de aplicaciones en áreas tales como automoción \cite{hsin2007}, defensa, medicina, agricultura \cite{tao2008}\cite{hwang2010}, etc. 

Algunas aplicaciones de este tipo de sistemas relacionadas con los riesgos ambientales incluyen, por ejemplo, el control de movimientos de personas y ganado, monitorización de factores ambientales que afectan a la calidad de los cultivos, detección de incendios forestales, mediciones meteorológicas y detección de inundaciones etc. \cite{Akyildiz2002}

Las redes de sensores inalámbricos están diseñadas para monitorizar y controlar eventos en lugares que presentan riesgos ambientales tales como bosques, terrenos montañosos etc. Por esta razón se diseña este tipo de sistemas de forma que sean lo más autónomos posible. Esto incluye, por ejemplo, la utilización de energías renovables (típicamente energía solar), que posibilitan que los sensores desplegados en el terreno sean totalmente auto-suficientes.

El diseño y despliegue de una red de sensores inalámbricos queda totalmente fuera del alcance de este proyecto y por tanto no se va a entrar en ningún tipo de detalle técnico sobre estos sistemas. El objetivo de este proyecto es evaluar algoritmos de inferencia difusa sobre los valores que serían entregados al sistema de decisión en una aplicación real . 

Las mediciones realizadas por la red de sensores constituyen las entradas del sistema de inferencia y decisión. Estas entradas son, originalmente, valores escalares que son transformados por el sistema de lógica difusa en valores difusos (por medio de variables lingüísticas). Esto permite realizar, sobre estas mediciones, un proceso de razonamiento aproximado que puede tratar mejor las imprecisiones, en comparación con utilizar dichos valores escalares directamente.


\section{Magnitudes medidas y variables lingüísticas}

Las entradas del sistema difuso son valores escalares de magnitudes que describen el incendio forestal, tales como la temperatura o la luminosidad, medidas por la red de sensores. Para cada una de estas magnitudes se define una variable lingüística con tres valores posibles (definidas por sus correspondientes funciones de pertenencia). Así pues, los valores escalares de entrada serán transformados en conjuntos difusos que posteriormente serán comparados con las variables lingüísticas según las reglas definidas, para obtener como salida el riesgo de incendio (que es también una variable lingüística). Estas variables lingüísticas así como sus correspondientes universos (rangos de valores posibles) son:

\begin{enumerate}[label=($\chi_\arabic*$),ref=(X\arabic*)]
   \item \bfseries Temperatura: \normalfont medida en grados centígrados (0ºC a 120ºC).
   \item \bfseries Humo: \normalfont medida en partes por millón (0 a 100ppm).
   \item \bfseries Luz: \normalfont medida en lux (0 a 1000 lux).
   \item \bfseries Humedad: \normalfont medida en partes por millón (0 a 100ppm).
   \item \bfseries Distancia: \normalfont medida en metros (0 a 80m).
\end{enumerate}

Para cada una de estas variables lingüísticas se definen tres valores posibles: \emph{Baja} (L,\emph{Low}), \emph{Media} (M, \emph{Medium}) y \emph{Alta} (H, \emph{High}) que los valores tienen el sentido de \emph{Cerca} (C, \emph{Close}), \emph{Media} (M, \emph{Medium}) y \emph{Lejos} (F, \emph{Far}) respectivamente.

La salida del sistema viene dada por la variable lingüística $y = $ \emph{``Riesgo de incendio''}, que determina el riesgo de incendio forestal en una escala del 0 al 100 (\%). Esta variable lingüística puede tomar los valores: \emph{Muy bajo} (VL, \emph{Very Low}), \emph{Bajo} (L, \emph{Low}), \emph{Medio} (M, \emph{Medium}), \emph{Alto} (H, \emph{High}) y \emph{Muy Alto} (VH, \emph{Very High}). 

En la figura \ref{fig:fire-detection-lang-variables} se representan las variables lingüística descritas anteriormente. Como se puede ver, se han elegido funciones de pertenencia lineales para modelar los valores de las variables lingüísticas.

\begin{figure}[t]
	\centering
	\begin{subfigure}[b]{0.45\textwidth}
		\setlength\figureheight{2.5cm}
		\setlength\figurewidth{6cm}
		% This file was created by matlab2tikz v0.4.7 (commit aa5a2b3a09d5219f83b03f452820de20dfb2a2eb) running on MATLAB 8.0.
% Copyright (c) 2008--2014, Nico Schlömer <nico.schloemer@gmail.com>
% All rights reserved.
% Minimal pgfplots version: 1.3
% 
\begin{tikzpicture}

\begin{axis}[%
width=\figurewidth,
height=\figureheight,
clip=false,
scale only axis,
xmin=0,
xmax=120,
ymin=0,
ymax=1
]
\addplot [color=black,solid,line width=1.4pt,forget plot]
  table[row sep=crcr]{0	1\\
1	0.98\\
2	0.96\\
3	0.94\\
4	0.92\\
5	0.9\\
6	0.88\\
7	0.86\\
8	0.84\\
9	0.82\\
10	0.8\\
11	0.78\\
12	0.76\\
13	0.74\\
14	0.72\\
15	0.7\\
16	0.68\\
17	0.66\\
18	0.64\\
19	0.62\\
20	0.6\\
21	0.58\\
22	0.56\\
23	0.54\\
24	0.52\\
25	0.5\\
26	0.48\\
27	0.46\\
28	0.44\\
29	0.42\\
30	0.4\\
31	0.38\\
32	0.36\\
33	0.34\\
34	0.32\\
35	0.3\\
36	0.28\\
37	0.26\\
38	0.24\\
39	0.22\\
40	0.2\\
41	0.18\\
42	0.16\\
43	0.14\\
44	0.12\\
45	0.1\\
46	0.08\\
47	0.0600000000000001\\
48	0.04\\
49	0.02\\
50	0\\
51	0\\
52	0\\
53	0\\
54	0\\
55	0\\
56	0\\
57	0\\
58	0\\
59	0\\
60	0\\
61	0\\
62	0\\
63	0\\
64	0\\
65	0\\
66	0\\
67	0\\
68	0\\
69	0\\
70	0\\
71	0\\
72	0\\
73	0\\
74	0\\
75	0\\
76	0\\
77	0\\
78	0\\
79	0\\
80	0\\
81	0\\
82	0\\
83	0\\
84	0\\
85	0\\
86	0\\
87	0\\
88	0\\
89	0\\
90	0\\
91	0\\
92	0\\
93	0\\
94	0\\
95	0\\
96	0\\
97	0\\
98	0\\
99	0\\
100	0\\
101	0\\
102	0\\
103	0\\
104	0\\
105	0\\
106	0\\
107	0\\
108	0\\
109	0\\
110	0\\
111	0\\
112	0\\
113	0\\
114	0\\
115	0\\
116	0\\
117	0\\
118	0\\
119	0\\
120	0\\
};
\addplot [color=black,solid,line width=1.4pt,forget plot]
  table[row sep=crcr]{0	0\\
1	0\\
2	0\\
3	0\\
4	0\\
5	0\\
6	0\\
7	0\\
8	0\\
9	0\\
10	0\\
11	0.02\\
12	0.04\\
13	0.06\\
14	0.08\\
15	0.1\\
16	0.12\\
17	0.14\\
18	0.16\\
19	0.18\\
20	0.2\\
21	0.22\\
22	0.24\\
23	0.26\\
24	0.28\\
25	0.3\\
26	0.32\\
27	0.34\\
28	0.36\\
29	0.38\\
30	0.4\\
31	0.42\\
32	0.44\\
33	0.46\\
34	0.48\\
35	0.5\\
36	0.52\\
37	0.54\\
38	0.56\\
39	0.58\\
40	0.6\\
41	0.62\\
42	0.64\\
43	0.66\\
44	0.68\\
45	0.7\\
46	0.72\\
47	0.74\\
48	0.76\\
49	0.78\\
50	0.8\\
51	0.82\\
52	0.84\\
53	0.86\\
54	0.88\\
55	0.9\\
56	0.92\\
57	0.94\\
58	0.96\\
59	0.98\\
60	1\\
61	0.98\\
62	0.96\\
63	0.94\\
64	0.92\\
65	0.9\\
66	0.88\\
67	0.86\\
68	0.84\\
69	0.82\\
70	0.8\\
71	0.78\\
72	0.76\\
73	0.74\\
74	0.72\\
75	0.7\\
76	0.68\\
77	0.66\\
78	0.64\\
79	0.62\\
80	0.6\\
81	0.58\\
82	0.56\\
83	0.54\\
84	0.52\\
85	0.5\\
86	0.48\\
87	0.46\\
88	0.44\\
89	0.42\\
90	0.4\\
91	0.38\\
92	0.36\\
93	0.34\\
94	0.32\\
95	0.3\\
96	0.28\\
97	0.26\\
98	0.24\\
99	0.22\\
100	0.2\\
101	0.18\\
102	0.16\\
103	0.14\\
104	0.12\\
105	0.1\\
106	0.08\\
107	0.0600000000000001\\
108	0.04\\
109	0.02\\
110	0\\
111	0\\
112	0\\
113	0\\
114	0\\
115	0\\
116	0\\
117	0\\
118	0\\
119	0\\
120	0\\
};
\addplot [color=black,solid,line width=1.4pt,forget plot]
  table[row sep=crcr]{0	0\\
1	0\\
2	0\\
3	0\\
4	0\\
5	0\\
6	0\\
7	0\\
8	0\\
9	0\\
10	0\\
11	0\\
12	0\\
13	0\\
14	0\\
15	0\\
16	0\\
17	0\\
18	0\\
19	0\\
20	0\\
21	0\\
22	0\\
23	0\\
24	0\\
25	0\\
26	0\\
27	0\\
28	0\\
29	0\\
30	0\\
31	0\\
32	0\\
33	0\\
34	0\\
35	0\\
36	0\\
37	0\\
38	0\\
39	0\\
40	0\\
41	0\\
42	0\\
43	0\\
44	0\\
45	0\\
46	0\\
47	0\\
48	0\\
49	0\\
50	0\\
51	0\\
52	0\\
53	0\\
54	0\\
55	0\\
56	0\\
57	0\\
58	0\\
59	0\\
60	0\\
61	0\\
62	0\\
63	0\\
64	0\\
65	0\\
66	0\\
67	0\\
68	0\\
69	0\\
70	0\\
71	0.02\\
72	0.04\\
73	0.06\\
74	0.08\\
75	0.1\\
76	0.12\\
77	0.14\\
78	0.16\\
79	0.18\\
80	0.2\\
81	0.22\\
82	0.24\\
83	0.26\\
84	0.28\\
85	0.3\\
86	0.32\\
87	0.34\\
88	0.36\\
89	0.38\\
90	0.4\\
91	0.42\\
92	0.44\\
93	0.46\\
94	0.48\\
95	0.5\\
96	0.52\\
97	0.54\\
98	0.56\\
99	0.58\\
100	0.6\\
101	0.62\\
102	0.64\\
103	0.66\\
104	0.68\\
105	0.7\\
106	0.72\\
107	0.74\\
108	0.76\\
109	0.78\\
110	0.8\\
111	0.82\\
112	0.84\\
113	0.86\\
114	0.88\\
115	0.9\\
116	0.92\\
117	0.94\\
118	0.96\\
119	0.98\\
120	1\\
};
\node[right, inner sep=0mm, text=black]
at (axis cs:1,1.1,0) {Baja (L)};
\node[right, inner sep=0mm, text=black]
at (axis cs:50,1.1,0) {Media (M)};
\node[right, inner sep=0mm, text=black]
at (axis cs:105,1.1,0) {Alta (H)};
\end{axis}
\end{tikzpicture}%
		\caption{$\chi_1$ - Temperatura (ºC)}
		\label{fig:temp-lang-variable}
	\end{subfigure}
	\qquad
	\begin{subfigure}[b]{0.45\textwidth}
		\setlength\figureheight{2.5cm}
		\setlength\figurewidth{6cm}
		% This file was created by matlab2tikz v0.4.7 (commit aa5a2b3a09d5219f83b03f452820de20dfb2a2eb) running on MATLAB 8.0.
% Copyright (c) 2008--2014, Nico Schlömer <nico.schloemer@gmail.com>
% All rights reserved.
% Minimal pgfplots version: 1.3
% 
\begin{tikzpicture}

\begin{axis}[%
width=\figurewidth,
height=\figureheight,
clip=false,
scale only axis,
xmin=0,
xmax=100,
ymin=0,
ymax=1
]
\addplot [color=black,solid,line width=1.4pt,forget plot]
  table[row sep=crcr]{0	1\\
1	0.975\\
2	0.95\\
3	0.925\\
4	0.9\\
5	0.875\\
6	0.85\\
7	0.825\\
8	0.8\\
9	0.775\\
10	0.75\\
11	0.725\\
12	0.7\\
13	0.675\\
14	0.65\\
15	0.625\\
16	0.6\\
17	0.575\\
18	0.55\\
19	0.525\\
20	0.5\\
21	0.475\\
22	0.45\\
23	0.425\\
24	0.4\\
25	0.375\\
26	0.35\\
27	0.325\\
28	0.3\\
29	0.275\\
30	0.25\\
31	0.225\\
32	0.2\\
33	0.175\\
34	0.15\\
35	0.125\\
36	0.1\\
37	0.075\\
38	0.05\\
39	0.025\\
40	0\\
41	0\\
42	0\\
43	0\\
44	0\\
45	0\\
46	0\\
47	0\\
48	0\\
49	0\\
50	0\\
51	0\\
52	0\\
53	0\\
54	0\\
55	0\\
56	0\\
57	0\\
58	0\\
59	0\\
60	0\\
61	0\\
62	0\\
63	0\\
64	0\\
65	0\\
66	0\\
67	0\\
68	0\\
69	0\\
70	0\\
71	0\\
72	0\\
73	0\\
74	0\\
75	0\\
76	0\\
77	0\\
78	0\\
79	0\\
80	0\\
81	0\\
82	0\\
83	0\\
84	0\\
85	0\\
86	0\\
87	0\\
88	0\\
89	0\\
90	0\\
91	0\\
92	0\\
93	0\\
94	0\\
95	0\\
96	0\\
97	0\\
98	0\\
99	0\\
100	0\\
};
\addplot [color=black,solid,line width=1.4pt,forget plot]
  table[row sep=crcr]{0	0\\
1	0\\
2	0\\
3	0\\
4	0\\
5	0\\
6	0\\
7	0\\
8	0\\
9	0\\
10	0\\
11	0.025\\
12	0.05\\
13	0.075\\
14	0.1\\
15	0.125\\
16	0.15\\
17	0.175\\
18	0.2\\
19	0.225\\
20	0.25\\
21	0.275\\
22	0.3\\
23	0.325\\
24	0.35\\
25	0.375\\
26	0.4\\
27	0.425\\
28	0.45\\
29	0.475\\
30	0.5\\
31	0.525\\
32	0.55\\
33	0.575\\
34	0.6\\
35	0.625\\
36	0.65\\
37	0.675\\
38	0.7\\
39	0.725\\
40	0.75\\
41	0.775\\
42	0.8\\
43	0.825\\
44	0.85\\
45	0.875\\
46	0.9\\
47	0.925\\
48	0.95\\
49	0.975\\
50	1\\
51	0.975\\
52	0.95\\
53	0.925\\
54	0.9\\
55	0.875\\
56	0.85\\
57	0.825\\
58	0.8\\
59	0.775\\
60	0.75\\
61	0.725\\
62	0.7\\
63	0.675\\
64	0.65\\
65	0.625\\
66	0.6\\
67	0.575\\
68	0.55\\
69	0.525\\
70	0.5\\
71	0.475\\
72	0.45\\
73	0.425\\
74	0.4\\
75	0.375\\
76	0.35\\
77	0.325\\
78	0.3\\
79	0.275\\
80	0.25\\
81	0.225\\
82	0.2\\
83	0.175\\
84	0.15\\
85	0.125\\
86	0.1\\
87	0.075\\
88	0.05\\
89	0.025\\
90	0\\
91	0\\
92	0\\
93	0\\
94	0\\
95	0\\
96	0\\
97	0\\
98	0\\
99	0\\
100	0\\
};
\addplot [color=black,solid,line width=1.4pt,forget plot]
  table[row sep=crcr]{0	0\\
1	0\\
2	0\\
3	0\\
4	0\\
5	0\\
6	0\\
7	0\\
8	0\\
9	0\\
10	0\\
11	0\\
12	0\\
13	0\\
14	0\\
15	0\\
16	0\\
17	0\\
18	0\\
19	0\\
20	0\\
21	0\\
22	0\\
23	0\\
24	0\\
25	0\\
26	0\\
27	0\\
28	0\\
29	0\\
30	0\\
31	0\\
32	0\\
33	0\\
34	0\\
35	0\\
36	0\\
37	0\\
38	0\\
39	0\\
40	0\\
41	0\\
42	0\\
43	0\\
44	0\\
45	0\\
46	0\\
47	0\\
48	0\\
49	0\\
50	0\\
51	0\\
52	0\\
53	0\\
54	0\\
55	0\\
56	0\\
57	0\\
58	0\\
59	0\\
60	0\\
61	0.025\\
62	0.05\\
63	0.075\\
64	0.1\\
65	0.125\\
66	0.15\\
67	0.175\\
68	0.2\\
69	0.225\\
70	0.25\\
71	0.275\\
72	0.3\\
73	0.325\\
74	0.35\\
75	0.375\\
76	0.4\\
77	0.425\\
78	0.45\\
79	0.475\\
80	0.5\\
81	0.525\\
82	0.55\\
83	0.575\\
84	0.6\\
85	0.625\\
86	0.65\\
87	0.675\\
88	0.7\\
89	0.725\\
90	0.75\\
91	0.775\\
92	0.8\\
93	0.825\\
94	0.85\\
95	0.875\\
96	0.9\\
97	0.925\\
98	0.95\\
99	0.975\\
100	1\\
};
\node[right, inner sep=0mm, text=black]
at (axis cs:1,1.1,0) {Bajo (L)};
\node[right, inner sep=0mm, text=black]
at (axis cs:45,1.1,0) {Medio (M)};
\node[right, inner sep=0mm, text=black]
at (axis cs:90,1.1,0) {Alto (H)};
\end{axis}
\end{tikzpicture}%
		\caption{$\chi_2$ - Humo (ppm)}
		\label{fig:smoke-lang-variable}
	\end{subfigure}
	
	\vspace{1 cm}
	\begin{subfigure}[b]{0.45\textwidth}
		\setlength\figureheight{2.5cm}
		\setlength\figurewidth{6cm}
		% This file was created by matlab2tikz v0.4.7 (commit aa5a2b3a09d5219f83b03f452820de20dfb2a2eb) running on MATLAB 8.0.
% Copyright (c) 2008--2014, Nico Schlömer <nico.schloemer@gmail.com>
% All rights reserved.
% Minimal pgfplots version: 1.3
% 
\begin{tikzpicture}

\begin{axis}[%
width=\figurewidth,
height=\figureheight,
clip=false,
scale only axis,
xmin=0,
xmax=1000,
ymin=0,
ymax=1
]
\addplot [color=black,solid,line width=1.4pt,forget plot]
  table[row sep=crcr]{0	1\\
1	0.9975\\
2	0.995\\
3	0.9925\\
4	0.99\\
5	0.9875\\
6	0.985\\
7	0.9825\\
8	0.98\\
9	0.9775\\
10	0.975\\
11	0.9725\\
12	0.97\\
13	0.9675\\
14	0.965\\
15	0.9625\\
16	0.96\\
17	0.9575\\
18	0.955\\
19	0.9525\\
20	0.95\\
21	0.9475\\
22	0.945\\
23	0.9425\\
24	0.94\\
25	0.9375\\
26	0.935\\
27	0.9325\\
28	0.93\\
29	0.9275\\
30	0.925\\
31	0.9225\\
32	0.92\\
33	0.9175\\
34	0.915\\
35	0.9125\\
36	0.91\\
37	0.9075\\
38	0.905\\
39	0.9025\\
40	0.9\\
41	0.8975\\
42	0.895\\
43	0.8925\\
44	0.89\\
45	0.8875\\
46	0.885\\
47	0.8825\\
48	0.88\\
49	0.8775\\
50	0.875\\
51	0.8725\\
52	0.87\\
53	0.8675\\
54	0.865\\
55	0.8625\\
56	0.86\\
57	0.8575\\
58	0.855\\
59	0.8525\\
60	0.85\\
61	0.8475\\
62	0.845\\
63	0.8425\\
64	0.84\\
65	0.8375\\
66	0.835\\
67	0.8325\\
68	0.83\\
69	0.8275\\
70	0.825\\
71	0.8225\\
72	0.82\\
73	0.8175\\
74	0.815\\
75	0.8125\\
76	0.81\\
77	0.8075\\
78	0.805\\
79	0.8025\\
80	0.8\\
81	0.7975\\
82	0.795\\
83	0.7925\\
84	0.79\\
85	0.7875\\
86	0.785\\
87	0.7825\\
88	0.78\\
89	0.7775\\
90	0.775\\
91	0.7725\\
92	0.77\\
93	0.7675\\
94	0.765\\
95	0.7625\\
96	0.76\\
97	0.7575\\
98	0.755\\
99	0.7525\\
100	0.75\\
101	0.7475\\
102	0.745\\
103	0.7425\\
104	0.74\\
105	0.7375\\
106	0.735\\
107	0.7325\\
108	0.73\\
109	0.7275\\
110	0.725\\
111	0.7225\\
112	0.72\\
113	0.7175\\
114	0.715\\
115	0.7125\\
116	0.71\\
117	0.7075\\
118	0.705\\
119	0.7025\\
120	0.7\\
121	0.6975\\
122	0.695\\
123	0.6925\\
124	0.69\\
125	0.6875\\
126	0.685\\
127	0.6825\\
128	0.68\\
129	0.6775\\
130	0.675\\
131	0.6725\\
132	0.67\\
133	0.6675\\
134	0.665\\
135	0.6625\\
136	0.66\\
137	0.6575\\
138	0.655\\
139	0.6525\\
140	0.65\\
141	0.6475\\
142	0.645\\
143	0.6425\\
144	0.64\\
145	0.6375\\
146	0.635\\
147	0.6325\\
148	0.63\\
149	0.6275\\
150	0.625\\
151	0.6225\\
152	0.62\\
153	0.6175\\
154	0.615\\
155	0.6125\\
156	0.61\\
157	0.6075\\
158	0.605\\
159	0.6025\\
160	0.6\\
161	0.5975\\
162	0.595\\
163	0.5925\\
164	0.59\\
165	0.5875\\
166	0.585\\
167	0.5825\\
168	0.58\\
169	0.5775\\
170	0.575\\
171	0.5725\\
172	0.57\\
173	0.5675\\
174	0.565\\
175	0.5625\\
176	0.56\\
177	0.5575\\
178	0.555\\
179	0.5525\\
180	0.55\\
181	0.5475\\
182	0.545\\
183	0.5425\\
184	0.54\\
185	0.5375\\
186	0.535\\
187	0.5325\\
188	0.53\\
189	0.5275\\
190	0.525\\
191	0.5225\\
192	0.52\\
193	0.5175\\
194	0.515\\
195	0.5125\\
196	0.51\\
197	0.5075\\
198	0.505\\
199	0.5025\\
200	0.5\\
201	0.4975\\
202	0.495\\
203	0.4925\\
204	0.49\\
205	0.4875\\
206	0.485\\
207	0.4825\\
208	0.48\\
209	0.4775\\
210	0.475\\
211	0.4725\\
212	0.47\\
213	0.4675\\
214	0.465\\
215	0.4625\\
216	0.46\\
217	0.4575\\
218	0.455\\
219	0.4525\\
220	0.45\\
221	0.4475\\
222	0.445\\
223	0.4425\\
224	0.44\\
225	0.4375\\
226	0.435\\
227	0.4325\\
228	0.43\\
229	0.4275\\
230	0.425\\
231	0.4225\\
232	0.42\\
233	0.4175\\
234	0.415\\
235	0.4125\\
236	0.41\\
237	0.4075\\
238	0.405\\
239	0.4025\\
240	0.4\\
241	0.3975\\
242	0.395\\
243	0.3925\\
244	0.39\\
245	0.3875\\
246	0.385\\
247	0.3825\\
248	0.38\\
249	0.3775\\
250	0.375\\
251	0.3725\\
252	0.37\\
253	0.3675\\
254	0.365\\
255	0.3625\\
256	0.36\\
257	0.3575\\
258	0.355\\
259	0.3525\\
260	0.35\\
261	0.3475\\
262	0.345\\
263	0.3425\\
264	0.34\\
265	0.3375\\
266	0.335\\
267	0.3325\\
268	0.33\\
269	0.3275\\
270	0.325\\
271	0.3225\\
272	0.32\\
273	0.3175\\
274	0.315\\
275	0.3125\\
276	0.31\\
277	0.3075\\
278	0.305\\
279	0.3025\\
280	0.3\\
281	0.2975\\
282	0.295\\
283	0.2925\\
284	0.29\\
285	0.2875\\
286	0.285\\
287	0.2825\\
288	0.28\\
289	0.2775\\
290	0.275\\
291	0.2725\\
292	0.27\\
293	0.2675\\
294	0.265\\
295	0.2625\\
296	0.26\\
297	0.2575\\
298	0.255\\
299	0.2525\\
300	0.25\\
301	0.2475\\
302	0.245\\
303	0.2425\\
304	0.24\\
305	0.2375\\
306	0.235\\
307	0.2325\\
308	0.23\\
309	0.2275\\
310	0.225\\
311	0.2225\\
312	0.22\\
313	0.2175\\
314	0.215\\
315	0.2125\\
316	0.21\\
317	0.2075\\
318	0.205\\
319	0.2025\\
320	0.2\\
321	0.1975\\
322	0.195\\
323	0.1925\\
324	0.19\\
325	0.1875\\
326	0.185\\
327	0.1825\\
328	0.18\\
329	0.1775\\
330	0.175\\
331	0.1725\\
332	0.17\\
333	0.1675\\
334	0.165\\
335	0.1625\\
336	0.16\\
337	0.1575\\
338	0.155\\
339	0.1525\\
340	0.15\\
341	0.1475\\
342	0.145\\
343	0.1425\\
344	0.14\\
345	0.1375\\
346	0.135\\
347	0.1325\\
348	0.13\\
349	0.1275\\
350	0.125\\
351	0.1225\\
352	0.12\\
353	0.1175\\
354	0.115\\
355	0.1125\\
356	0.11\\
357	0.1075\\
358	0.105\\
359	0.1025\\
360	0.1\\
361	0.0975\\
362	0.095\\
363	0.0925\\
364	0.09\\
365	0.0875\\
366	0.085\\
367	0.0825\\
368	0.08\\
369	0.0775\\
370	0.075\\
371	0.0725\\
372	0.07\\
373	0.0675\\
374	0.0649999999999999\\
375	0.0625\\
376	0.0600000000000001\\
377	0.0575\\
378	0.055\\
379	0.0525\\
380	0.05\\
381	0.0475\\
382	0.045\\
383	0.0425\\
384	0.04\\
385	0.0375\\
386	0.035\\
387	0.0325\\
388	0.03\\
389	0.0275\\
390	0.025\\
391	0.0225\\
392	0.02\\
393	0.0175\\
394	0.015\\
395	0.0125\\
396	0.01\\
397	0.00749999999999995\\
398	0.005\\
399	0.00249999999999995\\
400	0\\
401	0\\
402	0\\
403	0\\
404	0\\
405	0\\
406	0\\
407	0\\
408	0\\
409	0\\
410	0\\
411	0\\
412	0\\
413	0\\
414	0\\
415	0\\
416	0\\
417	0\\
418	0\\
419	0\\
420	0\\
421	0\\
422	0\\
423	0\\
424	0\\
425	0\\
426	0\\
427	0\\
428	0\\
429	0\\
430	0\\
431	0\\
432	0\\
433	0\\
434	0\\
435	0\\
436	0\\
437	0\\
438	0\\
439	0\\
440	0\\
441	0\\
442	0\\
443	0\\
444	0\\
445	0\\
446	0\\
447	0\\
448	0\\
449	0\\
450	0\\
451	0\\
452	0\\
453	0\\
454	0\\
455	0\\
456	0\\
457	0\\
458	0\\
459	0\\
460	0\\
461	0\\
462	0\\
463	0\\
464	0\\
465	0\\
466	0\\
467	0\\
468	0\\
469	0\\
470	0\\
471	0\\
472	0\\
473	0\\
474	0\\
475	0\\
476	0\\
477	0\\
478	0\\
479	0\\
480	0\\
481	0\\
482	0\\
483	0\\
484	0\\
485	0\\
486	0\\
487	0\\
488	0\\
489	0\\
490	0\\
491	0\\
492	0\\
493	0\\
494	0\\
495	0\\
496	0\\
497	0\\
498	0\\
499	0\\
500	0\\
501	0\\
502	0\\
503	0\\
504	0\\
505	0\\
506	0\\
507	0\\
508	0\\
509	0\\
510	0\\
511	0\\
512	0\\
513	0\\
514	0\\
515	0\\
516	0\\
517	0\\
518	0\\
519	0\\
520	0\\
521	0\\
522	0\\
523	0\\
524	0\\
525	0\\
526	0\\
527	0\\
528	0\\
529	0\\
530	0\\
531	0\\
532	0\\
533	0\\
534	0\\
535	0\\
536	0\\
537	0\\
538	0\\
539	0\\
540	0\\
541	0\\
542	0\\
543	0\\
544	0\\
545	0\\
546	0\\
547	0\\
548	0\\
549	0\\
550	0\\
551	0\\
552	0\\
553	0\\
554	0\\
555	0\\
556	0\\
557	0\\
558	0\\
559	0\\
560	0\\
561	0\\
562	0\\
563	0\\
564	0\\
565	0\\
566	0\\
567	0\\
568	0\\
569	0\\
570	0\\
571	0\\
572	0\\
573	0\\
574	0\\
575	0\\
576	0\\
577	0\\
578	0\\
579	0\\
580	0\\
581	0\\
582	0\\
583	0\\
584	0\\
585	0\\
586	0\\
587	0\\
588	0\\
589	0\\
590	0\\
591	0\\
592	0\\
593	0\\
594	0\\
595	0\\
596	0\\
597	0\\
598	0\\
599	0\\
600	0\\
601	0\\
602	0\\
603	0\\
604	0\\
605	0\\
606	0\\
607	0\\
608	0\\
609	0\\
610	0\\
611	0\\
612	0\\
613	0\\
614	0\\
615	0\\
616	0\\
617	0\\
618	0\\
619	0\\
620	0\\
621	0\\
622	0\\
623	0\\
624	0\\
625	0\\
626	0\\
627	0\\
628	0\\
629	0\\
630	0\\
631	0\\
632	0\\
633	0\\
634	0\\
635	0\\
636	0\\
637	0\\
638	0\\
639	0\\
640	0\\
641	0\\
642	0\\
643	0\\
644	0\\
645	0\\
646	0\\
647	0\\
648	0\\
649	0\\
650	0\\
651	0\\
652	0\\
653	0\\
654	0\\
655	0\\
656	0\\
657	0\\
658	0\\
659	0\\
660	0\\
661	0\\
662	0\\
663	0\\
664	0\\
665	0\\
666	0\\
667	0\\
668	0\\
669	0\\
670	0\\
671	0\\
672	0\\
673	0\\
674	0\\
675	0\\
676	0\\
677	0\\
678	0\\
679	0\\
680	0\\
681	0\\
682	0\\
683	0\\
684	0\\
685	0\\
686	0\\
687	0\\
688	0\\
689	0\\
690	0\\
691	0\\
692	0\\
693	0\\
694	0\\
695	0\\
696	0\\
697	0\\
698	0\\
699	0\\
700	0\\
701	0\\
702	0\\
703	0\\
704	0\\
705	0\\
706	0\\
707	0\\
708	0\\
709	0\\
710	0\\
711	0\\
712	0\\
713	0\\
714	0\\
715	0\\
716	0\\
717	0\\
718	0\\
719	0\\
720	0\\
721	0\\
722	0\\
723	0\\
724	0\\
725	0\\
726	0\\
727	0\\
728	0\\
729	0\\
730	0\\
731	0\\
732	0\\
733	0\\
734	0\\
735	0\\
736	0\\
737	0\\
738	0\\
739	0\\
740	0\\
741	0\\
742	0\\
743	0\\
744	0\\
745	0\\
746	0\\
747	0\\
748	0\\
749	0\\
750	0\\
751	0\\
752	0\\
753	0\\
754	0\\
755	0\\
756	0\\
757	0\\
758	0\\
759	0\\
760	0\\
761	0\\
762	0\\
763	0\\
764	0\\
765	0\\
766	0\\
767	0\\
768	0\\
769	0\\
770	0\\
771	0\\
772	0\\
773	0\\
774	0\\
775	0\\
776	0\\
777	0\\
778	0\\
779	0\\
780	0\\
781	0\\
782	0\\
783	0\\
784	0\\
785	0\\
786	0\\
787	0\\
788	0\\
789	0\\
790	0\\
791	0\\
792	0\\
793	0\\
794	0\\
795	0\\
796	0\\
797	0\\
798	0\\
799	0\\
800	0\\
801	0\\
802	0\\
803	0\\
804	0\\
805	0\\
806	0\\
807	0\\
808	0\\
809	0\\
810	0\\
811	0\\
812	0\\
813	0\\
814	0\\
815	0\\
816	0\\
817	0\\
818	0\\
819	0\\
820	0\\
821	0\\
822	0\\
823	0\\
824	0\\
825	0\\
826	0\\
827	0\\
828	0\\
829	0\\
830	0\\
831	0\\
832	0\\
833	0\\
834	0\\
835	0\\
836	0\\
837	0\\
838	0\\
839	0\\
840	0\\
841	0\\
842	0\\
843	0\\
844	0\\
845	0\\
846	0\\
847	0\\
848	0\\
849	0\\
850	0\\
851	0\\
852	0\\
853	0\\
854	0\\
855	0\\
856	0\\
857	0\\
858	0\\
859	0\\
860	0\\
861	0\\
862	0\\
863	0\\
864	0\\
865	0\\
866	0\\
867	0\\
868	0\\
869	0\\
870	0\\
871	0\\
872	0\\
873	0\\
874	0\\
875	0\\
876	0\\
877	0\\
878	0\\
879	0\\
880	0\\
881	0\\
882	0\\
883	0\\
884	0\\
885	0\\
886	0\\
887	0\\
888	0\\
889	0\\
890	0\\
891	0\\
892	0\\
893	0\\
894	0\\
895	0\\
896	0\\
897	0\\
898	0\\
899	0\\
900	0\\
901	0\\
902	0\\
903	0\\
904	0\\
905	0\\
906	0\\
907	0\\
908	0\\
909	0\\
910	0\\
911	0\\
912	0\\
913	0\\
914	0\\
915	0\\
916	0\\
917	0\\
918	0\\
919	0\\
920	0\\
921	0\\
922	0\\
923	0\\
924	0\\
925	0\\
926	0\\
927	0\\
928	0\\
929	0\\
930	0\\
931	0\\
932	0\\
933	0\\
934	0\\
935	0\\
936	0\\
937	0\\
938	0\\
939	0\\
940	0\\
941	0\\
942	0\\
943	0\\
944	0\\
945	0\\
946	0\\
947	0\\
948	0\\
949	0\\
950	0\\
951	0\\
952	0\\
953	0\\
954	0\\
955	0\\
956	0\\
957	0\\
958	0\\
959	0\\
960	0\\
961	0\\
962	0\\
963	0\\
964	0\\
965	0\\
966	0\\
967	0\\
968	0\\
969	0\\
970	0\\
971	0\\
972	0\\
973	0\\
974	0\\
975	0\\
976	0\\
977	0\\
978	0\\
979	0\\
980	0\\
981	0\\
982	0\\
983	0\\
984	0\\
985	0\\
986	0\\
987	0\\
988	0\\
989	0\\
990	0\\
991	0\\
992	0\\
993	0\\
994	0\\
995	0\\
996	0\\
997	0\\
998	0\\
999	0\\
1000	0\\
};
\addplot [color=black,solid,line width=1.4pt,forget plot]
  table[row sep=crcr]{0	0\\
1	0\\
2	0\\
3	0\\
4	0\\
5	0\\
6	0\\
7	0\\
8	0\\
9	0\\
10	0\\
11	0\\
12	0\\
13	0\\
14	0\\
15	0\\
16	0\\
17	0\\
18	0\\
19	0\\
20	0\\
21	0\\
22	0\\
23	0\\
24	0\\
25	0\\
26	0\\
27	0\\
28	0\\
29	0\\
30	0\\
31	0\\
32	0\\
33	0\\
34	0\\
35	0\\
36	0\\
37	0\\
38	0\\
39	0\\
40	0\\
41	0\\
42	0\\
43	0\\
44	0\\
45	0\\
46	0\\
47	0\\
48	0\\
49	0\\
50	0\\
51	0\\
52	0\\
53	0\\
54	0\\
55	0\\
56	0\\
57	0\\
58	0\\
59	0\\
60	0\\
61	0\\
62	0\\
63	0\\
64	0\\
65	0\\
66	0\\
67	0\\
68	0\\
69	0\\
70	0\\
71	0\\
72	0\\
73	0\\
74	0\\
75	0\\
76	0\\
77	0\\
78	0\\
79	0\\
80	0\\
81	0\\
82	0\\
83	0\\
84	0\\
85	0\\
86	0\\
87	0\\
88	0\\
89	0\\
90	0\\
91	0\\
92	0\\
93	0\\
94	0\\
95	0\\
96	0\\
97	0\\
98	0\\
99	0\\
100	0\\
101	0.0025\\
102	0.005\\
103	0.0075\\
104	0.01\\
105	0.0125\\
106	0.015\\
107	0.0175\\
108	0.02\\
109	0.0225\\
110	0.025\\
111	0.0275\\
112	0.03\\
113	0.0325\\
114	0.035\\
115	0.0375\\
116	0.04\\
117	0.0425\\
118	0.045\\
119	0.0475\\
120	0.05\\
121	0.0525\\
122	0.055\\
123	0.0575\\
124	0.06\\
125	0.0625\\
126	0.065\\
127	0.0675\\
128	0.07\\
129	0.0725\\
130	0.075\\
131	0.0775\\
132	0.08\\
133	0.0825\\
134	0.085\\
135	0.0875\\
136	0.09\\
137	0.0925\\
138	0.095\\
139	0.0975\\
140	0.1\\
141	0.1025\\
142	0.105\\
143	0.1075\\
144	0.11\\
145	0.1125\\
146	0.115\\
147	0.1175\\
148	0.12\\
149	0.1225\\
150	0.125\\
151	0.1275\\
152	0.13\\
153	0.1325\\
154	0.135\\
155	0.1375\\
156	0.14\\
157	0.1425\\
158	0.145\\
159	0.1475\\
160	0.15\\
161	0.1525\\
162	0.155\\
163	0.1575\\
164	0.16\\
165	0.1625\\
166	0.165\\
167	0.1675\\
168	0.17\\
169	0.1725\\
170	0.175\\
171	0.1775\\
172	0.18\\
173	0.1825\\
174	0.185\\
175	0.1875\\
176	0.19\\
177	0.1925\\
178	0.195\\
179	0.1975\\
180	0.2\\
181	0.2025\\
182	0.205\\
183	0.2075\\
184	0.21\\
185	0.2125\\
186	0.215\\
187	0.2175\\
188	0.22\\
189	0.2225\\
190	0.225\\
191	0.2275\\
192	0.23\\
193	0.2325\\
194	0.235\\
195	0.2375\\
196	0.24\\
197	0.2425\\
198	0.245\\
199	0.2475\\
200	0.25\\
201	0.2525\\
202	0.255\\
203	0.2575\\
204	0.26\\
205	0.2625\\
206	0.265\\
207	0.2675\\
208	0.27\\
209	0.2725\\
210	0.275\\
211	0.2775\\
212	0.28\\
213	0.2825\\
214	0.285\\
215	0.2875\\
216	0.29\\
217	0.2925\\
218	0.295\\
219	0.2975\\
220	0.3\\
221	0.3025\\
222	0.305\\
223	0.3075\\
224	0.31\\
225	0.3125\\
226	0.315\\
227	0.3175\\
228	0.32\\
229	0.3225\\
230	0.325\\
231	0.3275\\
232	0.33\\
233	0.3325\\
234	0.335\\
235	0.3375\\
236	0.34\\
237	0.3425\\
238	0.345\\
239	0.3475\\
240	0.35\\
241	0.3525\\
242	0.355\\
243	0.3575\\
244	0.36\\
245	0.3625\\
246	0.365\\
247	0.3675\\
248	0.37\\
249	0.3725\\
250	0.375\\
251	0.3775\\
252	0.38\\
253	0.3825\\
254	0.385\\
255	0.3875\\
256	0.39\\
257	0.3925\\
258	0.395\\
259	0.3975\\
260	0.4\\
261	0.4025\\
262	0.405\\
263	0.4075\\
264	0.41\\
265	0.4125\\
266	0.415\\
267	0.4175\\
268	0.42\\
269	0.4225\\
270	0.425\\
271	0.4275\\
272	0.43\\
273	0.4325\\
274	0.435\\
275	0.4375\\
276	0.44\\
277	0.4425\\
278	0.445\\
279	0.4475\\
280	0.45\\
281	0.4525\\
282	0.455\\
283	0.4575\\
284	0.46\\
285	0.4625\\
286	0.465\\
287	0.4675\\
288	0.47\\
289	0.4725\\
290	0.475\\
291	0.4775\\
292	0.48\\
293	0.4825\\
294	0.485\\
295	0.4875\\
296	0.49\\
297	0.4925\\
298	0.495\\
299	0.4975\\
300	0.5\\
301	0.5025\\
302	0.505\\
303	0.5075\\
304	0.51\\
305	0.5125\\
306	0.515\\
307	0.5175\\
308	0.52\\
309	0.5225\\
310	0.525\\
311	0.5275\\
312	0.53\\
313	0.5325\\
314	0.535\\
315	0.5375\\
316	0.54\\
317	0.5425\\
318	0.545\\
319	0.5475\\
320	0.55\\
321	0.5525\\
322	0.555\\
323	0.5575\\
324	0.56\\
325	0.5625\\
326	0.565\\
327	0.5675\\
328	0.57\\
329	0.5725\\
330	0.575\\
331	0.5775\\
332	0.58\\
333	0.5825\\
334	0.585\\
335	0.5875\\
336	0.59\\
337	0.5925\\
338	0.595\\
339	0.5975\\
340	0.6\\
341	0.6025\\
342	0.605\\
343	0.6075\\
344	0.61\\
345	0.6125\\
346	0.615\\
347	0.6175\\
348	0.62\\
349	0.6225\\
350	0.625\\
351	0.6275\\
352	0.63\\
353	0.6325\\
354	0.635\\
355	0.6375\\
356	0.64\\
357	0.6425\\
358	0.645\\
359	0.6475\\
360	0.65\\
361	0.6525\\
362	0.655\\
363	0.6575\\
364	0.66\\
365	0.6625\\
366	0.665\\
367	0.6675\\
368	0.67\\
369	0.6725\\
370	0.675\\
371	0.6775\\
372	0.68\\
373	0.6825\\
374	0.685\\
375	0.6875\\
376	0.69\\
377	0.6925\\
378	0.695\\
379	0.6975\\
380	0.7\\
381	0.7025\\
382	0.705\\
383	0.7075\\
384	0.71\\
385	0.7125\\
386	0.715\\
387	0.7175\\
388	0.72\\
389	0.7225\\
390	0.725\\
391	0.7275\\
392	0.73\\
393	0.7325\\
394	0.735\\
395	0.7375\\
396	0.74\\
397	0.7425\\
398	0.745\\
399	0.7475\\
400	0.75\\
401	0.7525\\
402	0.755\\
403	0.7575\\
404	0.76\\
405	0.7625\\
406	0.765\\
407	0.7675\\
408	0.77\\
409	0.7725\\
410	0.775\\
411	0.7775\\
412	0.78\\
413	0.7825\\
414	0.785\\
415	0.7875\\
416	0.79\\
417	0.7925\\
418	0.795\\
419	0.7975\\
420	0.8\\
421	0.8025\\
422	0.805\\
423	0.8075\\
424	0.81\\
425	0.8125\\
426	0.815\\
427	0.8175\\
428	0.82\\
429	0.8225\\
430	0.825\\
431	0.8275\\
432	0.83\\
433	0.8325\\
434	0.835\\
435	0.8375\\
436	0.84\\
437	0.8425\\
438	0.845\\
439	0.8475\\
440	0.85\\
441	0.8525\\
442	0.855\\
443	0.8575\\
444	0.86\\
445	0.8625\\
446	0.865\\
447	0.8675\\
448	0.87\\
449	0.8725\\
450	0.875\\
451	0.8775\\
452	0.88\\
453	0.8825\\
454	0.885\\
455	0.8875\\
456	0.89\\
457	0.8925\\
458	0.895\\
459	0.8975\\
460	0.9\\
461	0.9025\\
462	0.905\\
463	0.9075\\
464	0.91\\
465	0.9125\\
466	0.915\\
467	0.9175\\
468	0.92\\
469	0.9225\\
470	0.925\\
471	0.9275\\
472	0.93\\
473	0.9325\\
474	0.935\\
475	0.9375\\
476	0.94\\
477	0.9425\\
478	0.945\\
479	0.9475\\
480	0.95\\
481	0.9525\\
482	0.955\\
483	0.9575\\
484	0.96\\
485	0.9625\\
486	0.965\\
487	0.9675\\
488	0.97\\
489	0.9725\\
490	0.975\\
491	0.9775\\
492	0.98\\
493	0.9825\\
494	0.985\\
495	0.9875\\
496	0.99\\
497	0.9925\\
498	0.995\\
499	0.9975\\
500	1\\
501	0.9975\\
502	0.995\\
503	0.9925\\
504	0.99\\
505	0.9875\\
506	0.985\\
507	0.9825\\
508	0.98\\
509	0.9775\\
510	0.975\\
511	0.9725\\
512	0.97\\
513	0.9675\\
514	0.965\\
515	0.9625\\
516	0.96\\
517	0.9575\\
518	0.955\\
519	0.9525\\
520	0.95\\
521	0.9475\\
522	0.945\\
523	0.9425\\
524	0.94\\
525	0.9375\\
526	0.935\\
527	0.9325\\
528	0.93\\
529	0.9275\\
530	0.925\\
531	0.9225\\
532	0.92\\
533	0.9175\\
534	0.915\\
535	0.9125\\
536	0.91\\
537	0.9075\\
538	0.905\\
539	0.9025\\
540	0.9\\
541	0.8975\\
542	0.895\\
543	0.8925\\
544	0.89\\
545	0.8875\\
546	0.885\\
547	0.8825\\
548	0.88\\
549	0.8775\\
550	0.875\\
551	0.8725\\
552	0.87\\
553	0.8675\\
554	0.865\\
555	0.8625\\
556	0.86\\
557	0.8575\\
558	0.855\\
559	0.8525\\
560	0.85\\
561	0.8475\\
562	0.845\\
563	0.8425\\
564	0.84\\
565	0.8375\\
566	0.835\\
567	0.8325\\
568	0.83\\
569	0.8275\\
570	0.825\\
571	0.8225\\
572	0.82\\
573	0.8175\\
574	0.815\\
575	0.8125\\
576	0.81\\
577	0.8075\\
578	0.805\\
579	0.8025\\
580	0.8\\
581	0.7975\\
582	0.795\\
583	0.7925\\
584	0.79\\
585	0.7875\\
586	0.785\\
587	0.7825\\
588	0.78\\
589	0.7775\\
590	0.775\\
591	0.7725\\
592	0.77\\
593	0.7675\\
594	0.765\\
595	0.7625\\
596	0.76\\
597	0.7575\\
598	0.755\\
599	0.7525\\
600	0.75\\
601	0.7475\\
602	0.745\\
603	0.7425\\
604	0.74\\
605	0.7375\\
606	0.735\\
607	0.7325\\
608	0.73\\
609	0.7275\\
610	0.725\\
611	0.7225\\
612	0.72\\
613	0.7175\\
614	0.715\\
615	0.7125\\
616	0.71\\
617	0.7075\\
618	0.705\\
619	0.7025\\
620	0.7\\
621	0.6975\\
622	0.695\\
623	0.6925\\
624	0.69\\
625	0.6875\\
626	0.685\\
627	0.6825\\
628	0.68\\
629	0.6775\\
630	0.675\\
631	0.6725\\
632	0.67\\
633	0.6675\\
634	0.665\\
635	0.6625\\
636	0.66\\
637	0.6575\\
638	0.655\\
639	0.6525\\
640	0.65\\
641	0.6475\\
642	0.645\\
643	0.6425\\
644	0.64\\
645	0.6375\\
646	0.635\\
647	0.6325\\
648	0.63\\
649	0.6275\\
650	0.625\\
651	0.6225\\
652	0.62\\
653	0.6175\\
654	0.615\\
655	0.6125\\
656	0.61\\
657	0.6075\\
658	0.605\\
659	0.6025\\
660	0.6\\
661	0.5975\\
662	0.595\\
663	0.5925\\
664	0.59\\
665	0.5875\\
666	0.585\\
667	0.5825\\
668	0.58\\
669	0.5775\\
670	0.575\\
671	0.5725\\
672	0.57\\
673	0.5675\\
674	0.565\\
675	0.5625\\
676	0.56\\
677	0.5575\\
678	0.555\\
679	0.5525\\
680	0.55\\
681	0.5475\\
682	0.545\\
683	0.5425\\
684	0.54\\
685	0.5375\\
686	0.535\\
687	0.5325\\
688	0.53\\
689	0.5275\\
690	0.525\\
691	0.5225\\
692	0.52\\
693	0.5175\\
694	0.515\\
695	0.5125\\
696	0.51\\
697	0.5075\\
698	0.505\\
699	0.5025\\
700	0.5\\
701	0.4975\\
702	0.495\\
703	0.4925\\
704	0.49\\
705	0.4875\\
706	0.485\\
707	0.4825\\
708	0.48\\
709	0.4775\\
710	0.475\\
711	0.4725\\
712	0.47\\
713	0.4675\\
714	0.465\\
715	0.4625\\
716	0.46\\
717	0.4575\\
718	0.455\\
719	0.4525\\
720	0.45\\
721	0.4475\\
722	0.445\\
723	0.4425\\
724	0.44\\
725	0.4375\\
726	0.435\\
727	0.4325\\
728	0.43\\
729	0.4275\\
730	0.425\\
731	0.4225\\
732	0.42\\
733	0.4175\\
734	0.415\\
735	0.4125\\
736	0.41\\
737	0.4075\\
738	0.405\\
739	0.4025\\
740	0.4\\
741	0.3975\\
742	0.395\\
743	0.3925\\
744	0.39\\
745	0.3875\\
746	0.385\\
747	0.3825\\
748	0.38\\
749	0.3775\\
750	0.375\\
751	0.3725\\
752	0.37\\
753	0.3675\\
754	0.365\\
755	0.3625\\
756	0.36\\
757	0.3575\\
758	0.355\\
759	0.3525\\
760	0.35\\
761	0.3475\\
762	0.345\\
763	0.3425\\
764	0.34\\
765	0.3375\\
766	0.335\\
767	0.3325\\
768	0.33\\
769	0.3275\\
770	0.325\\
771	0.3225\\
772	0.32\\
773	0.3175\\
774	0.315\\
775	0.3125\\
776	0.31\\
777	0.3075\\
778	0.305\\
779	0.3025\\
780	0.3\\
781	0.2975\\
782	0.295\\
783	0.2925\\
784	0.29\\
785	0.2875\\
786	0.285\\
787	0.2825\\
788	0.28\\
789	0.2775\\
790	0.275\\
791	0.2725\\
792	0.27\\
793	0.2675\\
794	0.265\\
795	0.2625\\
796	0.26\\
797	0.2575\\
798	0.255\\
799	0.2525\\
800	0.25\\
801	0.2475\\
802	0.245\\
803	0.2425\\
804	0.24\\
805	0.2375\\
806	0.235\\
807	0.2325\\
808	0.23\\
809	0.2275\\
810	0.225\\
811	0.2225\\
812	0.22\\
813	0.2175\\
814	0.215\\
815	0.2125\\
816	0.21\\
817	0.2075\\
818	0.205\\
819	0.2025\\
820	0.2\\
821	0.1975\\
822	0.195\\
823	0.1925\\
824	0.19\\
825	0.1875\\
826	0.185\\
827	0.1825\\
828	0.18\\
829	0.1775\\
830	0.175\\
831	0.1725\\
832	0.17\\
833	0.1675\\
834	0.165\\
835	0.1625\\
836	0.16\\
837	0.1575\\
838	0.155\\
839	0.1525\\
840	0.15\\
841	0.1475\\
842	0.145\\
843	0.1425\\
844	0.14\\
845	0.1375\\
846	0.135\\
847	0.1325\\
848	0.13\\
849	0.1275\\
850	0.125\\
851	0.1225\\
852	0.12\\
853	0.1175\\
854	0.115\\
855	0.1125\\
856	0.11\\
857	0.1075\\
858	0.105\\
859	0.1025\\
860	0.1\\
861	0.0975\\
862	0.095\\
863	0.0925\\
864	0.09\\
865	0.0875\\
866	0.085\\
867	0.0825\\
868	0.08\\
869	0.0775\\
870	0.075\\
871	0.0725\\
872	0.07\\
873	0.0675\\
874	0.0649999999999999\\
875	0.0625\\
876	0.0600000000000001\\
877	0.0575\\
878	0.055\\
879	0.0525\\
880	0.05\\
881	0.0475\\
882	0.045\\
883	0.0425\\
884	0.04\\
885	0.0375\\
886	0.035\\
887	0.0325\\
888	0.03\\
889	0.0275\\
890	0.025\\
891	0.0225\\
892	0.02\\
893	0.0175\\
894	0.015\\
895	0.0125\\
896	0.01\\
897	0.00749999999999995\\
898	0.005\\
899	0.00249999999999995\\
900	0\\
901	0\\
902	0\\
903	0\\
904	0\\
905	0\\
906	0\\
907	0\\
908	0\\
909	0\\
910	0\\
911	0\\
912	0\\
913	0\\
914	0\\
915	0\\
916	0\\
917	0\\
918	0\\
919	0\\
920	0\\
921	0\\
922	0\\
923	0\\
924	0\\
925	0\\
926	0\\
927	0\\
928	0\\
929	0\\
930	0\\
931	0\\
932	0\\
933	0\\
934	0\\
935	0\\
936	0\\
937	0\\
938	0\\
939	0\\
940	0\\
941	0\\
942	0\\
943	0\\
944	0\\
945	0\\
946	0\\
947	0\\
948	0\\
949	0\\
950	0\\
951	0\\
952	0\\
953	0\\
954	0\\
955	0\\
956	0\\
957	0\\
958	0\\
959	0\\
960	0\\
961	0\\
962	0\\
963	0\\
964	0\\
965	0\\
966	0\\
967	0\\
968	0\\
969	0\\
970	0\\
971	0\\
972	0\\
973	0\\
974	0\\
975	0\\
976	0\\
977	0\\
978	0\\
979	0\\
980	0\\
981	0\\
982	0\\
983	0\\
984	0\\
985	0\\
986	0\\
987	0\\
988	0\\
989	0\\
990	0\\
991	0\\
992	0\\
993	0\\
994	0\\
995	0\\
996	0\\
997	0\\
998	0\\
999	0\\
1000	0\\
};
\addplot [color=black,solid,line width=1.4pt,forget plot]
  table[row sep=crcr]{0	0\\
1	0\\
2	0\\
3	0\\
4	0\\
5	0\\
6	0\\
7	0\\
8	0\\
9	0\\
10	0\\
11	0\\
12	0\\
13	0\\
14	0\\
15	0\\
16	0\\
17	0\\
18	0\\
19	0\\
20	0\\
21	0\\
22	0\\
23	0\\
24	0\\
25	0\\
26	0\\
27	0\\
28	0\\
29	0\\
30	0\\
31	0\\
32	0\\
33	0\\
34	0\\
35	0\\
36	0\\
37	0\\
38	0\\
39	0\\
40	0\\
41	0\\
42	0\\
43	0\\
44	0\\
45	0\\
46	0\\
47	0\\
48	0\\
49	0\\
50	0\\
51	0\\
52	0\\
53	0\\
54	0\\
55	0\\
56	0\\
57	0\\
58	0\\
59	0\\
60	0\\
61	0\\
62	0\\
63	0\\
64	0\\
65	0\\
66	0\\
67	0\\
68	0\\
69	0\\
70	0\\
71	0\\
72	0\\
73	0\\
74	0\\
75	0\\
76	0\\
77	0\\
78	0\\
79	0\\
80	0\\
81	0\\
82	0\\
83	0\\
84	0\\
85	0\\
86	0\\
87	0\\
88	0\\
89	0\\
90	0\\
91	0\\
92	0\\
93	0\\
94	0\\
95	0\\
96	0\\
97	0\\
98	0\\
99	0\\
100	0\\
101	0\\
102	0\\
103	0\\
104	0\\
105	0\\
106	0\\
107	0\\
108	0\\
109	0\\
110	0\\
111	0\\
112	0\\
113	0\\
114	0\\
115	0\\
116	0\\
117	0\\
118	0\\
119	0\\
120	0\\
121	0\\
122	0\\
123	0\\
124	0\\
125	0\\
126	0\\
127	0\\
128	0\\
129	0\\
130	0\\
131	0\\
132	0\\
133	0\\
134	0\\
135	0\\
136	0\\
137	0\\
138	0\\
139	0\\
140	0\\
141	0\\
142	0\\
143	0\\
144	0\\
145	0\\
146	0\\
147	0\\
148	0\\
149	0\\
150	0\\
151	0\\
152	0\\
153	0\\
154	0\\
155	0\\
156	0\\
157	0\\
158	0\\
159	0\\
160	0\\
161	0\\
162	0\\
163	0\\
164	0\\
165	0\\
166	0\\
167	0\\
168	0\\
169	0\\
170	0\\
171	0\\
172	0\\
173	0\\
174	0\\
175	0\\
176	0\\
177	0\\
178	0\\
179	0\\
180	0\\
181	0\\
182	0\\
183	0\\
184	0\\
185	0\\
186	0\\
187	0\\
188	0\\
189	0\\
190	0\\
191	0\\
192	0\\
193	0\\
194	0\\
195	0\\
196	0\\
197	0\\
198	0\\
199	0\\
200	0\\
201	0\\
202	0\\
203	0\\
204	0\\
205	0\\
206	0\\
207	0\\
208	0\\
209	0\\
210	0\\
211	0\\
212	0\\
213	0\\
214	0\\
215	0\\
216	0\\
217	0\\
218	0\\
219	0\\
220	0\\
221	0\\
222	0\\
223	0\\
224	0\\
225	0\\
226	0\\
227	0\\
228	0\\
229	0\\
230	0\\
231	0\\
232	0\\
233	0\\
234	0\\
235	0\\
236	0\\
237	0\\
238	0\\
239	0\\
240	0\\
241	0\\
242	0\\
243	0\\
244	0\\
245	0\\
246	0\\
247	0\\
248	0\\
249	0\\
250	0\\
251	0\\
252	0\\
253	0\\
254	0\\
255	0\\
256	0\\
257	0\\
258	0\\
259	0\\
260	0\\
261	0\\
262	0\\
263	0\\
264	0\\
265	0\\
266	0\\
267	0\\
268	0\\
269	0\\
270	0\\
271	0\\
272	0\\
273	0\\
274	0\\
275	0\\
276	0\\
277	0\\
278	0\\
279	0\\
280	0\\
281	0\\
282	0\\
283	0\\
284	0\\
285	0\\
286	0\\
287	0\\
288	0\\
289	0\\
290	0\\
291	0\\
292	0\\
293	0\\
294	0\\
295	0\\
296	0\\
297	0\\
298	0\\
299	0\\
300	0\\
301	0\\
302	0\\
303	0\\
304	0\\
305	0\\
306	0\\
307	0\\
308	0\\
309	0\\
310	0\\
311	0\\
312	0\\
313	0\\
314	0\\
315	0\\
316	0\\
317	0\\
318	0\\
319	0\\
320	0\\
321	0\\
322	0\\
323	0\\
324	0\\
325	0\\
326	0\\
327	0\\
328	0\\
329	0\\
330	0\\
331	0\\
332	0\\
333	0\\
334	0\\
335	0\\
336	0\\
337	0\\
338	0\\
339	0\\
340	0\\
341	0\\
342	0\\
343	0\\
344	0\\
345	0\\
346	0\\
347	0\\
348	0\\
349	0\\
350	0\\
351	0\\
352	0\\
353	0\\
354	0\\
355	0\\
356	0\\
357	0\\
358	0\\
359	0\\
360	0\\
361	0\\
362	0\\
363	0\\
364	0\\
365	0\\
366	0\\
367	0\\
368	0\\
369	0\\
370	0\\
371	0\\
372	0\\
373	0\\
374	0\\
375	0\\
376	0\\
377	0\\
378	0\\
379	0\\
380	0\\
381	0\\
382	0\\
383	0\\
384	0\\
385	0\\
386	0\\
387	0\\
388	0\\
389	0\\
390	0\\
391	0\\
392	0\\
393	0\\
394	0\\
395	0\\
396	0\\
397	0\\
398	0\\
399	0\\
400	0\\
401	0\\
402	0\\
403	0\\
404	0\\
405	0\\
406	0\\
407	0\\
408	0\\
409	0\\
410	0\\
411	0\\
412	0\\
413	0\\
414	0\\
415	0\\
416	0\\
417	0\\
418	0\\
419	0\\
420	0\\
421	0\\
422	0\\
423	0\\
424	0\\
425	0\\
426	0\\
427	0\\
428	0\\
429	0\\
430	0\\
431	0\\
432	0\\
433	0\\
434	0\\
435	0\\
436	0\\
437	0\\
438	0\\
439	0\\
440	0\\
441	0\\
442	0\\
443	0\\
444	0\\
445	0\\
446	0\\
447	0\\
448	0\\
449	0\\
450	0\\
451	0\\
452	0\\
453	0\\
454	0\\
455	0\\
456	0\\
457	0\\
458	0\\
459	0\\
460	0\\
461	0\\
462	0\\
463	0\\
464	0\\
465	0\\
466	0\\
467	0\\
468	0\\
469	0\\
470	0\\
471	0\\
472	0\\
473	0\\
474	0\\
475	0\\
476	0\\
477	0\\
478	0\\
479	0\\
480	0\\
481	0\\
482	0\\
483	0\\
484	0\\
485	0\\
486	0\\
487	0\\
488	0\\
489	0\\
490	0\\
491	0\\
492	0\\
493	0\\
494	0\\
495	0\\
496	0\\
497	0\\
498	0\\
499	0\\
500	0\\
501	0\\
502	0\\
503	0\\
504	0\\
505	0\\
506	0\\
507	0\\
508	0\\
509	0\\
510	0\\
511	0\\
512	0\\
513	0\\
514	0\\
515	0\\
516	0\\
517	0\\
518	0\\
519	0\\
520	0\\
521	0\\
522	0\\
523	0\\
524	0\\
525	0\\
526	0\\
527	0\\
528	0\\
529	0\\
530	0\\
531	0\\
532	0\\
533	0\\
534	0\\
535	0\\
536	0\\
537	0\\
538	0\\
539	0\\
540	0\\
541	0\\
542	0\\
543	0\\
544	0\\
545	0\\
546	0\\
547	0\\
548	0\\
549	0\\
550	0\\
551	0\\
552	0\\
553	0\\
554	0\\
555	0\\
556	0\\
557	0\\
558	0\\
559	0\\
560	0\\
561	0\\
562	0\\
563	0\\
564	0\\
565	0\\
566	0\\
567	0\\
568	0\\
569	0\\
570	0\\
571	0\\
572	0\\
573	0\\
574	0\\
575	0\\
576	0\\
577	0\\
578	0\\
579	0\\
580	0\\
581	0\\
582	0\\
583	0\\
584	0\\
585	0\\
586	0\\
587	0\\
588	0\\
589	0\\
590	0\\
591	0\\
592	0\\
593	0\\
594	0\\
595	0\\
596	0\\
597	0\\
598	0\\
599	0\\
600	0\\
601	0.0025\\
602	0.005\\
603	0.0075\\
604	0.01\\
605	0.0125\\
606	0.015\\
607	0.0175\\
608	0.02\\
609	0.0225\\
610	0.025\\
611	0.0275\\
612	0.03\\
613	0.0325\\
614	0.035\\
615	0.0375\\
616	0.04\\
617	0.0425\\
618	0.045\\
619	0.0475\\
620	0.05\\
621	0.0525\\
622	0.055\\
623	0.0575\\
624	0.06\\
625	0.0625\\
626	0.065\\
627	0.0675\\
628	0.07\\
629	0.0725\\
630	0.075\\
631	0.0775\\
632	0.08\\
633	0.0825\\
634	0.085\\
635	0.0875\\
636	0.09\\
637	0.0925\\
638	0.095\\
639	0.0975\\
640	0.1\\
641	0.1025\\
642	0.105\\
643	0.1075\\
644	0.11\\
645	0.1125\\
646	0.115\\
647	0.1175\\
648	0.12\\
649	0.1225\\
650	0.125\\
651	0.1275\\
652	0.13\\
653	0.1325\\
654	0.135\\
655	0.1375\\
656	0.14\\
657	0.1425\\
658	0.145\\
659	0.1475\\
660	0.15\\
661	0.1525\\
662	0.155\\
663	0.1575\\
664	0.16\\
665	0.1625\\
666	0.165\\
667	0.1675\\
668	0.17\\
669	0.1725\\
670	0.175\\
671	0.1775\\
672	0.18\\
673	0.1825\\
674	0.185\\
675	0.1875\\
676	0.19\\
677	0.1925\\
678	0.195\\
679	0.1975\\
680	0.2\\
681	0.2025\\
682	0.205\\
683	0.2075\\
684	0.21\\
685	0.2125\\
686	0.215\\
687	0.2175\\
688	0.22\\
689	0.2225\\
690	0.225\\
691	0.2275\\
692	0.23\\
693	0.2325\\
694	0.235\\
695	0.2375\\
696	0.24\\
697	0.2425\\
698	0.245\\
699	0.2475\\
700	0.25\\
701	0.2525\\
702	0.255\\
703	0.2575\\
704	0.26\\
705	0.2625\\
706	0.265\\
707	0.2675\\
708	0.27\\
709	0.2725\\
710	0.275\\
711	0.2775\\
712	0.28\\
713	0.2825\\
714	0.285\\
715	0.2875\\
716	0.29\\
717	0.2925\\
718	0.295\\
719	0.2975\\
720	0.3\\
721	0.3025\\
722	0.305\\
723	0.3075\\
724	0.31\\
725	0.3125\\
726	0.315\\
727	0.3175\\
728	0.32\\
729	0.3225\\
730	0.325\\
731	0.3275\\
732	0.33\\
733	0.3325\\
734	0.335\\
735	0.3375\\
736	0.34\\
737	0.3425\\
738	0.345\\
739	0.3475\\
740	0.35\\
741	0.3525\\
742	0.355\\
743	0.3575\\
744	0.36\\
745	0.3625\\
746	0.365\\
747	0.3675\\
748	0.37\\
749	0.3725\\
750	0.375\\
751	0.3775\\
752	0.38\\
753	0.3825\\
754	0.385\\
755	0.3875\\
756	0.39\\
757	0.3925\\
758	0.395\\
759	0.3975\\
760	0.4\\
761	0.4025\\
762	0.405\\
763	0.4075\\
764	0.41\\
765	0.4125\\
766	0.415\\
767	0.4175\\
768	0.42\\
769	0.4225\\
770	0.425\\
771	0.4275\\
772	0.43\\
773	0.4325\\
774	0.435\\
775	0.4375\\
776	0.44\\
777	0.4425\\
778	0.445\\
779	0.4475\\
780	0.45\\
781	0.4525\\
782	0.455\\
783	0.4575\\
784	0.46\\
785	0.4625\\
786	0.465\\
787	0.4675\\
788	0.47\\
789	0.4725\\
790	0.475\\
791	0.4775\\
792	0.48\\
793	0.4825\\
794	0.485\\
795	0.4875\\
796	0.49\\
797	0.4925\\
798	0.495\\
799	0.4975\\
800	0.5\\
801	0.5025\\
802	0.505\\
803	0.5075\\
804	0.51\\
805	0.5125\\
806	0.515\\
807	0.5175\\
808	0.52\\
809	0.5225\\
810	0.525\\
811	0.5275\\
812	0.53\\
813	0.5325\\
814	0.535\\
815	0.5375\\
816	0.54\\
817	0.5425\\
818	0.545\\
819	0.5475\\
820	0.55\\
821	0.5525\\
822	0.555\\
823	0.5575\\
824	0.56\\
825	0.5625\\
826	0.565\\
827	0.5675\\
828	0.57\\
829	0.5725\\
830	0.575\\
831	0.5775\\
832	0.58\\
833	0.5825\\
834	0.585\\
835	0.5875\\
836	0.59\\
837	0.5925\\
838	0.595\\
839	0.5975\\
840	0.6\\
841	0.6025\\
842	0.605\\
843	0.6075\\
844	0.61\\
845	0.6125\\
846	0.615\\
847	0.6175\\
848	0.62\\
849	0.6225\\
850	0.625\\
851	0.6275\\
852	0.63\\
853	0.6325\\
854	0.635\\
855	0.6375\\
856	0.64\\
857	0.6425\\
858	0.645\\
859	0.6475\\
860	0.65\\
861	0.6525\\
862	0.655\\
863	0.6575\\
864	0.66\\
865	0.6625\\
866	0.665\\
867	0.6675\\
868	0.67\\
869	0.6725\\
870	0.675\\
871	0.6775\\
872	0.68\\
873	0.6825\\
874	0.685\\
875	0.6875\\
876	0.69\\
877	0.6925\\
878	0.695\\
879	0.6975\\
880	0.7\\
881	0.7025\\
882	0.705\\
883	0.7075\\
884	0.71\\
885	0.7125\\
886	0.715\\
887	0.7175\\
888	0.72\\
889	0.7225\\
890	0.725\\
891	0.7275\\
892	0.73\\
893	0.7325\\
894	0.735\\
895	0.7375\\
896	0.74\\
897	0.7425\\
898	0.745\\
899	0.7475\\
900	0.75\\
901	0.7525\\
902	0.755\\
903	0.7575\\
904	0.76\\
905	0.7625\\
906	0.765\\
907	0.7675\\
908	0.77\\
909	0.7725\\
910	0.775\\
911	0.7775\\
912	0.78\\
913	0.7825\\
914	0.785\\
915	0.7875\\
916	0.79\\
917	0.7925\\
918	0.795\\
919	0.7975\\
920	0.8\\
921	0.8025\\
922	0.805\\
923	0.8075\\
924	0.81\\
925	0.8125\\
926	0.815\\
927	0.8175\\
928	0.82\\
929	0.8225\\
930	0.825\\
931	0.8275\\
932	0.83\\
933	0.8325\\
934	0.835\\
935	0.8375\\
936	0.84\\
937	0.8425\\
938	0.845\\
939	0.8475\\
940	0.85\\
941	0.8525\\
942	0.855\\
943	0.8575\\
944	0.86\\
945	0.8625\\
946	0.865\\
947	0.8675\\
948	0.87\\
949	0.8725\\
950	0.875\\
951	0.8775\\
952	0.88\\
953	0.8825\\
954	0.885\\
955	0.8875\\
956	0.89\\
957	0.8925\\
958	0.895\\
959	0.8975\\
960	0.9\\
961	0.9025\\
962	0.905\\
963	0.9075\\
964	0.91\\
965	0.9125\\
966	0.915\\
967	0.9175\\
968	0.92\\
969	0.9225\\
970	0.925\\
971	0.9275\\
972	0.93\\
973	0.9325\\
974	0.935\\
975	0.9375\\
976	0.94\\
977	0.9425\\
978	0.945\\
979	0.9475\\
980	0.95\\
981	0.9525\\
982	0.955\\
983	0.9575\\
984	0.96\\
985	0.9625\\
986	0.965\\
987	0.9675\\
988	0.97\\
989	0.9725\\
990	0.975\\
991	0.9775\\
992	0.98\\
993	0.9825\\
994	0.985\\
995	0.9875\\
996	0.99\\
997	0.9925\\
998	0.995\\
999	0.9975\\
1000	1\\
};
\node[right, inner sep=0mm, text=black]
at (axis cs:1,1.1,0) {Baja (L)};
\node[right, inner sep=0mm, text=black]
at (axis cs:450,1.1,0) {Media (M)};
\node[right, inner sep=0mm, text=black]
at (axis cs:900,1.1,0) {Alta (H)};
\end{axis}
\end{tikzpicture}%
		\caption{$\chi_3$ - Luz (lux)}
		\label{fig:light-lang-variable}
	\end{subfigure}
	\qquad
	\begin{subfigure}[b]{0.45\textwidth}
		\setlength\figureheight{2.5cm}
		\setlength\figurewidth{6cm}
		% This file was created by matlab2tikz v0.4.7 (commit aa5a2b3a09d5219f83b03f452820de20dfb2a2eb) running on MATLAB 8.0.
% Copyright (c) 2008--2014, Nico Schlömer <nico.schloemer@gmail.com>
% All rights reserved.
% Minimal pgfplots version: 1.3
% 
\begin{tikzpicture}

\begin{axis}[%
width=\figurewidth,
height=\figureheight,
clip=false,
scale only axis,
xmin=0,
xmax=100,
ymin=0,
ymax=1
]
\addplot [color=black,solid,line width=1.4pt,forget plot]
  table[row sep=crcr]{0	1\\
1	0.975\\
2	0.95\\
3	0.925\\
4	0.9\\
5	0.875\\
6	0.85\\
7	0.825\\
8	0.8\\
9	0.775\\
10	0.75\\
11	0.725\\
12	0.7\\
13	0.675\\
14	0.65\\
15	0.625\\
16	0.6\\
17	0.575\\
18	0.55\\
19	0.525\\
20	0.5\\
21	0.475\\
22	0.45\\
23	0.425\\
24	0.4\\
25	0.375\\
26	0.35\\
27	0.325\\
28	0.3\\
29	0.275\\
30	0.25\\
31	0.225\\
32	0.2\\
33	0.175\\
34	0.15\\
35	0.125\\
36	0.1\\
37	0.075\\
38	0.05\\
39	0.025\\
40	0\\
41	0\\
42	0\\
43	0\\
44	0\\
45	0\\
46	0\\
47	0\\
48	0\\
49	0\\
50	0\\
51	0\\
52	0\\
53	0\\
54	0\\
55	0\\
56	0\\
57	0\\
58	0\\
59	0\\
60	0\\
61	0\\
62	0\\
63	0\\
64	0\\
65	0\\
66	0\\
67	0\\
68	0\\
69	0\\
70	0\\
71	0\\
72	0\\
73	0\\
74	0\\
75	0\\
76	0\\
77	0\\
78	0\\
79	0\\
80	0\\
81	0\\
82	0\\
83	0\\
84	0\\
85	0\\
86	0\\
87	0\\
88	0\\
89	0\\
90	0\\
91	0\\
92	0\\
93	0\\
94	0\\
95	0\\
96	0\\
97	0\\
98	0\\
99	0\\
100	0\\
};
\addplot [color=black,solid,line width=1.4pt,forget plot]
  table[row sep=crcr]{0	0\\
1	0\\
2	0\\
3	0\\
4	0\\
5	0\\
6	0\\
7	0\\
8	0\\
9	0\\
10	0\\
11	0.025\\
12	0.05\\
13	0.075\\
14	0.1\\
15	0.125\\
16	0.15\\
17	0.175\\
18	0.2\\
19	0.225\\
20	0.25\\
21	0.275\\
22	0.3\\
23	0.325\\
24	0.35\\
25	0.375\\
26	0.4\\
27	0.425\\
28	0.45\\
29	0.475\\
30	0.5\\
31	0.525\\
32	0.55\\
33	0.575\\
34	0.6\\
35	0.625\\
36	0.65\\
37	0.675\\
38	0.7\\
39	0.725\\
40	0.75\\
41	0.775\\
42	0.8\\
43	0.825\\
44	0.85\\
45	0.875\\
46	0.9\\
47	0.925\\
48	0.95\\
49	0.975\\
50	1\\
51	0.975\\
52	0.95\\
53	0.925\\
54	0.9\\
55	0.875\\
56	0.85\\
57	0.825\\
58	0.8\\
59	0.775\\
60	0.75\\
61	0.725\\
62	0.7\\
63	0.675\\
64	0.65\\
65	0.625\\
66	0.6\\
67	0.575\\
68	0.55\\
69	0.525\\
70	0.5\\
71	0.475\\
72	0.45\\
73	0.425\\
74	0.4\\
75	0.375\\
76	0.35\\
77	0.325\\
78	0.3\\
79	0.275\\
80	0.25\\
81	0.225\\
82	0.2\\
83	0.175\\
84	0.15\\
85	0.125\\
86	0.1\\
87	0.075\\
88	0.05\\
89	0.025\\
90	0\\
91	0\\
92	0\\
93	0\\
94	0\\
95	0\\
96	0\\
97	0\\
98	0\\
99	0\\
100	0\\
};
\addplot [color=black,solid,line width=1.4pt,forget plot]
  table[row sep=crcr]{0	0\\
1	0\\
2	0\\
3	0\\
4	0\\
5	0\\
6	0\\
7	0\\
8	0\\
9	0\\
10	0\\
11	0\\
12	0\\
13	0\\
14	0\\
15	0\\
16	0\\
17	0\\
18	0\\
19	0\\
20	0\\
21	0\\
22	0\\
23	0\\
24	0\\
25	0\\
26	0\\
27	0\\
28	0\\
29	0\\
30	0\\
31	0\\
32	0\\
33	0\\
34	0\\
35	0\\
36	0\\
37	0\\
38	0\\
39	0\\
40	0\\
41	0\\
42	0\\
43	0\\
44	0\\
45	0\\
46	0\\
47	0\\
48	0\\
49	0\\
50	0\\
51	0\\
52	0\\
53	0\\
54	0\\
55	0\\
56	0\\
57	0\\
58	0\\
59	0\\
60	0\\
61	0.025\\
62	0.05\\
63	0.075\\
64	0.1\\
65	0.125\\
66	0.15\\
67	0.175\\
68	0.2\\
69	0.225\\
70	0.25\\
71	0.275\\
72	0.3\\
73	0.325\\
74	0.35\\
75	0.375\\
76	0.4\\
77	0.425\\
78	0.45\\
79	0.475\\
80	0.5\\
81	0.525\\
82	0.55\\
83	0.575\\
84	0.6\\
85	0.625\\
86	0.65\\
87	0.675\\
88	0.7\\
89	0.725\\
90	0.75\\
91	0.775\\
92	0.8\\
93	0.825\\
94	0.85\\
95	0.875\\
96	0.9\\
97	0.925\\
98	0.95\\
99	0.975\\
100	1\\
};
\node[right, inner sep=0mm, text=black]
at (axis cs:1,1.1,0) {Baja (L)};
\node[right, inner sep=0mm, text=black]
at (axis cs:45,1.1,0) {Media (M)};
\node[right, inner sep=0mm, text=black]
at (axis cs:90,1.1,0) {Alta (H)};
\end{axis}
\end{tikzpicture}%
		\caption{$\chi_4$ - Humedad (ppm)}
		\label{fig:humidity-lang-variable}
	\end{subfigure}
	
	\vspace{1 cm}
	\begin{subfigure}[b]{0.45\textwidth}
		\setlength\figureheight{2.5cm}
		\setlength\figurewidth{6cm}
		% This file was created by matlab2tikz v0.4.7 (commit aa5a2b3a09d5219f83b03f452820de20dfb2a2eb) running on MATLAB 8.0.
% Copyright (c) 2008--2014, Nico Schlömer <nico.schloemer@gmail.com>
% All rights reserved.
% Minimal pgfplots version: 1.3
% 
\begin{tikzpicture}

\begin{axis}[%
width=\figurewidth,
height=\figureheight,
clip=false,
scale only axis,
xmin=0,
xmax=80,
ymin=0,
ymax=1
]
\addplot [color=black,solid,line width=1.4pt,forget plot]
  table[row sep=crcr]{0	1\\
1	0.966666666666667\\
2	0.933333333333333\\
3	0.9\\
4	0.866666666666667\\
5	0.833333333333333\\
6	0.8\\
7	0.766666666666667\\
8	0.733333333333333\\
9	0.7\\
10	0.666666666666667\\
11	0.633333333333333\\
12	0.6\\
13	0.566666666666667\\
14	0.533333333333333\\
15	0.5\\
16	0.466666666666667\\
17	0.433333333333333\\
18	0.4\\
19	0.366666666666667\\
20	0.333333333333333\\
21	0.3\\
22	0.266666666666667\\
23	0.233333333333333\\
24	0.2\\
25	0.166666666666667\\
26	0.133333333333333\\
27	0.1\\
28	0.0666666666666667\\
29	0.0333333333333333\\
30	0\\
31	0\\
32	0\\
33	0\\
34	0\\
35	0\\
36	0\\
37	0\\
38	0\\
39	0\\
40	0\\
41	0\\
42	0\\
43	0\\
44	0\\
45	0\\
46	0\\
47	0\\
48	0\\
49	0\\
50	0\\
51	0\\
52	0\\
53	0\\
54	0\\
55	0\\
56	0\\
57	0\\
58	0\\
59	0\\
60	0\\
61	0\\
62	0\\
63	0\\
64	0\\
65	0\\
66	0\\
67	0\\
68	0\\
69	0\\
70	0\\
71	0\\
72	0\\
73	0\\
74	0\\
75	0\\
76	0\\
77	0\\
78	0\\
79	0\\
80	0\\
};
\addplot [color=black,solid,line width=1.4pt,forget plot]
  table[row sep=crcr]{0	0\\
1	0\\
2	0\\
3	0\\
4	0\\
5	0\\
6	0\\
7	0\\
8	0\\
9	0\\
10	0\\
11	0.0333333333333333\\
12	0.0666666666666667\\
13	0.1\\
14	0.133333333333333\\
15	0.166666666666667\\
16	0.2\\
17	0.233333333333333\\
18	0.266666666666667\\
19	0.3\\
20	0.333333333333333\\
21	0.366666666666667\\
22	0.4\\
23	0.433333333333333\\
24	0.466666666666667\\
25	0.5\\
26	0.533333333333333\\
27	0.566666666666667\\
28	0.6\\
29	0.633333333333333\\
30	0.666666666666667\\
31	0.7\\
32	0.733333333333333\\
33	0.766666666666667\\
34	0.8\\
35	0.833333333333333\\
36	0.866666666666667\\
37	0.9\\
38	0.933333333333333\\
39	0.966666666666667\\
40	1\\
41	0.966666666666667\\
42	0.933333333333333\\
43	0.9\\
44	0.866666666666667\\
45	0.833333333333333\\
46	0.8\\
47	0.766666666666667\\
48	0.733333333333333\\
49	0.7\\
50	0.666666666666667\\
51	0.633333333333333\\
52	0.6\\
53	0.566666666666667\\
54	0.533333333333333\\
55	0.5\\
56	0.466666666666667\\
57	0.433333333333333\\
58	0.4\\
59	0.366666666666667\\
60	0.333333333333333\\
61	0.3\\
62	0.266666666666667\\
63	0.233333333333333\\
64	0.2\\
65	0.166666666666667\\
66	0.133333333333333\\
67	0.1\\
68	0.0666666666666667\\
69	0.0333333333333333\\
70	0\\
71	0\\
72	0\\
73	0\\
74	0\\
75	0\\
76	0\\
77	0\\
78	0\\
79	0\\
80	0\\
};
\addplot [color=black,solid,line width=1.4pt,forget plot]
  table[row sep=crcr]{0	0\\
1	0\\
2	0\\
3	0\\
4	0\\
5	0\\
6	0\\
7	0\\
8	0\\
9	0\\
10	0\\
11	0\\
12	0\\
13	0\\
14	0\\
15	0\\
16	0\\
17	0\\
18	0\\
19	0\\
20	0\\
21	0\\
22	0\\
23	0\\
24	0\\
25	0\\
26	0\\
27	0\\
28	0\\
29	0\\
30	0\\
31	0\\
32	0\\
33	0\\
34	0\\
35	0\\
36	0\\
37	0\\
38	0\\
39	0\\
40	0\\
41	0\\
42	0\\
43	0\\
44	0\\
45	0\\
46	0\\
47	0\\
48	0\\
49	0\\
50	0\\
51	0.0333333333333333\\
52	0.0666666666666667\\
53	0.1\\
54	0.133333333333333\\
55	0.166666666666667\\
56	0.2\\
57	0.233333333333333\\
58	0.266666666666667\\
59	0.3\\
60	0.333333333333333\\
61	0.366666666666667\\
62	0.4\\
63	0.433333333333333\\
64	0.466666666666667\\
65	0.5\\
66	0.533333333333333\\
67	0.566666666666667\\
68	0.6\\
69	0.633333333333333\\
70	0.666666666666667\\
71	0.7\\
72	0.733333333333333\\
73	0.766666666666667\\
74	0.8\\
75	0.833333333333333\\
76	0.866666666666667\\
77	0.9\\
78	0.933333333333333\\
79	0.966666666666667\\
80	1\\
};
\node[right, inner sep=0mm, text=black]
at (axis cs:1,1.1,0) {Cerca (L)};
\node[right, inner sep=0mm, text=black]
at (axis cs:35,1.1,0) {Media (M)};
\node[right, inner sep=0mm, text=black]
at (axis cs:70,1.1,0) {Lejos (H)};
\end{axis}
\end{tikzpicture}%
		\caption{$\chi_5$ - Distancia (m)}
		\label{fig:distance-lang-variable}
	\end{subfigure}
	\qquad
	\begin{subfigure}[b]{0.45\textwidth}
		\setlength\figureheight{2.5cm}
		\setlength\figurewidth{6cm}
		% This file was created by matlab2tikz v0.4.7 (commit aa5a2b3a09d5219f83b03f452820de20dfb2a2eb) running on MATLAB 8.0.
% Copyright (c) 2008--2014, Nico Schlömer <nico.schloemer@gmail.com>
% All rights reserved.
% Minimal pgfplots version: 1.3
% 
\begin{tikzpicture}

\begin{axis}[%
width=\figurewidth,
height=\figureheight,
clip=false,
scale only axis,
xmin=0,
xmax=100,
ymin=0,
ymax=1
]
\addplot [color=black,solid,line width=1.4pt,forget plot]
  table[row sep=crcr]{0	1\\
1	0.96\\
2	0.92\\
3	0.88\\
4	0.84\\
5	0.8\\
6	0.76\\
7	0.72\\
8	0.68\\
9	0.64\\
10	0.6\\
11	0.56\\
12	0.52\\
13	0.48\\
14	0.44\\
15	0.4\\
16	0.36\\
17	0.32\\
18	0.28\\
19	0.24\\
20	0.2\\
21	0.16\\
22	0.12\\
23	0.08\\
24	0.04\\
25	0\\
26	0\\
27	0\\
28	0\\
29	0\\
30	0\\
31	0\\
32	0\\
33	0\\
34	0\\
35	0\\
36	0\\
37	0\\
38	0\\
39	0\\
40	0\\
41	0\\
42	0\\
43	0\\
44	0\\
45	0\\
46	0\\
47	0\\
48	0\\
49	0\\
50	0\\
51	0\\
52	0\\
53	0\\
54	0\\
55	0\\
56	0\\
57	0\\
58	0\\
59	0\\
60	0\\
61	0\\
62	0\\
63	0\\
64	0\\
65	0\\
66	0\\
67	0\\
68	0\\
69	0\\
70	0\\
71	0\\
72	0\\
73	0\\
74	0\\
75	0\\
76	0\\
77	0\\
78	0\\
79	0\\
80	0\\
81	0\\
82	0\\
83	0\\
84	0\\
85	0\\
86	0\\
87	0\\
88	0\\
89	0\\
90	0\\
91	0\\
92	0\\
93	0\\
94	0\\
95	0\\
96	0\\
97	0\\
98	0\\
99	0\\
100	0\\
};
\addplot [color=black,solid,line width=1.4pt,forget plot]
  table[row sep=crcr]{0	0\\
1	0.04\\
2	0.08\\
3	0.12\\
4	0.16\\
5	0.2\\
6	0.24\\
7	0.28\\
8	0.32\\
9	0.36\\
10	0.4\\
11	0.44\\
12	0.48\\
13	0.52\\
14	0.56\\
15	0.6\\
16	0.64\\
17	0.68\\
18	0.72\\
19	0.76\\
20	0.8\\
21	0.84\\
22	0.88\\
23	0.92\\
24	0.96\\
25	1\\
26	0.96\\
27	0.92\\
28	0.88\\
29	0.84\\
30	0.8\\
31	0.76\\
32	0.72\\
33	0.68\\
34	0.64\\
35	0.6\\
36	0.56\\
37	0.52\\
38	0.48\\
39	0.44\\
40	0.4\\
41	0.36\\
42	0.32\\
43	0.28\\
44	0.24\\
45	0.2\\
46	0.16\\
47	0.12\\
48	0.08\\
49	0.04\\
50	0\\
51	0\\
52	0\\
53	0\\
54	0\\
55	0\\
56	0\\
57	0\\
58	0\\
59	0\\
60	0\\
61	0\\
62	0\\
63	0\\
64	0\\
65	0\\
66	0\\
67	0\\
68	0\\
69	0\\
70	0\\
71	0\\
72	0\\
73	0\\
74	0\\
75	0\\
76	0\\
77	0\\
78	0\\
79	0\\
80	0\\
81	0\\
82	0\\
83	0\\
84	0\\
85	0\\
86	0\\
87	0\\
88	0\\
89	0\\
90	0\\
91	0\\
92	0\\
93	0\\
94	0\\
95	0\\
96	0\\
97	0\\
98	0\\
99	0\\
100	0\\
};
\addplot [color=black,solid,line width=1.4pt,forget plot]
  table[row sep=crcr]{0	0\\
1	0\\
2	0\\
3	0\\
4	0\\
5	0\\
6	0\\
7	0\\
8	0\\
9	0\\
10	0\\
11	0\\
12	0\\
13	0\\
14	0\\
15	0\\
16	0\\
17	0\\
18	0\\
19	0\\
20	0\\
21	0\\
22	0\\
23	0\\
24	0\\
25	0\\
26	0.04\\
27	0.08\\
28	0.12\\
29	0.16\\
30	0.2\\
31	0.24\\
32	0.28\\
33	0.32\\
34	0.36\\
35	0.4\\
36	0.44\\
37	0.48\\
38	0.52\\
39	0.56\\
40	0.6\\
41	0.64\\
42	0.68\\
43	0.72\\
44	0.76\\
45	0.8\\
46	0.84\\
47	0.88\\
48	0.92\\
49	0.96\\
50	1\\
51	0.96\\
52	0.92\\
53	0.88\\
54	0.84\\
55	0.8\\
56	0.76\\
57	0.72\\
58	0.68\\
59	0.64\\
60	0.6\\
61	0.56\\
62	0.52\\
63	0.48\\
64	0.44\\
65	0.4\\
66	0.36\\
67	0.32\\
68	0.28\\
69	0.24\\
70	0.2\\
71	0.16\\
72	0.12\\
73	0.08\\
74	0.04\\
75	0\\
76	0\\
77	0\\
78	0\\
79	0\\
80	0\\
81	0\\
82	0\\
83	0\\
84	0\\
85	0\\
86	0\\
87	0\\
88	0\\
89	0\\
90	0\\
91	0\\
92	0\\
93	0\\
94	0\\
95	0\\
96	0\\
97	0\\
98	0\\
99	0\\
100	0\\
};
\addplot [color=black,solid,line width=1.4pt,forget plot]
  table[row sep=crcr]{0	0\\
1	0\\
2	0\\
3	0\\
4	0\\
5	0\\
6	0\\
7	0\\
8	0\\
9	0\\
10	0\\
11	0\\
12	0\\
13	0\\
14	0\\
15	0\\
16	0\\
17	0\\
18	0\\
19	0\\
20	0\\
21	0\\
22	0\\
23	0\\
24	0\\
25	0\\
26	0\\
27	0\\
28	0\\
29	0\\
30	0\\
31	0\\
32	0\\
33	0\\
34	0\\
35	0\\
36	0\\
37	0\\
38	0\\
39	0\\
40	0\\
41	0\\
42	0\\
43	0\\
44	0\\
45	0\\
46	0\\
47	0\\
48	0\\
49	0\\
50	0\\
51	0.04\\
52	0.08\\
53	0.12\\
54	0.16\\
55	0.2\\
56	0.24\\
57	0.28\\
58	0.32\\
59	0.36\\
60	0.4\\
61	0.44\\
62	0.48\\
63	0.52\\
64	0.56\\
65	0.6\\
66	0.64\\
67	0.68\\
68	0.72\\
69	0.76\\
70	0.8\\
71	0.84\\
72	0.88\\
73	0.92\\
74	0.96\\
75	1\\
76	0.96\\
77	0.92\\
78	0.88\\
79	0.84\\
80	0.8\\
81	0.76\\
82	0.72\\
83	0.68\\
84	0.64\\
85	0.6\\
86	0.56\\
87	0.52\\
88	0.48\\
89	0.44\\
90	0.4\\
91	0.36\\
92	0.32\\
93	0.28\\
94	0.24\\
95	0.2\\
96	0.16\\
97	0.12\\
98	0.08\\
99	0.04\\
100	0\\
};
\addplot [color=black,solid,line width=1.4pt,forget plot]
  table[row sep=crcr]{0	0\\
1	0\\
2	0\\
3	0\\
4	0\\
5	0\\
6	0\\
7	0\\
8	0\\
9	0\\
10	0\\
11	0\\
12	0\\
13	0\\
14	0\\
15	0\\
16	0\\
17	0\\
18	0\\
19	0\\
20	0\\
21	0\\
22	0\\
23	0\\
24	0\\
25	0\\
26	0\\
27	0\\
28	0\\
29	0\\
30	0\\
31	0\\
32	0\\
33	0\\
34	0\\
35	0\\
36	0\\
37	0\\
38	0\\
39	0\\
40	0\\
41	0\\
42	0\\
43	0\\
44	0\\
45	0\\
46	0\\
47	0\\
48	0\\
49	0\\
50	0\\
51	0\\
52	0\\
53	0\\
54	0\\
55	0\\
56	0\\
57	0\\
58	0\\
59	0\\
60	0\\
61	0\\
62	0\\
63	0\\
64	0\\
65	0\\
66	0\\
67	0\\
68	0\\
69	0\\
70	0\\
71	0\\
72	0\\
73	0\\
74	0\\
75	0\\
76	0.04\\
77	0.08\\
78	0.12\\
79	0.16\\
80	0.2\\
81	0.24\\
82	0.28\\
83	0.32\\
84	0.36\\
85	0.4\\
86	0.44\\
87	0.48\\
88	0.52\\
89	0.56\\
90	0.6\\
91	0.64\\
92	0.68\\
93	0.72\\
94	0.76\\
95	0.8\\
96	0.84\\
97	0.88\\
98	0.92\\
99	0.96\\
100	1\\
};
\node[right, inner sep=0mm, text=black]
at (axis cs:1,1.1,0) {VL};
\node[right, inner sep=0mm, text=black]
at (axis cs:25,1.1,0) {L};
\node[right, inner sep=0mm, text=black]
at (axis cs:50,1.1,0) {M};
\node[right, inner sep=0mm, text=black]
at (axis cs:75,1.1,0) {H};
\node[right, inner sep=0mm, text=black]
at (axis cs:95,1.1,0) {VH};
\end{axis}
\end{tikzpicture}%
		\caption{$y$ - Riesgo de incendio (\%)}
		\label{fig:threat-lang-variable}
	\end{subfigure}
	\caption{Variables lingüísticas utilizadas en la determinación de riesgo de incendios.}
	\label{fig:fire-detection-lang-variables}
\end{figure}

\section{Conjunto de reglas}

Una vez definidas las variables lingüísticas, es necesario construir un conjunto de reglas, que forman la base de conocimiento sobre el problema y permiten, mediante el mecanismo de inferencia, transformar las entradas en salidas. En este caso particular las reglas tienen la forma:

\begin{equation}
\text{IF }\chi_1  \text{ es }A_1\text{ AND }\chi_2 \text{ es }A_2\text{ AND }\chi_3 \text{ es }A_3\text{ AND }\chi_4  \text{ es }A_4\text{ AND }\chi_5 \text{ es }A_5\text{ THEN }y\text{ es }B
\end{equation}

Donde $\chi_1,\ldots,\chi_n$ son las variables de entrada del sistema difuso (definidas en la sección anterior) e $y$ es la salida (riesgo de incendio). Todas las reglas definidas en este sistema harán uso de las 5 variables lingüísticas conectadas mediante el operador de conjunción (AND).

\begin{center}
	\def\arraystretch{1.5}
    \begin{longtable}{| l | l | l | l | l | l |}
    \hline
    $\chi_1=$ Temp. &  $\chi_2=$ Humo &  $\chi_3=$ Luz &  $\chi_4=$ Humedad &  $\chi_5=$ Distancia &  $y=$ Riesgo \\ \hline
    L & L & L & H & H & VL \\ \hline 
    L & L & L & H & M & VL \\ \hline 
    L & L & L & H & L & VL \\ \hline 
    L & L & L & M & H & VL \\ \hline 
    L & L & L & M & M & VL \\ \hline 
    L & L & L & M & L & L \\ \hline 
    L & L & L & L & H & VL \\ \hline 
    L & L & L & L & M & L \\ \hline 
    L & L & L & L & L & L \\ \hline 
    L & L & M & H & H & VL \\ \hline 
    L & L & M & H & M & L \\ \hline 
    L & L & M & H & L & L \\ \hline 
    L & L & M & M & H & VL \\ \hline 
    L & L & M & M & M & L \\ \hline 
    L & L & M & M & L & L \\ \hline 
    L & L & M & L & H & VL \\ \hline 
    L & L & M & L & M & L \\ \hline 
    L & L & M & L & L & L \\ \hline 
    L & L & H & H & H & L \\ \hline 
    L & L & H & H & M & L \\ \hline 
    L & L & H & H & L & L \\ \hline 
    L & L & H & M & H & L \\ \hline 
    L & L & H & M & M & L \\ \hline 
    L & L & H & M & L & M \\ \hline 
    L & L & H & L & H & L \\ \hline 
    L & L & H & L & M & M \\ \hline 
    L & L & H & L & L & M \\ \hline 
    L & M & L & H & H & VL \\ \hline 
    L & M & L & H & M & VL \\ \hline 
    L & M & L & H & L & L \\ \hline 
    L & M & L & M & H & VL \\ \hline 
    L & M & L & M & M & L \\ \hline 
    L & M & L & M & L & L \\ \hline 
    L & M & L & L & H & L \\ \hline 
    L & M & L & L & M & L \\ \hline 
    L & M & L & L & L & M \\ \hline 
    L & M & M & H & H & L \\ \hline 
    L & M & M & H & M & L \\ \hline 
    L & M & M & H & L & M \\ \hline 
    L & M & M & M & H & L \\ \hline 
    L & M & M & M & M & L \\ \hline 
    L & M & M & M & L & M \\ \hline 
    L & M & M & L & H & L \\ \hline 
    L & M & M & L & M & L \\ \hline 
    L & M & M & L & L & M \\ \hline 
    L & M & H & H & H & L \\ \hline 
    L & M & H & H & M & M \\ \hline 
    L & M & H & H & L & M \\ \hline 
    L & M & H & M & H & L \\ \hline 
    L & M & H & M & M & M \\ \hline 
    L & M & H & M & L & M \\ \hline 
    L & M & H & L & H & L \\ \hline 
    L & M & H & L & M & M \\ \hline 
    L & M & H & L & L & H \\ \hline 
    L & H & L & H & H & L \\ \hline 
    L & H & L & H & M & L \\ \hline 
    L & H & L & H & L & M \\ \hline 
    L & H & L & M & H & L \\ \hline 
    L & H & L & M & M & L \\ \hline 
    L & H & L & M & L & M \\ \hline 
    L & H & L & L & H & L \\ \hline 
    L & H & L & L & M & M \\ \hline 
    L & H & L & L & L & M \\ \hline 
    L & H & M & H & H & L \\ \hline 
    L & H & M & H & M & M \\ \hline 
    L & H & M & H & L & M \\ \hline 
    L & H & M & M & H & L \\ \hline 
    L & H & M & M & M & M \\ \hline 
    L & H & M & M & L & M \\ \hline 
    L & H & M & L & H & L \\ \hline 
    L & H & M & L & M & M \\ \hline 
    L & H & M & L & L & M \\ \hline 
    L & H & H & H & H & M \\ \hline 
    L & H & H & H & M & M \\ \hline 
    L & H & H & H & L & M \\ \hline 
    L & H & H & M & H & M \\ \hline 
    L & H & H & M & M & M \\ \hline 
    L & H & H & M & L & H \\ \hline 
    L & H & H & L & H & M \\ \hline 
    L & H & H & L & M & H \\ \hline 
    L & H & H & L & L & H \\ \hline 
    M & L & L & H & H & VL \\ \hline 
    M & L & L & H & M & VL \\ \hline 
    M & L & L & H & L & L \\ \hline 
    M & L & L & M & H & VL \\ \hline 
    M & L & L & M & M & L \\ \hline 
    M & L & L & M & L & L \\ \hline 
    M & L & L & L & H & L \\ \hline 
    M & L & L & L & M & L \\ \hline 
    M & L & L & L & L & M \\ \hline 
    M & L & M & H & H & L \\ \hline 
    M & L & M & H & M & L \\ \hline 
    M & L & M & H & L & M \\ \hline 
    M & L & M & M & H & L \\ \hline 
    M & L & M & M & M & L \\ \hline 
    M & L & M & M & L & M \\ \hline 
    M & L & M & L & H & L \\ \hline 
    M & L & M & L & M & L \\ \hline 
    M & L & M & L & L & M \\ \hline 
    M & L & H & H & H & L \\ \hline 
    M & L & H & H & M & M \\ \hline 
    M & L & H & H & L & M \\ \hline 
    M & L & H & M & H & L \\ \hline 
    M & L & H & M & M & M \\ \hline 
    M & L & H & M & L & M \\ \hline 
    M & L & H & L & H & L \\ \hline 
    M & L & H & L & M & M \\ \hline 
    M & L & H & L & L & H \\ \hline 
    M & M & L & H & H & L \\ \hline 
    M & M & L & H & M & L \\ \hline 
    M & M & L & H & L & M \\ \hline 
    M & M & L & M & H & L \\ \hline 
    M & M & L & M & M & L \\ \hline 
    M & M & L & M & L & M \\ \hline 
    M & M & L & L & H & L \\ \hline 
    M & M & L & L & M & M \\ \hline 
    M & M & L & L & L & M \\ \hline 
    M & M & M & H & H & L \\ \hline 
    M & M & M & H & M & M \\ \hline 
    M & M & M & H & L & M \\ \hline 
    M & M & M & M & H & L \\ \hline 
    M & M & M & M & M & M \\ \hline 
    M & M & M & M & L & M \\ \hline 
    M & M & M & L & H & M \\ \hline 
    M & M & M & L & M & M \\ \hline 
    M & M & M & L & L & M \\ \hline 
    M & M & H & H & H & M \\ \hline 
    M & M & H & H & M & M \\ \hline 
    M & M & H & H & L & H \\ \hline 
    M & M & H & M & H & M \\ \hline 
    M & M & H & M & M & M \\ \hline 
    M & M & H & M & L & H \\ \hline 
    M & M & H & L & H & M \\ \hline 
    M & M & H & L & M & H \\ \hline 
    M & M & H & L & L & H \\ \hline 
    M & H & L & H & H & L \\ \hline 
    M & H & L & H & M & M \\ \hline 
    M & H & L & H & L & M \\ \hline 
    M & H & L & M & H & L \\ \hline 
    M & H & L & M & M & M \\ \hline 
    M & H & L & M & L & M \\ \hline 
    M & H & L & L & H & M \\ \hline 
    M & H & L & L & M & M \\ \hline 
    M & H & L & L & L & H \\ \hline 
    M & H & M & H & H & M \\ \hline 
    M & H & M & H & M & M \\ \hline 
    M & H & M & H & L & M \\ \hline 
    M & H & M & M & H & M \\ \hline 
    M & H & M & M & M & M \\ \hline 
    M & H & M & M & L & H \\ \hline 
    M & H & M & L & H & M \\ \hline 
    M & H & M & L & M & H \\ \hline 
    M & H & M & L & L & H \\ \hline 
    M & H & H & H & H & M \\ \hline 
    M & H & H & H & M & M \\ \hline 
    M & H & H & H & L & H \\ \hline 
    M & H & H & M & H & M \\ \hline 
    M & H & H & M & M & H \\ \hline 
    M & H & H & M & L & H \\ \hline 
    M & H & H & L & H & H \\ \hline 
    M & H & H & L & M & H \\ \hline 
    M & H & H & L & L & VH \\ \hline 
    H & L & L & H & H & L \\ \hline 
    H & L & L & H & M & L \\ \hline 
    H & L & L & H & L & M \\ \hline 
    H & L & L & M & H & L \\ \hline 
    H & L & L & M & M & L \\ \hline 
    H & L & L & M & L & M \\ \hline 
    H & L & L & L & H & L \\ \hline 
    H & L & L & L & M & M \\ \hline 
    H & L & L & L & L & M \\ \hline 
    H & L & M & H & H & L \\ \hline 
    H & L & M & H & M & M \\ \hline 
    H & L & M & H & L & M \\ \hline 
    H & L & M & M & H & L \\ \hline 
    H & L & M & M & M & M \\ \hline 
    H & L & M & M & L & M \\ \hline 
    H & L & M & L & H & L \\ \hline 
    H & L & M & L & M & M \\ \hline 
    H & L & M & L & L & M \\ \hline 
    H & L & H & H & H & M \\ \hline 
    H & L & H & H & M & M \\ \hline 
    H & L & H & H & L & M \\ \hline 
    H & L & H & M & H & M \\ \hline 
    H & L & H & M & M & M \\ \hline 
    H & L & H & M & L & H \\ \hline 
    H & L & H & L & H & M \\ \hline 
    H & L & H & L & M & H \\ \hline 
    H & L & H & L & L & H \\ \hline 
    H & M & L & H & H & L \\ \hline 
    H & M & L & H & M & M \\ \hline 
    H & M & L & H & L & M \\ \hline 
    H & M & L & M & H & L \\ \hline 
    H & M & L & M & M & M \\ \hline 
    H & M & L & M & L & M \\ \hline 
    H & M & L & L & H & M \\ \hline 
    H & M & L & L & M & M \\ \hline 
    H & M & L & L & L & H \\ \hline 
    H & M & M & H & H & M \\ \hline 
    H & M & M & H & M & M \\ \hline 
    H & M & M & H & L & H \\ \hline 
    H & M & M & M & H & M \\ \hline 
    H & M & M & M & M & H \\ \hline 
    H & M & M & M & L & H \\ \hline 
    H & M & M & L & H & M \\ \hline 
    H & M & M & L & M & H \\ \hline 
    H & M & M & L & L & H \\ \hline 
    H & M & H & H & H & H \\ \hline 
    H & M & H & H & M & H \\ \hline 
    H & M & H & H & L & H \\ \hline 
    H & M & H & M & H & H \\ \hline 
    H & M & H & M & M & H \\ \hline 
    H & M & H & M & L & VH \\ \hline 
    H & M & H & L & H & H \\ \hline 
    H & M & H & L & M & VH \\ \hline 
    H & M & H & L & L & VH \\ \hline 
    H & H & L & H & H & M \\ \hline 
    H & H & L & H & M & M \\ \hline 
    H & H & L & H & L & H \\ \hline 
    H & H & L & M & H & M \\ \hline 
    H & H & L & M & M & M \\ \hline 
    H & H & L & M & L & H \\ \hline 
    H & H & L & L & H & H \\ \hline 
    H & H & L & L & M & H \\ \hline 
    H & H & L & L & L & H \\ \hline 
    H & H & M & H & H & M \\ \hline 
    H & H & M & H & M & H \\ \hline 
    H & H & M & H & L & H \\ \hline 
    H & H & M & M & H & H \\ \hline 
    H & H & M & M & M & H \\ \hline 
    H & H & M & M & L & H \\ \hline 
    H & H & M & L & H & H \\ \hline 
    H & H & M & L & M & H \\ \hline 
    H & H & M & L & L & VH \\ \hline 
    H & H & H & H & H & H \\ \hline 
    H & H & H & H & M & H \\ \hline 
    H & H & H & H & L & VH \\ \hline 
    H & H & H & M & H & H \\ \hline 
    H & H & H & M & M & VH \\ \hline 
    H & H & H & M & L & VH \\ \hline 
    H & H & H & L & H & VH \\ \hline 
    H & H & H & L & M & VH \\ \hline 
    H & H & H & L & L & VH \\
    \hline
    \caption{Conjunto de reglas del sistema difuso de detección de incendios forestales.}
    \label{tab:fire-detection-rule-set}
    \end{longtable}
\end{center}

\section{Resultados y comparación de métodos}
\label{sec:fire-detection-results}
En esta sección se presentan los resultados obtenidos al aplicar los métodos descritos en secciones anteriores y se realiza una comparativa de los mismos, usando diferentes t-normas e índices de solapamiento (en el caso del algoritmo \ref{algo:multi-overlap-interpolation-method}). 

\subsection{Método de Mamdani}
El primer método estudiado es el \emph{método de Mamdani}, presentado en la sección \ref{sec:mamdani-controller} (algoritmo \ref{algo:mamdani}). En la figura \ref{fig:mamdani-fire-detection-example} y en la tabla \ref{tab:fire-detection-mamdani-example} se presentan los resultados de aplicar el método de Mamdani, utilizando el conjunto de reglas definidas en la tabla \ref{tab:fire-detection-rule-set} a diferentes entradas. Las entradas del sistema son valores escalares para cada una de las variables lingüísticas:

\begin{center}
	%\def\arraystretch{1.5}
	\begin{longtable}{| c | c | c | c | c | c | c | c | c  | c  | c |}
\hline
\multirow{2}{*}{\textbf{Fig.}} & \multirow{2}{*}{\textbf{Temp.}} & \multirow{2}{*}{\textbf{Humo}} & \multirow{2}{*}{\textbf{Luz}}& \multirow{2}{*}{\textbf{Hum.}} & \multirow{2}{*}{\textbf{Dist.}} &  \multicolumn{5}{|c|}{\textbf{Riesgo (\%)}} \\ 
\cline{7-11}
& & & & & & \textbf{cent.  (\textasteriskcentered)} & \textbf{bis. (+)} & \textbf{som ($\triangledown$)} & \textbf{mom ($\square$)} & \textbf{lom ($\vartriangle$)}  \\ 
\hline
\ref{fig:fire-detection-mamdani-low} & 25 & 0 & 200 & 20 & 70 & 20 & 18 & 0 & 6 & 12 \\ 
\hline
\ref{fig:fire-detection-mamdani-low-medium} & 30 & 20 & 500 & 50 & 40 & 35 & 33 & 10 & 25 & 40 \\
\hline
\ref{fig:fire-detection-mamdani-medium} & 40 & 50 & 500 & 30 & 40 & 43 & 45 & 38 & 50 & 62 \\
\hline
\ref{fig:fire-detection-mamdani-medium-high} & 80 & 80 & 700 & 20 & 30 & 67 & 69 & 63 & 75 & 87 \\
\hline
\ref{fig:fire-detection-mamdani-high} & 100 & 90 & 900 & 10 & 20 & 80 & 82 & 84 & 92 & 100 \\
\hline
\ref{fig:fire-detection-mamdani-very-high} & 120 & 100 & 1000 & 10 & 10 & 91 & 92 & 92 & 96 & 100 \\
\hline
\caption{Resultados de aplicar el método de Mamdani a la detección de incendios forestales. Las gráficas correspondientes se muestran en la figura \ref{fig:mamdani-fire-detection-example}.}
\label{tab:fire-detection-mamdani-example}
\end{longtable}
\end{center}

Los valores de salida (riesgo (\%)) se obtienen aplicando los desdifusificadores definidos en la sección \ref{sec:desdifusificadores} al conjunto difuso obtenido al aplicar el método de Mamdani sobre las entradas (representados en la figura \ref{fig:mamdani-fire-detection-example}). 

Se puede comprobar que los resultados obtenidos con el método de Mamdani parecen ser intuitivamente correctos. En el primer ejemplo (fig. \ref{fig:fire-detection-mamdani-low}) las entradas sugieren un riesgo de incendio bajo, debido a la baja temperatura medida, el humo inexistente y la baja luminosidad. Aplicando el método de Mamdani efectivamente se obtiene un nivel de riesgo bajo (entre 0\% y 20\%). 

Como se puede ver, en general, los valores obtenidos con el desdifusificador \emph{centroide} y el \emph{bisector} son similares y pueden diferir ampliamente de los obtenidos con los métodos basados en el máximo. En este ejemplo se puede comprobar además que los desdifusificadores basados en el máximo pueden no dar resultados intuitivamente acertados, ya que el método del menor de los máximos da un resultado del 0\% de riesgo en el primer ejemplo, a pesar de que las entradas no son mínimas.

A medida que aumentan la temperatura, el humo y la luz y disminuyen la humedad y la distancia el riesgo de incendio aumenta (figuras \ref{fig:fire-detection-mamdani-medium} y \ref{fig:fire-detection-mamdani-high}), lo cual es coherente con lo expresado en el conjunto de reglas.

\begin{figure}[t]
	\centering
	\begin{subfigure}[b]{0.45\textwidth}
		\setlength\figureheight{3cm}
		\setlength\figurewidth{6cm}
		% This file was created by matlab2tikz v0.4.7 running on MATLAB 8.0.
% Copyright (c) 2008--2014, Nico Schlömer <nico.schloemer@gmail.com>
% All rights reserved.
% Minimal pgfplots version: 1.3
% 
%
% defining custom colors
\definecolor{mycolor1}{rgb}{0.00000,1.00000,1.00000}%
\definecolor{mycolor2}{rgb}{1.00000,0.00000,1.00000}%
%
\begin{tikzpicture}

\begin{axis}[%
width=\figurewidth,
height=\figureheight,
scale only axis,
xmin=0,
xmax=100,
xlabel={Riesgo (\%)},
ymin=0,
ymax=0.5,
axis x line*=bottom,
axis y line*=left
]
\addplot [color=black,solid,line width=1.2pt,forget plot]
  table[row sep=crcr]{0	0.5\\
1	0.5\\
2	0.5\\
3	0.5\\
4	0.5\\
5	0.5\\
6	0.5\\
7	0.5\\
8	0.5\\
9	0.5\\
10	0.5\\
11	0.5\\
12	0.5\\
13	0.48\\
14	0.44\\
15	0.4\\
16	0.36\\
17	0.32\\
18	0.3\\
19	0.3\\
20	0.3\\
21	0.3\\
22	0.3\\
23	0.3\\
24	0.3\\
25	0.3\\
26	0.3\\
27	0.3\\
28	0.3\\
29	0.3\\
30	0.3\\
31	0.3\\
32	0.3\\
33	0.3\\
34	0.3\\
35	0.3\\
36	0.3\\
37	0.3\\
38	0.3\\
39	0.3\\
40	0.3\\
41	0.3\\
42	0.3\\
43	0.28\\
44	0.24\\
45	0.2\\
46	0.16\\
47	0.12\\
48	0.08\\
49	0.04\\
50	0\\
51	0\\
52	0\\
53	0\\
54	0\\
55	0\\
56	0\\
57	0\\
58	0\\
59	0\\
60	0\\
61	0\\
62	0\\
63	0\\
64	0\\
65	0\\
66	0\\
67	0\\
68	0\\
69	0\\
70	0\\
71	0\\
72	0\\
73	0\\
74	0\\
75	0\\
76	0\\
77	0\\
78	0\\
79	0\\
80	0\\
81	0\\
82	0\\
83	0\\
84	0\\
85	0\\
86	0\\
87	0\\
88	0\\
89	0\\
90	0\\
91	0\\
92	0\\
93	0\\
94	0\\
95	0\\
96	0\\
97	0\\
98	0\\
99	0\\
100	0\\
};
\addplot [color=blue,line width=1.2pt,mark size=4.5pt,only marks,mark=asterisk,mark options={solid},forget plot]
  table[row sep=crcr]{20	0.3\\
};
\addplot [color=green,line width=1.2pt,mark size=4.5pt,only marks,mark=+,mark options={solid},forget plot]
  table[row sep=crcr]{18	0.3\\
};
\addplot [color=red,line width=1.2pt,mark size=3.2pt,only marks,mark=square,mark options={solid},forget plot]
  table[row sep=crcr]{6	0.5\\
};
\addplot [color=mycolor1,line width=1.2pt,mark size=3.0pt,only marks,mark=triangle,mark options={solid,rotate=180},forget plot]
  table[row sep=crcr]{0	0.5\\
};
\addplot [color=mycolor2,line width=1.2pt,mark size=3.0pt,only marks,mark=triangle,mark options={solid},forget plot]
  table[row sep=crcr]{12	0.5\\
};
\end{axis}
\end{tikzpicture}%
		\caption{Riesgo bajo.}\label{fig:fire-detection-mamdani-low}
	\end{subfigure}
	\qquad
	\begin{subfigure}[b]{0.45\textwidth}
		\setlength\figureheight{3cm}
		\setlength\figurewidth{6cm}
		% This file was created by matlab2tikz v0.4.7 running on MATLAB 8.0.
% Copyright (c) 2008--2014, Nico Schlömer <nico.schloemer@gmail.com>
% All rights reserved.
% Minimal pgfplots version: 1.3
% 
%
% defining custom colors
\definecolor{mycolor1}{rgb}{0.00000,1.00000,1.00000}%
\definecolor{mycolor2}{rgb}{1.00000,0.00000,1.00000}%
%
\begin{tikzpicture}

\begin{axis}[%
width=\figurewidth,
height=\figureheight,
scale only axis,
xmin=0,
xmax=100,
xlabel={Riesgo (\%)},
ymin=0,
ymax=0.4,
axis x line*=bottom,
axis y line*=left
]
\addplot [color=black,solid,line width=1.2pt,forget plot]
  table[row sep=crcr]{0	0\\
1	0.04\\
2	0.08\\
3	0.12\\
4	0.16\\
5	0.2\\
6	0.24\\
7	0.28\\
8	0.32\\
9	0.36\\
10	0.4\\
11	0.4\\
12	0.4\\
13	0.4\\
14	0.4\\
15	0.4\\
16	0.4\\
17	0.4\\
18	0.4\\
19	0.4\\
20	0.4\\
21	0.4\\
22	0.4\\
23	0.4\\
24	0.4\\
25	0.4\\
26	0.4\\
27	0.4\\
28	0.4\\
29	0.4\\
30	0.4\\
31	0.4\\
32	0.4\\
33	0.4\\
34	0.4\\
35	0.4\\
36	0.4\\
37	0.4\\
38	0.4\\
39	0.4\\
40	0.4\\
41	0.36\\
42	0.32\\
43	0.28\\
44	0.25\\
45	0.25\\
46	0.25\\
47	0.25\\
48	0.25\\
49	0.25\\
50	0.25\\
51	0.25\\
52	0.25\\
53	0.25\\
54	0.25\\
55	0.25\\
56	0.25\\
57	0.25\\
58	0.25\\
59	0.25\\
60	0.25\\
61	0.25\\
62	0.25\\
63	0.25\\
64	0.25\\
65	0.25\\
66	0.25\\
67	0.25\\
68	0.25\\
69	0.24\\
70	0.2\\
71	0.16\\
72	0.12\\
73	0.08\\
74	0.04\\
75	0\\
76	0\\
77	0\\
78	0\\
79	0\\
80	0\\
81	0\\
82	0\\
83	0\\
84	0\\
85	0\\
86	0\\
87	0\\
88	0\\
89	0\\
90	0\\
91	0\\
92	0\\
93	0\\
94	0\\
95	0\\
96	0\\
97	0\\
98	0\\
99	0\\
100	0\\
};
\addplot [color=blue,line width=1.2pt,mark size=4.5pt,only marks,mark=asterisk,mark options={solid},forget plot]
  table[row sep=crcr]{35	0.4\\
};
\addplot [color=green,line width=1.2pt,mark size=4.5pt,only marks,mark=+,mark options={solid},forget plot]
  table[row sep=crcr]{33	0.4\\
};
\addplot [color=red,line width=1.2pt,mark size=3.2pt,only marks,mark=square,mark options={solid},forget plot]
  table[row sep=crcr]{25	0.4\\
};
\addplot [color=mycolor1,line width=1.2pt,mark size=3.0pt,only marks,mark=triangle,mark options={solid,rotate=180},forget plot]
  table[row sep=crcr]{10	0.4\\
};
\addplot [color=mycolor2,line width=1.2pt,mark size=3.0pt,only marks,mark=triangle,mark options={solid},forget plot]
  table[row sep=crcr]{40	0.4\\
};
\end{axis}
\end{tikzpicture}%
		\caption{Riesgo bajo-medio.}\label{fig:fire-detection-mamdani-low-medium}
	\end{subfigure}
	
	\vspace{1cm}
	\begin{subfigure}[b]{0.45\textwidth}
		\setlength\figureheight{3cm}
		\setlength\figurewidth{6cm}
		% This file was created by matlab2tikz v0.4.7 running on MATLAB 8.0.
% Copyright (c) 2008--2014, Nico Schlömer <nico.schloemer@gmail.com>
% All rights reserved.
% Minimal pgfplots version: 1.3
% 
%
% defining custom colors
\definecolor{mycolor1}{rgb}{0.00000,1.00000,1.00000}%
\definecolor{mycolor2}{rgb}{1.00000,0.00000,1.00000}%
%
\begin{tikzpicture}

\begin{axis}[%
width=\figurewidth,
height=\figureheight,
scale only axis,
xmin=0,
xmax=100,
xlabel={Riesgo (\%)},
ymin=0,
ymax=0.5,
axis x line*=bottom,
axis y line*=left
]
\addplot [color=black,solid,line width=1.2pt,forget plot]
  table[row sep=crcr]{0	0\\
1	0.04\\
2	0.08\\
3	0.12\\
4	0.16\\
5	0.2\\
6	0.2\\
7	0.2\\
8	0.2\\
9	0.2\\
10	0.2\\
11	0.2\\
12	0.2\\
13	0.2\\
14	0.2\\
15	0.2\\
16	0.2\\
17	0.2\\
18	0.2\\
19	0.2\\
20	0.2\\
21	0.2\\
22	0.2\\
23	0.2\\
24	0.2\\
25	0.2\\
26	0.2\\
27	0.2\\
28	0.2\\
29	0.2\\
30	0.2\\
31	0.24\\
32	0.28\\
33	0.32\\
34	0.36\\
35	0.4\\
36	0.44\\
37	0.48\\
38	0.5\\
39	0.5\\
40	0.5\\
41	0.5\\
42	0.5\\
43	0.5\\
44	0.5\\
45	0.5\\
46	0.5\\
47	0.5\\
48	0.5\\
49	0.5\\
50	0.5\\
51	0.5\\
52	0.5\\
53	0.5\\
54	0.5\\
55	0.5\\
56	0.5\\
57	0.5\\
58	0.5\\
59	0.5\\
60	0.5\\
61	0.5\\
62	0.5\\
63	0.48\\
64	0.44\\
65	0.4\\
66	0.36\\
67	0.32\\
68	0.28\\
69	0.24\\
70	0.2\\
71	0.16\\
72	0.12\\
73	0.08\\
74	0.04\\
75	0\\
76	0\\
77	0\\
78	0\\
79	0\\
80	0\\
81	0\\
82	0\\
83	0\\
84	0\\
85	0\\
86	0\\
87	0\\
88	0\\
89	0\\
90	0\\
91	0\\
92	0\\
93	0\\
94	0\\
95	0\\
96	0\\
97	0\\
98	0\\
99	0\\
100	0\\
};
\addplot [color=blue,line width=1.2pt,mark size=4.5pt,only marks,mark=asterisk,mark options={solid},forget plot]
  table[row sep=crcr]{43	0.5\\
};
\addplot [color=green,line width=1.2pt,mark size=4.5pt,only marks,mark=+,mark options={solid},forget plot]
  table[row sep=crcr]{45	0.5\\
};
\addplot [color=red,line width=1.2pt,mark size=3.2pt,only marks,mark=square,mark options={solid},forget plot]
  table[row sep=crcr]{50	0.5\\
};
\addplot [color=mycolor1,line width=1.2pt,mark size=3.0pt,only marks,mark=triangle,mark options={solid,rotate=180},forget plot]
  table[row sep=crcr]{38	0.5\\
};
\addplot [color=mycolor2,line width=1.2pt,mark size=3.0pt,only marks,mark=triangle,mark options={solid},forget plot]
  table[row sep=crcr]{62	0.5\\
};
\end{axis}
\end{tikzpicture}%
		\caption{Riesgo medio.}\label{fig:fire-detection-mamdani-medium}
	\end{subfigure}
	\qquad
	\begin{subfigure}[b]{0.45\textwidth}
		\setlength\figureheight{3cm}
		\setlength\figurewidth{6cm}
		% This file was created by matlab2tikz v0.4.7 running on MATLAB 8.0.
% Copyright (c) 2008--2014, Nico Schlömer <nico.schloemer@gmail.com>
% All rights reserved.
% Minimal pgfplots version: 1.3
% 
%
% defining custom colors
\definecolor{mycolor1}{rgb}{0.00000,1.00000,1.00000}%
\definecolor{mycolor2}{rgb}{1.00000,0.00000,1.00000}%
%
\begin{tikzpicture}

\begin{axis}[%
width=\figurewidth,
height=\figureheight,
scale only axis,
xmin=0,
xmax=100,
xlabel={Riesgo (\%)},
ymin=0,
ymax=0.5,
axis x line*=bottom,
axis y line*=left
]
\addplot [color=black,solid,line width=1.2pt,forget plot]
  table[row sep=crcr]{0	0\\
1	0\\
2	0\\
3	0\\
4	0\\
5	0\\
6	0\\
7	0\\
8	0\\
9	0\\
10	0\\
11	0\\
12	0\\
13	0\\
14	0\\
15	0\\
16	0\\
17	0\\
18	0\\
19	0\\
20	0\\
21	0\\
22	0\\
23	0\\
24	0\\
25	0\\
26	0.04\\
27	0.08\\
28	0.12\\
29	0.16\\
30	0.2\\
31	0.24\\
32	0.25\\
33	0.25\\
34	0.25\\
35	0.25\\
36	0.25\\
37	0.25\\
38	0.25\\
39	0.25\\
40	0.25\\
41	0.25\\
42	0.25\\
43	0.25\\
44	0.25\\
45	0.25\\
46	0.25\\
47	0.25\\
48	0.25\\
49	0.25\\
50	0.25\\
51	0.25\\
52	0.25\\
53	0.25\\
54	0.25\\
55	0.25\\
56	0.25\\
57	0.28\\
58	0.32\\
59	0.36\\
60	0.4\\
61	0.44\\
62	0.48\\
63	0.5\\
64	0.5\\
65	0.5\\
66	0.5\\
67	0.5\\
68	0.5\\
69	0.5\\
70	0.5\\
71	0.5\\
72	0.5\\
73	0.5\\
74	0.5\\
75	0.5\\
76	0.5\\
77	0.5\\
78	0.5\\
79	0.5\\
80	0.5\\
81	0.5\\
82	0.5\\
83	0.5\\
84	0.5\\
85	0.5\\
86	0.5\\
87	0.5\\
88	0.48\\
89	0.44\\
90	0.4\\
91	0.36\\
92	0.32\\
93	0.28\\
94	0.24\\
95	0.2\\
96	0.2\\
97	0.2\\
98	0.2\\
99	0.2\\
100	0.2\\
};
\addplot [color=blue,line width=1.2pt,mark size=4.5pt,only marks,mark=asterisk,mark options={solid},forget plot]
  table[row sep=crcr]{67	0.5\\
};
\addplot [color=green,line width=1.2pt,mark size=4.5pt,only marks,mark=+,mark options={solid},forget plot]
  table[row sep=crcr]{69	0.5\\
};
\addplot [color=red,line width=1.2pt,mark size=3.2pt,only marks,mark=square,mark options={solid},forget plot]
  table[row sep=crcr]{75	0.5\\
};
\addplot [color=mycolor1,line width=1.2pt,mark size=3.0pt,only marks,mark=triangle,mark options={solid,rotate=180},forget plot]
  table[row sep=crcr]{63	0.5\\
};
\addplot [color=mycolor2,line width=1.2pt,mark size=3.0pt,only marks,mark=triangle,mark options={solid},forget plot]
  table[row sep=crcr]{87	0.5\\
};
\end{axis}
\end{tikzpicture}%
		\caption{Riesgo medio-alto.}\label{fig:fire-detection-mamdani-medium-high}
	\end{subfigure}
	
		\vspace{1cm}
		\begin{subfigure}[b]{0.45\textwidth}
			\setlength\figureheight{3cm}
			\setlength\figurewidth{6cm}
			% This file was created by matlab2tikz v0.4.7 running on MATLAB 8.0.
% Copyright (c) 2008--2014, Nico Schlömer <nico.schloemer@gmail.com>
% All rights reserved.
% Minimal pgfplots version: 1.3
% 
%
% defining custom colors
\definecolor{mycolor1}{rgb}{0.00000,1.00000,1.00000}%
\definecolor{mycolor2}{rgb}{1.00000,0.00000,1.00000}%
%
\begin{tikzpicture}

\begin{axis}[%
width=\figurewidth,
height=\figureheight,
scale only axis,
xmin=0,
xmax=100,
xlabel={Riesgo (\%)},
ymin=0,
ymax=0.35,
axis x line*=bottom,
axis y line*=left
]
\addplot [color=black,solid,line width=1.2pt,forget plot]
  table[row sep=crcr]{0	0\\
1	0\\
2	0\\
3	0\\
4	0\\
5	0\\
6	0\\
7	0\\
8	0\\
9	0\\
10	0\\
11	0\\
12	0\\
13	0\\
14	0\\
15	0\\
16	0\\
17	0\\
18	0\\
19	0\\
20	0\\
21	0\\
22	0\\
23	0\\
24	0\\
25	0\\
26	0\\
27	0\\
28	0\\
29	0\\
30	0\\
31	0\\
32	0\\
33	0\\
34	0\\
35	0\\
36	0\\
37	0\\
38	0\\
39	0\\
40	0\\
41	0\\
42	0\\
43	0\\
44	0\\
45	0\\
46	0\\
47	0\\
48	0\\
49	0\\
50	0\\
51	0.04\\
52	0.08\\
53	0.12\\
54	0.16\\
55	0.2\\
56	0.2\\
57	0.2\\
58	0.2\\
59	0.2\\
60	0.2\\
61	0.2\\
62	0.2\\
63	0.2\\
64	0.2\\
65	0.2\\
66	0.2\\
67	0.2\\
68	0.2\\
69	0.2\\
70	0.2\\
71	0.2\\
72	0.2\\
73	0.2\\
74	0.2\\
75	0.2\\
76	0.2\\
77	0.2\\
78	0.2\\
79	0.2\\
80	0.2\\
81	0.24\\
82	0.28\\
83	0.32\\
84	0.333333333333333\\
85	0.333333333333333\\
86	0.333333333333333\\
87	0.333333333333333\\
88	0.333333333333333\\
89	0.333333333333333\\
90	0.333333333333333\\
91	0.333333333333333\\
92	0.333333333333333\\
93	0.333333333333333\\
94	0.333333333333333\\
95	0.333333333333333\\
96	0.333333333333333\\
97	0.333333333333333\\
98	0.333333333333333\\
99	0.333333333333333\\
100	0.333333333333333\\
};
\addplot [color=blue,line width=1.2pt,mark size=4.5pt,only marks,mark=asterisk,mark options={solid},forget plot]
  table[row sep=crcr]{80	0.2\\
};
\addplot [color=green,line width=1.2pt,mark size=4.5pt,only marks,mark=+,mark options={solid},forget plot]
  table[row sep=crcr]{82	0.28\\
};
\addplot [color=red,line width=1.2pt,mark size=3.2pt,only marks,mark=square,mark options={solid},forget plot]
  table[row sep=crcr]{92	0.333333333333333\\
};
\addplot [color=mycolor1,line width=1.2pt,mark size=3.0pt,only marks,mark=triangle,mark options={solid,rotate=180},forget plot]
  table[row sep=crcr]{84	0.333333333333333\\
};
\addplot [color=mycolor2,line width=1.2pt,mark size=3.0pt,only marks,mark=triangle,mark options={solid},forget plot]
  table[row sep=crcr]{100	0.333333333333333\\
};
\end{axis}
\end{tikzpicture}%
			\caption{Riesgo alto.}\label{fig:fire-detection-mamdani-high}
		\end{subfigure}
		\qquad
		\begin{subfigure}[b]{0.45\textwidth}
			\setlength\figureheight{3cm}
			\setlength\figurewidth{6cm}
			% This file was created by matlab2tikz v0.4.7 running on MATLAB 8.0.
% Copyright (c) 2008--2014, Nico Schlömer <nico.schloemer@gmail.com>
% All rights reserved.
% Minimal pgfplots version: 1.3
% 
%
% defining custom colors
\definecolor{mycolor1}{rgb}{0.00000,1.00000,1.00000}%
\definecolor{mycolor2}{rgb}{1.00000,0.00000,1.00000}%
%
\begin{tikzpicture}

\begin{axis}[%
width=\figurewidth,
height=\figureheight,
scale only axis,
xmin=0,
xmax=100,
xlabel={Riesgo (\%)},
ymin=0,
ymax=0.7,
axis x line*=bottom,
axis y line*=left
]
\addplot [color=black,solid,line width=1.2pt,forget plot]
  table[row sep=crcr]{0	0\\
1	0\\
2	0\\
3	0\\
4	0\\
5	0\\
6	0\\
7	0\\
8	0\\
9	0\\
10	0\\
11	0\\
12	0\\
13	0\\
14	0\\
15	0\\
16	0\\
17	0\\
18	0\\
19	0\\
20	0\\
21	0\\
22	0\\
23	0\\
24	0\\
25	0\\
26	0\\
27	0\\
28	0\\
29	0\\
30	0\\
31	0\\
32	0\\
33	0\\
34	0\\
35	0\\
36	0\\
37	0\\
38	0\\
39	0\\
40	0\\
41	0\\
42	0\\
43	0\\
44	0\\
45	0\\
46	0\\
47	0\\
48	0\\
49	0\\
50	0\\
51	0\\
52	0\\
53	0\\
54	0\\
55	0\\
56	0\\
57	0\\
58	0\\
59	0\\
60	0\\
61	0\\
62	0\\
63	0\\
64	0\\
65	0\\
66	0\\
67	0\\
68	0\\
69	0\\
70	0\\
71	0\\
72	0\\
73	0\\
74	0\\
75	0\\
76	0.04\\
77	0.08\\
78	0.12\\
79	0.16\\
80	0.2\\
81	0.24\\
82	0.28\\
83	0.32\\
84	0.36\\
85	0.4\\
86	0.44\\
87	0.48\\
88	0.52\\
89	0.56\\
90	0.6\\
91	0.64\\
92	0.666666666666667\\
93	0.666666666666667\\
94	0.666666666666667\\
95	0.666666666666667\\
96	0.666666666666667\\
97	0.666666666666667\\
98	0.666666666666667\\
99	0.666666666666667\\
100	0.666666666666667\\
};
\addplot [color=blue,line width=1.2pt,mark size=4.5pt,only marks,mark=asterisk,mark options={solid},forget plot]
  table[row sep=crcr]{91	0.64\\
};
\addplot [color=green,line width=1.2pt,mark size=4.5pt,only marks,mark=+,mark options={solid},forget plot]
  table[row sep=crcr]{92	0.666666666666667\\
};
\addplot [color=red,line width=1.2pt,mark size=3.2pt,only marks,mark=square,mark options={solid},forget plot]
  table[row sep=crcr]{96	0.666666666666667\\
};
\addplot [color=mycolor1,line width=1.2pt,mark size=3.0pt,only marks,mark=triangle,mark options={solid,rotate=180},forget plot]
  table[row sep=crcr]{92	0.666666666666667\\
};
\addplot [color=mycolor2,line width=1.2pt,mark size=3.0pt,only marks,mark=triangle,mark options={solid},forget plot]
  table[row sep=crcr]{100	0.666666666666667\\
};
\end{axis}
\end{tikzpicture}%
			\caption{Riesgo muy alto.}\label{fig:fire-detection-mamdani-very-high}
		\end{subfigure}
	\caption{Método de Mamdani aplicado a la detección de incendios forestales.}
	\label{fig:mamdani-fire-detection-example}
\end{figure}

\subsection{Método de interpolación basado en índices de solapamiento}\label{sec:fire-detection-interpolation-method}
Para ejecutar el algoritmo \ref{algo:multi-overlap-interpolation-method} es necesario seleccionar una t-norma(\emph{T}), un operador de agregación(\emph{M}) y un índice de solapamiento (\emph{O}). En este estudio se va a utilizar la media aritmética como operador de agregación. En cuanto a las t-normas, se van a utilizar las siguientes:

\begin{itemize}
	\item \bfseries Mínimo: $T_{min}(x_{1},x_{2},\cdots,x_{n}) = \min\{x_{1},x_{2},\cdots,x_{n}\}$
	\item \bfseries Producto: $T_{prod}(x_{1},x_{2},\cdots,x_{n}) = x_{1}*x_{2} * \cdots * x_{n}$
	\item \bfseries Media geométrica: $T_{geo}(x_{1},x_{2},\cdots,x_{n}) = (\prod\limits_{i=1}^{n}x_{i})^{\frac{1}{n}}$
	\item \bfseries Media armónica: $T_{harm}(x_{1},x_{2},\cdots,x_{n}) = \frac{n}{\sum_{i=1}^{n}\frac{1}{x_i}}$
	\item \bfseries Seno: $T_{sin}(x_{1},x_{2},\cdots,x_{n}) = \frac{1}{n}\sin(\frac{\pi}{2}(x_1 * x_2 * \ldots * x_n))^{\frac{1}{4}}$
	\item \bfseries Einstein: $T_{einstein}(x_{1},x_{2},\cdots,x_{n}) = \frac{(x_{1}*x_{2} * \cdots * x_{n})}{1 + (1 - x_{1} )* \ldots * (1 - x_{n} )}$
\end{itemize}

Además se van a utilizar los siguientes índices de solapamiento:

\begin{itemize}
	\item \bfseries Media de productos: $O_{\pi}(A,B) = \frac{1}{n}\sum_{i=1}^{n}\mu_A(x_i)*\mu_B(x_i)$
	\item \bfseries Media de mínimos: $O_{avgmin}(A,B) = \frac{1}{n}\sum_{i=1}^{n}\min(\mu_A(x_i),\mu_B(x_i))$
	\item \bfseries Máximo de mínimos: $O_{Z}(A,B) = \max\limits_{i=1}^{n}(\min(\mu_A(x_i),\mu_B(x_i)))$
	\item \bfseries Media de raíces cuadradas: $O_{\sqrt{\text{ }}}(A,B) =  \frac{1}{n}\sum_{i=1}^{n}\sqrt{\mu_A(x_i)*\mu_B(x_i)}$
	\item \bfseries Media del seno:  $O_{sin}(A, B) = \frac{1}{n}\sum_{i=1}^{n}\sin(\frac{\pi}{2}(\mu_A(x_i)*\mu_B(x_i)))^{\frac{1}{4}}$
\end{itemize}

Así por ejemplo, si tomamos $M = \text{Media aritmética}$, $T = T_{min}$ y $O = O_Z$ y aplicamos el método de interpolación(algoritmo \ref{algo:multi-overlap-interpolation-method}) a las entradas de la tabla \ref{tab:fire-detection-mamdani-example} se obtienen los resultados mostrados en la tabla \ref{tab:fire-detection-interpolation-example} y la figura \ref{fig:interpolation-fire-detection-example}:

\begin{center}
	%\def\arraystretch{1.5}
    \begin{longtable}{| c | c | c | c | c | c | c | c | c  | c  | c |}
    \hline
    \multirow{2}{*}{\textbf{Fig.}} & \multirow{2}{*}{\textbf{Temp.}} & \multirow{2}{*}{\textbf{Humo}} & \multirow{2}{*}{\textbf{Luz}}& \multirow{2}{*}{\textbf{Hum.}} & \multirow{2}{*}{\textbf{Dist.}} &  \multicolumn{5}{|c|}{\textbf{Riesgo (\%)}} \\ 
    \cline{7-11}
    & & & & & & \textbf{cent.  (\textasteriskcentered)} & \textbf{bis. (+)} & \textbf{som ($\triangledown$)} & \textbf{mom ($\square$)} & \textbf{lom ($\vartriangle$)}  \\ 
    \hline
    \ref{fig:fire-detection-interpolation-low} & 25 & 0 & 200 & 20 & 70 & 18 & 15 & 8 & 10 & 12  \\ 
    \hline
    \ref{fig:fire-detection-interpolation-low-medium} & 30 & 20 & 500 & 50 & 40 & 30 & 30 & 32 & 36 & 40 \\
    \hline
    \ref{fig:fire-detection-interpolation-medium} & 40 & 50 & 500 & 30 & 40 & 41 & 41 & 38 & 42 & 45 \\
    \hline
    \ref{fig:fire-detection-interpolation-medium-high} & 80 & 80 & 700 & 20 & 30 & 69 & 69 & 63 & 66 & 68 \\
    \hline
    \ref{fig:fire-detection-interpolation-high} & 100 & 90 & 900 & 10 & 20 & 85 & 87 & 84 & 90 & 95 \\
    \hline
    \ref{fig:fire-detection-interpolation-very-high} & 120 & 100 & 1000 & 10 & 10 & 91 & 92 & 92 & 96 & 100 \\
    \hline
    \caption{Resultados de aplicar el método de interpolación con $M = \text{Media aritmética}$, $T = T_{min}$ y $O = O_Z$  a la detección de incendios forestales. Las gráficas correspondientes se muestran en la figura \ref{fig:interpolation-fire-detection-example}.}
    \end{longtable}
\end{center}

\begin{figure}[t]
	\centering
	\begin{subfigure}[b]{0.45\textwidth}
		\setlength\figureheight{3cm}
		\setlength\figurewidth{6cm}
		% This file was created by matlab2tikz v0.4.7 running on MATLAB 8.0.
% Copyright (c) 2008--2014, Nico Schlömer <nico.schloemer@gmail.com>
% All rights reserved.
% Minimal pgfplots version: 1.3
% 
%
% defining custom colors
\definecolor{mycolor1}{rgb}{0.00000,1.00000,1.00000}%
\definecolor{mycolor2}{rgb}{1.00000,0.00000,1.00000}%
%
\begin{tikzpicture}

\begin{axis}[%
width=\figurewidth,
height=\figureheight,
scale only axis,
xmin=0,
xmax=100,
xlabel={Riesgo (\%)},
ymin=0,
ymax=0.01,
axis x line*=bottom,
axis y line*=left
]
\addplot [color=black,solid,line width=1.2pt,forget plot]
  table[row sep=crcr]{0	0.00617283950617284\\
1	0.00666666666666667\\
2	0.0071604938271605\\
3	0.00765432098765432\\
4	0.00814814814814815\\
5	0.00864197530864198\\
6	0.0091358024691358\\
7	0.00938271604938272\\
8	0.00946502057613169\\
9	0.00946502057613169\\
10	0.00946502057613169\\
11	0.00946502057613169\\
12	0.00946502057613169\\
13	0.00938271604938272\\
14	0.00921810699588477\\
15	0.00905349794238683\\
16	0.00888888888888889\\
17	0.00872427983539095\\
18	0.008559670781893\\
19	0.00823045267489712\\
20	0.00740740740740741\\
21	0.0065843621399177\\
22	0.00576131687242798\\
23	0.00493827160493827\\
24	0.00411522633744856\\
25	0.00329218106995885\\
26	0.00329218106995885\\
27	0.00329218106995885\\
28	0.00329218106995885\\
29	0.00329218106995885\\
30	0.00329218106995885\\
31	0.00329218106995885\\
32	0.00329218106995885\\
33	0.00329218106995885\\
34	0.00329218106995885\\
35	0.00329218106995885\\
36	0.00329218106995885\\
37	0.00329218106995885\\
38	0.00329218106995885\\
39	0.00329218106995885\\
40	0.00329218106995885\\
41	0.00329218106995885\\
42	0.00329218106995885\\
43	0.00320987654320988\\
44	0.00296296296296296\\
45	0.00246913580246914\\
46	0.00197530864197531\\
47	0.00148148148148148\\
48	0.000987654320987654\\
49	0.000493827160493828\\
50	0\\
51	0\\
52	0\\
53	0\\
54	0\\
55	0\\
56	0\\
57	0\\
58	0\\
59	0\\
60	0\\
61	0\\
62	0\\
63	0\\
64	0\\
65	0\\
66	0\\
67	0\\
68	0\\
69	0\\
70	0\\
71	0\\
72	0\\
73	0\\
74	0\\
75	0\\
76	0\\
77	0\\
78	0\\
79	0\\
80	0\\
81	0\\
82	0\\
83	0\\
84	0\\
85	0\\
86	0\\
87	0\\
88	0\\
89	0\\
90	0\\
91	0\\
92	0\\
93	0\\
94	0\\
95	0\\
96	0\\
97	0\\
98	0\\
99	0\\
100	0\\
};
\addplot [color=blue,line width=1.2pt,mark size=4.5pt,only marks,mark=asterisk,mark options={solid},forget plot]
  table[row sep=crcr]{18	0.008559670781893\\
};
\addplot [color=green,line width=1.2pt,mark size=4.5pt,only marks,mark=+,mark options={solid},forget plot]
  table[row sep=crcr]{15	0.00905349794238683\\
};
\addplot [color=red,line width=1.2pt,mark size=3.2pt,only marks,mark=square,mark options={solid},forget plot]
  table[row sep=crcr]{10	0.00946502057613169\\
};
\addplot [color=mycolor1,line width=1.2pt,mark size=3.0pt,only marks,mark=triangle,mark options={solid,rotate=180},forget plot]
  table[row sep=crcr]{8	0.00946502057613169\\
};
\addplot [color=mycolor2,line width=1.2pt,mark size=3.0pt,only marks,mark=triangle,mark options={solid},forget plot]
  table[row sep=crcr]{12	0.00946502057613169\\
};
\end{axis}
\end{tikzpicture}%
		\caption{Riesgo bajo.}\label{fig:fire-detection-interpolation-low}
	\end{subfigure}
	\qquad
	\begin{subfigure}[b]{0.45\textwidth}
		\setlength\figureheight{3cm}
		\setlength\figurewidth{6cm}
		% This file was created by matlab2tikz v0.4.7 running on MATLAB 8.0.
% Copyright (c) 2008--2014, Nico Schlömer <nico.schloemer@gmail.com>
% All rights reserved.
% Minimal pgfplots version: 1.3
% 
%
% defining custom colors
\definecolor{mycolor1}{rgb}{0.00000,1.00000,1.00000}%
\definecolor{mycolor2}{rgb}{1.00000,0.00000,1.00000}%
%
\begin{tikzpicture}

\begin{axis}[%
width=\figurewidth,
height=\figureheight,
scale only axis,
xmin=0,
xmax=100,
xlabel={Riesgo (\%)},
ymin=0,
ymax=0.006,
axis x line*=bottom,
axis y line*=left
]
\addplot [color=black,solid,line width=1.2pt,forget plot]
  table[row sep=crcr]{0	0\\
1	0.000493827160493827\\
2	0.000987654320987654\\
3	0.00148148148148148\\
4	0.00197530864197531\\
5	0.00246913580246914\\
6	0.00296296296296296\\
7	0.00333333333333333\\
8	0.00366255144032922\\
9	0.0039917695473251\\
10	0.00432098765432099\\
11	0.00432098765432099\\
12	0.00432098765432099\\
13	0.00432098765432099\\
14	0.00432098765432099\\
15	0.00432098765432099\\
16	0.00432098765432099\\
17	0.00432098765432099\\
18	0.00432098765432099\\
19	0.00432098765432099\\
20	0.00432098765432099\\
21	0.00432098765432099\\
22	0.00432098765432099\\
23	0.00432098765432099\\
24	0.00432098765432099\\
25	0.00432098765432099\\
26	0.00448559670781893\\
27	0.00465020576131687\\
28	0.00481481481481481\\
29	0.00497942386831276\\
30	0.0051440329218107\\
31	0.00530864197530864\\
32	0.00534979423868313\\
33	0.00534979423868313\\
34	0.00534979423868313\\
35	0.00534979423868313\\
36	0.00534979423868313\\
37	0.00534979423868313\\
38	0.00534979423868313\\
39	0.00534979423868313\\
40	0.00534979423868313\\
41	0.00502057613168724\\
42	0.00469135802469136\\
43	0.00436213991769547\\
44	0.0039917695473251\\
45	0.00349794238683128\\
46	0.00300411522633745\\
47	0.00251028806584362\\
48	0.00201646090534979\\
49	0.00152263374485597\\
50	0.00102880658436214\\
51	0.00102880658436214\\
52	0.00102880658436214\\
53	0.00102880658436214\\
54	0.00102880658436214\\
55	0.00102880658436214\\
56	0.00102880658436214\\
57	0.00102880658436214\\
58	0.00102880658436214\\
59	0.00102880658436214\\
60	0.00102880658436214\\
61	0.00102880658436214\\
62	0.00102880658436214\\
63	0.00102880658436214\\
64	0.00102880658436214\\
65	0.00102880658436214\\
66	0.00102880658436214\\
67	0.00102880658436214\\
68	0.00102880658436214\\
69	0.000987654320987654\\
70	0.000823045267489712\\
71	0.00065843621399177\\
72	0.000493827160493827\\
73	0.000329218106995885\\
74	0.000164609053497943\\
75	0\\
76	0\\
77	0\\
78	0\\
79	0\\
80	0\\
81	0\\
82	0\\
83	0\\
84	0\\
85	0\\
86	0\\
87	0\\
88	0\\
89	0\\
90	0\\
91	0\\
92	0\\
93	0\\
94	0\\
95	0\\
96	0\\
97	0\\
98	0\\
99	0\\
100	0\\
};
\addplot [color=blue,line width=1.2pt,mark size=4.5pt,only marks,mark=asterisk,mark options={solid},forget plot]
  table[row sep=crcr]{30	0.0051440329218107\\
};
\addplot [color=green,line width=1.2pt,mark size=4.5pt,only marks,mark=+,mark options={solid},forget plot]
  table[row sep=crcr]{30	0.0051440329218107\\
};
\addplot [color=red,line width=1.2pt,mark size=3.2pt,only marks,mark=square,mark options={solid},forget plot]
  table[row sep=crcr]{36	0.00534979423868313\\
};
\addplot [color=mycolor1,line width=1.2pt,mark size=3.0pt,only marks,mark=triangle,mark options={solid,rotate=180},forget plot]
  table[row sep=crcr]{32	0.00534979423868313\\
};
\addplot [color=mycolor2,line width=1.2pt,mark size=3.0pt,only marks,mark=triangle,mark options={solid},forget plot]
  table[row sep=crcr]{40	0.00534979423868313\\
};
\end{axis}
\end{tikzpicture}%
		\caption{Riesgo bajo-medio.}\label{fig:fire-detection-interpolation-low-medium}
	\end{subfigure}
	
	\vspace{1cm}
	\begin{subfigure}[b]{0.45\textwidth}
		\setlength\figureheight{3cm}
		\setlength\figurewidth{6cm}
		% This file was created by matlab2tikz v0.4.7 running on MATLAB 8.0.
% Copyright (c) 2008--2014, Nico Schlömer <nico.schloemer@gmail.com>
% All rights reserved.
% Minimal pgfplots version: 1.3
% 
%
% defining custom colors
\definecolor{mycolor1}{rgb}{0.00000,1.00000,1.00000}%
\definecolor{mycolor2}{rgb}{1.00000,0.00000,1.00000}%
%
\begin{tikzpicture}

\begin{axis}[%
width=\figurewidth,
height=\figureheight,
scale only axis,
xmin=0,
xmax=100,
xlabel={Riesgo (\%)},
ymin=0,
ymax=0.005,
axis x line*=bottom,
axis y line*=left
]
\addplot [color=black,solid,line width=1.2pt,forget plot]
  table[row sep=crcr]{0	0\\
1	0.000329218106995885\\
2	0.00065843621399177\\
3	0.000987654320987654\\
4	0.00131687242798354\\
5	0.00164609053497942\\
6	0.00164609053497942\\
7	0.00164609053497942\\
8	0.00164609053497942\\
9	0.00164609053497942\\
10	0.00164609053497942\\
11	0.00164609053497942\\
12	0.00164609053497942\\
13	0.00164609053497942\\
14	0.00164609053497942\\
15	0.00164609053497942\\
16	0.00164609053497942\\
17	0.00164609053497942\\
18	0.00164609053497942\\
19	0.00164609053497942\\
20	0.00164609053497942\\
21	0.00164609053497942\\
22	0.00164609053497942\\
23	0.00164609053497942\\
24	0.00164609053497942\\
25	0.00164609053497942\\
26	0.00197530864197531\\
27	0.00230452674897119\\
28	0.00263374485596708\\
29	0.00296296296296296\\
30	0.00329218106995885\\
31	0.00362139917695473\\
32	0.00382716049382716\\
33	0.0039917695473251\\
34	0.00415637860082304\\
35	0.00432098765432099\\
36	0.00448559670781893\\
37	0.00465020576131687\\
38	0.00473251028806584\\
39	0.00473251028806584\\
40	0.00473251028806584\\
41	0.00473251028806584\\
42	0.00473251028806584\\
43	0.00473251028806584\\
44	0.00473251028806584\\
45	0.00473251028806584\\
46	0.00440329218106996\\
47	0.00407407407407407\\
48	0.00374485596707819\\
49	0.0034156378600823\\
50	0.00308641975308642\\
51	0.00308641975308642\\
52	0.00308641975308642\\
53	0.00308641975308642\\
54	0.00308641975308642\\
55	0.00308641975308642\\
56	0.00308641975308642\\
57	0.00308641975308642\\
58	0.00308641975308642\\
59	0.00308641975308642\\
60	0.00308641975308642\\
61	0.00308641975308642\\
62	0.00308641975308642\\
63	0.00300411522633745\\
64	0.00283950617283951\\
65	0.00267489711934156\\
66	0.00251028806584362\\
67	0.00234567901234568\\
68	0.00218106995884774\\
69	0.00197530864197531\\
70	0.00164609053497942\\
71	0.00131687242798354\\
72	0.000987654320987654\\
73	0.000658436213991769\\
74	0.000329218106995885\\
75	0\\
76	0\\
77	0\\
78	0\\
79	0\\
80	0\\
81	0\\
82	0\\
83	0\\
84	0\\
85	0\\
86	0\\
87	0\\
88	0\\
89	0\\
90	0\\
91	0\\
92	0\\
93	0\\
94	0\\
95	0\\
96	0\\
97	0\\
98	0\\
99	0\\
100	0\\
};
\addplot [color=blue,line width=1.2pt,mark size=4.5pt,only marks,mark=asterisk,mark options={solid},forget plot]
  table[row sep=crcr]{41	0.00473251028806584\\
};
\addplot [color=green,line width=1.2pt,mark size=4.5pt,only marks,mark=+,mark options={solid},forget plot]
  table[row sep=crcr]{41	0.00473251028806584\\
};
\addplot [color=red,line width=1.2pt,mark size=3.2pt,only marks,mark=square,mark options={solid},forget plot]
  table[row sep=crcr]{42	0.00473251028806584\\
};
\addplot [color=mycolor1,line width=1.2pt,mark size=3.0pt,only marks,mark=triangle,mark options={solid,rotate=180},forget plot]
  table[row sep=crcr]{38	0.00473251028806584\\
};
\addplot [color=mycolor2,line width=1.2pt,mark size=3.0pt,only marks,mark=triangle,mark options={solid},forget plot]
  table[row sep=crcr]{45	0.00473251028806584\\
};
\end{axis}
\end{tikzpicture}%
		\caption{Riesgo medio.}\label{fig:fire-detection-interpolation-medium}
	\end{subfigure}
	\qquad
	\begin{subfigure}[b]{0.45\textwidth}
		\setlength\figureheight{3cm}
		\setlength\figurewidth{6cm}
		% This file was created by matlab2tikz v0.4.7 running on MATLAB 8.0.
% Copyright (c) 2008--2014, Nico Schlömer <nico.schloemer@gmail.com>
% All rights reserved.
% Minimal pgfplots version: 1.3
% 
%
% defining custom colors
\definecolor{mycolor1}{rgb}{0.00000,1.00000,1.00000}%
\definecolor{mycolor2}{rgb}{1.00000,0.00000,1.00000}%
%
\begin{tikzpicture}

\begin{axis}[%
width=\figurewidth,
height=\figureheight,
scale only axis,
xmin=0,
xmax=100,
xlabel={Riesgo (\%)},
ymin=0,
ymax=0.014,
axis x line*=bottom,
axis y line*=left
]
\addplot [color=black,solid,line width=1.2pt,forget plot]
  table[row sep=crcr]{0	0\\
1	0\\
2	0\\
3	0\\
4	0\\
5	0\\
6	0\\
7	0\\
8	0\\
9	0\\
10	0\\
11	0\\
12	0\\
13	0\\
14	0\\
15	0\\
16	0\\
17	0\\
18	0\\
19	0\\
20	0\\
21	0\\
22	0\\
23	0\\
24	0\\
25	0\\
26	0.00065843621399177\\
27	0.00131687242798354\\
28	0.00197530864197531\\
29	0.00263374485596708\\
30	0.00329218106995885\\
31	0.00395061728395062\\
32	0.00411522633744856\\
33	0.00411522633744856\\
34	0.00411522633744856\\
35	0.00411522633744856\\
36	0.00411522633744856\\
37	0.00411522633744856\\
38	0.00411522633744856\\
39	0.00411522633744856\\
40	0.00411522633744856\\
41	0.00411522633744856\\
42	0.00411522633744856\\
43	0.00411522633744856\\
44	0.00411522633744856\\
45	0.00411522633744856\\
46	0.00411522633744856\\
47	0.00411522633744856\\
48	0.00411522633744856\\
49	0.00411522633744856\\
50	0.00411522633744856\\
51	0.00559670781893004\\
52	0.00707818930041153\\
53	0.00855967078189301\\
54	0.0100411522633745\\
55	0.011522633744856\\
56	0.0121810699588477\\
57	0.0124691358024691\\
58	0.0126337448559671\\
59	0.012798353909465\\
60	0.012962962962963\\
61	0.0131275720164609\\
62	0.0132921810699589\\
63	0.0133744855967078\\
64	0.0133744855967078\\
65	0.0133744855967078\\
66	0.0133744855967078\\
67	0.0133744855967078\\
68	0.0133744855967078\\
69	0.0132098765432099\\
70	0.0125514403292181\\
71	0.0118930041152263\\
72	0.0112345679012346\\
73	0.0105761316872428\\
74	0.00991769547325103\\
75	0.00925925925925926\\
76	0.00975308641975309\\
77	0.0102469135802469\\
78	0.0107407407407407\\
79	0.0112345679012346\\
80	0.0117283950617284\\
81	0.0117283950617284\\
82	0.0117283950617284\\
83	0.0117283950617284\\
84	0.0117283950617284\\
85	0.0117283950617284\\
86	0.0117283950617284\\
87	0.0117283950617284\\
88	0.0116460905349794\\
89	0.0114814814814815\\
90	0.0113168724279835\\
91	0.0111522633744856\\
92	0.0109876543209877\\
93	0.0108230452674897\\
94	0.0105349794238683\\
95	0.00987654320987654\\
96	0.00839506172839506\\
97	0.00691358024691358\\
98	0.0054320987654321\\
99	0.00395061728395062\\
100	0.00246913580246914\\
};
\addplot [color=blue,line width=1.2pt,mark size=4.5pt,only marks,mark=asterisk,mark options={solid},forget plot]
  table[row sep=crcr]{69	0.0132098765432099\\
};
\addplot [color=green,line width=1.2pt,mark size=4.5pt,only marks,mark=+,mark options={solid},forget plot]
  table[row sep=crcr]{69	0.0132098765432099\\
};
\addplot [color=red,line width=1.2pt,mark size=3.2pt,only marks,mark=square,mark options={solid},forget plot]
  table[row sep=crcr]{66	0.0133744855967078\\
};
\addplot [color=mycolor1,line width=1.2pt,mark size=3.0pt,only marks,mark=triangle,mark options={solid,rotate=180},forget plot]
  table[row sep=crcr]{63	0.0133744855967078\\
};
\addplot [color=mycolor2,line width=1.2pt,mark size=3.0pt,only marks,mark=triangle,mark options={solid},forget plot]
  table[row sep=crcr]{68	0.0133744855967078\\
};
\end{axis}
\end{tikzpicture}%
		\caption{Riesgo medio-alto.}\label{fig:fire-detection-interpolation-medium-high}
	\end{subfigure}
	
		\vspace{1cm}
		\begin{subfigure}[b]{0.45\textwidth}
			\setlength\figureheight{3cm}
			\setlength\figurewidth{6cm}
			% This file was created by matlab2tikz v0.4.7 running on MATLAB 8.0.
% Copyright (c) 2008--2014, Nico Schlömer <nico.schloemer@gmail.com>
% All rights reserved.
% Minimal pgfplots version: 1.3
% 
%
% defining custom colors
\definecolor{mycolor1}{rgb}{0.00000,1.00000,1.00000}%
\definecolor{mycolor2}{rgb}{1.00000,0.00000,1.00000}%
%
\begin{tikzpicture}

\begin{axis}[%
width=\figurewidth,
height=\figureheight,
scale only axis,
xmin=0,
xmax=100,
xlabel={Riesgo (\%)},
ymin=0,
ymax=0.0045,
axis x line*=bottom,
axis y line*=left
]
\addplot [color=black,solid,line width=1.2pt,forget plot]
  table[row sep=crcr]{0	0\\
1	0\\
2	0\\
3	0\\
4	0\\
5	0\\
6	0\\
7	0\\
8	0\\
9	0\\
10	0\\
11	0\\
12	0\\
13	0\\
14	0\\
15	0\\
16	0\\
17	0\\
18	0\\
19	0\\
20	0\\
21	0\\
22	0\\
23	0\\
24	0\\
25	0\\
26	0\\
27	0\\
28	0\\
29	0\\
30	0\\
31	0\\
32	0\\
33	0\\
34	0\\
35	0\\
36	0\\
37	0\\
38	0\\
39	0\\
40	0\\
41	0\\
42	0\\
43	0\\
44	0\\
45	0\\
46	0\\
47	0\\
48	0\\
49	0\\
50	0\\
51	0.000164609053497942\\
52	0.000329218106995885\\
53	0.000493827160493827\\
54	0.00065843621399177\\
55	0.000823045267489712\\
56	0.000823045267489712\\
57	0.000823045267489712\\
58	0.000823045267489712\\
59	0.000823045267489712\\
60	0.000823045267489712\\
61	0.000823045267489712\\
62	0.000823045267489712\\
63	0.000823045267489712\\
64	0.000823045267489712\\
65	0.000823045267489712\\
66	0.000823045267489712\\
67	0.000823045267489712\\
68	0.000823045267489712\\
69	0.000823045267489712\\
70	0.000823045267489712\\
71	0.000823045267489712\\
72	0.000823045267489712\\
73	0.000823045267489712\\
74	0.000823045267489712\\
75	0.000823045267489712\\
76	0.00131687242798354\\
77	0.00181069958847737\\
78	0.00230452674897119\\
79	0.00279835390946502\\
80	0.00329218106995885\\
81	0.00362139917695473\\
82	0.00395061728395062\\
83	0.0042798353909465\\
84	0.00438957475994513\\
85	0.00438957475994513\\
86	0.00438957475994513\\
87	0.00438957475994513\\
88	0.00438957475994513\\
89	0.00438957475994513\\
90	0.00438957475994513\\
91	0.00438957475994513\\
92	0.00438957475994513\\
93	0.00438957475994513\\
94	0.00438957475994513\\
95	0.00438957475994513\\
96	0.00422496570644719\\
97	0.00406035665294925\\
98	0.0038957475994513\\
99	0.00373113854595336\\
100	0.00356652949245542\\
};
\addplot [color=blue,line width=1.2pt,mark size=4.5pt,only marks,mark=asterisk,mark options={solid},forget plot]
  table[row sep=crcr]{85	0.00438957475994513\\
};
\addplot [color=green,line width=1.2pt,mark size=4.5pt,only marks,mark=+,mark options={solid},forget plot]
  table[row sep=crcr]{87	0.00438957475994513\\
};
\addplot [color=red,line width=1.2pt,mark size=3.2pt,only marks,mark=square,mark options={solid},forget plot]
  table[row sep=crcr]{90	0.00438957475994513\\
};
\addplot [color=mycolor1,line width=1.2pt,mark size=3.0pt,only marks,mark=triangle,mark options={solid,rotate=180},forget plot]
  table[row sep=crcr]{84	0.00438957475994513\\
};
\addplot [color=mycolor2,line width=1.2pt,mark size=3.0pt,only marks,mark=triangle,mark options={solid},forget plot]
  table[row sep=crcr]{95	0.00438957475994513\\
};
\end{axis}
\end{tikzpicture}%
			\caption{Riesgo alto.}\label{fig:fire-detection-interpolation-high}
		\end{subfigure}
		\qquad
		\begin{subfigure}[b]{0.45\textwidth}
			\setlength\figureheight{3cm}
			\setlength\figurewidth{6cm}
			% This file was created by matlab2tikz v0.4.7 running on MATLAB 8.0.
% Copyright (c) 2008--2014, Nico Schlömer <nico.schloemer@gmail.com>
% All rights reserved.
% Minimal pgfplots version: 1.3
% 
%
% defining custom colors
\definecolor{mycolor1}{rgb}{0.00000,1.00000,1.00000}%
\definecolor{mycolor2}{rgb}{1.00000,0.00000,1.00000}%
%
\begin{tikzpicture}

\begin{axis}[%
width=\figurewidth,
height=\figureheight,
scale only axis,
xmin=0,
xmax=100,
xlabel={Riesgo (\%)},
ymin=0,
ymax=0.003,
axis x line*=bottom,
axis y line*=left
]
\addplot [color=black,solid,line width=1.2pt,forget plot]
  table[row sep=crcr]{0	0\\
1	0\\
2	0\\
3	0\\
4	0\\
5	0\\
6	0\\
7	0\\
8	0\\
9	0\\
10	0\\
11	0\\
12	0\\
13	0\\
14	0\\
15	0\\
16	0\\
17	0\\
18	0\\
19	0\\
20	0\\
21	0\\
22	0\\
23	0\\
24	0\\
25	0\\
26	0\\
27	0\\
28	0\\
29	0\\
30	0\\
31	0\\
32	0\\
33	0\\
34	0\\
35	0\\
36	0\\
37	0\\
38	0\\
39	0\\
40	0\\
41	0\\
42	0\\
43	0\\
44	0\\
45	0\\
46	0\\
47	0\\
48	0\\
49	0\\
50	0\\
51	0\\
52	0\\
53	0\\
54	0\\
55	0\\
56	0\\
57	0\\
58	0\\
59	0\\
60	0\\
61	0\\
62	0\\
63	0\\
64	0\\
65	0\\
66	0\\
67	0\\
68	0\\
69	0\\
70	0\\
71	0\\
72	0\\
73	0\\
74	0\\
75	0\\
76	0.000164609053497942\\
77	0.000329218106995885\\
78	0.000493827160493827\\
79	0.00065843621399177\\
80	0.000823045267489712\\
81	0.000987654320987654\\
82	0.0011522633744856\\
83	0.00131687242798354\\
84	0.00148148148148148\\
85	0.00164609053497942\\
86	0.00181069958847737\\
87	0.00197530864197531\\
88	0.00213991769547325\\
89	0.00230452674897119\\
90	0.00246913580246914\\
91	0.00263374485596708\\
92	0.00274348422496571\\
93	0.00274348422496571\\
94	0.00274348422496571\\
95	0.00274348422496571\\
96	0.00274348422496571\\
97	0.00274348422496571\\
98	0.00274348422496571\\
99	0.00274348422496571\\
100	0.00274348422496571\\
};
\addplot [color=blue,line width=1.2pt,mark size=4.5pt,only marks,mark=asterisk,mark options={solid},forget plot]
  table[row sep=crcr]{91	0.00263374485596708\\
};
\addplot [color=green,line width=1.2pt,mark size=4.5pt,only marks,mark=+,mark options={solid},forget plot]
  table[row sep=crcr]{92	0.00274348422496571\\
};
\addplot [color=red,line width=1.2pt,mark size=3.2pt,only marks,mark=square,mark options={solid},forget plot]
  table[row sep=crcr]{96	0.00274348422496571\\
};
\addplot [color=mycolor1,line width=1.2pt,mark size=3.0pt,only marks,mark=triangle,mark options={solid,rotate=180},forget plot]
  table[row sep=crcr]{92	0.00274348422496571\\
};
\addplot [color=mycolor2,line width=1.2pt,mark size=3.0pt,only marks,mark=triangle,mark options={solid},forget plot]
  table[row sep=crcr]{100	0.00274348422496571\\
};
\end{axis}
\end{tikzpicture}%
			\caption{Riesgo muy alto.}\label{fig:fire-detection-interpolation-very-high}
		\end{subfigure}
	\caption{Método de interpolación aplicado a la detección de incendios forestales.}
	\label{fig:interpolation-fire-detection-example}
\end{figure}

En este caso, los resultados son bastante similares a los obtenidos con el método de Mamdani, a pesar de que se obtienen conjuntos de salida con características diferentes. Los valores obtenidos con los desdifusificadores \emph{centroide} y \emph{bisector} se asemejan bastante a los obtenidos con el método de Mamdani, mientras que en los desdifusificadores basados en el máximo(\emph{som}, \emph{mom} y \emph{lom}) las diferencias son más notables.

La ventaja de este método es que permite seleccionar diferentes combinaciones de t-normas e índices de solapamiento y evaluar las salidas obtenidas con cada uno de ellos. En la tabla \ref{tab:fire-detection-interpolation-comparison} y en las figuras \ref{fig:fire-detection-interpolation-comparison-tnorm-prod}-\ref{fig:fire-detection-interpolation-comparison-tnorm-einstein} se presentan los resultados obtenidos al aplicar el método de interpolación utilizando diferentes combinaciones de t-normas e índices de solapamiento a las entradas: $Temperatura = 30 \text{ ºC}$, $Humo = 20 \text{ ppm}$, $Luz = 500 \text{ lux}$, $Humedad = 50 \text{ ppm}$ y $Distancia = 40 \text{ m}$.

\begin{center}
	\def\arraystretch{1.1}
    \begin{longtable}{| c | c | c | c | c | c | c |}
\hline
 \multirow{2}{*}{\textbf{T-norma}} & \multirow{2}{*}{\textbf{Índice de solapamiento}} &  \multicolumn{5}{|c|}{\textbf{Riesgo (\%)}} \\ \cline{3-7}
& & \textbf{cent.  (\textasteriskcentered)} & \textbf{bis. (+)} & \textbf{som ($\triangledown$)} & \textbf{mom ($\square$)} & \textbf{lom ($\vartriangle$)}  \\ \hline
\multirow{5}{*}{$T_{prod}$}&$O_{\pi}$ &29 &29 &26 &38 &49 \\ \cline{2-7}
&$O_{avgmin}$ &29 &29 &26 &38 &49 \\ \cline{2-7}
&$O_{Z}$ &29 &29 &28 &36 &44 \\ \cline{2-7}
&$O_{\sqrt{\text{ }}}$ &30 &30 &26 &38 &49 \\ \cline{2-7}
&$O_{sin}$ &31 &31 &26 &38 &49 \\ \hhline{|=|=|=|=|=|=|=|}
\multirow{5}{*}{$T_{min}$}&$O_{\pi}$ &31 &31 &26 &38 &49 \\ \cline{2-7}
&$O_{avgmin}$ &31 &31 &26 &38 &49 \\ \cline{2-7}
&$O_{Z}$ &30 &30 &32 &36 &40 \\ \cline{2-7}
&$O_{\sqrt{\text{ }}}$ &31 &31 &26 &38 &49 \\ \cline{2-7}
&$O_{sin}$ &31 &31 &26 &38 &49 \\ \hhline{|=|=|=|=|=|=|=|}
\multirow{5}{*}{$T_{geo}$}&$O_{\pi}$ &31 &31 &26 &38 &49 \\ \cline{2-7}
&$O_{avgmin}$ &31 &31 &26 &38 &49 \\ \cline{2-7}
&$O_{Z}$ &31 &30 &32 &32 &32 \\ \cline{2-7}
&$O_{\sqrt{\text{ }}}$ &31 &31 &26 &38 &49 \\ \cline{2-7}
&$O_{sin}$ &31 &31 &26 &38 &49 \\ \hhline{|=|=|=|=|=|=|=|}
\multirow{5}{*}{$T_{harm}$}&$O_{\pi}$ &31 &31 &26 &38 &49 \\ \cline{2-7}
&$O_{avgmin}$ &31 &31 &26 &38 &49 \\ \cline{2-7}
&$O_{Z}$ &31 &30 &33 &33 &33 \\ \cline{2-7}
&$O_{\sqrt{\text{ }}}$ &31 &31 &26 &38 &49 \\ \cline{2-7}
&$O_{sin}$ &31 &31 &26 &38 &49 \\ \hhline{|=|=|=|=|=|=|=|}
\multirow{5}{*}{$T_{sin}$}&$O_{\pi}$ &31 &31 &26 &38 &49 \\ \cline{2-7}
&$O_{avgmin}$ &31 &31 &26 &38 &49 \\ \cline{2-7}
&$O_{Z}$ &31 &31 &28 &37 &46 \\ \cline{2-7}
&$O_{\sqrt{\text{ }}}$ &31 &31 &26 &38 &49 \\ \cline{2-7}
&$O_{sin}$ &31 &31 &26 &38 &49 \\ \hhline{|=|=|=|=|=|=|=|}
\multirow{5}{*}{$T_{einstein}$}&$O_{\pi}$ &29 &29 &26 &38 &49 \\ \cline{2-7}
&$O_{avgmin}$ &29 &29 &26 &38 &49 \\ \cline{2-7}
&$O_{Z}$ &29 &29 &28 &36 &44 \\ \cline{2-7}
&$O_{\sqrt{\text{ }}}$ &30 &30 &26 &38 &49 \\ \cline{2-7}
&$O_{sin}$ &31 &31 &26 &38 &49 \\ \hline
\caption{Resultados de aplicar el método de interpolación con diferentes combinaciones de t-normas e índices de solapamiento.}
\label{tab:fire-detection-interpolation-comparison}
\end{longtable}

\end{center}

\begin{figure}[h!]
\centering
\begin{subfigure}[b]{\textwidth}
	\textbf{T-norma: Producto ($T_{prod}$)} \\
	\setlength\figureheight{2cm}
	\setlength\figurewidth{7cm}
	% This file was created by matlab2tikz v0.4.7 running on MATLAB 8.0.
% Copyright (c) 2008--2014, Nico Schlömer <nico.schloemer@gmail.com>
% All rights reserved.
% Minimal pgfplots version: 1.3
% 
%
% defining custom colors
\definecolor{mycolor1}{rgb}{0.00000,1.00000,1.00000}%
\definecolor{mycolor2}{rgb}{1.00000,0.00000,1.00000}%
%
\begin{tikzpicture}

\begin{axis}[%
width=\figurewidth,
height=\figureheight,
scale only axis,
xmin=0,
xmax=100,
ymin=0,
ymax=0.003,
name=plot3,
title={$O_Z$}
]
\addplot [color=black,solid,line width=1.2pt,forget plot]
  table[row sep=crcr]{0	0\\
1	0.000493827160493827\\
2	0.000987654320987654\\
3	0.00139917695473251\\
4	0.0017283950617284\\
5	0.00205761316872428\\
6	0.00205761316872428\\
7	0.00205761316872428\\
8	0.00205761316872428\\
9	0.00205761316872428\\
10	0.00205761316872428\\
11	0.00205761316872428\\
12	0.00205761316872428\\
13	0.00205761316872428\\
14	0.00205761316872428\\
15	0.00205761316872428\\
16	0.00205761316872428\\
17	0.00205761316872428\\
18	0.00205761316872428\\
19	0.00205761316872428\\
20	0.00205761316872428\\
21	0.00205761316872428\\
22	0.00205761316872428\\
23	0.00205761316872428\\
24	0.00205761316872428\\
25	0.00205761316872428\\
26	0.00222222222222222\\
27	0.00238683127572016\\
28	0.00246913580246914\\
29	0.00246913580246914\\
30	0.00246913580246914\\
31	0.00246913580246914\\
32	0.00246913580246914\\
33	0.00246913580246914\\
34	0.00246913580246914\\
35	0.00246913580246914\\
36	0.00246913580246914\\
37	0.00246913580246914\\
38	0.00246913580246914\\
39	0.00246913580246914\\
40	0.00246913580246914\\
41	0.00246913580246914\\
42	0.00246913580246914\\
43	0.00246913580246914\\
44	0.00246913580246914\\
45	0.00246913580246914\\
46	0.00213991769547325\\
47	0.00181069958847737\\
48	0.00139917695473251\\
49	0.000905349794238684\\
50	0.000411522633744856\\
51	0.000411522633744856\\
52	0.000411522633744856\\
53	0.000411522633744856\\
54	0.000411522633744856\\
55	0.000411522633744856\\
56	0.000411522633744856\\
57	0.000411522633744856\\
58	0.000411522633744856\\
59	0.000411522633744856\\
60	0.000411522633744856\\
61	0.000411522633744856\\
62	0.000411522633744856\\
63	0.000411522633744856\\
64	0.000411522633744856\\
65	0.000411522633744856\\
66	0.000411522633744856\\
67	0.000411522633744856\\
68	0.000411522633744856\\
69	0.000411522633744856\\
70	0.000411522633744856\\
71	0.000411522633744856\\
72	0.000411522633744856\\
73	0.000329218106995885\\
74	0.000164609053497943\\
75	0\\
76	0\\
77	0\\
78	0\\
79	0\\
80	0\\
81	0\\
82	0\\
83	0\\
84	0\\
85	0\\
86	0\\
87	0\\
88	0\\
89	0\\
90	0\\
91	0\\
92	0\\
93	0\\
94	0\\
95	0\\
96	0\\
97	0\\
98	0\\
99	0\\
100	0\\
};
\addplot [color=blue,line width=1.2pt,only marks,mark=asterisk,mark options={solid},forget plot]
  table[row sep=crcr]{29	0.00246913580246914\\
};
\addplot [color=green,line width=1.2pt,only marks,mark=+,mark options={solid},forget plot]
  table[row sep=crcr]{29	0.00246913580246914\\
};
\addplot [color=red,line width=1.2pt,only marks,mark=square,mark options={solid},forget plot]
  table[row sep=crcr]{36	0.00246913580246914\\
};
\addplot [color=mycolor1,line width=1.2pt,only marks,mark=triangle,mark options={solid,rotate=180},forget plot]
  table[row sep=crcr]{28	0.00246913580246914\\
};
\addplot [color=mycolor2,line width=1.2pt,only marks,mark=triangle,mark options={solid},forget plot]
  table[row sep=crcr]{44	0.00246913580246914\\
};
\end{axis}

\begin{axis}[%
width=\figurewidth,
height=\figureheight,
scale only axis,
xmin=0,
xmax=100,
ymin=0,
ymax=3e-14,
name=plot1,
at=(plot3.above north west),
anchor=below south west,
title={$O_{\pi}$}
]
\addplot [color=black,solid,line width=1.2pt,forget plot]
  table[row sep=crcr]{0	0\\
1	2.05596878431034e-14\\
2	2.05596878431034e-14\\
3	2.05596878431034e-14\\
4	2.05596878431034e-14\\
5	2.05596878431034e-14\\
6	2.05596878431034e-14\\
7	2.05596878431034e-14\\
8	2.05596878431034e-14\\
9	2.05596878431034e-14\\
10	2.05596878431034e-14\\
11	2.05596878431034e-14\\
12	2.05596878431034e-14\\
13	2.05596878431034e-14\\
14	2.05596878431034e-14\\
15	2.05596878431034e-14\\
16	2.05596878431034e-14\\
17	2.05596878431034e-14\\
18	2.05596878431034e-14\\
19	2.05596878431034e-14\\
20	2.05596878431034e-14\\
21	2.05596878431034e-14\\
22	2.05596878431034e-14\\
23	2.05596878431034e-14\\
24	2.05596878431034e-14\\
25	2.05596878431034e-14\\
26	2.46716254117241e-14\\
27	2.46716254117241e-14\\
28	2.46716254117241e-14\\
29	2.46716254117241e-14\\
30	2.46716254117241e-14\\
31	2.46716254117241e-14\\
32	2.46716254117241e-14\\
33	2.46716254117241e-14\\
34	2.46716254117241e-14\\
35	2.46716254117241e-14\\
36	2.46716254117241e-14\\
37	2.46716254117241e-14\\
38	2.46716254117241e-14\\
39	2.46716254117241e-14\\
40	2.46716254117241e-14\\
41	2.46716254117241e-14\\
42	2.46716254117241e-14\\
43	2.46716254117241e-14\\
44	2.46716254117241e-14\\
45	2.46716254117241e-14\\
46	2.46716254117241e-14\\
47	2.46716254117241e-14\\
48	2.46716254117241e-14\\
49	2.46716254117241e-14\\
50	4.11193756862068e-15\\
51	4.11193756862068e-15\\
52	4.11193756862068e-15\\
53	4.11193756862068e-15\\
54	4.11193756862068e-15\\
55	4.11193756862068e-15\\
56	4.11193756862068e-15\\
57	4.11193756862068e-15\\
58	4.11193756862068e-15\\
59	4.11193756862068e-15\\
60	4.11193756862068e-15\\
61	4.11193756862068e-15\\
62	4.11193756862068e-15\\
63	4.11193756862068e-15\\
64	4.11193756862068e-15\\
65	4.11193756862068e-15\\
66	4.11193756862068e-15\\
67	4.11193756862068e-15\\
68	4.11193756862068e-15\\
69	4.11193756862068e-15\\
70	4.11193756862068e-15\\
71	4.11193756862068e-15\\
72	4.11193756862068e-15\\
73	4.11193756862068e-15\\
74	4.11193756862068e-15\\
75	0\\
76	0\\
77	0\\
78	0\\
79	0\\
80	0\\
81	0\\
82	0\\
83	0\\
84	0\\
85	0\\
86	0\\
87	0\\
88	0\\
89	0\\
90	0\\
91	0\\
92	0\\
93	0\\
94	0\\
95	0\\
96	0\\
97	0\\
98	0\\
99	0\\
100	0\\
};
\addplot [color=blue,line width=1.2pt,only marks,mark=asterisk,mark options={solid},forget plot]
  table[row sep=crcr]{29	2.46716254117241e-14\\
};
\addplot [color=green,line width=1.2pt,only marks,mark=+,mark options={solid},forget plot]
  table[row sep=crcr]{29	2.46716254117241e-14\\
};
\addplot [color=red,line width=1.2pt,only marks,mark=square,mark options={solid},forget plot]
  table[row sep=crcr]{38	2.46716254117241e-14\\
};
\addplot [color=mycolor1,line width=1.2pt,only marks,mark=triangle,mark options={solid,rotate=180},forget plot]
  table[row sep=crcr]{26	2.46716254117241e-14\\
};
\addplot [color=mycolor2,line width=1.2pt,only marks,mark=triangle,mark options={solid},forget plot]
  table[row sep=crcr]{49	2.46716254117241e-14\\
};
\end{axis}

\begin{axis}[%
width=\figurewidth,
height=\figureheight,
scale only axis,
xmin=0,
xmax=100,
ymin=0,
ymax=3e-14,
name=plot2,
at=(plot1.right of south east),
anchor=left of south west,
title={$O_{avgmin}$}
]
\addplot [color=black,solid,line width=1.2pt,forget plot]
  table[row sep=crcr]{0	0\\
1	2.05596878431034e-14\\
2	2.05596878431034e-14\\
3	2.05596878431034e-14\\
4	2.05596878431034e-14\\
5	2.05596878431034e-14\\
6	2.05596878431034e-14\\
7	2.05596878431034e-14\\
8	2.05596878431034e-14\\
9	2.05596878431034e-14\\
10	2.05596878431034e-14\\
11	2.05596878431034e-14\\
12	2.05596878431034e-14\\
13	2.05596878431034e-14\\
14	2.05596878431034e-14\\
15	2.05596878431034e-14\\
16	2.05596878431034e-14\\
17	2.05596878431034e-14\\
18	2.05596878431034e-14\\
19	2.05596878431034e-14\\
20	2.05596878431034e-14\\
21	2.05596878431034e-14\\
22	2.05596878431034e-14\\
23	2.05596878431034e-14\\
24	2.05596878431034e-14\\
25	2.05596878431034e-14\\
26	2.46716254117241e-14\\
27	2.46716254117241e-14\\
28	2.46716254117241e-14\\
29	2.46716254117241e-14\\
30	2.46716254117241e-14\\
31	2.46716254117241e-14\\
32	2.46716254117241e-14\\
33	2.46716254117241e-14\\
34	2.46716254117241e-14\\
35	2.46716254117241e-14\\
36	2.46716254117241e-14\\
37	2.46716254117241e-14\\
38	2.46716254117241e-14\\
39	2.46716254117241e-14\\
40	2.46716254117241e-14\\
41	2.46716254117241e-14\\
42	2.46716254117241e-14\\
43	2.46716254117241e-14\\
44	2.46716254117241e-14\\
45	2.46716254117241e-14\\
46	2.46716254117241e-14\\
47	2.46716254117241e-14\\
48	2.46716254117241e-14\\
49	2.46716254117241e-14\\
50	4.11193756862068e-15\\
51	4.11193756862068e-15\\
52	4.11193756862068e-15\\
53	4.11193756862068e-15\\
54	4.11193756862068e-15\\
55	4.11193756862068e-15\\
56	4.11193756862068e-15\\
57	4.11193756862068e-15\\
58	4.11193756862068e-15\\
59	4.11193756862068e-15\\
60	4.11193756862068e-15\\
61	4.11193756862068e-15\\
62	4.11193756862068e-15\\
63	4.11193756862068e-15\\
64	4.11193756862068e-15\\
65	4.11193756862068e-15\\
66	4.11193756862068e-15\\
67	4.11193756862068e-15\\
68	4.11193756862068e-15\\
69	4.11193756862068e-15\\
70	4.11193756862068e-15\\
71	4.11193756862068e-15\\
72	4.11193756862068e-15\\
73	4.11193756862068e-15\\
74	4.11193756862068e-15\\
75	0\\
76	0\\
77	0\\
78	0\\
79	0\\
80	0\\
81	0\\
82	0\\
83	0\\
84	0\\
85	0\\
86	0\\
87	0\\
88	0\\
89	0\\
90	0\\
91	0\\
92	0\\
93	0\\
94	0\\
95	0\\
96	0\\
97	0\\
98	0\\
99	0\\
100	0\\
};
\addplot [color=blue,line width=1.2pt,only marks,mark=asterisk,mark options={solid},forget plot]
  table[row sep=crcr]{29	2.46716254117241e-14\\
};
\addplot [color=green,line width=1.2pt,only marks,mark=+,mark options={solid},forget plot]
  table[row sep=crcr]{29	2.46716254117241e-14\\
};
\addplot [color=red,line width=1.2pt,only marks,mark=square,mark options={solid},forget plot]
  table[row sep=crcr]{38	2.46716254117241e-14\\
};
\addplot [color=mycolor1,line width=1.2pt,only marks,mark=triangle,mark options={solid,rotate=180},forget plot]
  table[row sep=crcr]{26	2.46716254117241e-14\\
};
\addplot [color=mycolor2,line width=1.2pt,only marks,mark=triangle,mark options={solid},forget plot]
  table[row sep=crcr]{49	2.46716254117241e-14\\
};
\end{axis}

\begin{axis}[%
width=\figurewidth,
height=\figureheight,
scale only axis,
xmin=0,
xmax=100,
ymin=0,
ymax=1e-13,
at=(plot2.below south west),
anchor=above north west,
title={$O_{\sqrt{\text{ }}}$}
]
\addplot [color=black,solid,line width=1.2pt,forget plot]
  table[row sep=crcr]{0	0\\
1	4.97813760039405e-14\\
2	4.97813760039405e-14\\
3	4.97813760039405e-14\\
4	4.97813760039405e-14\\
5	4.97813760039405e-14\\
6	4.97813760039405e-14\\
7	4.97813760039405e-14\\
8	4.97813760039405e-14\\
9	4.97813760039405e-14\\
10	4.97813760039405e-14\\
11	4.97813760039405e-14\\
12	4.97813760039405e-14\\
13	4.97813760039405e-14\\
14	4.97813760039405e-14\\
15	4.97813760039405e-14\\
16	4.97813760039405e-14\\
17	4.97813760039405e-14\\
18	4.97813760039405e-14\\
19	4.97813760039405e-14\\
20	4.97813760039405e-14\\
21	4.97813760039405e-14\\
22	4.97813760039405e-14\\
23	4.97813760039405e-14\\
24	4.97813760039405e-14\\
25	4.97813760039405e-14\\
26	6.27844643171967e-14\\
27	6.27844643171967e-14\\
28	6.27844643171967e-14\\
29	6.27844643171967e-14\\
30	6.27844643171967e-14\\
31	6.27844643171967e-14\\
32	6.27844643171967e-14\\
33	6.27844643171967e-14\\
34	6.27844643171967e-14\\
35	6.27844643171967e-14\\
36	6.27844643171967e-14\\
37	6.27844643171967e-14\\
38	6.27844643171967e-14\\
39	6.27844643171967e-14\\
40	6.27844643171967e-14\\
41	6.27844643171967e-14\\
42	6.27844643171967e-14\\
43	6.27844643171967e-14\\
44	6.27844643171967e-14\\
45	6.27844643171967e-14\\
46	6.27844643171967e-14\\
47	6.27844643171967e-14\\
48	6.27844643171967e-14\\
49	6.27844643171967e-14\\
50	1.30030883132563e-14\\
51	1.30030883132563e-14\\
52	1.30030883132563e-14\\
53	1.30030883132563e-14\\
54	1.30030883132563e-14\\
55	1.30030883132563e-14\\
56	1.30030883132563e-14\\
57	1.30030883132563e-14\\
58	1.30030883132563e-14\\
59	1.30030883132563e-14\\
60	1.30030883132563e-14\\
61	1.30030883132563e-14\\
62	1.30030883132563e-14\\
63	1.30030883132563e-14\\
64	1.30030883132563e-14\\
65	1.30030883132563e-14\\
66	1.30030883132563e-14\\
67	1.30030883132563e-14\\
68	1.30030883132563e-14\\
69	1.30030883132563e-14\\
70	1.30030883132563e-14\\
71	1.30030883132563e-14\\
72	1.30030883132563e-14\\
73	1.30030883132563e-14\\
74	1.30030883132563e-14\\
75	0\\
76	0\\
77	0\\
78	0\\
79	0\\
80	0\\
81	0\\
82	0\\
83	0\\
84	0\\
85	0\\
86	0\\
87	0\\
88	0\\
89	0\\
90	0\\
91	0\\
92	0\\
93	0\\
94	0\\
95	0\\
96	0\\
97	0\\
98	0\\
99	0\\
100	0\\
};
\addplot [color=blue,line width=1.2pt,only marks,mark=asterisk,mark options={solid},forget plot]
  table[row sep=crcr]{30	6.27844643171967e-14\\
};
\addplot [color=green,line width=1.2pt,only marks,mark=+,mark options={solid},forget plot]
  table[row sep=crcr]{30	6.27844643171967e-14\\
};
\addplot [color=red,line width=1.2pt,only marks,mark=square,mark options={solid},forget plot]
  table[row sep=crcr]{38	6.27844643171967e-14\\
};
\addplot [color=mycolor1,line width=1.2pt,only marks,mark=triangle,mark options={solid,rotate=180},forget plot]
  table[row sep=crcr]{26	6.27844643171967e-14\\
};
\addplot [color=mycolor2,line width=1.2pt,only marks,mark=triangle,mark options={solid},forget plot]
  table[row sep=crcr]{49	6.27844643171967e-14\\
};
\end{axis}

\begin{axis}[%
width=\figurewidth,
height=\figureheight,
scale only axis,
xmin=0,
xmax=100,
ymin=0,
ymax=2e-13,
at=(plot3.below south west),
anchor=above north west,
title={$O_{sin}$}
]
\addplot [color=black,solid,line width=1.2pt,forget plot]
  table[row sep=crcr]{0	0\\
1	1.10553708723286e-13\\
2	1.10553708723286e-13\\
3	1.10553708723286e-13\\
4	1.10553708723286e-13\\
5	1.10553708723286e-13\\
6	1.10553708723286e-13\\
7	1.10553708723286e-13\\
8	1.10553708723286e-13\\
9	1.10553708723286e-13\\
10	1.10553708723286e-13\\
11	1.10553708723286e-13\\
12	1.10553708723286e-13\\
13	1.10553708723286e-13\\
14	1.10553708723286e-13\\
15	1.10553708723286e-13\\
16	1.10553708723286e-13\\
17	1.10553708723286e-13\\
18	1.10553708723286e-13\\
19	1.10553708723286e-13\\
20	1.10553708723286e-13\\
21	1.10553708723286e-13\\
22	1.10553708723286e-13\\
23	1.10553708723286e-13\\
24	1.10553708723286e-13\\
25	1.10553708723286e-13\\
26	1.45508642396137e-13\\
27	1.45508642396137e-13\\
28	1.45508642396137e-13\\
29	1.45508642396137e-13\\
30	1.45508642396137e-13\\
31	1.45508642396137e-13\\
32	1.45508642396137e-13\\
33	1.45508642396137e-13\\
34	1.45508642396137e-13\\
35	1.45508642396137e-13\\
36	1.45508642396137e-13\\
37	1.45508642396137e-13\\
38	1.45508642396137e-13\\
39	1.45508642396137e-13\\
40	1.45508642396137e-13\\
41	1.45508642396137e-13\\
42	1.45508642396137e-13\\
43	1.45508642396137e-13\\
44	1.45508642396137e-13\\
45	1.45508642396137e-13\\
46	1.45508642396137e-13\\
47	1.45508642396137e-13\\
48	1.45508642396137e-13\\
49	1.45508642396137e-13\\
50	3.4954933672851e-14\\
51	3.4954933672851e-14\\
52	3.4954933672851e-14\\
53	3.4954933672851e-14\\
54	3.4954933672851e-14\\
55	3.4954933672851e-14\\
56	3.4954933672851e-14\\
57	3.4954933672851e-14\\
58	3.4954933672851e-14\\
59	3.4954933672851e-14\\
60	3.4954933672851e-14\\
61	3.4954933672851e-14\\
62	3.4954933672851e-14\\
63	3.4954933672851e-14\\
64	3.4954933672851e-14\\
65	3.4954933672851e-14\\
66	3.4954933672851e-14\\
67	3.4954933672851e-14\\
68	3.4954933672851e-14\\
69	3.4954933672851e-14\\
70	3.4954933672851e-14\\
71	3.4954933672851e-14\\
72	3.4954933672851e-14\\
73	3.4954933672851e-14\\
74	3.4954933672851e-14\\
75	0\\
76	0\\
77	0\\
78	0\\
79	0\\
80	0\\
81	0\\
82	0\\
83	0\\
84	0\\
85	0\\
86	0\\
87	0\\
88	0\\
89	0\\
90	0\\
91	0\\
92	0\\
93	0\\
94	0\\
95	0\\
96	0\\
97	0\\
98	0\\
99	0\\
100	0\\
};
\addplot [color=blue,line width=1.2pt,only marks,mark=asterisk,mark options={solid},forget plot]
  table[row sep=crcr]{31	1.45508642396137e-13\\
};
\addplot [color=green,line width=1.2pt,only marks,mark=+,mark options={solid},forget plot]
  table[row sep=crcr]{31	1.45508642396137e-13\\
};
\addplot [color=red,line width=1.2pt,only marks,mark=square,mark options={solid},forget plot]
  table[row sep=crcr]{38	1.45508642396137e-13\\
};
\addplot [color=mycolor1,line width=1.2pt,only marks,mark=triangle,mark options={solid,rotate=180},forget plot]
  table[row sep=crcr]{26	1.45508642396137e-13\\
};
\addplot [color=mycolor2,line width=1.2pt,only marks,mark=triangle,mark options={solid},forget plot]
  table[row sep=crcr]{49	1.45508642396137e-13\\
};
\end{axis}
\end{tikzpicture}%
\end{subfigure}
\label{fig:fire-detection-interpolation-comparison-tnorm-prod}
\caption{T-norma: Producto ($T_{prod}$)}
\end{figure}

\begin{figure}[h!]
\centering
\begin{subfigure}[b]{\textwidth}
	\textbf{T-norma: Mínimo ($T_{min}$)} \\
	\setlength\figureheight{2cm}
	\setlength\figurewidth{7cm}
	% This file was created by matlab2tikz v0.4.7 running on MATLAB 8.0.
% Copyright (c) 2008--2014, Nico Schlömer <nico.schloemer@gmail.com>
% All rights reserved.
% Minimal pgfplots version: 1.3
% 
%
% defining custom colors
\definecolor{mycolor1}{rgb}{0.00000,1.00000,1.00000}%
\definecolor{mycolor2}{rgb}{1.00000,0.00000,1.00000}%
%
\begin{tikzpicture}

\begin{axis}[%
width=\figurewidth,
height=\figureheight,
scale only axis,
xmin=0,
xmax=100,
ymin=0,
ymax=0.006,
name=plot3,
title={Omaxmin}
]
\addplot [color=black,solid,line width=1.2pt,forget plot]
  table[row sep=crcr]{0	0\\
1	0.000493827160493827\\
2	0.000987654320987654\\
3	0.00148148148148148\\
4	0.00197530864197531\\
5	0.00246913580246914\\
6	0.00296296296296296\\
7	0.00333333333333333\\
8	0.00366255144032922\\
9	0.0039917695473251\\
10	0.00432098765432099\\
11	0.00432098765432099\\
12	0.00432098765432099\\
13	0.00432098765432099\\
14	0.00432098765432099\\
15	0.00432098765432099\\
16	0.00432098765432099\\
17	0.00432098765432099\\
18	0.00432098765432099\\
19	0.00432098765432099\\
20	0.00432098765432099\\
21	0.00432098765432099\\
22	0.00432098765432099\\
23	0.00432098765432099\\
24	0.00432098765432099\\
25	0.00432098765432099\\
26	0.00448559670781893\\
27	0.00465020576131687\\
28	0.00481481481481481\\
29	0.00497942386831276\\
30	0.0051440329218107\\
31	0.00530864197530864\\
32	0.00534979423868313\\
33	0.00534979423868313\\
34	0.00534979423868313\\
35	0.00534979423868313\\
36	0.00534979423868313\\
37	0.00534979423868313\\
38	0.00534979423868313\\
39	0.00534979423868313\\
40	0.00534979423868313\\
41	0.00502057613168724\\
42	0.00469135802469136\\
43	0.00436213991769547\\
44	0.0039917695473251\\
45	0.00349794238683128\\
46	0.00300411522633745\\
47	0.00251028806584362\\
48	0.00201646090534979\\
49	0.00152263374485597\\
50	0.00102880658436214\\
51	0.00102880658436214\\
52	0.00102880658436214\\
53	0.00102880658436214\\
54	0.00102880658436214\\
55	0.00102880658436214\\
56	0.00102880658436214\\
57	0.00102880658436214\\
58	0.00102880658436214\\
59	0.00102880658436214\\
60	0.00102880658436214\\
61	0.00102880658436214\\
62	0.00102880658436214\\
63	0.00102880658436214\\
64	0.00102880658436214\\
65	0.00102880658436214\\
66	0.00102880658436214\\
67	0.00102880658436214\\
68	0.00102880658436214\\
69	0.000987654320987654\\
70	0.000823045267489712\\
71	0.00065843621399177\\
72	0.000493827160493827\\
73	0.000329218106995885\\
74	0.000164609053497943\\
75	0\\
76	0\\
77	0\\
78	0\\
79	0\\
80	0\\
81	0\\
82	0\\
83	0\\
84	0\\
85	0\\
86	0\\
87	0\\
88	0\\
89	0\\
90	0\\
91	0\\
92	0\\
93	0\\
94	0\\
95	0\\
96	0\\
97	0\\
98	0\\
99	0\\
100	0\\
};
\addplot [color=blue,line width=1.2pt,only marks,mark=asterisk,mark options={solid},forget plot]
  table[row sep=crcr]{30	0.0051440329218107\\
};
\addplot [color=green,line width=1.2pt,only marks,mark=+,mark options={solid},forget plot]
  table[row sep=crcr]{30	0.0051440329218107\\
};
\addplot [color=red,line width=1.2pt,only marks,mark=square,mark options={solid},forget plot]
  table[row sep=crcr]{36	0.00534979423868313\\
};
\addplot [color=mycolor1,line width=1.2pt,only marks,mark=triangle,mark options={solid,rotate=180},forget plot]
  table[row sep=crcr]{32	0.00534979423868313\\
};
\addplot [color=mycolor2,line width=1.2pt,only marks,mark=triangle,mark options={solid},forget plot]
  table[row sep=crcr]{40	0.00534979423868313\\
};
\end{axis}

\begin{axis}[%
width=\figurewidth,
height=\figureheight,
scale only axis,
xmin=0,
xmax=100,
ymin=0,
ymax=2e-05,
name=plot1,
at=(plot3.above north west),
anchor=below south west,
title={Opi}
]
\addplot [color=black,solid,line width=1.2pt,forget plot]
  table[row sep=crcr]{0	0\\
1	1.2333345666679e-05\\
2	1.2333345666679e-05\\
3	1.2333345666679e-05\\
4	1.2333345666679e-05\\
5	1.2333345666679e-05\\
6	1.2333345666679e-05\\
7	1.2333345666679e-05\\
8	1.2333345666679e-05\\
9	1.2333345666679e-05\\
10	1.2333345666679e-05\\
11	1.2333345666679e-05\\
12	1.2333345666679e-05\\
13	1.2333345666679e-05\\
14	1.2333345666679e-05\\
15	1.2333345666679e-05\\
16	1.2333345666679e-05\\
17	1.2333345666679e-05\\
18	1.2333345666679e-05\\
19	1.2333345666679e-05\\
20	1.2333345666679e-05\\
21	1.2333345666679e-05\\
22	1.2333345666679e-05\\
23	1.2333345666679e-05\\
24	1.2333345666679e-05\\
25	1.2333345666679e-05\\
26	1.64444608889053e-05\\
27	1.64444608889053e-05\\
28	1.64444608889053e-05\\
29	1.64444608889053e-05\\
30	1.64444608889053e-05\\
31	1.64444608889053e-05\\
32	1.64444608889053e-05\\
33	1.64444608889053e-05\\
34	1.64444608889053e-05\\
35	1.64444608889053e-05\\
36	1.64444608889053e-05\\
37	1.64444608889053e-05\\
38	1.64444608889053e-05\\
39	1.64444608889053e-05\\
40	1.64444608889053e-05\\
41	1.64444608889053e-05\\
42	1.64444608889053e-05\\
43	1.64444608889053e-05\\
44	1.64444608889053e-05\\
45	1.64444608889053e-05\\
46	1.64444608889053e-05\\
47	1.64444608889053e-05\\
48	1.64444608889053e-05\\
49	1.64444608889053e-05\\
50	4.11111522222633e-06\\
51	4.11111522222633e-06\\
52	4.11111522222633e-06\\
53	4.11111522222633e-06\\
54	4.11111522222633e-06\\
55	4.11111522222633e-06\\
56	4.11111522222633e-06\\
57	4.11111522222633e-06\\
58	4.11111522222633e-06\\
59	4.11111522222633e-06\\
60	4.11111522222633e-06\\
61	4.11111522222633e-06\\
62	4.11111522222633e-06\\
63	4.11111522222633e-06\\
64	4.11111522222633e-06\\
65	4.11111522222633e-06\\
66	4.11111522222633e-06\\
67	4.11111522222633e-06\\
68	4.11111522222633e-06\\
69	4.11111522222633e-06\\
70	4.11111522222633e-06\\
71	4.11111522222633e-06\\
72	4.11111522222633e-06\\
73	4.11111522222633e-06\\
74	4.11111522222633e-06\\
75	0\\
76	0\\
77	0\\
78	0\\
79	0\\
80	0\\
81	0\\
82	0\\
83	0\\
84	0\\
85	0\\
86	0\\
87	0\\
88	0\\
89	0\\
90	0\\
91	0\\
92	0\\
93	0\\
94	0\\
95	0\\
96	0\\
97	0\\
98	0\\
99	0\\
100	0\\
};
\addplot [color=blue,line width=1.2pt,only marks,mark=asterisk,mark options={solid},forget plot]
  table[row sep=crcr]{31	1.64444608889053e-05\\
};
\addplot [color=green,line width=1.2pt,only marks,mark=+,mark options={solid},forget plot]
  table[row sep=crcr]{31	1.64444608889053e-05\\
};
\addplot [color=red,line width=1.2pt,only marks,mark=square,mark options={solid},forget plot]
  table[row sep=crcr]{38	1.64444608889053e-05\\
};
\addplot [color=mycolor1,line width=1.2pt,only marks,mark=triangle,mark options={solid,rotate=180},forget plot]
  table[row sep=crcr]{26	1.64444608889053e-05\\
};
\addplot [color=mycolor2,line width=1.2pt,only marks,mark=triangle,mark options={solid},forget plot]
  table[row sep=crcr]{49	1.64444608889053e-05\\
};
\end{axis}

\begin{axis}[%
width=\figurewidth,
height=\figureheight,
scale only axis,
xmin=0,
xmax=100,
ymin=0,
ymax=2e-05,
name=plot2,
at=(plot1.right of south east),
anchor=left of south west,
title={Oavgmin}
]
\addplot [color=black,solid,line width=1.2pt,forget plot]
  table[row sep=crcr]{0	0\\
1	1.2333345666679e-05\\
2	1.2333345666679e-05\\
3	1.2333345666679e-05\\
4	1.2333345666679e-05\\
5	1.2333345666679e-05\\
6	1.2333345666679e-05\\
7	1.2333345666679e-05\\
8	1.2333345666679e-05\\
9	1.2333345666679e-05\\
10	1.2333345666679e-05\\
11	1.2333345666679e-05\\
12	1.2333345666679e-05\\
13	1.2333345666679e-05\\
14	1.2333345666679e-05\\
15	1.2333345666679e-05\\
16	1.2333345666679e-05\\
17	1.2333345666679e-05\\
18	1.2333345666679e-05\\
19	1.2333345666679e-05\\
20	1.2333345666679e-05\\
21	1.2333345666679e-05\\
22	1.2333345666679e-05\\
23	1.2333345666679e-05\\
24	1.2333345666679e-05\\
25	1.2333345666679e-05\\
26	1.64444608889053e-05\\
27	1.64444608889053e-05\\
28	1.64444608889053e-05\\
29	1.64444608889053e-05\\
30	1.64444608889053e-05\\
31	1.64444608889053e-05\\
32	1.64444608889053e-05\\
33	1.64444608889053e-05\\
34	1.64444608889053e-05\\
35	1.64444608889053e-05\\
36	1.64444608889053e-05\\
37	1.64444608889053e-05\\
38	1.64444608889053e-05\\
39	1.64444608889053e-05\\
40	1.64444608889053e-05\\
41	1.64444608889053e-05\\
42	1.64444608889053e-05\\
43	1.64444608889053e-05\\
44	1.64444608889053e-05\\
45	1.64444608889053e-05\\
46	1.64444608889053e-05\\
47	1.64444608889053e-05\\
48	1.64444608889053e-05\\
49	1.64444608889053e-05\\
50	4.11111522222633e-06\\
51	4.11111522222633e-06\\
52	4.11111522222633e-06\\
53	4.11111522222633e-06\\
54	4.11111522222633e-06\\
55	4.11111522222633e-06\\
56	4.11111522222633e-06\\
57	4.11111522222633e-06\\
58	4.11111522222633e-06\\
59	4.11111522222633e-06\\
60	4.11111522222633e-06\\
61	4.11111522222633e-06\\
62	4.11111522222633e-06\\
63	4.11111522222633e-06\\
64	4.11111522222633e-06\\
65	4.11111522222633e-06\\
66	4.11111522222633e-06\\
67	4.11111522222633e-06\\
68	4.11111522222633e-06\\
69	4.11111522222633e-06\\
70	4.11111522222633e-06\\
71	4.11111522222633e-06\\
72	4.11111522222633e-06\\
73	4.11111522222633e-06\\
74	4.11111522222633e-06\\
75	0\\
76	0\\
77	0\\
78	0\\
79	0\\
80	0\\
81	0\\
82	0\\
83	0\\
84	0\\
85	0\\
86	0\\
87	0\\
88	0\\
89	0\\
90	0\\
91	0\\
92	0\\
93	0\\
94	0\\
95	0\\
96	0\\
97	0\\
98	0\\
99	0\\
100	0\\
};
\addplot [color=blue,line width=1.2pt,only marks,mark=asterisk,mark options={solid},forget plot]
  table[row sep=crcr]{31	1.64444608889053e-05\\
};
\addplot [color=green,line width=1.2pt,only marks,mark=+,mark options={solid},forget plot]
  table[row sep=crcr]{31	1.64444608889053e-05\\
};
\addplot [color=red,line width=1.2pt,only marks,mark=square,mark options={solid},forget plot]
  table[row sep=crcr]{38	1.64444608889053e-05\\
};
\addplot [color=mycolor1,line width=1.2pt,only marks,mark=triangle,mark options={solid,rotate=180},forget plot]
  table[row sep=crcr]{26	1.64444608889053e-05\\
};
\addplot [color=mycolor2,line width=1.2pt,only marks,mark=triangle,mark options={solid},forget plot]
  table[row sep=crcr]{49	1.64444608889053e-05\\
};
\end{axis}

\begin{axis}[%
width=\figurewidth,
height=\figureheight,
scale only axis,
xmin=0,
xmax=100,
ymin=0,
ymax=2e-05,
at=(plot2.below south west),
anchor=above north west,
title={Osqrt}
]
\addplot [color=black,solid,line width=1.2pt,forget plot]
  table[row sep=crcr]{0	0\\
1	1.2333345666679e-05\\
2	1.2333345666679e-05\\
3	1.2333345666679e-05\\
4	1.2333345666679e-05\\
5	1.2333345666679e-05\\
6	1.2333345666679e-05\\
7	1.2333345666679e-05\\
8	1.2333345666679e-05\\
9	1.2333345666679e-05\\
10	1.2333345666679e-05\\
11	1.2333345666679e-05\\
12	1.2333345666679e-05\\
13	1.2333345666679e-05\\
14	1.2333345666679e-05\\
15	1.2333345666679e-05\\
16	1.2333345666679e-05\\
17	1.2333345666679e-05\\
18	1.2333345666679e-05\\
19	1.2333345666679e-05\\
20	1.2333345666679e-05\\
21	1.2333345666679e-05\\
22	1.2333345666679e-05\\
23	1.2333345666679e-05\\
24	1.2333345666679e-05\\
25	1.2333345666679e-05\\
26	1.64444608889053e-05\\
27	1.64444608889053e-05\\
28	1.64444608889053e-05\\
29	1.64444608889053e-05\\
30	1.64444608889053e-05\\
31	1.64444608889053e-05\\
32	1.64444608889053e-05\\
33	1.64444608889053e-05\\
34	1.64444608889053e-05\\
35	1.64444608889053e-05\\
36	1.64444608889053e-05\\
37	1.64444608889053e-05\\
38	1.64444608889053e-05\\
39	1.64444608889053e-05\\
40	1.64444608889053e-05\\
41	1.64444608889053e-05\\
42	1.64444608889053e-05\\
43	1.64444608889053e-05\\
44	1.64444608889053e-05\\
45	1.64444608889053e-05\\
46	1.64444608889053e-05\\
47	1.64444608889053e-05\\
48	1.64444608889053e-05\\
49	1.64444608889053e-05\\
50	4.11111522222633e-06\\
51	4.11111522222633e-06\\
52	4.11111522222633e-06\\
53	4.11111522222633e-06\\
54	4.11111522222633e-06\\
55	4.11111522222633e-06\\
56	4.11111522222633e-06\\
57	4.11111522222633e-06\\
58	4.11111522222633e-06\\
59	4.11111522222633e-06\\
60	4.11111522222633e-06\\
61	4.11111522222633e-06\\
62	4.11111522222633e-06\\
63	4.11111522222633e-06\\
64	4.11111522222633e-06\\
65	4.11111522222633e-06\\
66	4.11111522222633e-06\\
67	4.11111522222633e-06\\
68	4.11111522222633e-06\\
69	4.11111522222633e-06\\
70	4.11111522222633e-06\\
71	4.11111522222633e-06\\
72	4.11111522222633e-06\\
73	4.11111522222633e-06\\
74	4.11111522222633e-06\\
75	0\\
76	0\\
77	0\\
78	0\\
79	0\\
80	0\\
81	0\\
82	0\\
83	0\\
84	0\\
85	0\\
86	0\\
87	0\\
88	0\\
89	0\\
90	0\\
91	0\\
92	0\\
93	0\\
94	0\\
95	0\\
96	0\\
97	0\\
98	0\\
99	0\\
100	0\\
};
\addplot [color=blue,line width=1.2pt,only marks,mark=asterisk,mark options={solid},forget plot]
  table[row sep=crcr]{31	1.64444608889053e-05\\
};
\addplot [color=green,line width=1.2pt,only marks,mark=+,mark options={solid},forget plot]
  table[row sep=crcr]{31	1.64444608889053e-05\\
};
\addplot [color=red,line width=1.2pt,only marks,mark=square,mark options={solid},forget plot]
  table[row sep=crcr]{38	1.64444608889053e-05\\
};
\addplot [color=mycolor1,line width=1.2pt,only marks,mark=triangle,mark options={solid,rotate=180},forget plot]
  table[row sep=crcr]{26	1.64444608889053e-05\\
};
\addplot [color=mycolor2,line width=1.2pt,only marks,mark=triangle,mark options={solid},forget plot]
  table[row sep=crcr]{49	1.64444608889053e-05\\
};
\end{axis}

\begin{axis}[%
width=\figurewidth,
height=\figureheight,
scale only axis,
xmin=0,
xmax=100,
ymin=0,
ymax=2e-05,
at=(plot3.below south west),
anchor=above north west,
title={Osin}
]
\addplot [color=black,solid,line width=1.2pt,forget plot]
  table[row sep=crcr]{0	0\\
1	1.2333345666679e-05\\
2	1.2333345666679e-05\\
3	1.2333345666679e-05\\
4	1.2333345666679e-05\\
5	1.2333345666679e-05\\
6	1.2333345666679e-05\\
7	1.2333345666679e-05\\
8	1.2333345666679e-05\\
9	1.2333345666679e-05\\
10	1.2333345666679e-05\\
11	1.2333345666679e-05\\
12	1.2333345666679e-05\\
13	1.2333345666679e-05\\
14	1.2333345666679e-05\\
15	1.2333345666679e-05\\
16	1.2333345666679e-05\\
17	1.2333345666679e-05\\
18	1.2333345666679e-05\\
19	1.2333345666679e-05\\
20	1.2333345666679e-05\\
21	1.2333345666679e-05\\
22	1.2333345666679e-05\\
23	1.2333345666679e-05\\
24	1.2333345666679e-05\\
25	1.2333345666679e-05\\
26	1.64444608889053e-05\\
27	1.64444608889053e-05\\
28	1.64444608889053e-05\\
29	1.64444608889053e-05\\
30	1.64444608889053e-05\\
31	1.64444608889053e-05\\
32	1.64444608889053e-05\\
33	1.64444608889053e-05\\
34	1.64444608889053e-05\\
35	1.64444608889053e-05\\
36	1.64444608889053e-05\\
37	1.64444608889053e-05\\
38	1.64444608889053e-05\\
39	1.64444608889053e-05\\
40	1.64444608889053e-05\\
41	1.64444608889053e-05\\
42	1.64444608889053e-05\\
43	1.64444608889053e-05\\
44	1.64444608889053e-05\\
45	1.64444608889053e-05\\
46	1.64444608889053e-05\\
47	1.64444608889053e-05\\
48	1.64444608889053e-05\\
49	1.64444608889053e-05\\
50	4.11111522222633e-06\\
51	4.11111522222633e-06\\
52	4.11111522222633e-06\\
53	4.11111522222633e-06\\
54	4.11111522222633e-06\\
55	4.11111522222633e-06\\
56	4.11111522222633e-06\\
57	4.11111522222633e-06\\
58	4.11111522222633e-06\\
59	4.11111522222633e-06\\
60	4.11111522222633e-06\\
61	4.11111522222633e-06\\
62	4.11111522222633e-06\\
63	4.11111522222633e-06\\
64	4.11111522222633e-06\\
65	4.11111522222633e-06\\
66	4.11111522222633e-06\\
67	4.11111522222633e-06\\
68	4.11111522222633e-06\\
69	4.11111522222633e-06\\
70	4.11111522222633e-06\\
71	4.11111522222633e-06\\
72	4.11111522222633e-06\\
73	4.11111522222633e-06\\
74	4.11111522222633e-06\\
75	0\\
76	0\\
77	0\\
78	0\\
79	0\\
80	0\\
81	0\\
82	0\\
83	0\\
84	0\\
85	0\\
86	0\\
87	0\\
88	0\\
89	0\\
90	0\\
91	0\\
92	0\\
93	0\\
94	0\\
95	0\\
96	0\\
97	0\\
98	0\\
99	0\\
100	0\\
};
\addplot [color=blue,line width=1.2pt,only marks,mark=asterisk,mark options={solid},forget plot]
  table[row sep=crcr]{31	1.64444608889053e-05\\
};
\addplot [color=green,line width=1.2pt,only marks,mark=+,mark options={solid},forget plot]
  table[row sep=crcr]{31	1.64444608889053e-05\\
};
\addplot [color=red,line width=1.2pt,only marks,mark=square,mark options={solid},forget plot]
  table[row sep=crcr]{38	1.64444608889053e-05\\
};
\addplot [color=mycolor1,line width=1.2pt,only marks,mark=triangle,mark options={solid,rotate=180},forget plot]
  table[row sep=crcr]{26	1.64444608889053e-05\\
};
\addplot [color=mycolor2,line width=1.2pt,only marks,mark=triangle,mark options={solid},forget plot]
  table[row sep=crcr]{49	1.64444608889053e-05\\
};
\end{axis}
\end{tikzpicture}%
\end{subfigure}
\label{fig:fire-detection-interpolation-comparison-tnorm-min}
\caption{T-norma: Mínimo ($T_{min}$)}
\end{figure}
\begin{figure}[h!]
\centering
\begin{subfigure}[b]{\textwidth}
	\textbf{T-norma: Media geométrica ($T_{geo}$)} \\
	\setlength\figureheight{2cm}
	\setlength\figurewidth{7cm}
	% This file was created by matlab2tikz v0.4.7 running on MATLAB 8.0.
% Copyright (c) 2008--2014, Nico Schlömer <nico.schloemer@gmail.com>
% All rights reserved.
% Minimal pgfplots version: 1.3
% 
%
% defining custom colors
\definecolor{mycolor1}{rgb}{0.00000,1.00000,1.00000}%
\definecolor{mycolor2}{rgb}{1.00000,0.00000,1.00000}%
%
\begin{tikzpicture}

\begin{axis}[%
width=\figurewidth,
height=\figureheight,
scale only axis,
xmin=0,
xmax=100,
ymin=0,
ymax=0.01,
name=plot3,
title={Omaxmin}
]
\addplot [color=black,solid,line width=1.2pt,forget plot]
  table[row sep=crcr]{0	0\\
1	0.000493827160493827\\
2	0.000987654320987654\\
3	0.00148148148148148\\
4	0.00197530864197531\\
5	0.00246913580246914\\
6	0.00296296296296296\\
7	0.00345679012345679\\
8	0.00395061728395062\\
9	0.00444444444444444\\
10	0.00493827160493827\\
11	0.0054320987654321\\
12	0.00592592592592593\\
13	0.00641975308641975\\
14	0.00691358024691358\\
15	0.00740740740740741\\
16	0.00786402199374565\\
17	0.00819324010074154\\
18	0.00852245820773742\\
19	0.00856179700343862\\
20	0.00856179700343862\\
21	0.00856179700343862\\
22	0.00856179700343862\\
23	0.00856179700343862\\
24	0.00856179700343862\\
25	0.00856179700343862\\
26	0.00872640605693656\\
27	0.0088910151104345\\
28	0.00905562416393245\\
29	0.00922023321743039\\
30	0.00938484227092833\\
31	0.00954945132442627\\
32	0.00967472158222302\\
33	0.00951011252872507\\
34	0.00934550347522713\\
35	0.00905349794238683\\
36	0.00872427983539095\\
37	0.00839506172839506\\
38	0.00806584362139918\\
39	0.00773662551440329\\
40	0.00740740740740741\\
41	0.00704097672625594\\
42	0.00654714956576211\\
43	0.00605332240526828\\
44	0.00555949524477446\\
45	0.00506566808428063\\
46	0.0045718409237868\\
47	0.00407801376329298\\
48	0.00358418660279915\\
49	0.00309035944230532\\
50	0.0025965322818115\\
51	0.0025965322818115\\
52	0.0025965322818115\\
53	0.0025965322818115\\
54	0.0025965322818115\\
55	0.0025965322818115\\
56	0.0025965322818115\\
57	0.0025965322818115\\
58	0.0025965322818115\\
59	0.0025965322818115\\
60	0.00246913580246914\\
61	0.00230452674897119\\
62	0.00213991769547325\\
63	0.00197530864197531\\
64	0.00181069958847737\\
65	0.00164609053497942\\
66	0.00148148148148148\\
67	0.00131687242798354\\
68	0.0011522633744856\\
69	0.000987654320987654\\
70	0.000823045267489712\\
71	0.00065843621399177\\
72	0.000493827160493827\\
73	0.000329218106995885\\
74	0.000164609053497943\\
75	0\\
76	0\\
77	0\\
78	0\\
79	0\\
80	0\\
81	0\\
82	0\\
83	0\\
84	0\\
85	0\\
86	0\\
87	0\\
88	0\\
89	0\\
90	0\\
91	0\\
92	0\\
93	0\\
94	0\\
95	0\\
96	0\\
97	0\\
98	0\\
99	0\\
100	0\\
};
\addplot [color=blue,line width=1.2pt,only marks,mark=asterisk,mark options={solid},forget plot]
  table[row sep=crcr]{31	0.00954945132442627\\
};
\addplot [color=green,line width=1.2pt,only marks,mark=+,mark options={solid},forget plot]
  table[row sep=crcr]{30	0.00938484227092833\\
};
\addplot [color=red,line width=1.2pt,only marks,mark=square,mark options={solid},forget plot]
  table[row sep=crcr]{32	0.00967472158222302\\
};
\addplot [color=mycolor1,line width=1.2pt,only marks,mark=triangle,mark options={solid,rotate=180},forget plot]
  table[row sep=crcr]{32	0.00967472158222302\\
};
\addplot [color=mycolor2,line width=1.2pt,only marks,mark=triangle,mark options={solid},forget plot]
  table[row sep=crcr]{32	0.00967472158222302\\
};
\end{axis}

\begin{axis}[%
width=\figurewidth,
height=\figureheight,
scale only axis,
xmin=0,
xmax=100,
ymin=0,
ymax=0.0001,
name=plot1,
at=(plot3.above north west),
anchor=below south west,
title={Opi}
]
\addplot [color=black,solid,line width=1.2pt,forget plot]
  table[row sep=crcr]{0	0\\
1	5.40126498038937e-05\\
2	5.40126498038937e-05\\
3	5.40126498038937e-05\\
4	5.40126498038937e-05\\
5	5.40126498038937e-05\\
6	5.40126498038937e-05\\
7	5.40126498038937e-05\\
8	5.40126498038937e-05\\
9	5.40126498038937e-05\\
10	5.40126498038937e-05\\
11	5.40126498038937e-05\\
12	5.40126498038937e-05\\
13	5.40126498038937e-05\\
14	5.40126498038937e-05\\
15	5.40126498038937e-05\\
16	5.40126498038937e-05\\
17	5.40126498038937e-05\\
18	5.40126498038937e-05\\
19	5.40126498038937e-05\\
20	5.40126498038937e-05\\
21	5.40126498038937e-05\\
22	5.40126498038937e-05\\
23	5.40126498038937e-05\\
24	5.40126498038937e-05\\
25	5.40126498038937e-05\\
26	7.03930415354032e-05\\
27	7.03930415354032e-05\\
28	7.03930415354032e-05\\
29	7.03930415354032e-05\\
30	7.03930415354032e-05\\
31	7.03930415354032e-05\\
32	7.03930415354032e-05\\
33	7.03930415354032e-05\\
34	7.03930415354032e-05\\
35	7.03930415354032e-05\\
36	7.03930415354032e-05\\
37	7.03930415354032e-05\\
38	7.03930415354032e-05\\
39	7.03930415354032e-05\\
40	7.03930415354032e-05\\
41	7.03930415354032e-05\\
42	7.03930415354032e-05\\
43	7.03930415354032e-05\\
44	7.03930415354032e-05\\
45	7.03930415354032e-05\\
46	7.03930415354032e-05\\
47	7.03930415354032e-05\\
48	7.03930415354032e-05\\
49	7.03930415354032e-05\\
50	1.63803917315096e-05\\
51	1.63803917315096e-05\\
52	1.63803917315096e-05\\
53	1.63803917315096e-05\\
54	1.63803917315096e-05\\
55	1.63803917315096e-05\\
56	1.63803917315096e-05\\
57	1.63803917315096e-05\\
58	1.63803917315096e-05\\
59	1.63803917315096e-05\\
60	1.63803917315096e-05\\
61	1.63803917315096e-05\\
62	1.63803917315096e-05\\
63	1.63803917315096e-05\\
64	1.63803917315096e-05\\
65	1.63803917315096e-05\\
66	1.63803917315096e-05\\
67	1.63803917315096e-05\\
68	1.63803917315096e-05\\
69	1.63803917315096e-05\\
70	1.63803917315096e-05\\
71	1.63803917315096e-05\\
72	1.63803917315096e-05\\
73	1.63803917315096e-05\\
74	1.63803917315096e-05\\
75	0\\
76	0\\
77	0\\
78	0\\
79	0\\
80	0\\
81	0\\
82	0\\
83	0\\
84	0\\
85	0\\
86	0\\
87	0\\
88	0\\
89	0\\
90	0\\
91	0\\
92	0\\
93	0\\
94	0\\
95	0\\
96	0\\
97	0\\
98	0\\
99	0\\
100	0\\
};
\addplot [color=blue,line width=1.2pt,only marks,mark=asterisk,mark options={solid},forget plot]
  table[row sep=crcr]{31	7.03930415354032e-05\\
};
\addplot [color=green,line width=1.2pt,only marks,mark=+,mark options={solid},forget plot]
  table[row sep=crcr]{31	7.03930415354032e-05\\
};
\addplot [color=red,line width=1.2pt,only marks,mark=square,mark options={solid},forget plot]
  table[row sep=crcr]{38	7.03930415354032e-05\\
};
\addplot [color=mycolor1,line width=1.2pt,only marks,mark=triangle,mark options={solid,rotate=180},forget plot]
  table[row sep=crcr]{26	7.03930415354032e-05\\
};
\addplot [color=mycolor2,line width=1.2pt,only marks,mark=triangle,mark options={solid},forget plot]
  table[row sep=crcr]{49	7.03930415354032e-05\\
};
\end{axis}

\begin{axis}[%
width=\figurewidth,
height=\figureheight,
scale only axis,
xmin=0,
xmax=100,
ymin=0,
ymax=0.0001,
name=plot2,
at=(plot1.right of south east),
anchor=left of south west,
title={Oavgmin}
]
\addplot [color=black,solid,line width=1.2pt,forget plot]
  table[row sep=crcr]{0	0\\
1	5.40126498038937e-05\\
2	5.40126498038937e-05\\
3	5.40126498038937e-05\\
4	5.40126498038937e-05\\
5	5.40126498038937e-05\\
6	5.40126498038937e-05\\
7	5.40126498038937e-05\\
8	5.40126498038937e-05\\
9	5.40126498038937e-05\\
10	5.40126498038937e-05\\
11	5.40126498038937e-05\\
12	5.40126498038937e-05\\
13	5.40126498038937e-05\\
14	5.40126498038937e-05\\
15	5.40126498038937e-05\\
16	5.40126498038937e-05\\
17	5.40126498038937e-05\\
18	5.40126498038937e-05\\
19	5.40126498038937e-05\\
20	5.40126498038937e-05\\
21	5.40126498038937e-05\\
22	5.40126498038937e-05\\
23	5.40126498038937e-05\\
24	5.40126498038937e-05\\
25	5.40126498038937e-05\\
26	7.03930415354032e-05\\
27	7.03930415354032e-05\\
28	7.03930415354032e-05\\
29	7.03930415354032e-05\\
30	7.03930415354032e-05\\
31	7.03930415354032e-05\\
32	7.03930415354032e-05\\
33	7.03930415354032e-05\\
34	7.03930415354032e-05\\
35	7.03930415354032e-05\\
36	7.03930415354032e-05\\
37	7.03930415354032e-05\\
38	7.03930415354032e-05\\
39	7.03930415354032e-05\\
40	7.03930415354032e-05\\
41	7.03930415354032e-05\\
42	7.03930415354032e-05\\
43	7.03930415354032e-05\\
44	7.03930415354032e-05\\
45	7.03930415354032e-05\\
46	7.03930415354032e-05\\
47	7.03930415354032e-05\\
48	7.03930415354032e-05\\
49	7.03930415354032e-05\\
50	1.63803917315096e-05\\
51	1.63803917315096e-05\\
52	1.63803917315096e-05\\
53	1.63803917315096e-05\\
54	1.63803917315096e-05\\
55	1.63803917315096e-05\\
56	1.63803917315096e-05\\
57	1.63803917315096e-05\\
58	1.63803917315096e-05\\
59	1.63803917315096e-05\\
60	1.63803917315096e-05\\
61	1.63803917315096e-05\\
62	1.63803917315096e-05\\
63	1.63803917315096e-05\\
64	1.63803917315096e-05\\
65	1.63803917315096e-05\\
66	1.63803917315096e-05\\
67	1.63803917315096e-05\\
68	1.63803917315096e-05\\
69	1.63803917315096e-05\\
70	1.63803917315096e-05\\
71	1.63803917315096e-05\\
72	1.63803917315096e-05\\
73	1.63803917315096e-05\\
74	1.63803917315096e-05\\
75	0\\
76	0\\
77	0\\
78	0\\
79	0\\
80	0\\
81	0\\
82	0\\
83	0\\
84	0\\
85	0\\
86	0\\
87	0\\
88	0\\
89	0\\
90	0\\
91	0\\
92	0\\
93	0\\
94	0\\
95	0\\
96	0\\
97	0\\
98	0\\
99	0\\
100	0\\
};
\addplot [color=blue,line width=1.2pt,only marks,mark=asterisk,mark options={solid},forget plot]
  table[row sep=crcr]{31	7.03930415354032e-05\\
};
\addplot [color=green,line width=1.2pt,only marks,mark=+,mark options={solid},forget plot]
  table[row sep=crcr]{31	7.03930415354032e-05\\
};
\addplot [color=red,line width=1.2pt,only marks,mark=square,mark options={solid},forget plot]
  table[row sep=crcr]{38	7.03930415354032e-05\\
};
\addplot [color=mycolor1,line width=1.2pt,only marks,mark=triangle,mark options={solid,rotate=180},forget plot]
  table[row sep=crcr]{26	7.03930415354032e-05\\
};
\addplot [color=mycolor2,line width=1.2pt,only marks,mark=triangle,mark options={solid},forget plot]
  table[row sep=crcr]{49	7.03930415354032e-05\\
};
\end{axis}

\begin{axis}[%
width=\figurewidth,
height=\figureheight,
scale only axis,
xmin=0,
xmax=100,
ymin=0,
ymax=0.0001,
at=(plot2.below south west),
anchor=above north west,
title={Osqrt}
]
\addplot [color=black,solid,line width=1.2pt,forget plot]
  table[row sep=crcr]{0	0\\
1	6.48252546089399e-05\\
2	6.48252546089399e-05\\
3	6.48252546089399e-05\\
4	6.48252546089399e-05\\
5	6.48252546089399e-05\\
6	6.48252546089399e-05\\
7	6.48252546089399e-05\\
8	6.48252546089399e-05\\
9	6.48252546089399e-05\\
10	6.48252546089399e-05\\
11	6.48252546089399e-05\\
12	6.48252546089399e-05\\
13	6.48252546089399e-05\\
14	6.48252546089399e-05\\
15	6.48252546089399e-05\\
16	6.48252546089399e-05\\
17	6.48252546089399e-05\\
18	6.48252546089399e-05\\
19	6.48252546089399e-05\\
20	6.48252546089399e-05\\
21	6.48252546089399e-05\\
22	6.48252546089399e-05\\
23	6.48252546089399e-05\\
24	6.48252546089399e-05\\
25	6.48252546089399e-05\\
26	8.54469460148803e-05\\
27	8.54469460148803e-05\\
28	8.54469460148803e-05\\
29	8.54469460148803e-05\\
30	8.54469460148803e-05\\
31	8.54469460148803e-05\\
32	8.54469460148803e-05\\
33	8.54469460148803e-05\\
34	8.54469460148803e-05\\
35	8.54469460148803e-05\\
36	8.54469460148803e-05\\
37	8.54469460148803e-05\\
38	8.54469460148803e-05\\
39	8.54469460148803e-05\\
40	8.54469460148803e-05\\
41	8.54469460148803e-05\\
42	8.54469460148803e-05\\
43	8.54469460148803e-05\\
44	8.54469460148803e-05\\
45	8.54469460148803e-05\\
46	8.54469460148803e-05\\
47	8.54469460148803e-05\\
48	8.54469460148803e-05\\
49	8.54469460148803e-05\\
50	2.06216914059404e-05\\
51	2.06216914059404e-05\\
52	2.06216914059404e-05\\
53	2.06216914059404e-05\\
54	2.06216914059404e-05\\
55	2.06216914059404e-05\\
56	2.06216914059404e-05\\
57	2.06216914059404e-05\\
58	2.06216914059404e-05\\
59	2.06216914059404e-05\\
60	2.06216914059404e-05\\
61	2.06216914059404e-05\\
62	2.06216914059404e-05\\
63	2.06216914059404e-05\\
64	2.06216914059404e-05\\
65	2.06216914059404e-05\\
66	2.06216914059404e-05\\
67	2.06216914059404e-05\\
68	2.06216914059404e-05\\
69	2.06216914059404e-05\\
70	2.06216914059404e-05\\
71	2.06216914059404e-05\\
72	2.06216914059404e-05\\
73	2.06216914059404e-05\\
74	2.06216914059404e-05\\
75	0\\
76	0\\
77	0\\
78	0\\
79	0\\
80	0\\
81	0\\
82	0\\
83	0\\
84	0\\
85	0\\
86	0\\
87	0\\
88	0\\
89	0\\
90	0\\
91	0\\
92	0\\
93	0\\
94	0\\
95	0\\
96	0\\
97	0\\
98	0\\
99	0\\
100	0\\
};
\addplot [color=blue,line width=1.2pt,only marks,mark=asterisk,mark options={solid},forget plot]
  table[row sep=crcr]{31	8.54469460148803e-05\\
};
\addplot [color=green,line width=1.2pt,only marks,mark=+,mark options={solid},forget plot]
  table[row sep=crcr]{31	8.54469460148803e-05\\
};
\addplot [color=red,line width=1.2pt,only marks,mark=square,mark options={solid},forget plot]
  table[row sep=crcr]{38	8.54469460148803e-05\\
};
\addplot [color=mycolor1,line width=1.2pt,only marks,mark=triangle,mark options={solid,rotate=180},forget plot]
  table[row sep=crcr]{26	8.54469460148803e-05\\
};
\addplot [color=mycolor2,line width=1.2pt,only marks,mark=triangle,mark options={solid},forget plot]
  table[row sep=crcr]{49	8.54469460148803e-05\\
};
\end{axis}

\begin{axis}[%
width=\figurewidth,
height=\figureheight,
scale only axis,
xmin=0,
xmax=100,
ymin=0,
ymax=0.00015,
at=(plot3.below south west),
anchor=above north west,
title={Osin}
]
\addplot [color=black,solid,line width=1.2pt,forget plot]
  table[row sep=crcr]{0	0\\
1	7.61868286641183e-05\\
2	7.61868286641183e-05\\
3	7.61868286641183e-05\\
4	7.61868286641183e-05\\
5	7.61868286641183e-05\\
6	7.61868286641183e-05\\
7	7.61868286641183e-05\\
8	7.61868286641183e-05\\
9	7.61868286641183e-05\\
10	7.61868286641183e-05\\
11	7.61868286641183e-05\\
12	7.61868286641183e-05\\
13	7.61868286641183e-05\\
14	7.61868286641183e-05\\
15	7.61868286641183e-05\\
16	7.61868286641183e-05\\
17	7.61868286641183e-05\\
18	7.61868286641183e-05\\
19	7.61868286641183e-05\\
20	7.61868286641183e-05\\
21	7.61868286641183e-05\\
22	7.61868286641183e-05\\
23	7.61868286641183e-05\\
24	7.61868286641183e-05\\
25	7.61868286641183e-05\\
26	0.0001013182283749\\
27	0.0001013182283749\\
28	0.0001013182283749\\
29	0.0001013182283749\\
30	0.0001013182283749\\
31	0.0001013182283749\\
32	0.0001013182283749\\
33	0.0001013182283749\\
34	0.0001013182283749\\
35	0.0001013182283749\\
36	0.0001013182283749\\
37	0.0001013182283749\\
38	0.0001013182283749\\
39	0.0001013182283749\\
40	0.0001013182283749\\
41	0.0001013182283749\\
42	0.0001013182283749\\
43	0.0001013182283749\\
44	0.0001013182283749\\
45	0.0001013182283749\\
46	0.0001013182283749\\
47	0.0001013182283749\\
48	0.0001013182283749\\
49	0.0001013182283749\\
50	2.51313997107821e-05\\
51	2.51313997107821e-05\\
52	2.51313997107821e-05\\
53	2.51313997107821e-05\\
54	2.51313997107821e-05\\
55	2.51313997107821e-05\\
56	2.51313997107821e-05\\
57	2.51313997107821e-05\\
58	2.51313997107821e-05\\
59	2.51313997107821e-05\\
60	2.51313997107821e-05\\
61	2.51313997107821e-05\\
62	2.51313997107821e-05\\
63	2.51313997107821e-05\\
64	2.51313997107821e-05\\
65	2.51313997107821e-05\\
66	2.51313997107821e-05\\
67	2.51313997107821e-05\\
68	2.51313997107821e-05\\
69	2.51313997107821e-05\\
70	2.51313997107821e-05\\
71	2.51313997107821e-05\\
72	2.51313997107821e-05\\
73	2.51313997107821e-05\\
74	2.51313997107821e-05\\
75	0\\
76	0\\
77	0\\
78	0\\
79	0\\
80	0\\
81	0\\
82	0\\
83	0\\
84	0\\
85	0\\
86	0\\
87	0\\
88	0\\
89	0\\
90	0\\
91	0\\
92	0\\
93	0\\
94	0\\
95	0\\
96	0\\
97	0\\
98	0\\
99	0\\
100	0\\
};
\addplot [color=blue,line width=1.2pt,only marks,mark=asterisk,mark options={solid},forget plot]
  table[row sep=crcr]{31	0.0001013182283749\\
};
\addplot [color=green,line width=1.2pt,only marks,mark=+,mark options={solid},forget plot]
  table[row sep=crcr]{31	0.0001013182283749\\
};
\addplot [color=red,line width=1.2pt,only marks,mark=square,mark options={solid},forget plot]
  table[row sep=crcr]{38	0.0001013182283749\\
};
\addplot [color=mycolor1,line width=1.2pt,only marks,mark=triangle,mark options={solid,rotate=180},forget plot]
  table[row sep=crcr]{26	0.0001013182283749\\
};
\addplot [color=mycolor2,line width=1.2pt,only marks,mark=triangle,mark options={solid},forget plot]
  table[row sep=crcr]{49	0.0001013182283749\\
};
\end{axis}
\end{tikzpicture}%
\end{subfigure}
\caption{T-norma: Media geométrica ($T_{geo}$)}
\label{fig:fire-detection-interpolation-comparison-tnorm-geomean}
\end{figure}
\begin{figure}[h!]
\centering
\begin{subfigure}[b]{\textwidth}
	\textbf{T-norma: Media armónica ($T_{harm}$)} \\
	\setlength\figureheight{2cm}
	\setlength\figurewidth{7cm}
	% This file was created by matlab2tikz v0.4.7 running on MATLAB 8.0.
% Copyright (c) 2008--2014, Nico Schlömer <nico.schloemer@gmail.com>
% All rights reserved.
% Minimal pgfplots version: 1.3
% 
%
% defining custom colors
\definecolor{mycolor1}{rgb}{0.00000,1.00000,1.00000}%
\definecolor{mycolor2}{rgb}{1.00000,0.00000,1.00000}%
%
\begin{tikzpicture}

\begin{axis}[%
width=\figurewidth,
height=\figureheight,
scale only axis,
xmin=0,
xmax=100,
ymin=0,
ymax=0.01,
name=plot3,
title={$O_Z$}
]
\addplot [color=black,solid,line width=1.2pt,forget plot]
  table[row sep=crcr]{0	0\\
1	0.000493827160493827\\
2	0.000987654320987654\\
3	0.00148148148148148\\
4	0.00197530864197531\\
5	0.00246913580246914\\
6	0.00296296296296296\\
7	0.00345679012345679\\
8	0.00395061728395062\\
9	0.00444444444444444\\
10	0.00493827160493827\\
11	0.0054320987654321\\
12	0.00592592592592593\\
13	0.00641975308641975\\
14	0.00677496209659952\\
15	0.00710418020359541\\
16	0.00743339831059129\\
17	0.00765287704858855\\
18	0.00765287704858855\\
19	0.00765287704858855\\
20	0.00765287704858855\\
21	0.00765287704858855\\
22	0.00765287704858855\\
23	0.00765287704858855\\
24	0.00765287704858855\\
25	0.00765287704858855\\
26	0.00781748610208649\\
27	0.00798209515558443\\
28	0.00814670420908238\\
29	0.00831131326258032\\
30	0.00847592231607826\\
31	0.0086405313695762\\
32	0.00880514042307414\\
33	0.00896974947657209\\
34	0.00891487979207277\\
35	0.00875027073857483\\
36	0.00858566168507689\\
37	0.00839506172839506\\
38	0.00806584362139918\\
39	0.00759800736408923\\
40	0.00710418020359541\\
41	0.00661035304310158\\
42	0.00611652588260775\\
43	0.00562269872211393\\
44	0.0051288715616201\\
45	0.00463504440112627\\
46	0.00414121724063245\\
47	0.00364739008013862\\
48	0.00315356291964479\\
49	0.00265973575915096\\
50	0.00216590859865714\\
51	0.00216590859865714\\
52	0.00216590859865714\\
53	0.00216590859865714\\
54	0.00216590859865714\\
55	0.00216590859865714\\
56	0.00216590859865714\\
57	0.00216590859865714\\
58	0.00216590859865714\\
59	0.00216590859865714\\
60	0.00216590859865714\\
61	0.00216590859865714\\
62	0.00213991769547325\\
63	0.00197530864197531\\
64	0.00181069958847737\\
65	0.00164609053497942\\
66	0.00148148148148148\\
67	0.00131687242798354\\
68	0.0011522633744856\\
69	0.000987654320987654\\
70	0.000823045267489712\\
71	0.00065843621399177\\
72	0.000493827160493827\\
73	0.000329218106995885\\
74	0.000164609053497943\\
75	0\\
76	0\\
77	0\\
78	0\\
79	0\\
80	0\\
81	0\\
82	0\\
83	0\\
84	0\\
85	0\\
86	0\\
87	0\\
88	0\\
89	0\\
90	0\\
91	0\\
92	0\\
93	0\\
94	0\\
95	0\\
96	0\\
97	0\\
98	0\\
99	0\\
100	0\\
};
\addplot [color=blue,line width=1.2pt,only marks,mark=asterisk,mark options={solid},forget plot]
  table[row sep=crcr]{31	0.0086405313695762\\
};
\addplot [color=green,line width=1.2pt,only marks,mark=+,mark options={solid},forget plot]
  table[row sep=crcr]{30	0.00847592231607826\\
};
\addplot [color=red,line width=1.2pt,only marks,mark=square,mark options={solid},forget plot]
  table[row sep=crcr]{33	0.00896974947657209\\
};
\addplot [color=mycolor1,line width=1.2pt,only marks,mark=triangle,mark options={solid,rotate=180},forget plot]
  table[row sep=crcr]{33	0.00896974947657209\\
};
\addplot [color=mycolor2,line width=1.2pt,only marks,mark=triangle,mark options={solid},forget plot]
  table[row sep=crcr]{33	0.00896974947657209\\
};
\end{axis}

\begin{axis}[%
width=\figurewidth,
height=\figureheight,
scale only axis,
xmin=0,
xmax=100,
ymin=0,
ymax=6e-05,
name=plot1,
at=(plot3.above north west),
anchor=below south west,
title={$O_{\pi}$}
]
\addplot [color=black,solid,line width=1.2pt,forget plot]
  table[row sep=crcr]{0	0\\
1	3.52762494858135e-05\\
2	3.52762494858135e-05\\
3	3.52762494858135e-05\\
4	3.52762494858135e-05\\
5	3.52762494858135e-05\\
6	3.52762494858135e-05\\
7	3.52762494858135e-05\\
8	3.52762494858135e-05\\
9	3.52762494858135e-05\\
10	3.52762494858135e-05\\
11	3.52762494858135e-05\\
12	3.52762494858135e-05\\
13	3.52762494858135e-05\\
14	3.52762494858135e-05\\
15	3.52762494858135e-05\\
16	3.52762494858135e-05\\
17	3.52762494858135e-05\\
18	3.52762494858135e-05\\
19	3.52762494858135e-05\\
20	3.52762494858135e-05\\
21	3.52762494858135e-05\\
22	3.52762494858135e-05\\
23	3.52762494858135e-05\\
24	3.52762494858135e-05\\
25	3.52762494858135e-05\\
26	4.61659725274874e-05\\
27	4.61659725274874e-05\\
28	4.61659725274874e-05\\
29	4.61659725274874e-05\\
30	4.61659725274874e-05\\
31	4.61659725274874e-05\\
32	4.61659725274874e-05\\
33	4.61659725274874e-05\\
34	4.61659725274874e-05\\
35	4.61659725274874e-05\\
36	4.61659725274874e-05\\
37	4.61659725274874e-05\\
38	4.61659725274874e-05\\
39	4.61659725274874e-05\\
40	4.61659725274874e-05\\
41	4.61659725274874e-05\\
42	4.61659725274874e-05\\
43	4.61659725274874e-05\\
44	4.61659725274874e-05\\
45	4.61659725274874e-05\\
46	4.61659725274874e-05\\
47	4.61659725274874e-05\\
48	4.61659725274874e-05\\
49	4.61659725274874e-05\\
50	1.08897230416739e-05\\
51	1.08897230416739e-05\\
52	1.08897230416739e-05\\
53	1.08897230416739e-05\\
54	1.08897230416739e-05\\
55	1.08897230416739e-05\\
56	1.08897230416739e-05\\
57	1.08897230416739e-05\\
58	1.08897230416739e-05\\
59	1.08897230416739e-05\\
60	1.08897230416739e-05\\
61	1.08897230416739e-05\\
62	1.08897230416739e-05\\
63	1.08897230416739e-05\\
64	1.08897230416739e-05\\
65	1.08897230416739e-05\\
66	1.08897230416739e-05\\
67	1.08897230416739e-05\\
68	1.08897230416739e-05\\
69	1.08897230416739e-05\\
70	1.08897230416739e-05\\
71	1.08897230416739e-05\\
72	1.08897230416739e-05\\
73	1.08897230416739e-05\\
74	1.08897230416739e-05\\
75	0\\
76	0\\
77	0\\
78	0\\
79	0\\
80	0\\
81	0\\
82	0\\
83	0\\
84	0\\
85	0\\
86	0\\
87	0\\
88	0\\
89	0\\
90	0\\
91	0\\
92	0\\
93	0\\
94	0\\
95	0\\
96	0\\
97	0\\
98	0\\
99	0\\
100	0\\
};
\addplot [color=blue,line width=1.2pt,only marks,mark=asterisk,mark options={solid},forget plot]
  table[row sep=crcr]{31	4.61659725274874e-05\\
};
\addplot [color=green,line width=1.2pt,only marks,mark=+,mark options={solid},forget plot]
  table[row sep=crcr]{31	4.61659725274874e-05\\
};
\addplot [color=red,line width=1.2pt,only marks,mark=square,mark options={solid},forget plot]
  table[row sep=crcr]{38	4.61659725274874e-05\\
};
\addplot [color=mycolor1,line width=1.2pt,only marks,mark=triangle,mark options={solid,rotate=180},forget plot]
  table[row sep=crcr]{26	4.61659725274874e-05\\
};
\addplot [color=mycolor2,line width=1.2pt,only marks,mark=triangle,mark options={solid},forget plot]
  table[row sep=crcr]{49	4.61659725274874e-05\\
};
\end{axis}

\begin{axis}[%
width=\figurewidth,
height=\figureheight,
scale only axis,
xmin=0,
xmax=100,
ymin=0,
ymax=6e-05,
name=plot2,
at=(plot1.right of south east),
anchor=left of south west,
title={$O_{avgmin}$}
]
\addplot [color=black,solid,line width=1.2pt,forget plot]
  table[row sep=crcr]{0	0\\
1	3.52762494858135e-05\\
2	3.52762494858135e-05\\
3	3.52762494858135e-05\\
4	3.52762494858135e-05\\
5	3.52762494858135e-05\\
6	3.52762494858135e-05\\
7	3.52762494858135e-05\\
8	3.52762494858135e-05\\
9	3.52762494858135e-05\\
10	3.52762494858135e-05\\
11	3.52762494858135e-05\\
12	3.52762494858135e-05\\
13	3.52762494858135e-05\\
14	3.52762494858135e-05\\
15	3.52762494858135e-05\\
16	3.52762494858135e-05\\
17	3.52762494858135e-05\\
18	3.52762494858135e-05\\
19	3.52762494858135e-05\\
20	3.52762494858135e-05\\
21	3.52762494858135e-05\\
22	3.52762494858135e-05\\
23	3.52762494858135e-05\\
24	3.52762494858135e-05\\
25	3.52762494858135e-05\\
26	4.61659725274874e-05\\
27	4.61659725274874e-05\\
28	4.61659725274874e-05\\
29	4.61659725274874e-05\\
30	4.61659725274874e-05\\
31	4.61659725274874e-05\\
32	4.61659725274874e-05\\
33	4.61659725274874e-05\\
34	4.61659725274874e-05\\
35	4.61659725274874e-05\\
36	4.61659725274874e-05\\
37	4.61659725274874e-05\\
38	4.61659725274874e-05\\
39	4.61659725274874e-05\\
40	4.61659725274874e-05\\
41	4.61659725274874e-05\\
42	4.61659725274874e-05\\
43	4.61659725274874e-05\\
44	4.61659725274874e-05\\
45	4.61659725274874e-05\\
46	4.61659725274874e-05\\
47	4.61659725274874e-05\\
48	4.61659725274874e-05\\
49	4.61659725274874e-05\\
50	1.08897230416739e-05\\
51	1.08897230416739e-05\\
52	1.08897230416739e-05\\
53	1.08897230416739e-05\\
54	1.08897230416739e-05\\
55	1.08897230416739e-05\\
56	1.08897230416739e-05\\
57	1.08897230416739e-05\\
58	1.08897230416739e-05\\
59	1.08897230416739e-05\\
60	1.08897230416739e-05\\
61	1.08897230416739e-05\\
62	1.08897230416739e-05\\
63	1.08897230416739e-05\\
64	1.08897230416739e-05\\
65	1.08897230416739e-05\\
66	1.08897230416739e-05\\
67	1.08897230416739e-05\\
68	1.08897230416739e-05\\
69	1.08897230416739e-05\\
70	1.08897230416739e-05\\
71	1.08897230416739e-05\\
72	1.08897230416739e-05\\
73	1.08897230416739e-05\\
74	1.08897230416739e-05\\
75	0\\
76	0\\
77	0\\
78	0\\
79	0\\
80	0\\
81	0\\
82	0\\
83	0\\
84	0\\
85	0\\
86	0\\
87	0\\
88	0\\
89	0\\
90	0\\
91	0\\
92	0\\
93	0\\
94	0\\
95	0\\
96	0\\
97	0\\
98	0\\
99	0\\
100	0\\
};
\addplot [color=blue,line width=1.2pt,only marks,mark=asterisk,mark options={solid},forget plot]
  table[row sep=crcr]{31	4.61659725274874e-05\\
};
\addplot [color=green,line width=1.2pt,only marks,mark=+,mark options={solid},forget plot]
  table[row sep=crcr]{31	4.61659725274874e-05\\
};
\addplot [color=red,line width=1.2pt,only marks,mark=square,mark options={solid},forget plot]
  table[row sep=crcr]{38	4.61659725274874e-05\\
};
\addplot [color=mycolor1,line width=1.2pt,only marks,mark=triangle,mark options={solid,rotate=180},forget plot]
  table[row sep=crcr]{26	4.61659725274874e-05\\
};
\addplot [color=mycolor2,line width=1.2pt,only marks,mark=triangle,mark options={solid},forget plot]
  table[row sep=crcr]{49	4.61659725274874e-05\\
};
\end{axis}

\begin{axis}[%
width=\figurewidth,
height=\figureheight,
scale only axis,
xmin=0,
xmax=100,
ymin=0,
ymax=6e-05,
at=(plot2.below south west),
anchor=above north west,
title={$O_{\sqrt{\text{ }}}$}
]
\addplot [color=black,solid,line width=1.2pt,forget plot]
  table[row sep=crcr]{0	0\\
1	4.01779450484702e-05\\
2	4.01779450484702e-05\\
3	4.01779450484702e-05\\
4	4.01779450484702e-05\\
5	4.01779450484702e-05\\
6	4.01779450484702e-05\\
7	4.01779450484702e-05\\
8	4.01779450484702e-05\\
9	4.01779450484702e-05\\
10	4.01779450484702e-05\\
11	4.01779450484702e-05\\
12	4.01779450484702e-05\\
13	4.01779450484702e-05\\
14	4.01779450484702e-05\\
15	4.01779450484702e-05\\
16	4.01779450484702e-05\\
17	4.01779450484702e-05\\
18	4.01779450484702e-05\\
19	4.01779450484702e-05\\
20	4.01779450484702e-05\\
21	4.01779450484702e-05\\
22	4.01779450484702e-05\\
23	4.01779450484702e-05\\
24	4.01779450484702e-05\\
25	4.01779450484702e-05\\
26	5.32312339903342e-05\\
27	5.32312339903342e-05\\
28	5.32312339903342e-05\\
29	5.32312339903342e-05\\
30	5.32312339903342e-05\\
31	5.32312339903342e-05\\
32	5.32312339903342e-05\\
33	5.32312339903342e-05\\
34	5.32312339903342e-05\\
35	5.32312339903342e-05\\
36	5.32312339903342e-05\\
37	5.32312339903342e-05\\
38	5.32312339903342e-05\\
39	5.32312339903342e-05\\
40	5.32312339903342e-05\\
41	5.32312339903342e-05\\
42	5.32312339903342e-05\\
43	5.32312339903342e-05\\
44	5.32312339903342e-05\\
45	5.32312339903342e-05\\
46	5.32312339903342e-05\\
47	5.32312339903342e-05\\
48	5.32312339903342e-05\\
49	5.32312339903342e-05\\
50	1.30532889418641e-05\\
51	1.30532889418641e-05\\
52	1.30532889418641e-05\\
53	1.30532889418641e-05\\
54	1.30532889418641e-05\\
55	1.30532889418641e-05\\
56	1.30532889418641e-05\\
57	1.30532889418641e-05\\
58	1.30532889418641e-05\\
59	1.30532889418641e-05\\
60	1.30532889418641e-05\\
61	1.30532889418641e-05\\
62	1.30532889418641e-05\\
63	1.30532889418641e-05\\
64	1.30532889418641e-05\\
65	1.30532889418641e-05\\
66	1.30532889418641e-05\\
67	1.30532889418641e-05\\
68	1.30532889418641e-05\\
69	1.30532889418641e-05\\
70	1.30532889418641e-05\\
71	1.30532889418641e-05\\
72	1.30532889418641e-05\\
73	1.30532889418641e-05\\
74	1.30532889418641e-05\\
75	0\\
76	0\\
77	0\\
78	0\\
79	0\\
80	0\\
81	0\\
82	0\\
83	0\\
84	0\\
85	0\\
86	0\\
87	0\\
88	0\\
89	0\\
90	0\\
91	0\\
92	0\\
93	0\\
94	0\\
95	0\\
96	0\\
97	0\\
98	0\\
99	0\\
100	0\\
};
\addplot [color=blue,line width=1.2pt,only marks,mark=asterisk,mark options={solid},forget plot]
  table[row sep=crcr]{31	5.32312339903342e-05\\
};
\addplot [color=green,line width=1.2pt,only marks,mark=+,mark options={solid},forget plot]
  table[row sep=crcr]{31	5.32312339903342e-05\\
};
\addplot [color=red,line width=1.2pt,only marks,mark=square,mark options={solid},forget plot]
  table[row sep=crcr]{38	5.32312339903342e-05\\
};
\addplot [color=mycolor1,line width=1.2pt,only marks,mark=triangle,mark options={solid,rotate=180},forget plot]
  table[row sep=crcr]{26	5.32312339903342e-05\\
};
\addplot [color=mycolor2,line width=1.2pt,only marks,mark=triangle,mark options={solid},forget plot]
  table[row sep=crcr]{49	5.32312339903342e-05\\
};
\end{axis}

\begin{axis}[%
width=\figurewidth,
height=\figureheight,
scale only axis,
xmin=0,
xmax=100,
ymin=0,
ymax=6e-05,
at=(plot3.below south west),
anchor=above north west,
title={$O_{sin}$}
]
\addplot [color=black,solid,line width=1.2pt,forget plot]
  table[row sep=crcr]{0	0\\
1	4.35445136869088e-05\\
2	4.35445136869088e-05\\
3	4.35445136869088e-05\\
4	4.35445136869088e-05\\
5	4.35445136869088e-05\\
6	4.35445136869088e-05\\
7	4.35445136869088e-05\\
8	4.35445136869088e-05\\
9	4.35445136869088e-05\\
10	4.35445136869088e-05\\
11	4.35445136869088e-05\\
12	4.35445136869088e-05\\
13	4.35445136869088e-05\\
14	4.35445136869088e-05\\
15	4.35445136869088e-05\\
16	4.35445136869088e-05\\
17	4.35445136869088e-05\\
18	4.35445136869088e-05\\
19	4.35445136869088e-05\\
20	4.35445136869088e-05\\
21	4.35445136869088e-05\\
22	4.35445136869088e-05\\
23	4.35445136869088e-05\\
24	4.35445136869088e-05\\
25	4.35445136869088e-05\\
26	5.80015660788597e-05\\
27	5.80015660788597e-05\\
28	5.80015660788597e-05\\
29	5.80015660788597e-05\\
30	5.80015660788597e-05\\
31	5.80015660788597e-05\\
32	5.80015660788597e-05\\
33	5.80015660788597e-05\\
34	5.80015660788597e-05\\
35	5.80015660788597e-05\\
36	5.80015660788597e-05\\
37	5.80015660788597e-05\\
38	5.80015660788597e-05\\
39	5.80015660788597e-05\\
40	5.80015660788597e-05\\
41	5.80015660788597e-05\\
42	5.80015660788597e-05\\
43	5.80015660788597e-05\\
44	5.80015660788597e-05\\
45	5.80015660788597e-05\\
46	5.80015660788597e-05\\
47	5.80015660788597e-05\\
48	5.80015660788597e-05\\
49	5.80015660788597e-05\\
50	1.44570523919509e-05\\
51	1.44570523919509e-05\\
52	1.44570523919509e-05\\
53	1.44570523919509e-05\\
54	1.44570523919509e-05\\
55	1.44570523919509e-05\\
56	1.44570523919509e-05\\
57	1.44570523919509e-05\\
58	1.44570523919509e-05\\
59	1.44570523919509e-05\\
60	1.44570523919509e-05\\
61	1.44570523919509e-05\\
62	1.44570523919509e-05\\
63	1.44570523919509e-05\\
64	1.44570523919509e-05\\
65	1.44570523919509e-05\\
66	1.44570523919509e-05\\
67	1.44570523919509e-05\\
68	1.44570523919509e-05\\
69	1.44570523919509e-05\\
70	1.44570523919509e-05\\
71	1.44570523919509e-05\\
72	1.44570523919509e-05\\
73	1.44570523919509e-05\\
74	1.44570523919509e-05\\
75	0\\
76	0\\
77	0\\
78	0\\
79	0\\
80	0\\
81	0\\
82	0\\
83	0\\
84	0\\
85	0\\
86	0\\
87	0\\
88	0\\
89	0\\
90	0\\
91	0\\
92	0\\
93	0\\
94	0\\
95	0\\
96	0\\
97	0\\
98	0\\
99	0\\
100	0\\
};
\addplot [color=blue,line width=1.2pt,only marks,mark=asterisk,mark options={solid},forget plot]
  table[row sep=crcr]{31	5.80015660788597e-05\\
};
\addplot [color=green,line width=1.2pt,only marks,mark=+,mark options={solid},forget plot]
  table[row sep=crcr]{31	5.80015660788597e-05\\
};
\addplot [color=red,line width=1.2pt,only marks,mark=square,mark options={solid},forget plot]
  table[row sep=crcr]{38	5.80015660788597e-05\\
};
\addplot [color=mycolor1,line width=1.2pt,only marks,mark=triangle,mark options={solid,rotate=180},forget plot]
  table[row sep=crcr]{26	5.80015660788597e-05\\
};
\addplot [color=mycolor2,line width=1.2pt,only marks,mark=triangle,mark options={solid},forget plot]
  table[row sep=crcr]{49	5.80015660788597e-05\\
};
\end{axis}
\end{tikzpicture}%
\end{subfigure}
\caption{T-norma: Media armónica ($T_{harm}$)}
\label{fig:fire-detection-interpolation-comparison-tnorm-harmmean}
\end{figure}
\begin{figure}[h!]
\centering
\begin{subfigure}[b]{\textwidth}
	\textbf{T-norma: Seno ($T_{sin}$)} \\
	\setlength\figureheight{2cm}
	\setlength\figurewidth{7cm}
	% This file was created by matlab2tikz v0.4.7 running on MATLAB 8.0.
% Copyright (c) 2008--2014, Nico Schlömer <nico.schloemer@gmail.com>
% All rights reserved.
% Minimal pgfplots version: 1.3
% 
%
% defining custom colors
\definecolor{mycolor1}{rgb}{0.00000,1.00000,1.00000}%
\definecolor{mycolor2}{rgb}{1.00000,0.00000,1.00000}%
%
\begin{tikzpicture}

\begin{axis}[%
width=\figurewidth,
height=\figureheight,
scale only axis,
xmin=0,
xmax=100,
ymin=0,
ymax=0.003,
name=plot3,
title={Omaxmin}
]
\addplot [color=black,solid,line width=1.2pt,forget plot]
  table[row sep=crcr]{0	0\\
1	0.000493827160493827\\
2	0.000987654320987654\\
3	0.00147224755756767\\
4	0.00160502331013939\\
5	0.00160502331013939\\
6	0.00160502331013939\\
7	0.00160502331013939\\
8	0.00160502331013939\\
9	0.00160502331013939\\
10	0.00160502331013939\\
11	0.00160502331013939\\
12	0.00160502331013939\\
13	0.00160502331013939\\
14	0.00160502331013939\\
15	0.00160502331013939\\
16	0.00160502331013939\\
17	0.00160502331013939\\
18	0.00160502331013939\\
19	0.00160502331013939\\
20	0.00160502331013939\\
21	0.00160502331013939\\
22	0.00160502331013939\\
23	0.00160502331013939\\
24	0.00160502331013939\\
25	0.00160502331013939\\
26	0.00176963236363733\\
27	0.00193424141713527\\
28	0.0020896165467194\\
29	0.0020896165467194\\
30	0.0020896165467194\\
31	0.0020896165467194\\
32	0.0020896165467194\\
33	0.0020896165467194\\
34	0.0020896165467194\\
35	0.0020896165467194\\
36	0.0020896165467194\\
37	0.0020896165467194\\
38	0.0020896165467194\\
39	0.0020896165467194\\
40	0.0020896165467194\\
41	0.0020896165467194\\
42	0.0020896165467194\\
43	0.0020896165467194\\
44	0.0020896165467194\\
45	0.0020896165467194\\
46	0.0020896165467194\\
47	0.00195684079414768\\
48	0.00147224755756767\\
49	0.000978420397073839\\
50	0.000484593236580011\\
51	0.000484593236580011\\
52	0.000484593236580011\\
53	0.000484593236580011\\
54	0.000484593236580011\\
55	0.000484593236580011\\
56	0.000484593236580011\\
57	0.000484593236580011\\
58	0.000484593236580011\\
59	0.000484593236580011\\
60	0.000484593236580011\\
61	0.000484593236580011\\
62	0.000484593236580011\\
63	0.000484593236580011\\
64	0.000484593236580011\\
65	0.000484593236580011\\
66	0.000484593236580011\\
67	0.000484593236580011\\
68	0.000484593236580011\\
69	0.000484593236580011\\
70	0.000484593236580011\\
71	0.000484593236580011\\
72	0.000484593236580011\\
73	0.000329218106995885\\
74	0.000164609053497943\\
75	0\\
76	0\\
77	0\\
78	0\\
79	0\\
80	0\\
81	0\\
82	0\\
83	0\\
84	0\\
85	0\\
86	0\\
87	0\\
88	0\\
89	0\\
90	0\\
91	0\\
92	0\\
93	0\\
94	0\\
95	0\\
96	0\\
97	0\\
98	0\\
99	0\\
100	0\\
};
\addplot [color=blue,line width=1.2pt,only marks,mark=asterisk,mark options={solid},forget plot]
  table[row sep=crcr]{31	0.0020896165467194\\
};
\addplot [color=green,line width=1.2pt,only marks,mark=+,mark options={solid},forget plot]
  table[row sep=crcr]{31	0.0020896165467194\\
};
\addplot [color=red,line width=1.2pt,only marks,mark=square,mark options={solid},forget plot]
  table[row sep=crcr]{37	0.0020896165467194\\
};
\addplot [color=mycolor1,line width=1.2pt,only marks,mark=triangle,mark options={solid,rotate=180},forget plot]
  table[row sep=crcr]{28	0.0020896165467194\\
};
\addplot [color=mycolor2,line width=1.2pt,only marks,mark=triangle,mark options={solid},forget plot]
  table[row sep=crcr]{46	0.0020896165467194\\
};
\end{axis}

\begin{axis}[%
width=\figurewidth,
height=\figureheight,
scale only axis,
xmin=0,
xmax=100,
ymin=0,
ymax=6e-06,
name=plot1,
at=(plot3.above north west),
anchor=below south west,
title={Opi}
]
\addplot [color=black,solid,line width=1.2pt,forget plot]
  table[row sep=crcr]{0	0\\
1	3.11228721233702e-06\\
2	3.11228721233702e-06\\
3	3.11228721233702e-06\\
4	3.11228721233702e-06\\
5	3.11228721233702e-06\\
6	3.11228721233702e-06\\
7	3.11228721233702e-06\\
8	3.11228721233702e-06\\
9	3.11228721233702e-06\\
10	3.11228721233702e-06\\
11	3.11228721233702e-06\\
12	3.11228721233702e-06\\
13	3.11228721233702e-06\\
14	3.11228721233702e-06\\
15	3.11228721233702e-06\\
16	3.11228721233702e-06\\
17	3.11228721233702e-06\\
18	3.11228721233702e-06\\
19	3.11228721233702e-06\\
20	3.11228721233702e-06\\
21	3.11228721233702e-06\\
22	3.11228721233702e-06\\
23	3.11228721233702e-06\\
24	3.11228721233702e-06\\
25	3.11228721233702e-06\\
26	4.03351450714213e-06\\
27	4.03351450714213e-06\\
28	4.03351450714213e-06\\
29	4.03351450714213e-06\\
30	4.03351450714213e-06\\
31	4.03351450714213e-06\\
32	4.03351450714213e-06\\
33	4.03351450714213e-06\\
34	4.03351450714213e-06\\
35	4.03351450714213e-06\\
36	4.03351450714213e-06\\
37	4.03351450714213e-06\\
38	4.03351450714213e-06\\
39	4.03351450714213e-06\\
40	4.03351450714213e-06\\
41	4.03351450714213e-06\\
42	4.03351450714213e-06\\
43	4.03351450714213e-06\\
44	4.03351450714213e-06\\
45	4.03351450714213e-06\\
46	4.03351450714213e-06\\
47	4.03351450714213e-06\\
48	4.03351450714213e-06\\
49	4.03351450714213e-06\\
50	9.21227294805105e-07\\
51	9.21227294805105e-07\\
52	9.21227294805105e-07\\
53	9.21227294805105e-07\\
54	9.21227294805105e-07\\
55	9.21227294805105e-07\\
56	9.21227294805105e-07\\
57	9.21227294805105e-07\\
58	9.21227294805105e-07\\
59	9.21227294805105e-07\\
60	9.21227294805105e-07\\
61	9.21227294805105e-07\\
62	9.21227294805105e-07\\
63	9.21227294805105e-07\\
64	9.21227294805105e-07\\
65	9.21227294805105e-07\\
66	9.21227294805105e-07\\
67	9.21227294805105e-07\\
68	9.21227294805105e-07\\
69	9.21227294805105e-07\\
70	9.21227294805105e-07\\
71	9.21227294805105e-07\\
72	9.21227294805105e-07\\
73	9.21227294805105e-07\\
74	9.21227294805105e-07\\
75	0\\
76	0\\
77	0\\
78	0\\
79	0\\
80	0\\
81	0\\
82	0\\
83	0\\
84	0\\
85	0\\
86	0\\
87	0\\
88	0\\
89	0\\
90	0\\
91	0\\
92	0\\
93	0\\
94	0\\
95	0\\
96	0\\
97	0\\
98	0\\
99	0\\
100	0\\
};
\addplot [color=blue,line width=1.2pt,only marks,mark=asterisk,mark options={solid},forget plot]
  table[row sep=crcr]{31	4.03351450714213e-06\\
};
\addplot [color=green,line width=1.2pt,only marks,mark=+,mark options={solid},forget plot]
  table[row sep=crcr]{31	4.03351450714213e-06\\
};
\addplot [color=red,line width=1.2pt,only marks,mark=square,mark options={solid},forget plot]
  table[row sep=crcr]{38	4.03351450714213e-06\\
};
\addplot [color=mycolor1,line width=1.2pt,only marks,mark=triangle,mark options={solid,rotate=180},forget plot]
  table[row sep=crcr]{26	4.03351450714213e-06\\
};
\addplot [color=mycolor2,line width=1.2pt,only marks,mark=triangle,mark options={solid},forget plot]
  table[row sep=crcr]{49	4.03351450714213e-06\\
};
\end{axis}

\begin{axis}[%
width=\figurewidth,
height=\figureheight,
scale only axis,
xmin=0,
xmax=100,
ymin=0,
ymax=6e-06,
name=plot2,
at=(plot1.right of south east),
anchor=left of south west,
title={Oavgmin}
]
\addplot [color=black,solid,line width=1.2pt,forget plot]
  table[row sep=crcr]{0	0\\
1	3.11228721233702e-06\\
2	3.11228721233702e-06\\
3	3.11228721233702e-06\\
4	3.11228721233702e-06\\
5	3.11228721233702e-06\\
6	3.11228721233702e-06\\
7	3.11228721233702e-06\\
8	3.11228721233702e-06\\
9	3.11228721233702e-06\\
10	3.11228721233702e-06\\
11	3.11228721233702e-06\\
12	3.11228721233702e-06\\
13	3.11228721233702e-06\\
14	3.11228721233702e-06\\
15	3.11228721233702e-06\\
16	3.11228721233702e-06\\
17	3.11228721233702e-06\\
18	3.11228721233702e-06\\
19	3.11228721233702e-06\\
20	3.11228721233702e-06\\
21	3.11228721233702e-06\\
22	3.11228721233702e-06\\
23	3.11228721233702e-06\\
24	3.11228721233702e-06\\
25	3.11228721233702e-06\\
26	4.03351450714213e-06\\
27	4.03351450714213e-06\\
28	4.03351450714213e-06\\
29	4.03351450714213e-06\\
30	4.03351450714213e-06\\
31	4.03351450714213e-06\\
32	4.03351450714213e-06\\
33	4.03351450714213e-06\\
34	4.03351450714213e-06\\
35	4.03351450714213e-06\\
36	4.03351450714213e-06\\
37	4.03351450714213e-06\\
38	4.03351450714213e-06\\
39	4.03351450714213e-06\\
40	4.03351450714213e-06\\
41	4.03351450714213e-06\\
42	4.03351450714213e-06\\
43	4.03351450714213e-06\\
44	4.03351450714213e-06\\
45	4.03351450714213e-06\\
46	4.03351450714213e-06\\
47	4.03351450714213e-06\\
48	4.03351450714213e-06\\
49	4.03351450714213e-06\\
50	9.21227294805105e-07\\
51	9.21227294805105e-07\\
52	9.21227294805105e-07\\
53	9.21227294805105e-07\\
54	9.21227294805105e-07\\
55	9.21227294805105e-07\\
56	9.21227294805105e-07\\
57	9.21227294805105e-07\\
58	9.21227294805105e-07\\
59	9.21227294805105e-07\\
60	9.21227294805105e-07\\
61	9.21227294805105e-07\\
62	9.21227294805105e-07\\
63	9.21227294805105e-07\\
64	9.21227294805105e-07\\
65	9.21227294805105e-07\\
66	9.21227294805105e-07\\
67	9.21227294805105e-07\\
68	9.21227294805105e-07\\
69	9.21227294805105e-07\\
70	9.21227294805105e-07\\
71	9.21227294805105e-07\\
72	9.21227294805105e-07\\
73	9.21227294805105e-07\\
74	9.21227294805105e-07\\
75	0\\
76	0\\
77	0\\
78	0\\
79	0\\
80	0\\
81	0\\
82	0\\
83	0\\
84	0\\
85	0\\
86	0\\
87	0\\
88	0\\
89	0\\
90	0\\
91	0\\
92	0\\
93	0\\
94	0\\
95	0\\
96	0\\
97	0\\
98	0\\
99	0\\
100	0\\
};
\addplot [color=blue,line width=1.2pt,only marks,mark=asterisk,mark options={solid},forget plot]
  table[row sep=crcr]{31	4.03351450714213e-06\\
};
\addplot [color=green,line width=1.2pt,only marks,mark=+,mark options={solid},forget plot]
  table[row sep=crcr]{31	4.03351450714213e-06\\
};
\addplot [color=red,line width=1.2pt,only marks,mark=square,mark options={solid},forget plot]
  table[row sep=crcr]{38	4.03351450714213e-06\\
};
\addplot [color=mycolor1,line width=1.2pt,only marks,mark=triangle,mark options={solid,rotate=180},forget plot]
  table[row sep=crcr]{26	4.03351450714213e-06\\
};
\addplot [color=mycolor2,line width=1.2pt,only marks,mark=triangle,mark options={solid},forget plot]
  table[row sep=crcr]{49	4.03351450714213e-06\\
};
\end{axis}

\begin{axis}[%
width=\figurewidth,
height=\figureheight,
scale only axis,
xmin=0,
xmax=100,
ymin=0,
ymax=6e-06,
at=(plot2.below south west),
anchor=above north west,
title={Osqrt}
]
\addplot [color=black,solid,line width=1.2pt,forget plot]
  table[row sep=crcr]{0	0\\
1	3.90780141787935e-06\\
2	3.90780141787935e-06\\
3	3.90780141787935e-06\\
4	3.90780141787935e-06\\
5	3.90780141787935e-06\\
6	3.90780141787935e-06\\
7	3.90780141787935e-06\\
8	3.90780141787935e-06\\
9	3.90780141787935e-06\\
10	3.90780141787935e-06\\
11	3.90780141787935e-06\\
12	3.90780141787935e-06\\
13	3.90780141787935e-06\\
14	3.90780141787935e-06\\
15	3.90780141787935e-06\\
16	3.90780141787935e-06\\
17	3.90780141787935e-06\\
18	3.90780141787935e-06\\
19	3.90780141787935e-06\\
20	3.90780141787935e-06\\
21	3.90780141787935e-06\\
22	3.90780141787935e-06\\
23	3.90780141787935e-06\\
24	3.90780141787935e-06\\
25	3.90780141787935e-06\\
26	5.13627755976046e-06\\
27	5.13627755976046e-06\\
28	5.13627755976046e-06\\
29	5.13627755976046e-06\\
30	5.13627755976046e-06\\
31	5.13627755976046e-06\\
32	5.13627755976046e-06\\
33	5.13627755976046e-06\\
34	5.13627755976046e-06\\
35	5.13627755976046e-06\\
36	5.13627755976046e-06\\
37	5.13627755976046e-06\\
38	5.13627755976046e-06\\
39	5.13627755976046e-06\\
40	5.13627755976046e-06\\
41	5.13627755976046e-06\\
42	5.13627755976046e-06\\
43	5.13627755976046e-06\\
44	5.13627755976046e-06\\
45	5.13627755976046e-06\\
46	5.13627755976046e-06\\
47	5.13627755976046e-06\\
48	5.13627755976046e-06\\
49	5.13627755976046e-06\\
50	1.2284761418811e-06\\
51	1.2284761418811e-06\\
52	1.2284761418811e-06\\
53	1.2284761418811e-06\\
54	1.2284761418811e-06\\
55	1.2284761418811e-06\\
56	1.2284761418811e-06\\
57	1.2284761418811e-06\\
58	1.2284761418811e-06\\
59	1.2284761418811e-06\\
60	1.2284761418811e-06\\
61	1.2284761418811e-06\\
62	1.2284761418811e-06\\
63	1.2284761418811e-06\\
64	1.2284761418811e-06\\
65	1.2284761418811e-06\\
66	1.2284761418811e-06\\
67	1.2284761418811e-06\\
68	1.2284761418811e-06\\
69	1.2284761418811e-06\\
70	1.2284761418811e-06\\
71	1.2284761418811e-06\\
72	1.2284761418811e-06\\
73	1.2284761418811e-06\\
74	1.2284761418811e-06\\
75	0\\
76	0\\
77	0\\
78	0\\
79	0\\
80	0\\
81	0\\
82	0\\
83	0\\
84	0\\
85	0\\
86	0\\
87	0\\
88	0\\
89	0\\
90	0\\
91	0\\
92	0\\
93	0\\
94	0\\
95	0\\
96	0\\
97	0\\
98	0\\
99	0\\
100	0\\
};
\addplot [color=blue,line width=1.2pt,only marks,mark=asterisk,mark options={solid},forget plot]
  table[row sep=crcr]{31	5.13627755976046e-06\\
};
\addplot [color=green,line width=1.2pt,only marks,mark=+,mark options={solid},forget plot]
  table[row sep=crcr]{31	5.13627755976046e-06\\
};
\addplot [color=red,line width=1.2pt,only marks,mark=square,mark options={solid},forget plot]
  table[row sep=crcr]{38	5.13627755976046e-06\\
};
\addplot [color=mycolor1,line width=1.2pt,only marks,mark=triangle,mark options={solid,rotate=180},forget plot]
  table[row sep=crcr]{26	5.13627755976046e-06\\
};
\addplot [color=mycolor2,line width=1.2pt,only marks,mark=triangle,mark options={solid},forget plot]
  table[row sep=crcr]{49	5.13627755976046e-06\\
};
\end{axis}

\begin{axis}[%
width=\figurewidth,
height=\figureheight,
scale only axis,
xmin=0,
xmax=100,
ymin=0,
ymax=1e-05,
at=(plot3.below south west),
anchor=above north west,
title={Osin}
]
\addplot [color=black,solid,line width=1.2pt,forget plot]
  table[row sep=crcr]{0	0\\
1	4.7811734098003e-06\\
2	4.7811734098003e-06\\
3	4.7811734098003e-06\\
4	4.7811734098003e-06\\
5	4.7811734098003e-06\\
6	4.7811734098003e-06\\
7	4.7811734098003e-06\\
8	4.7811734098003e-06\\
9	4.7811734098003e-06\\
10	4.7811734098003e-06\\
11	4.7811734098003e-06\\
12	4.7811734098003e-06\\
13	4.7811734098003e-06\\
14	4.7811734098003e-06\\
15	4.7811734098003e-06\\
16	4.7811734098003e-06\\
17	4.7811734098003e-06\\
18	4.7811734098003e-06\\
19	4.7811734098003e-06\\
20	4.7811734098003e-06\\
21	4.7811734098003e-06\\
22	4.7811734098003e-06\\
23	4.7811734098003e-06\\
24	4.7811734098003e-06\\
25	4.7811734098003e-06\\
26	6.35418570028017e-06\\
27	6.35418570028017e-06\\
28	6.35418570028017e-06\\
29	6.35418570028017e-06\\
30	6.35418570028017e-06\\
31	6.35418570028017e-06\\
32	6.35418570028017e-06\\
33	6.35418570028017e-06\\
34	6.35418570028017e-06\\
35	6.35418570028017e-06\\
36	6.35418570028017e-06\\
37	6.35418570028017e-06\\
38	6.35418570028017e-06\\
39	6.35418570028017e-06\\
40	6.35418570028017e-06\\
41	6.35418570028017e-06\\
42	6.35418570028017e-06\\
43	6.35418570028017e-06\\
44	6.35418570028017e-06\\
45	6.35418570028017e-06\\
46	6.35418570028017e-06\\
47	6.35418570028017e-06\\
48	6.35418570028017e-06\\
49	6.35418570028017e-06\\
50	1.57301229047987e-06\\
51	1.57301229047987e-06\\
52	1.57301229047987e-06\\
53	1.57301229047987e-06\\
54	1.57301229047987e-06\\
55	1.57301229047987e-06\\
56	1.57301229047987e-06\\
57	1.57301229047987e-06\\
58	1.57301229047987e-06\\
59	1.57301229047987e-06\\
60	1.57301229047987e-06\\
61	1.57301229047987e-06\\
62	1.57301229047987e-06\\
63	1.57301229047987e-06\\
64	1.57301229047987e-06\\
65	1.57301229047987e-06\\
66	1.57301229047987e-06\\
67	1.57301229047987e-06\\
68	1.57301229047987e-06\\
69	1.57301229047987e-06\\
70	1.57301229047987e-06\\
71	1.57301229047987e-06\\
72	1.57301229047987e-06\\
73	1.57301229047987e-06\\
74	1.57301229047987e-06\\
75	0\\
76	0\\
77	0\\
78	0\\
79	0\\
80	0\\
81	0\\
82	0\\
83	0\\
84	0\\
85	0\\
86	0\\
87	0\\
88	0\\
89	0\\
90	0\\
91	0\\
92	0\\
93	0\\
94	0\\
95	0\\
96	0\\
97	0\\
98	0\\
99	0\\
100	0\\
};
\addplot [color=blue,line width=1.2pt,only marks,mark=asterisk,mark options={solid},forget plot]
  table[row sep=crcr]{31	6.35418570028017e-06\\
};
\addplot [color=green,line width=1.2pt,only marks,mark=+,mark options={solid},forget plot]
  table[row sep=crcr]{31	6.35418570028017e-06\\
};
\addplot [color=red,line width=1.2pt,only marks,mark=square,mark options={solid},forget plot]
  table[row sep=crcr]{38	6.35418570028017e-06\\
};
\addplot [color=mycolor1,line width=1.2pt,only marks,mark=triangle,mark options={solid,rotate=180},forget plot]
  table[row sep=crcr]{26	6.35418570028017e-06\\
};
\addplot [color=mycolor2,line width=1.2pt,only marks,mark=triangle,mark options={solid},forget plot]
  table[row sep=crcr]{49	6.35418570028017e-06\\
};
\end{axis}
\end{tikzpicture}%
\end{subfigure}
\caption{T-norma: Seno ($T_{sin}$)}
\label{fig:fire-detection-interpolation-comparison-tnorm-sinmean}
\end{figure}
\begin{figure}[h!]
\centering
\begin{subfigure}[b]{\textwidth}
	\textbf{T-norma: Einstein ($T_{einstein}$)} \\
	\setlength\figureheight{2cm}
	\setlength\figurewidth{7cm}
	% This file was created by matlab2tikz v0.4.7 running on MATLAB 8.0.
% Copyright (c) 2008--2014, Nico Schlömer <nico.schloemer@gmail.com>
% All rights reserved.
% Minimal pgfplots version: 1.3
% 
%
% defining custom colors
\definecolor{mycolor1}{rgb}{0.00000,1.00000,1.00000}%
\definecolor{mycolor2}{rgb}{1.00000,0.00000,1.00000}%
%
\begin{tikzpicture}

\begin{axis}[%
width=\figurewidth,
height=\figureheight,
scale only axis,
xmin=0,
xmax=100,
ymin=0,
ymax=0.003,
name=plot3,
title={Omaxmin}
]
\addplot [color=black,solid,line width=1.2pt,forget plot]
  table[row sep=crcr]{0	0\\
1	0.000493827160493827\\
2	0.000987654320987654\\
3	0.00139917695473251\\
4	0.0017283950617284\\
5	0.00205761316872428\\
6	0.00205761316872428\\
7	0.00205761316872428\\
8	0.00205761316872428\\
9	0.00205761316872428\\
10	0.00205761316872428\\
11	0.00205761316872428\\
12	0.00205761316872428\\
13	0.00205761316872428\\
14	0.00205761316872428\\
15	0.00205761316872428\\
16	0.00205761316872428\\
17	0.00205761316872428\\
18	0.00205761316872428\\
19	0.00205761316872428\\
20	0.00205761316872428\\
21	0.00205761316872428\\
22	0.00205761316872428\\
23	0.00205761316872428\\
24	0.00205761316872428\\
25	0.00205761316872428\\
26	0.00222222222222222\\
27	0.00238683127572016\\
28	0.00246913580246914\\
29	0.00246913580246914\\
30	0.00246913580246914\\
31	0.00246913580246914\\
32	0.00246913580246914\\
33	0.00246913580246914\\
34	0.00246913580246914\\
35	0.00246913580246914\\
36	0.00246913580246914\\
37	0.00246913580246914\\
38	0.00246913580246914\\
39	0.00246913580246914\\
40	0.00246913580246914\\
41	0.00246913580246914\\
42	0.00246913580246914\\
43	0.00246913580246914\\
44	0.00246913580246914\\
45	0.00246913580246914\\
46	0.00213991769547325\\
47	0.00181069958847737\\
48	0.00139917695473251\\
49	0.000905349794238684\\
50	0.000411522633744856\\
51	0.000411522633744856\\
52	0.000411522633744856\\
53	0.000411522633744856\\
54	0.000411522633744856\\
55	0.000411522633744856\\
56	0.000411522633744856\\
57	0.000411522633744856\\
58	0.000411522633744856\\
59	0.000411522633744856\\
60	0.000411522633744856\\
61	0.000411522633744856\\
62	0.000411522633744856\\
63	0.000411522633744856\\
64	0.000411522633744856\\
65	0.000411522633744856\\
66	0.000411522633744856\\
67	0.000411522633744856\\
68	0.000411522633744856\\
69	0.000411522633744856\\
70	0.000411522633744856\\
71	0.000411522633744856\\
72	0.000411522633744856\\
73	0.000329218106995885\\
74	0.000164609053497943\\
75	0\\
76	0\\
77	0\\
78	0\\
79	0\\
80	0\\
81	0\\
82	0\\
83	0\\
84	0\\
85	0\\
86	0\\
87	0\\
88	0\\
89	0\\
90	0\\
91	0\\
92	0\\
93	0\\
94	0\\
95	0\\
96	0\\
97	0\\
98	0\\
99	0\\
100	0\\
};
\addplot [color=blue,line width=1.2pt,only marks,mark=asterisk,mark options={solid},forget plot]
  table[row sep=crcr]{29	0.00246913580246914\\
};
\addplot [color=green,line width=1.2pt,only marks,mark=+,mark options={solid},forget plot]
  table[row sep=crcr]{29	0.00246913580246914\\
};
\addplot [color=red,line width=1.2pt,only marks,mark=square,mark options={solid},forget plot]
  table[row sep=crcr]{36	0.00246913580246914\\
};
\addplot [color=mycolor1,line width=1.2pt,only marks,mark=triangle,mark options={solid,rotate=180},forget plot]
  table[row sep=crcr]{28	0.00246913580246914\\
};
\addplot [color=mycolor2,line width=1.2pt,only marks,mark=triangle,mark options={solid},forget plot]
  table[row sep=crcr]{44	0.00246913580246914\\
};
\end{axis}

\begin{axis}[%
width=\figurewidth,
height=\figureheight,
scale only axis,
xmin=0,
xmax=100,
ymin=0,
ymax=1.5e-14,
name=plot1,
at=(plot3.above north west),
anchor=below south west,
title={Opi}
]
\addplot [color=black,solid,line width=1.2pt,forget plot]
  table[row sep=crcr]{0	0\\
1	1.04399369740669e-14\\
2	1.04399369740669e-14\\
3	1.04399369740669e-14\\
4	1.04399369740669e-14\\
5	1.04399369740669e-14\\
6	1.04399369740669e-14\\
7	1.04399369740669e-14\\
8	1.04399369740669e-14\\
9	1.04399369740669e-14\\
10	1.04399369740669e-14\\
11	1.04399369740669e-14\\
12	1.04399369740669e-14\\
13	1.04399369740669e-14\\
14	1.04399369740669e-14\\
15	1.04399369740669e-14\\
16	1.04399369740669e-14\\
17	1.04399369740669e-14\\
18	1.04399369740669e-14\\
19	1.04399369740669e-14\\
20	1.04399369740669e-14\\
21	1.04399369740669e-14\\
22	1.04399369740669e-14\\
23	1.04399369740669e-14\\
24	1.04399369740669e-14\\
25	1.04399369740669e-14\\
26	1.25258816414174e-14\\
27	1.25258816414174e-14\\
28	1.25258816414174e-14\\
29	1.25258816414174e-14\\
30	1.25258816414174e-14\\
31	1.25258816414174e-14\\
32	1.25258816414174e-14\\
33	1.25258816414174e-14\\
34	1.25258816414174e-14\\
35	1.25258816414174e-14\\
36	1.25258816414174e-14\\
37	1.25258816414174e-14\\
38	1.25258816414174e-14\\
39	1.25258816414174e-14\\
40	1.25258816414174e-14\\
41	1.25258816414174e-14\\
42	1.25258816414174e-14\\
43	1.25258816414174e-14\\
44	1.25258816414174e-14\\
45	1.25258816414174e-14\\
46	1.25258816414174e-14\\
47	1.25258816414174e-14\\
48	1.25258816414174e-14\\
49	1.25258816414174e-14\\
50	2.08594466735055e-15\\
51	2.08594466735055e-15\\
52	2.08594466735055e-15\\
53	2.08594466735055e-15\\
54	2.08594466735055e-15\\
55	2.08594466735055e-15\\
56	2.08594466735055e-15\\
57	2.08594466735055e-15\\
58	2.08594466735055e-15\\
59	2.08594466735055e-15\\
60	2.08594466735055e-15\\
61	2.08594466735055e-15\\
62	2.08594466735055e-15\\
63	2.08594466735055e-15\\
64	2.08594466735055e-15\\
65	2.08594466735055e-15\\
66	2.08594466735055e-15\\
67	2.08594466735055e-15\\
68	2.08594466735055e-15\\
69	2.08594466735055e-15\\
70	2.08594466735055e-15\\
71	2.08594466735055e-15\\
72	2.08594466735055e-15\\
73	2.08594466735055e-15\\
74	2.08594466735055e-15\\
75	0\\
76	0\\
77	0\\
78	0\\
79	0\\
80	0\\
81	0\\
82	0\\
83	0\\
84	0\\
85	0\\
86	0\\
87	0\\
88	0\\
89	0\\
90	0\\
91	0\\
92	0\\
93	0\\
94	0\\
95	0\\
96	0\\
97	0\\
98	0\\
99	0\\
100	0\\
};
\addplot [color=blue,line width=1.2pt,only marks,mark=asterisk,mark options={solid},forget plot]
  table[row sep=crcr]{29	1.25258816414174e-14\\
};
\addplot [color=green,line width=1.2pt,only marks,mark=+,mark options={solid},forget plot]
  table[row sep=crcr]{29	1.25258816414174e-14\\
};
\addplot [color=red,line width=1.2pt,only marks,mark=square,mark options={solid},forget plot]
  table[row sep=crcr]{38	1.25258816414174e-14\\
};
\addplot [color=mycolor1,line width=1.2pt,only marks,mark=triangle,mark options={solid,rotate=180},forget plot]
  table[row sep=crcr]{26	1.25258816414174e-14\\
};
\addplot [color=mycolor2,line width=1.2pt,only marks,mark=triangle,mark options={solid},forget plot]
  table[row sep=crcr]{49	1.25258816414174e-14\\
};
\end{axis}

\begin{axis}[%
width=\figurewidth,
height=\figureheight,
scale only axis,
xmin=0,
xmax=100,
ymin=0,
ymax=1.5e-14,
name=plot2,
at=(plot1.right of south east),
anchor=left of south west,
title={Oavgmin}
]
\addplot [color=black,solid,line width=1.2pt,forget plot]
  table[row sep=crcr]{0	0\\
1	1.04399369740669e-14\\
2	1.04399369740669e-14\\
3	1.04399369740669e-14\\
4	1.04399369740669e-14\\
5	1.04399369740669e-14\\
6	1.04399369740669e-14\\
7	1.04399369740669e-14\\
8	1.04399369740669e-14\\
9	1.04399369740669e-14\\
10	1.04399369740669e-14\\
11	1.04399369740669e-14\\
12	1.04399369740669e-14\\
13	1.04399369740669e-14\\
14	1.04399369740669e-14\\
15	1.04399369740669e-14\\
16	1.04399369740669e-14\\
17	1.04399369740669e-14\\
18	1.04399369740669e-14\\
19	1.04399369740669e-14\\
20	1.04399369740669e-14\\
21	1.04399369740669e-14\\
22	1.04399369740669e-14\\
23	1.04399369740669e-14\\
24	1.04399369740669e-14\\
25	1.04399369740669e-14\\
26	1.25258816414174e-14\\
27	1.25258816414174e-14\\
28	1.25258816414174e-14\\
29	1.25258816414174e-14\\
30	1.25258816414174e-14\\
31	1.25258816414174e-14\\
32	1.25258816414174e-14\\
33	1.25258816414174e-14\\
34	1.25258816414174e-14\\
35	1.25258816414174e-14\\
36	1.25258816414174e-14\\
37	1.25258816414174e-14\\
38	1.25258816414174e-14\\
39	1.25258816414174e-14\\
40	1.25258816414174e-14\\
41	1.25258816414174e-14\\
42	1.25258816414174e-14\\
43	1.25258816414174e-14\\
44	1.25258816414174e-14\\
45	1.25258816414174e-14\\
46	1.25258816414174e-14\\
47	1.25258816414174e-14\\
48	1.25258816414174e-14\\
49	1.25258816414174e-14\\
50	2.08594466735055e-15\\
51	2.08594466735055e-15\\
52	2.08594466735055e-15\\
53	2.08594466735055e-15\\
54	2.08594466735055e-15\\
55	2.08594466735055e-15\\
56	2.08594466735055e-15\\
57	2.08594466735055e-15\\
58	2.08594466735055e-15\\
59	2.08594466735055e-15\\
60	2.08594466735055e-15\\
61	2.08594466735055e-15\\
62	2.08594466735055e-15\\
63	2.08594466735055e-15\\
64	2.08594466735055e-15\\
65	2.08594466735055e-15\\
66	2.08594466735055e-15\\
67	2.08594466735055e-15\\
68	2.08594466735055e-15\\
69	2.08594466735055e-15\\
70	2.08594466735055e-15\\
71	2.08594466735055e-15\\
72	2.08594466735055e-15\\
73	2.08594466735055e-15\\
74	2.08594466735055e-15\\
75	0\\
76	0\\
77	0\\
78	0\\
79	0\\
80	0\\
81	0\\
82	0\\
83	0\\
84	0\\
85	0\\
86	0\\
87	0\\
88	0\\
89	0\\
90	0\\
91	0\\
92	0\\
93	0\\
94	0\\
95	0\\
96	0\\
97	0\\
98	0\\
99	0\\
100	0\\
};
\addplot [color=blue,line width=1.2pt,only marks,mark=asterisk,mark options={solid},forget plot]
  table[row sep=crcr]{29	1.25258816414174e-14\\
};
\addplot [color=green,line width=1.2pt,only marks,mark=+,mark options={solid},forget plot]
  table[row sep=crcr]{29	1.25258816414174e-14\\
};
\addplot [color=red,line width=1.2pt,only marks,mark=square,mark options={solid},forget plot]
  table[row sep=crcr]{38	1.25258816414174e-14\\
};
\addplot [color=mycolor1,line width=1.2pt,only marks,mark=triangle,mark options={solid,rotate=180},forget plot]
  table[row sep=crcr]{26	1.25258816414174e-14\\
};
\addplot [color=mycolor2,line width=1.2pt,only marks,mark=triangle,mark options={solid},forget plot]
  table[row sep=crcr]{49	1.25258816414174e-14\\
};
\end{axis}

\begin{axis}[%
width=\figurewidth,
height=\figureheight,
scale only axis,
xmin=0,
xmax=100,
ymin=0,
ymax=4e-14,
at=(plot2.below south west),
anchor=above north west,
title={Osqrt}
]
\addplot [color=black,solid,line width=1.2pt,forget plot]
  table[row sep=crcr]{0	0\\
1	2.53274722703583e-14\\
2	2.53274722703583e-14\\
3	2.53274722703583e-14\\
4	2.53274722703583e-14\\
5	2.53274722703583e-14\\
6	2.53274722703583e-14\\
7	2.53274722703583e-14\\
8	2.53274722703583e-14\\
9	2.53274722703583e-14\\
10	2.53274722703583e-14\\
11	2.53274722703583e-14\\
12	2.53274722703583e-14\\
13	2.53274722703583e-14\\
14	2.53274722703583e-14\\
15	2.53274722703583e-14\\
16	2.53274722703583e-14\\
17	2.53274722703583e-14\\
18	2.53274722703583e-14\\
19	2.53274722703583e-14\\
20	2.53274722703583e-14\\
21	2.53274722703583e-14\\
22	2.53274722703583e-14\\
23	2.53274722703583e-14\\
24	2.53274722703583e-14\\
25	2.53274722703583e-14\\
26	3.19381532973811e-14\\
27	3.19381532973811e-14\\
28	3.19381532973811e-14\\
29	3.19381532973811e-14\\
30	3.19381532973811e-14\\
31	3.19381532973811e-14\\
32	3.19381532973811e-14\\
33	3.19381532973811e-14\\
34	3.19381532973811e-14\\
35	3.19381532973811e-14\\
36	3.19381532973811e-14\\
37	3.19381532973811e-14\\
38	3.19381532973811e-14\\
39	3.19381532973811e-14\\
40	3.19381532973811e-14\\
41	3.19381532973811e-14\\
42	3.19381532973811e-14\\
43	3.19381532973811e-14\\
44	3.19381532973811e-14\\
45	3.19381532973811e-14\\
46	3.19381532973811e-14\\
47	3.19381532973811e-14\\
48	3.19381532973811e-14\\
49	3.19381532973811e-14\\
50	6.61068102702279e-15\\
51	6.61068102702279e-15\\
52	6.61068102702279e-15\\
53	6.61068102702279e-15\\
54	6.61068102702279e-15\\
55	6.61068102702279e-15\\
56	6.61068102702279e-15\\
57	6.61068102702279e-15\\
58	6.61068102702279e-15\\
59	6.61068102702279e-15\\
60	6.61068102702279e-15\\
61	6.61068102702279e-15\\
62	6.61068102702279e-15\\
63	6.61068102702279e-15\\
64	6.61068102702279e-15\\
65	6.61068102702279e-15\\
66	6.61068102702279e-15\\
67	6.61068102702279e-15\\
68	6.61068102702279e-15\\
69	6.61068102702279e-15\\
70	6.61068102702279e-15\\
71	6.61068102702279e-15\\
72	6.61068102702279e-15\\
73	6.61068102702279e-15\\
74	6.61068102702279e-15\\
75	0\\
76	0\\
77	0\\
78	0\\
79	0\\
80	0\\
81	0\\
82	0\\
83	0\\
84	0\\
85	0\\
86	0\\
87	0\\
88	0\\
89	0\\
90	0\\
91	0\\
92	0\\
93	0\\
94	0\\
95	0\\
96	0\\
97	0\\
98	0\\
99	0\\
100	0\\
};
\addplot [color=blue,line width=1.2pt,only marks,mark=asterisk,mark options={solid},forget plot]
  table[row sep=crcr]{30	3.19381532973811e-14\\
};
\addplot [color=green,line width=1.2pt,only marks,mark=+,mark options={solid},forget plot]
  table[row sep=crcr]{30	3.19381532973811e-14\\
};
\addplot [color=red,line width=1.2pt,only marks,mark=square,mark options={solid},forget plot]
  table[row sep=crcr]{38	3.19381532973811e-14\\
};
\addplot [color=mycolor1,line width=1.2pt,only marks,mark=triangle,mark options={solid,rotate=180},forget plot]
  table[row sep=crcr]{26	3.19381532973811e-14\\
};
\addplot [color=mycolor2,line width=1.2pt,only marks,mark=triangle,mark options={solid},forget plot]
  table[row sep=crcr]{49	3.19381532973811e-14\\
};
\end{axis}

\begin{axis}[%
width=\figurewidth,
height=\figureheight,
scale only axis,
xmin=0,
xmax=100,
ymin=0,
ymax=1e-13,
at=(plot3.below south west),
anchor=above north west,
title={Osin}
]
\addplot [color=black,solid,line width=1.2pt,forget plot]
  table[row sep=crcr]{0	0\\
1	5.64002943904551e-14\\
2	5.64002943904551e-14\\
3	5.64002943904551e-14\\
4	5.64002943904551e-14\\
5	5.64002943904551e-14\\
6	5.64002943904551e-14\\
7	5.64002943904551e-14\\
8	5.64002943904551e-14\\
9	5.64002943904551e-14\\
10	5.64002943904551e-14\\
11	5.64002943904551e-14\\
12	5.64002943904551e-14\\
13	5.64002943904551e-14\\
14	5.64002943904551e-14\\
15	5.64002943904551e-14\\
16	5.64002943904551e-14\\
17	5.64002943904551e-14\\
18	5.64002943904551e-14\\
19	5.64002943904551e-14\\
20	5.64002943904551e-14\\
21	5.64002943904551e-14\\
22	5.64002943904551e-14\\
23	5.64002943904551e-14\\
24	5.64002943904551e-14\\
25	5.64002943904551e-14\\
26	7.42286187523654e-14\\
27	7.42286187523654e-14\\
28	7.42286187523654e-14\\
29	7.42286187523654e-14\\
30	7.42286187523654e-14\\
31	7.42286187523654e-14\\
32	7.42286187523654e-14\\
33	7.42286187523654e-14\\
34	7.42286187523654e-14\\
35	7.42286187523654e-14\\
36	7.42286187523654e-14\\
37	7.42286187523654e-14\\
38	7.42286187523654e-14\\
39	7.42286187523654e-14\\
40	7.42286187523654e-14\\
41	7.42286187523654e-14\\
42	7.42286187523654e-14\\
43	7.42286187523654e-14\\
44	7.42286187523654e-14\\
45	7.42286187523654e-14\\
46	7.42286187523654e-14\\
47	7.42286187523654e-14\\
48	7.42286187523654e-14\\
49	7.42286187523654e-14\\
50	1.78283243619103e-14\\
51	1.78283243619103e-14\\
52	1.78283243619103e-14\\
53	1.78283243619103e-14\\
54	1.78283243619103e-14\\
55	1.78283243619103e-14\\
56	1.78283243619103e-14\\
57	1.78283243619103e-14\\
58	1.78283243619103e-14\\
59	1.78283243619103e-14\\
60	1.78283243619103e-14\\
61	1.78283243619103e-14\\
62	1.78283243619103e-14\\
63	1.78283243619103e-14\\
64	1.78283243619103e-14\\
65	1.78283243619103e-14\\
66	1.78283243619103e-14\\
67	1.78283243619103e-14\\
68	1.78283243619103e-14\\
69	1.78283243619103e-14\\
70	1.78283243619103e-14\\
71	1.78283243619103e-14\\
72	1.78283243619103e-14\\
73	1.78283243619103e-14\\
74	1.78283243619103e-14\\
75	0\\
76	0\\
77	0\\
78	0\\
79	0\\
80	0\\
81	0\\
82	0\\
83	0\\
84	0\\
85	0\\
86	0\\
87	0\\
88	0\\
89	0\\
90	0\\
91	0\\
92	0\\
93	0\\
94	0\\
95	0\\
96	0\\
97	0\\
98	0\\
99	0\\
100	0\\
};
\addplot [color=blue,line width=1.2pt,only marks,mark=asterisk,mark options={solid},forget plot]
  table[row sep=crcr]{31	7.42286187523654e-14\\
};
\addplot [color=green,line width=1.2pt,only marks,mark=+,mark options={solid},forget plot]
  table[row sep=crcr]{31	7.42286187523654e-14\\
};
\addplot [color=red,line width=1.2pt,only marks,mark=square,mark options={solid},forget plot]
  table[row sep=crcr]{38	7.42286187523654e-14\\
};
\addplot [color=mycolor1,line width=1.2pt,only marks,mark=triangle,mark options={solid,rotate=180},forget plot]
  table[row sep=crcr]{26	7.42286187523654e-14\\
};
\addplot [color=mycolor2,line width=1.2pt,only marks,mark=triangle,mark options={solid},forget plot]
  table[row sep=crcr]{49	7.42286187523654e-14\\
};
\end{axis}
\end{tikzpicture}%
\end{subfigure}
\caption{T-norma: Einstein ($T_{einstein}$)}
\label{fig:fire-detection-interpolation-comparison-tnorm-einstein}
\end{figure}

Como se puede ver en la tabla, los resultados finales (desdifusificados) al aplicar diferentes t-normas y operadores de agregación no varían significativamente (al menos con estos datos de entrada). En las gráficas se puede apreciar que los conjuntos de salida sí presentan algunas diferencias entre sí, especialmente cuando se utiliza el índice de solapamiento del máximo de mínimos ($O_Z$).

También existen diferencias notables en la escala de la función de pertenencia de los conjuntos de salida. El ejemplo más claro se puede ver en las gráficas obtenidas con la t-norma producto. Los valores son del orden de $10^{-14}$, mientras que al utilizar otras t-normas como el mínimo se obtienen resultados ``más grandes'' (de órdenes entre $10^{-3}$ y $10^{-5}$). Esto se debe a que los índices de solapamiento entre las entradas y los antecedentes son valores en el rango [0,1] y por tanto al realizar el producto de estos valores se obtienen números cada vez más pequeños. 

Es importante tener esto en cuenta a la hora de elegir una t-norma puesto que en el caso de reglas con muchos antecedentes se podrían producir problemas de precisión al trabajar con números tan pequeños. Un caso similar es el de la t-norma de \L{}ukasiewicz, que no se puede aplicar de forma efectiva a reglas con más de dos antecedentes, puesto que el resultado tiende a cero muy rápidamente.


\chapter{Implementación en MATLAB}
En este capítulo se incluye el código fuente de la implementación en MATLAB de los métodos estudiados en este proyecto.

\section{General}
En esta sección se incluyen algunas funciones de propósito general que han sido implementadas en el desarrollo de este proyecto.

\subsection{T-normas y operadores de agregación}
MATLAB dispone de una colección importante de funciones predefinidas para calcular t-normas y operadores de agregación. En este proyecto se han utilizado las siguientes:

\begin{itemize}
\item Media aritmética: \lstinline|mean()|.
\item Máximo: \lstinline|max()|.
\item Producto: \lstinline|prod()|.
\item Mínimo: \lstinline|min()|.
\item Media geométrica: \lstinline|geomean()|.
\item Media armónica: \lstinline|harmmean()|.
\end{itemize}

Además se han implementado sendas funciones para calcular las t-normas del seno y de Einstein:

\lstinputlisting[label=lst:sinmean, caption=T-norma del seno (\lstinline|functions/sinmean.m|),inputencoding=cp1252]{./matlab/functions/sinmean.m}

\lstinputlisting[label=lst:einsteinmean, caption=T-norma de Einstein (\lstinline|functions/einsteinmean.m|),inputencoding=cp1252]{./matlab/functions/einsteinmean.m}

Como se puede ver, en la función \lstinline|einsteinmean()| se ha utilizado una función anónima (\lstinline|@(x)(1-x)|), que es una funcionalidad muy útil de MATLAB, inspirada en los lenguajes funcionales, que permite declarar funciones \emph{lambda} o \emph{anónimas} (sin nombre) y utilizarlas como argumentos de otras funciones (en este caso de \lstinline|arrayfun()|, que aplica dicha función anónima a todos los elementos de un vector). Esta característica del lenguaje de MATLAB se utilizará también a la hora de definir índices de solapamiento.

\subsection{Construcción de índices de solapamiento}
Los índices de solapamiento (\ref{sec:overlap-indexes}) juegan un papel crucial en el desarrollo de este proyecto, al ser la base del método de interpolación estudiado (\ref{sec:interpolation-method}). En la sección \ref{sec:overlap-index-construction} se ha presentado un método para construir índices de solapamiento a partir de operadores de agregación y funciones de solapamiento. 

Este mismo concepto de ``construir'' índices de solapamiento a partir de otras funciones se puede implementar en MATLAB. Para ello se va a utilizar una técnica de programación llamada \emph{clausura}\footnote{Conocida generalmente por su término inglés ``closure''.}, que junto con las funciones anónimas, es muy utilizada en los lenguajes funcionales. A grandes rasgos, una clausura es una función de primer orden (puede tratarse como cualquier otra variable de tipo básico) que mantiene el estado de las variables del ámbito en el que fue creada. Así pues, el objetivo es definir una función que tome como parámetros una función de agregación y una de solapamiento y devuelva una nueva función que sea el índice de solapamiento construido a partir de estas funciones recibidas como parámetros.

Para implementar la construcción de índices de solapamiento se ha definido la función \lstinline|make_overlap_index()| (\lstinline|functions/make_overlap_index.m|), que toma como parámetros una función de agregación \lstinline|M| y una de solapamiento \lstinline|Go| y devuelve una nueva función que es el índice de solapamiento construido a partir de \lstinline|M| y 
\lstinline|Go| (según el teorema \ref{sec:overlap-index-construction}):

\lstinputlisting[label=lst:make_overlap_index,caption=Construcción de índices de solapamiento (\lstinline|functions/make_overlap_index.m|).,inputencoding=cp1252]{./matlab/functions/make_overlap_index.m}

El resultado de llamar a esta función es una nueva función que calcula el índice de solapamiento entre dos conjuntos pasados como parámetros \lstinline|A| y \lstinline|B|. Por ejemplo, si queremos construir el índice de solapamiento de Zadeh $O_Z$ (ecuación \ref{eq:zadeh-consistency}) debemos pasar a la función \lstinline|make_overlap_index()| un puntero a la función máximo (\lstinline|@max|) como función de agregación y un puntero a la función mínimo (\lstinline|@min|) como función de solapamiento:

\begin{lstlisting}
Oz = make_overlap_index(@max,@min);
\end{lstlisting}

La función \lstinline|Oz()| ahora puede utilizarse para calcular el índice de solapamiento entre dos conjuntos \lstinline|A| y \lstinline|B| pasados como parámetros:

\begin{lstlisting}
overlap_index = Oz(A, B);
\end{lstlisting}
 
La construcción de los índices de solapamiento utilizados en la comparativa entre el método de interpolación y Mamdani (sec. \ref{sec:fire-detection-interpolation-method}) se ha implementado en la clase \lstinline|O| (\lstinline|functions/O.m|) utilizando métodos estáticos:

\lstinputlisting[label=lst:overlap_indexes,caption=Índices de solapamiento (\lstinline|functions/O.m|).,inputencoding=cp1252]{./matlab/functions/O.m}

De esta forma se puede obtener el índice de solapamiento ya construido llamando a su función correspondiente. Por ejemplo, para obtener la consistencia de Zadeh:

\begin{lstlisting}
Oz = O.maxmin();
\end{lstlisting}

\section{Sistemas difusos basados en reglas}
En esta sección se incluyen las funciones implementadas relacionadas con los sistemas difusos basados en reglas estudiados a lo largo del proyecto.

\subsection{Conjuntos de reglas}\label{sec:matlab-rule-base}
Los conjuntos de reglas juegan un papel fundamental en los sistemas difusos basados en ellas, dado que expresan el conocimiento sobre el dominio del problema y tienen la misión de definir cómo se transforman las entradas del sistema difuso en salidas. Los métodos de Mamdani e interpolación implementados en este proyecto reciben el conjunto de reglas como un parámetro, de forma que pueden ser utilizados con cualquier conjunto de reglas y aplicados a cualquier tipo de problema.

Existen multitud de formas de implementar el conjunto de reglas. En este caso se ha optado por utilizar un vector de estructuras, debido a que el código resultante es sencillo de entender y de modificar. Una estructura en MATLAB es un tipo de variable compuesta que puede tener un número arbitrario de campos con nombre (similar a las estructuras de lenguajes como C). En el caso del conjunto de reglas, cada regla es una estructura \lstinline|R| con los siguientes campos:

\begin{itemize}
\item \lstinline|R.n|: Nº de regla.
\item \lstinline|R.A|: Vector con las funciones de pertenencia de los valores (punteros a funciones) de cada uno de los antecedentes de la regla.
\item \lstinline|R.B|: Función de pertenencia del valor del consecuente.
\end{itemize}
 
Por ejemplo, para definir una regla para el caso de la detección de incendios forestales tal que:
\begin{multline}
\text{IF \emph{Temperatura} es ``Baja'' AND }  \text{ \emph{Humo} es ``Alto'' AND } \text{ \emph{Luz} es ``Baja''} \\
\text{AND \emph{Humedad} es ``Alta'' AND }  \text{ \emph{Distancia} es ``Media'' THEN }  \text{ \emph{Riesgo} es ``Bajo'' }
\end{multline}
hay que construir una estructura de la siguiente forma:
 
\begin{lstlisting}
R.n = 1;
R.A = {@temp.low, @smoke.high, @llight.low, @humidity.high, @distance.medium};
R.B = @threat.low;
\end{lstlisting}
 
Para definir un conjunto con varias reglas basta con añadir índices:
 
\begin{lstlisting}
R(1).n = 1;
R(1).A = {@temp.low, @smoke.high, @llight.low, @humidity.high, @distance.medium};
R(1).B = @threat.low;

R(2).n = 2;
R(2).A = {@temp.low, @smoke.high, @llight.low, @humidity.high, @distance.high};
R(2).B = @threat.low;
\end{lstlisting}

La variable \lstinline|R| resultante es un vector de estructuras y se puede obtener la regla \emph{i-ésima} haciendo uso del índice (como cualquier otro vector): \lstinline|R(i)|.

\subsection{Difusificadores}
Otro componente importante de los sistemas difusos basados en reglas son los difusificadores. Un difusificador es una función que transforma una entrada escalar del sistema en un conjunto difuso. Existen diferente tipos de difusificadores, tal y como se ha estudiado en la sección \ref{sec:difusificadores}. En este proyecto, sin embargo, sólo se ha utilizado el difusificador \emph{singleton} cuya implementación en MATLAB es la siguiente:

\lstinputlisting[caption=Difusificador singleton (\lstinline|functions/singleton_fuzzifier.m|),inputencoding=cp1252]{./matlab/functions/singleton_fuzzifier.m}

\subsection{Desdifusificadores}
El último componente de un sistema difuso basado en reglas es el \emph{desdifusificador}. Un desdifusificador es una función que transforma el conjunto difuso de salida en un valor escalar. En la sección \ref{sec:desdifusificadores} se han establecido las propiedades que debe tener un desdifusificador y se han definido algunos de los más utilizados.

MATLAB, en su toolbox \emph{Fuzzy Logic,} ya tiene implementados los desdifusificadores más utilizados por medio de la función \lstinline|defuzz(x, mf, type)|. Esta función toma como parámetros el universo de referencia del conjunto \lstinline|x|, los grados de pertenencia para los valores del universo de referencia \lstinline|mf| y el tipo de desdifusificador a utilizar \lstinline|type|. Los tipos de desdifusificadores posibles son:

\begin{itemize}
\item \lstinline|centroid|: Centroide.
\item \lstinline|bisector|: Bisector.
\item \lstinline|mom|: Media de máximos.
\item \lstinline|som|: Menor de máximos.
\item \lstinline|lom|: Mayor de máximos.
\end{itemize}


\subsection{Método de Mamdani}
El primero de los métodos estudiados en este proyecto es el método de Mamdani (algoritmo \ref{algo:mamdani}). La implementación de este método es una función que toma como parámetros:
\begin{itemize}
\item \lstinline|R|: Conjunto de reglas tal y como se ha definido en la sección \ref{sec:matlab-rule-base}.
\item \lstinline|fact|: Premisa o valores escalares de entrada al sistema. Un vector con tantas posiciones como antecedentes tienen las reglas (variables lingüísticas), de forma que el valor para el antecedente \emph{i-ésimo} debe estar en \lstinline|fact(i)|.
\item \lstinline|y|: Universo de salida (vector).
\end{itemize}

\lstinputlisting[caption=Método de Mamdani (\lstinline|functions/mamdani.m|),inputencoding=cp1252]{./matlab/functions/mamdani.m}

\subsection{Método de interpolación basado en índices de solapamiento}
El método más importante estudiado en este proyecto es el método de interpolación basado en índices de solapamiento (algoritmo \ref{algo:multi-overlap-interpolation-method}). Este método se ha implementado en la función \lstinline|interpolation(R, fact, y, O, T, M)| (\lstinline|functions/interpolation.m|) que toma como parámetros:

\begin{itemize}
\item \lstinline|R|: Conjunto de reglas tal y como se ha definido en la sección \ref{sec:matlab-rule-base}.
\item \lstinline|fact|: Premisa. Vector de conjuntos difusos con tantas posiciones como antecedentes tienen las reglas (variables lingüísticas), de forma que el valor para el antecedente \emph{i-ésimo} debe estar en \lstinline|fact(i)|.
\item \lstinline|y|: Universo de salida (vector).
\item \lstinline|O|: Índice de solapamiento. Función que puede construirse utilizando el método implementado en \ref{lst:make_overlap_index}.
\item \lstinline|T|: T-norma. Utilizada para calcular el grado de activación de cada regla agregando los índices de solapamiento de cada antecedente.
\item \lstinline|M|: Operador de agregación.
\end{itemize}

El método toma estos parámetros y devuelve el conjunto difuso de salida, resultado de aplicar el conjunto de reglas ( \lstinline|R|) a la entrada (\lstinline|fact|).

\lstinputlisting[label=lst:interpolation,caption=Método de interpolación basado en índices de solapamiento (\lstinline|functions/interpolation.m|),inputencoding=cp1252]{./matlab/functions/interpolation.m}

Como se puede, ver en la línea 23 del código \ref{lst:interpolation} se utiliza una función llamada \lstinline|matching_degree()| (\lstinline|functions/matching_degree.m|) para calcular el grado de activación de cada regla. Esta función toma como parámetros:

\begin{itemize}
\item \lstinline|R|: Regla.
\item \lstinline|fact|: Premisa.
\item \lstinline|O|: Índice de solapamiento.
\item \lstinline|T|: T-norma.
\end{itemize}

Esta función calcula el índice de solapamiento (\lstinline|O|) de cada antecedente de la regla \lstinline|R| con su correspondiente valor en la premisa \lstinline|fact|  y agrega los valores obtenidos utilizando la t-norma \lstinline|T|:

\lstinputlisting[label=lst:matching_degree,caption=Cálculo del grado de activación basado en índices de solapamiento (\lstinline|functions/matching_degree.m|),inputencoding=cp1252]{./matlab/functions/matching_degree.m}

Estos métodos, a pesar de ser correctos y devolver los resultados esperados son mejorables en lo que a rendimiento se refiere. El principal problema de esta implementación es que se recalculan los conjuntos difusos de los antecedentes y el consecuente en cada regla (línea 27 en \ref{lst:interpolation} y línea 16 en \ref{lst:matching_degree}). Este cálculo sólo es necesario hacerlo la primera vez para cada valor de cada variable lingüística y después reutilizarlo.

Por ello se ha implementado una nueva versión de estos métodos que utiliza una caché de conjuntos difusos ya calculados, que se almacenan en una especie de \emph{tabla hash}. La tabla hash se ha implementado utilizando una clausura, de forma que se devuelve una función que mantiene la estructura de la tabla hash en el ámbito y la utiliza en llamadas posteriores. La función \lstinline|fuzzyset_hash()| se ha implementado en \lstinline|functions/fuzzyset_hash.m| de la siguiente manera:

\lstinputlisting[label=lst:fuzzyset_hash,caption=Función para construir una tabla hash de conjuntos difusos,inputencoding=cp1252]{./matlab/functions/fuzzyset_hash.m}

El resultado de llamar a esta función es un puntero a una nueva función (\lstinline|get_fuzzy_set(mf, x)|) que calcula los grados de pertenencia de un conjunto difuso dados su función de pertenencia \lstinline|mf| y su universo de referencia \lstinline|x|. En caso de que este cálculo haya sido realizado previamente, devuelve el resultado anterior, para evitar trabajo innecesario. 

Esta función hace uso de la estructura \lstinline|hash| (línea 12), que es una variable que ha quedado en su ámbito al crearse la función. Esta variable permite guardar los cálculos realizados en llamadas anteriores a la función y mantener el estado de la tabla hash mientras la función exista. La tabla hash de conjuntos se indexa utilizando el nombre de la función de pertenencia (obtenido mediante la función \lstinline|func2str()|).

Una vez implementada esta función, se puede utilizar en los métodos \lstinline|interpolation()| y \lstinline|matching_degree()| para ganar rendimiento:

\lstinputlisting[label=lst:hashed_interpolation,caption=Método de interpolación basado en índices de solapamiento (\lstinline|functions/hashed_interpolation.m|),inputencoding=cp1252]{./matlab/functions/hashed_interpolation.m}

\lstinputlisting[label=lst:hashed_matching_degree,caption=Cálculo del grado de activación basado en índices de solapamiento (\lstinline|functions/hashed_matching_degree.m|),inputencoding=cp1252]{./matlab/functions/hashed_matching_degree.m}

A pesar de que esta función es muy simple, la ganancia de rendimiento al utilizar esta técnica es notable.
 
\section{Detección de incendios forestales}
\subsection{Variables lingüísticas}
Como se ha definido en \ref{def:formal-lang-variable}, una variable lingüística es una variable que puede tomar conjuntos difusos como valores. Así mismo, un conjunto difuso viene definido por su función de pertenencia. Por tanto, se debe definir una función para cada valor de cada variable lingüística utilizada. 

Generalmente, en MATLAB, cada función se define en un fichero separado que únicamente contiene el código de dicha función. En este caso y dado que se utilizan 5 variables lingüísticas para los antecedentes con 3 valores posibles cada una y una variable de salida para el consecuente con 5 valores posibles, habría que crear 20 ficheros para definir estas funciones. Aunque es perfectamente posible y válido hacerlo de esta manera, resulta tedioso manejar tantos ficheros. 

Una solución adecuada al problema es definir cada variable lingüística en un fichero en el que se implementan las funciones de pertenencia de sus posibles valores. Para ello se puede crear una clase para cada variable lingüística y definir la función de pertenencia de cada valor como un método estático de la misma. Esta forma de definir las variables lingüísticas tiene la ventaja de que todo el código asociado a la variable está en un mismo fichero. Además la forma de utilizar las funciones resulta muy cómoda y expresiva.

Por ejemplo, para la variable lingüística $\chi_1$ = \emph{Temperatura} se ha definido la clase \lstinline|temp| en el fichero \lstinline|lang_variables/temp.m|. Esta clase tiene los siguientes métodos estáticos:
\begin{itemize}
\item \lstinline|get_x()| : Devuelve el universo de referencia para la variable (en este caso el vector $[0,100]$).
\item \lstinline|low(t)| : Función de pertenencia para el valor ``Bajo''.
\item \lstinline|medium(t)| : Función de pertenencia para el valor ``Medio''.
\item \lstinline|high(t)| : Función de pertenencia para el valor ``Alto''.
\end{itemize}

De esta forma, para obtener el grado de pertenencia de un valor \lstinline|t| al conjunto \emph{``Temperatura Media''} basta con hacer:

\begin{lstlisting}
v = temp.medium(t);
\end{lstlisting}

Además es posible obtener un puntero a estas funciones utilizando el operador \lstinline|@|:

\begin{lstlisting}
fh = @temp.medium;
v = fh(t);
\end{lstlisting}

La variable \lstinline|fh| es un puntero a la función \lstinline|temp.medium()| y puede utilizarse de la misma forma que ésta (línea 2). Además la variable \lstinline|fh| se puede utilizar de forma similar a otros tipos de variables, es decir, se pueden crear vectores de punteros de funciones, pasarlos como parámetros a otras funciones etc. Esta propiedad se utilizará de forma intensiva a la hora de implementar los métodos estudiados en este proyecto y el conjunto de reglas.

A continuación se incluyen las definiciones de las variables lingüísticas utilizadas en la aplicación práctica de detección de incendios forestales:

\lstinputlisting[caption=$\chi_1$ - Temperatura (\lstinline|lang_variables/temp.m|),inputencoding=cp1252]{./matlab/lang_variables/temp.m}

\lstinputlisting[caption=$\chi_2$ - Humo (\lstinline|lang_variables/smoke.m|), inputencoding=cp1252]{./matlab/lang_variables/smoke.m}

\lstinputlisting[caption=$\chi_3$ - Luz (\lstinline|lang_variables/llight.m|),inputencoding=cp1252]{./matlab/lang_variables/llight.m}

\lstinputlisting[caption=$\chi_4$ - Humedad (\lstinline|lang_variables/humidity.m|),inputencoding=cp1252]{./matlab/lang_variables/humidity.m}

\lstinputlisting[caption=$\chi_5$ - Distancia (\lstinline|lang_variables/distance.m|),inputencoding=cp1252]{./matlab/lang_variables/distance.m}

\lstinputlisting[caption=$y$ - Riesgo (\lstinline|lang_variables/threat.m|),inputencoding=cp1252]{./matlab/lang_variables/threat.m}

La función \lstinline|lineal()| es una función muy simple que permite calcular el valor de una función de pertenencia lineal entre los puntos pasados en los dos primeros parámetros. Esta función toma los siguientes parámetros:

\begin{itemize}
\item \lstinline|first|: Primer punto del rango (eje x).
\item \lstinline|last|: Último punto del rango (eje x).
\item \lstinline|value|: Punto en el que se quiere calcular el valor de la función de pertenencia (eje x).
\item \lstinline|r|: Parámetro opcional. Si tiene valor 'r' se calcula la función de pertenencia ``inversa''.
\end{itemize}

Este método construye la función de pertenencia lineal que pasa entre los puntos (\lstinline|first|, 0) y (\lstinline|last|, 1) o (\lstinline|first|, 1) y (\lstinline|last|, 0) en caso de que se pase el último parámetro con valor \lstinline|'r'| y obtiene el grado de pertenencia del punto \lstinline|value| para esa función. En caso de que el valor a calcular esté fuera del rango [\lstinline|first|, \lstinline|last|], el grado de pertenencia devuelto es 0:

\lstinputlisting[caption=Función de pertenencia lineal (\lstinline|functions/lineal.m|),inputencoding=cp1252]{./matlab/functions/lineal.m}

\subsection{Conjunto de reglas}

El conjunto de reglas se ha definido en la función \lstinline|fire_detection_rules()| (\lstinline|fire_detection_rules.m|), de forma que sea sencillo obtener y utilizar este conjunto en otros scripts y funciones. Esta función devuelve un vector de estructuras tal y como se ha definido en la sección \ref{sec:matlab-rule-base}.

Este tipo de estructura tiene la ventaja de que es sencilla de implementar y de entender. Sin embargo, para conjuntos de reglas grandes (como es el caso) resulta tedioso tener que indicar el índice para cada regla y repetir una y otra vez el nombre del conjunto (\lstinline|R|) o de los campos. 

Por esta razón se ha optado por una forma de definir el conjunto de reglas más concisa y clara. El conjunto de reglas se define como una matriz, en la que cada fila corresponde a una regla y cada columna es un puntero a la función de pertenencia del valor de la variable lingüística que ocupa esa posición. Las primeras $n-1$ columnas de cada fila corresponden a los antecedentes de la regla, y la última al consecuente.

En la función  \lstinline|fire_detection_rules| (código \ref{matlab_fire_detection_rules}) se define el conjunto de reglas en la matriz \lstinline|rules| (líneas 2-335) de esta forma. Dado que el resto de funciones implementadas esperan un conjunto de reglas definido como un vector de estructuras, se transforma esta matriz en dicho vector en las líneas 340-344. De esta forma, el conjunto de reglas es muy sencillo de leer y modificar y se asemeja mucho a la tabla \ref{tab:fire-detection-rule-set}.

\lstset{linewidth=18cm}
\lstinputlisting[label=matlab_fire_detection_rules, caption=Conjunto de reglas (\lstinline|fire_detection_rules.m|) ,inputencoding=cp1252]{./matlab/fire_detection_rules.m}
\lstset{linewidth=\textwidth}
\subsection{Método de Mamdani aplicado a la detección de incendios forestales}
Con el código implementado hasta ahora ya es posible aplicar el método de Mamdani a la detección y determinación de riesgos de incendios forestales. Para ello se ha definido la función \lstinline|fire_detection_mamdani(temp, smoke, light, humidity, distance)| que recibe como parámetros las entradas escalares de temperatura, humo , etc. y devuelve el riesgo de incendio resultante al aplicar el método de Mamdani. Esta función además devuelve los valores desdifusificados de este conjunto de salida:

\lstinputlisting[caption=Método de Mamdani aplicado a la determinación de riesgo de incendios forestales (\lstinline|fire_detection_mamdani.m|),inputencoding=cp1252]{./matlab/fire_detection_mamdani.m}

Utilizando esta función se puede definir una nueva función \lstinline|fire_detection_mamdani_plot()| que acepta los mismos argumentos (y algunos parámetros opcionales más) y que además de aplicar el método de Mamdani sobre las entradas, genera una gráfica con los resultados obtenidos:

\lstinputlisting[caption=Método de Mamdani y gráfica con resultados (\lstinline|fire_detection_mamdani_plot.m|),inputencoding=cp1252]{./matlab/fire_detection_mamdani_plot.m}

Este método acepta los siguientes parámetros opcionales:

\begin{itemize}
\item \lstinline|'showTitle'| (true/false): Mostrar el título en la gráfica generada.
\item \lstinline|'showLegend'| (true/false): Mostrar la leyenda en la gráfica generada.
\item \lstinline|'exportTikz'| (true/false): Exportar la gráfica generada a tikz.
\end{itemize} 

Para obtener y validar estos parámetros opcionales se utiliza un parser de entradas, que es el método estándar de MATLAB. Este parser es creado por la función \lstinline|fire_detection_plot_input_parser()| (utilizada en la línea 26):

\lstinputlisting[label=lst:fire_detection_plot_input_parser, caption=Parser de parámetros opcionales (\lstinline|fire_detection_plot_input_parser.m|),inputencoding=cp1252]{./matlab/fire_detection_plot_input_parser.m}

Esta función define el parser y los parámetros opcionales que acepta, así como sus valores por defecto. 

La función \lstinline|fire_detection_mamdani_plot()| se puede utilizar por ejemplo para generar las gráficas de la figura \ref{fig:mamdani-fire-detection-example} ejecutando el script \lstinline|fire_detection_mamdani_examples.m|:

\lstinputlisting[caption=Ejemplos del método de Mamdani (\lstinline|fire_detection_mamdani_examples.m|),inputencoding=cp1252]{./matlab/fire_detection_mamdani_examples.m}

\subsection{Método de interpolación basado en índices de solapamiento aplicado a la detección de incendios forestales}
Con las funciones implementadas también se puede aplicar el método de interpolación basado en índices de solapamiento a la determinación de riesgo de incendios forestales. Para ello se ha definido la función \lstinline|fire_detection_interpolation(O, T, M, t, s, l, h, d)|, similar a la ya definida para Mamdani. Esta función, además de las entradas escalares de temperatura, humo, etc. (\lstinline|t, s, l, h, d|) recibe el índice de solapamiento \lstinline|O|, la t-norma \lstinline|T| y la función de agregación \lstinline|M| utilizadas en el método. Al igual que el método de Mamdani, esta función devuelve el valor de riesgo como un conjunto difuso y sus correspondientes valores desdifusificados:

\lstinputlisting[caption=Método de interpolación basado en índices de solapamiento aplicado a la determinación de riesgo de incendios forestales (\lstinline|fire_detection_interpolation.m|),inputencoding=cp1252]{./matlab/fire_detection_interpolation.m}

Al igual que en el método de Mamdani, se puede definir una función que utilice esta función para generar gráficas con los resultados. Esta función (\lstinline|fire_detection_interpolation_plot()|) se ha implementado en \lstinline|fire_detection_interpolation_plot.m| de la siguiente forma:

\lstinputlisting[caption=Método de interpolación y generación de gráficas con resultados (\lstinline|fire_detection_interpolation_plot.m|),inputencoding=cp1252]{./matlab/fire_detection_interpolation_plot.m}

Los parámetros opcionales de este método son los mismos que en el método de Mamdani ya que, de hecho, utiliza el mismo parser de entradas (código \ref{lst:fire_detection_plot_input_parser}). Esta función se puede utilizar por ejemplo para generar las gráficas de la figura \ref{fig:interpolation-fire-detection-example} ejecutando el script:

\lstinputlisting[caption=Ejemplos del método de interpolación (\lstinline|fire_detection_interpolation_examples.m|),inputencoding=cp1252]{./matlab/fire_detection_interpolation_examples.m}

\subsection{Comparativa de resultados obtenidos con el método de interpolación}
En la sección \ref{sec:fire-detection-interpolation-method} se ha presentado una comparativa de los resultados obtenidos al aplicar el método de interpolación basado en índices de solapamiento a unas entradas concretas utilizando diferentes combinaciones de t-normas e índices de solapamiento.

Para obtener estos resultados se ha implementado la función  \lstinline|fire_detection_interpolation_comparison(t, s, l, h, d)| que aplica el método de interpolación con cada combinación de t-normas e índices de solapamiento  a las entradas \lstinline|t, s, l, h, d|:

\lstinputlisting[caption=Comparativa de resultados obtenidos con el método de interpolación (\lstinline|fire_detection_interpolation_comparison.m|),inputencoding=cp1252]{./matlab/fire_detection_interpolation_comparison.m}

Esta función genera también gráficas con los resultados obtenidos, que se imprimen por pantalla y en ficheros .tikz que se pueden importar después en \LaTeX . Además vuelca los resultados en formato de tabla \LaTeX, tal y como se puede ver en la sección \ref{sec:fire-detection-interpolation-method}. De hecho, para generar esos resultados es tan sencillo como ejecutar la siguiente orden desde la línea de comandos:

\begin{lstlisting}
fire_detection_interpolation_comparison(30, 20, 500, 50, 40);
\end{lstlisting}


\chapter{Conclusiones}
\label{cha:conclusiones}
\begin{tabbing}
AAAA\=AAAA\=AAAA\kill\\
\>que\\
\>\>tal?\\
\end{tabbing}
\begin{theorem}
If \emph{proposition}\index{proposition|textbf} $P$ is a tautology then $\sim P$ is a contradiction, and conversely.
\end{theorem}
\begin{theorem}[Tautologies and Contradictions]
content...
\end{theorem}

Esto es una cita \cite{bolourchi2013}




\backmatter

% Bibliografía
\bibliographystyle{ieeetr}
\bibliography{bibliography}

% Glosario de términos
\printindex

\end{document}
